\documentclass[11pt]{article}

\usepackage[utf8]{inputenc}
% microtype 의 폰트 확장은 scalable 폰트를 요구한다. 기본 CM 이 비트맵(Type3)으로
% 잡히는 배포판에서는 lmodern 없이는 "auto expansion is only possible with
% scalable fonts" 로 컴파일이 죽는다.
\usepackage{lmodern}
\usepackage[T1]{fontenc}
\usepackage[margin=1in]{geometry}
\usepackage{amsmath,amssymb}
\usepackage{graphicx}
\usepackage{longtable,booktabs}
\usepackage{microtype}
% svgnames 를 줘야 teal 이 정의된다. 없으면 \cite 하나마다
% "Undefined color `teal'" 이 뜬다 (본편 산출물에서 732회).
\usepackage[svgnames]{xcolor}
% hyperref goes last; it makes \cite clickable through to the bibliography.
\usepackage[colorlinks=true,linkcolor=blue,citecolor=teal,urlcolor=magenta]{hyperref}

% pandoc emits \tightlist but expects the preamble to define it.
\providecommand{\tightlist}{%
  \setlength{\itemsep}{0pt}\setlength{\parskip}{0pt}}

% If the compile times out on Overleaf, comment out \tableofcontents below first.
\title{A Comprehensive Survey of In-Context Learning: Mechanisms, Methodologies, and Multimodal Frontiers}
\author{Generated by AutoSurvey}
\date{}

\begin{document}
\maketitle
\tableofcontents
\newpage

\section{Introduction and Foundations of In-Context Learning}

\subsection{Defining In-Context Learning: The Emergent Paradigm}

In-context learning (ICL) represents a fundamental shift in how artificial intelligence models, particularly large language models (LLMs), adapt to and perform new tasks, a shift that, as detailed in the following historical account, was largely precipitated by the scaling of model and data size. Formally, ICL is defined as a learning paradigm where a model conditions on a prompt, comprising a task specification and optionally a set of input-output demonstrations, to directly generate a desired output for a novel query without any gradient-based updates to its parameters. This process is entirely \emph{inference-time}: the model's weights remain frozen, and all ``learning'' is transient, existing only within the activations of the forward pass for a given prompt. This distinguishes it sharply from the ``pre-train and fine-tune'' paradigm that dominated the earlier NLP landscape, where learning is a training-time process that updates model weights. The core components of an ICL prompt decompose into three elements: a natural language instruction (optional but common), a set of \(k\) demonstrations \(\{(x_i, y_i)\}_{i=1}^k\), and a query input \(x_{query}\) for which the model predicts the output \(\hat{y}_{query}\). The model implicitly infers a latent function \(f: \mathcal{X} \to \mathcal{Y}\) mapping inputs to labels entirely from the structure and patterns presented in the few-shot demonstrations \cite{ref1,ref2}.

To establish precise terminology for this survey, it is crucial to delineate ICL from closely related and often conflated concepts: zero-shot transfer and few-shot prompting. Zero-shot transfer, in its purest form, involves a model performing a task solely based on a task instruction without any provided demonstrations. The model relies entirely on its parametric knowledge acquired during pretraining to interpret the instruction and generate an answer. ICL, in its canonical few-shot form, augments this instruction with a set of demonstrations, enabling the model to ``learn'' the task-specific input-label mapping from the prompt context. Drawing a sharp line, few-shot prompting is the \emph{mechanism} or technique of providing demonstrations, whereas ICL is the overarching \emph{paradigm} of learning from that context. This distinction became critical as the field evolved beyond the initial few-shot findings of GPT-3 into more sophisticated techniques. For instance, the concepts of ``zero-shot ICL'' have emerged, where models like Self-ICL \cite{ref3} or Universal Self-Adaptive Prompting \cite{ref4} generate their own pseudo-demonstrations, bootstrapping a few-shot learning signal from a zero-shot prompt. Similarly, SEC \cite{ref5} prompts LLMs to first internally construct examples before solving the final query. These methodologies demonstrate that the defining trait of ICL is not the \emph{source} of the demonstrations (human-curated, retrieved, or self-generated) but the \emph{mechanism} of conditioning on example-based input-output structure to adapt behavior without parameter updates.

The emergent nature of ICL is foundational to its definition. It is a capability that was not explicitly architected into models but rather arose unpredictably with scale in sufficiently large transformers trained on diverse, open-domain text corpora---a phenomenon that catalyzed the entire research trajectory detailed in the subsequent historical subsection. The survey posits that ICL is not a monolithic ability but can be conceptually decomposed into two synergistic sub-processes, a framework formalized in recent analyses \cite{ref6}: \emph{skill recognition} and \emph{skill learning}. Skill recognition refers to the model's ability to identify a known task from its parametric memory based on the prompt's demonstrations, effectively matching the pattern to a pre-learned function. Skill learning, in contrast, is the more profound and less understood ability to dynamically infer a new, unseen input-label mapping and apply it to a query, potentially learning a new function on the fly. This distinction is critical for analyzing ICL's failure modes, which can stem from either an inability to recognize a task or an inability to learn its novel mapping from context \cite{ref7}.

The learning signal in ICL is fundamentally structured by the demonstrations, and the model's objective is to analyze the relationship between the inputs \((x)\) and their labels \((y)\) embedded in this sequence. This process has been rigorously formulated as an algorithm learning problem, where a transformer implicitly constructs a hypothesis function \(h\) at inference time to predict \(y = h(x)\) \cite{ref2}. The effectiveness of ICL is thus measured by how well the model's implicit hypothesis generalizes from the few-shot examples to the novel query. This framing establishes a key piece of terminology for our discussion: the ``input-label mapping'' is the conceptual target of ICL. A model's success in ICL is contingent on accurately capturing this mapping, a process that can be disrupted by biases in the demonstrations. For example, Comparable Demonstrations (CDs), which minimally edit text to flip labels, have been shown to mitigate spurious correlations and better highlight the task's essential mapping \cite{ref8}. This underscores that the nature of the demonstrations directly governs the quality of the transient function learned in context, a principle that underpins the many-shot and prompt-engineering innovations chronicled in the following historical evolution.

Furthermore, the formal definition of ICL must account for the mechanism's implicit connection to more traditional optimization, despite the absence of weight updates. A compelling theoretical perspective, which we explore in depth in Section 2, posits that ICL can be interpreted as performing implicit fine-tuning or implicit gradient descent. This view is supported by empirical findings showing that the changes in a model's hidden states during an ICL forward pass closely resemble those that would occur if the demonstrations were used to explicitly update the model's parameters via instruction tuning \cite{ref9}. In this sense, ICL acts as a ``shortcut'' for a gradient-based learning step, a dual formulation that has been the subject of intense theoretical scrutiny \cite{ref10}. This perspective solidifies ICL's status not just as a prompting trick but as a genuine, if implicit, learning algorithm operationalized through the mechanics of a transformer architecture. Consequently, in this survey, the term ``in-context learner'' refers to any model capable of this inference-time adaptation, and ``demonstrations'' specifically denote the (x, y) pairs that constitute the primary learning signal within the prompt's context. This definitional groundwork sets the stage for tracing how this emergent paradigm was historically discovered, scaled, and extended across modalities and input regimes, as explored in the next subsection.

\subsection{Historical Origins and Evolution}

Having established a precise definition of ICL as an emergent, inference-time learning paradigm distinct from both zero-shot and fine-tuning approaches, we now trace its historical trajectory. This evolution is intrinsically linked to the scaling of language models and the serendipitous discovery of their emergent few-shot capabilities, a journey that transformed ICL from a curious observation into the universal adaptation interface formalized in the preceding subsection. Before the modern formulation of ICL, transfer learning in natural language processing (NLP) was dominated by the ``pre-train and fine-tune'' paradigm, where models like BERT were adapted to downstream tasks via gradient updates on task-specific data. The conceptual pivot toward learning from input context alone, without any parameter modification, can be traced back to early meta-learning ideas and phrasing-based prompting, but it was crystallized with the release of GPT-3. The landmark paper demonstrating GPT-3's capabilities empirically established that scaling up model size and data unlocks an ability to rapidly adapt to new tasks using only a handful of input-output examples provided in the prompt. This discovery marked the genesis of ICL as a formal object of study, shifting the focus from explicit fine-tuning to the implicit, transient learning through conditioning that we have defined as the core of the ICL paradigm. GPT-3 showed that, by simply prepending a query with a few demonstrations, the model could perform translation, question answering, and even novel word usage tasks at a competitive level, fundamentally redefining the landscape of few-shot learning and laying the groundwork for the algorithmic and mechanistic interpretations that now dominate the field.

In the immediate aftermath of GPT-3, the evolution of ICL was shaped by two parallel forces: the rapid scaling of model architectures and the methodological refinement of prompting. Subsequent flagship models demonstrated that increased parameter counts and training corpus diversity not only improved task performance but also stabilized the otherwise brittle nature of ICL. These models exhibited more robust abilities to reason over multiple demonstrations and follow complex instructions provided purely in-context, solidifying the distinction between skill recognition and skill learning that is central to our formal decomposition. A critical extension during this period was the introduction of chain-of-thought (CoT) prompting, which integrated intermediate reasoning steps into the few-shot demonstrations. This innovation unlocked advanced multi-step reasoning capabilities in LLMs, a domain where standard ICL with input-output pairs alone often faltered, and demonstrated how the implicit hypothesis function \(h(x)\), introduced in our definition, could be guided through more complex algorithmic steps. The concept of the model as a general-purpose learner, capable of absorbing a task specification on the fly, became a central premise, influencing both the architectural design of subsequent models and the research agendas aimed at understanding the phenomenon.

Crucially, the historical expansion of ICL was not confined to NLP, a broadening that directly informs the multimodal and cross-domain applications we will later explore. The paradigm rapidly transcended textual boundaries, giving rise to a vibrant subfield of multimodal in-context learning (M-ICL). This shift was fueled by the development of Vision-Language Models (VLMs) and Multimodal Large Language Models (MLLMs) such as Flamingo, Kosmos-1 \cite{ref11}, and GPT-4V(ision) \cite{ref12}. These models demonstrated that the ICL capability of large language models could be naturally extended to visual inputs, proving that the core mechanism of conditioning on structured demonstrations was not an artifact of text but a general property of large transformer architectures. Early analyses of GPT-4V probed its ability to process interleaved image and text demonstrations for complex visual reasoning tasks, a stark departure from the unimodal text prompts of the GPT-3 era \cite{ref12}. Kosmos-1 was explicitly trained from scratch on multimodal corpora to learn in context across modalities, evaluating zero-shot, few-shot, and multimodal chain-of-thought prompting, thereby cementing the concept of a general multimodal learner \cite{ref11}. This period saw the emergence of dedicated multimodal ICL benchmarks, such as VL-ICL Bench, which systematically moved beyond simple visual question answering (VQA) and image captioning to test the genuine cross-modal ICL limits of models like GPT-4, revealing strengths and weaknesses in perception-to-reasoning tasks \cite{ref13}. Concurrently, works like MMICL \cite{ref14} proposed novel context schemes and datasets to empower VLMs to handle complex prompts with multiple interlaced images, tackling the intricate challenges of cross-modal interaction that are not present in text-only settings and foreshadowing the formal decomposition of the prompt structure \(\mathcal{P}\) presented in the next subsection.

The evolution of ICL further progressed along a quantitative axis, moving from the few-shot regime, typically involving tens of examples, into the many-shot regime with the advent of models supporting drastically larger context windows. This shift directly addresses the central tension in our definition of ICL between the quality and quantity of the learning signal. The many-shot regime, scaling the demonstration set \(\mathcal{D}_K\) to hundreds or even thousands of examples, yielded significant and consistent performance gains across a wide array of tasks \cite{ref15}. It revealed that with sufficient demonstrations, ICL could override stubborn pretraining biases present in the model and learn high-dimensional functions with numerical inputs---feats that were previously unattainable and which directly test the model's capacity for genuine skill learning from context. To overcome the practical bottleneck of sourcing large quantities of human-generated examples, this era also saw the exploration of Reinforced and Unsupervised ICL, where model-generated rationales or entirely unlabeled domain-specific data could effectively replace hand-crafted demonstrations in the prompt, particularly for complex reasoning tasks \cite{ref15}. This evolution underscored a key insight that connects back to our framework: ICL performance is not solely a function of model scale but is dynamically intertwined with the scale and quality of the prompt's exemplar history, a principle that governs the demonstration selection and ordering strategies we formalize subsequently.

Simultaneously, there was a growing recognition that emergent ICL capabilities were not strictly tied to astronomical model sizes. Investigations into reduced-scale data and models revealed that by simplifying pre-training data to make the ``language'' more coherent and structured, zero-shot and few-shot capabilities could emerge in models with as few as a million parameters, challenging the narrative that ICL is purely a property of scale \cite{ref16}. This finding pointed toward data-centric determinants as equally crucial to the historical emergence and evolution of ICL, offering an early empirical anchor for the theoretical link between pre-training data distribution and the implicit learning algorithms we explore later. In the domain of speech processing, researchers also began initial explorations into transferring the ICL paradigm to audio modalities. Studies showed that while naïve speech LMs lack inherent ICL capability, a ``warmup training'' procedure could induce this emergent property, enabling a speech LM to perform few-shot classification on unseen spoken tasks by conditioning on utterance-label demonstrations, thereby opening the door for auditory in-context adaptation without text supervision \cite{ref17}. The cumulative historical arc thus traces a path from a surprising text-only phenomenon in a single oversized model, through a phase of multi-turn reasoning and cross-modal fusion, into an era of scalable many-shot learning and modality-agnostic adaptation. This trajectory establishes ICL as a universal interface for modern foundation models, whose constituent components---task specification, demonstration engineering, and output generation---we are now poised to formally deconstruct.

\subsection{Formal Problem Formulation and Core Components}

Building upon this historical trajectory, which established ICL as a universal inference-time adaptation interface, we now turn to a formal deconstruction of its operational components. In-context learning (ICL) can be formally conceptualized as a conditional text generation process where a pre-trained model approximates a target function based on a structured input sequence. Let \$ f\_\textbackslash theta \$ denote a large language model (LLM) parameterized by fixed weights \$ \textbackslash theta \$. The core objective of ICL is to predict an output \$ \textbackslash hat\{y\} \$ for a query input \$ x\_\{\textbackslash text\{query\}\} \$, conditioned not only on an optional natural language instruction \$ I \$, but critically on a set of \$ K \$ demonstrations, \$ \textbackslash mathcal\{D\}\emph{K = \{(x\_i, y\_i)\}}\{i=1\}\^{}\{K\} \$, where \$ x\_i \$ is an input example and \$ y\_i \$ is the corresponding ground-truth label or output. The prompt \$ \textbackslash mathcal\{P\} \$, which serves as the sole input to the frozen model, is typically constructed by concatenating these components according to a specific template structure: \$ \textbackslash mathcal\{P\} = I \textbackslash oplus x\_1 \textbackslash oplus y\_1 \textbackslash oplus \textbackslash dots \textbackslash oplus x\_K \textbackslash oplus y\_K \textbackslash oplus x\_\{\textbackslash text\{query\}\} \$. The model then generates the prediction by sampling from the conditional probability distribution \$ \textbackslash hat\{y\} \textbackslash sim f\_\textbackslash theta(\textbackslash cdot \textbackslash mid \textbackslash mathcal\{P\}) \$. This formulation underscores the fundamental departure from gradient-based learning that defines the contrasts explored in the following subsection: the model's behavior is modulated entirely through the information embedded in the prompt's structure and content, rather than through weight updates.

The effectiveness of this frozen forward pass hinges on the precise decomposition and design of its constituent components. The first critical component is the \textbf{task specification}, which can be conveyed either implicitly through the input-output patterns in the demonstrations or explicitly through a natural language instruction. The seminal analysis presented in \cite{ref18} revealed a counter-intuitive property of this specification process: ground-truth input-label mappings are not strictly necessary for effective ICL. By showing that randomly replacing labels in demonstrations barely degrades performance across a range of classification and multi-choice tasks, the authors demonstrated that a significant part of ICL's power comes from the demonstrations specifying the label space, the distribution of the input text, and the overall sequence format. This suggests that the model leverages demonstrations to delimit the output space and to calibrate its understanding of the task structure, a process further deconstructed in \cite{ref7}. This work decomposed ICL performance into three dimensions---label space, format, and discrimination---finding that the primary benefit of demonstrations lies in regulating the label space and format to help LLMs respond using desired label words, while the marginal impact on activating discriminative knowledge from the model's pre-trained priors was surprisingly limited. This finding directly informs the later discussion on why ICL can struggle to override strong pretrained priors when contrasted with explicit fine-tuning or meta-learning.

The second core component is \textbf{demonstration selection and ordering}, which directly constructs the input-label mapping context. The selection of the demonstration set \$ \textbackslash mathcal\{D\}\_K \$ is far from trivial, as ICL performance is notoriously sensitive to the choice of examples. Early work in \cite{ref19} formalized retrieval-based selection, demonstrating that even simple word-overlap similarity measures like BM25 could outperform a random selection of demonstrations. The process of selecting examples can be viewed as estimating the most informative context for a given query, where standard approaches rank individual examples independently. However, this independence assumption can lead to redundancy and a lack of coverage over the task's salient aspects. To address this, \cite{ref20} introduced a set-level optimization perspective, proposing BERTScore-Recall as a metric that selects subsets of demonstrations to better cover the reasoning patterns of the test input. This shifts the formulation from a simple top-K retrieval problem to a more complex set selection problem aimed at maximizing the contextual information accessible to the model, mirroring the broader historical shift toward optimizing prompt quality at scale.

Beyond mere selection, the \emph{ordering} of demonstrations within the prompt \$ \textbackslash mathcal\{P\} \$ introduces a sequential prior that biases model predictions. The autoregressive nature of causal LLMs means that demonstrations at the end of the prompt exert a greater influence on the query prediction due to the attention mask and recency bias. This has motivated strategies to structure the prompt sequence to optimize learning. The concept of a learning curriculum was elegantly applied in \cite{ref21}, which introduced in-context curriculum learning (ICCL). This method posits that ordering demonstrations by gradually increasing complexity---from simpler to harder examples---improves performance, mirroring effective human learning strategies. This reveals that the prompt is not just a set of information but a sequence through which a form of implicit, forward-forward learning occurs, a dynamic that contributes to the layer causality discrepancies between ICL and standard gradient descent explored in the next subsection.

An essential concept that bridges the input-label mapping in the prompt to the model's output is the \textbf{task vector}. A major theoretical advance in understanding the mechanics of this process was provided by \cite{ref22}, which proposed a structural interpretation of the learned function in ICL. The authors showed that ICL can be understood as the model compressing the entire demonstration set \$ \textbackslash mathcal\{D\}\_K \$ into a single, abstract ``task vector'' \$ \textbackslash boldsymbol\{\textbackslash theta\}(\textbackslash mathcal\{D\}\emph{K) \$. In this framework, the transformer's function becomes one where its only inputs are the query \$ x}\{\textbackslash text\{query\}\} \$ and this learned task vector, which is computed from the demonstrations and is used to modulate the transformer's operations to produce the final prediction. This compression hypothesis elegantly maps the ICL process onto a more standard machine learning framework where \$ \textbackslash mathcal\{D\}\_K \$ acts as a training set used to find a best-fitting function, albeit within a compressed, abstract space rather than through explicit gradient descent on the model's weights. This provides a conceptual bridge to the implicit meta-optimization perspective that positions ICL relative to explicit fine-tuning and meta-learning.

Finally, the \textbf{output generation} component involves decoding the model's internal state into a final prediction. This stage is vulnerable to biases introduced during the task specification. A significant issue is the ``Demonstration Shortcut,'' as identified in \cite{ref23}, where LLMs default to their pre-trained semantic priors about the labels rather than genuinely learning the novel input-label relationships presented in the prompt. This miscalibration is a direct failure of the formal ICL objective, where the likelihood \$ f\_\textbackslash theta(y \textbackslash mid x, \textbackslash mathcal\{D\}\_K) \$ is dominated by spurious correlations from pre-training rather than the conditional structure of the demonstrations. To rectify this, calibration methods must be integrated into the output generation process to recalibrate the prediction distribution, aligning it with the intended, context-specific mapping. Furthermore, \cite{ref24} revealed that specific lexical tokens, such as template words and stopwords, act as critical task-encoding tokens whose representations heavily determine task performance, adding another layer of complexity to the output generation mechanism. Collectively, these components---task specification, demonstration selection and ordering, task vector computation, and calibrated output generation---form the formal backbone of the ICL paradigm, outlining a complex inference-time computation graph that is only partially understood, and whose operational characteristics define its unique contrasts with traditional learning paradigms.

\subsection{Contrasting ICL with Traditional Learning Paradigms}

The formal decomposition of ICL in the previous subsection defined a complex inference-time computation graph, where task specification, demonstration design, task vector compression, and calibrated output generation occur entirely within a single frozen forward pass. This purely input-driven mechanism stands in stark opposition to the parameter-updating paradigms that dominate traditional machine learning, and it is this contrast that defines ICL's unique position in the landscape of model adaptation. To fully appreciate the nature and implications of ICL, it is essential to formally contrast it with three closely related but distinct paradigms: traditional fine-tuning, meta-learning, and instruction tuning. The core differentiator lies in the locus and explicitness of the adaptation process, which in turn dictates computational cost, data efficiency, and generalization behavior---properties that directly connect to the emergent challenges of sensitivity, calibration, and bias that are discussed in the following subsection.

The most straightforward contrast is between ICL and \textbf{traditional fine-tuning}. Fine-tuning is an explicit learning paradigm where a pre-trained model's parameters are updated via gradient descent on a task-specific, labeled dataset. This process is resource-intensive, requiring significant computation and memory for backpropagation, and it permanently alters the model's weights for a specialized purpose. ICL, conversely, keeps the model entirely frozen. The ``learning'' signal is provided entirely through the input prompt itself, which prepends a few demonstration examples to the query. A landmark study \cite{ref25} provided a theoretical bridge between these paradigms, framing ICL as a form of \emph{implicit fine-tuning}. The work established a dual form between transformer attention and gradient descent, suggesting that a forward pass during ICL effectively computes meta-gradients from the demonstrations and applies them to the model's original representations to construct an ``ICL model.'' Empirical analyses in this study showed that ICL and explicit fine-tuning behave similarly in multiple respects, such as in how they update the model's output distributions. However, this equivalence is not absolute. Subsequent research revisited this hypothesis and found gaps, particularly in realistic NLP settings. For instance, \cite{ref10} discovered that even untrained models could achieve high ICL-GD similarity scores, and identified a ``Layer Causality'' discrepancy in the flow of information between the two processes. When this layer causality was enforced in a GD-based procedure, similarity metrics improved, suggesting that standard GD optimization does not fully mirror the sequential, layer-by-layer dynamics of a transformer's forward pass during ICL. Further complicating the picture, \cite{ref26} directly probed the ICL-GD hypothesis in natural pre-trained LLMs like LLaMA-7B, finding that ICL and GD exhibit different sensitivities to demonstration order and modify output distributions in inconsistent ways, leading the authors to conclude that a strict equivalence remains an open hypothesis. Therefore, while the implicit gradient descent framework is a powerful conceptual tool, ICL in real-world models is a more complex phenomenon that deviates from standard gradient-based optimization in subtle but critical ways---deviations that contribute directly to the reliability and calibration challenges that await further exploration.

ICL is also deeply intertwined with \textbf{meta-learning}, an older paradigm often described as ``learning to learn.'' In conventional meta-learning, a model is trained over a distribution of tasks to learn an inner-loop algorithm that can rapidly adapt to new tasks with few examples. Gradient-based meta-learners like Model-Agnostic Meta-Learning (MAML) explicitly unroll this inner-loop gradient descent during meta-training, incurring substantial computational cost and memory overhead \cite{ref27}. ICL can be viewed as a highly efficient, emergent form of meta-learning where the inner-loop adaptation is performed \emph{implicitly} through the model's forward pass, a behavior that emerges from pre-training on vast and diverse corpora. This perspective aligns naturally with the task vector formalism introduced in Section 1.3, where the demonstration set \(\mathcal{D}_K\) is compressed into an abstract representation that modulates the model's behavior---a process that mirrors the inner-loop learning step but without any explicit optimization. The connection is clearly articulated in work that explains language models as meta-optimizers \cite{ref25}. From this perspective, the pre-trained LLM serves as a meta-learner that has internalized a generic optimization algorithm, and the in-context demonstrations act as the training data for an instantaneous learning episode. This implicit nature avoids the high memory and computational burdens of explicit meta-learning inner loops, as highlighted by the development of implicit MAML algorithms \cite{ref27}. However, a key distinction remains: explicit meta-learning typically meta-trains a model specifically for the objective of few-shot adaptation, whereas for LLMs, ICL is often an emergent byproduct of language modeling objectives \cite{ref28}. This difference in training objective leads to distinct failure modes that manifest as the bias-related challenges previewed in the next subsection. Models can struggle to override strong task priors ingrained during pre-training, a phenomenon particularly evident in subjective tasks like emotion recognition where prior knowledge can ``ossify'' predictions \cite{ref29}. In a traditional meta-learning setting, one could address this by adjusting the meta-training task distribution, but in emergent ICL, the designer is limited to crafting prompts that can overcome the model's intrinsic biases---a constraint that directly motivates the extensive research on demonstration selection and calibration techniques.

The boundary between ICL and \textbf{instruction tuning} (IT) is more nuanced, as both aim to steer a frozen or partially frozen model, but a series of recent investigations has begun to delineate their mechanistic relationship. Instruction tuning is an explicit form of meta-training where a model is fine-tuned on a diverse collection of tasks described via instructions, with the goal of improving zero-shot generalization to unseen tasks \cite{ref28}. ICL, in contrast, performs task adaptation at inference time without any parameter updates. This distinction directly echoes the formal framework of Section 1.3, where the task specification component can operate through either explicit instructions or implicit demonstration patterns---instruction tuning effectively internalizes the former during training, while ICL must rely on the latter at inference. The relationship can be further characterized by their impact on internal representations. Research exploring this connection \cite{ref9} found that ICL is essentially ``implicit instruction tuning.'' Their experiments on LLaMA models showed that ICL changes a model's hidden states as if the demonstrations were used to instructionally tune the model, and the convergence between the two is contingent on factors like demonstration quality. A more practical comparison is provided by \cite{ref30}, which examined the trade-offs between ICL and parameter-efficient fine-tuning (PEFT), a lightweight cousin of full instruction tuning. Their findings demonstrated that PEFT often offers superior accuracy and dramatically lower computational costs per prediction compared to ICL, which must process all training examples for every inference call. This positions ICL not as a universally superior method, but as a versatile tool with distinct advantages in scenarios requiring zero-weight-change adaptation, such as non-parametric knowledge retrieval or situations where model weights are inaccessible. The interplay is also practical: hybrid methods like ``instruction prompt tuning'' (IPT) combine learned prompt embeddings with natural language ICL demonstrations, revealing that ICL not only provides a task blueprint but can stabilize the high-variance nature of prompt tuning and facilitate positive transfer across tasks \cite{ref31}.

In summary, ICL occupies a unique point in the landscape of model adaptation. It is not mere retrieval, explicit fine-tuning, or formal meta-learning, but an emergent capability best described as \textbf{implicit meta-optimization conditioned on pretrained priors}. Its sole reliance on the forward pass for learning bypasses the gradient computation of fine-tuning, sidesteps the complexity of configuring explicit meta-learning inner-loops, and offers a distinct calibration and performance profile compared to instruction tuning. These contrasts are not merely academic taxonomies; they directly shape the practical challenges that define the frontier of ICL research. Understanding the precise ways in which ICL diverges from its traditional counterparts is crucial for diagnosing why it exhibits extreme sensitivity to prompt design, why it struggles with miscalibration and overconfidence, and why its reliance on pretrained priors can lead to ossified predictions. The following subsection systematically unpacks these emergent challenges, grounding them in the unique mechanistic and operational properties that this contrast has brought into focus.

\subsection{Key Challenges and Open Problems}

The previous section established that in-context learning (ICL) occupies a unique and valuable point in the landscape of model adaptation: a form of emergent implicit meta-optimization that sidesteps the gradient computations of fine-tuning and the design complexity of explicit meta-learning. However, this very nature---a frozen model conditioned entirely on prompt-based exemplars---gives rise to a distinct set of challenges that limit its practical deployment and theoretical understanding. These challenges span issues of stability, reliability, bias, efficiency, and a deep lack of explanatory depth, and they constitute the central motivating problems that the subsequent sections of this survey aim to address.

One of the most prominent and widely-studied consequences of this inference-only adaptation is the extreme sensitivity of ICL to the selection and ordering of demonstrations. Slight variations in the composition of the prompt can lead to wildly fluctuating performance, making ICL appear unreliable in practice \cite{ref32}. This instability undermines the paradigm's promise of low-cost adaptation. The choice of which specific demonstrations are included is a critical factor \cite{ref33,ref34}, and a significant body of work is dedicated to retrieval-based methods to find optimally informative examples \cite{ref19,ref35}. However, the problem extends beyond simple similarity, as the model's existing knowledge and the ambiguity of the test input must be considered \cite{ref36}. Furthermore, even with a fixed set of demonstrations, the \emph{order} in which they are presented can dramatically alter outcomes, a direct manifestation of the lack of parameter-based stabilization. While strategies like in-context curriculum learning attempt to structure this order by progressively increasing demonstration complexity \cite{ref21}, the inherent sensitivity to position remains a core instability that limits the robustness and consistency of ICL systems \cite{ref37}.

A second critical challenge, and one that directly threatens the safe deployment of ICL, is the miscalibration and overconfidence of models. As noted when contrasting ICL with instruction tuning and fine-tuning, the frozen model's outputs can exhibit a distinct calibration profile. Research has shown that ICL frequently disrupts the alignment between a model's confidence and its accuracy, with models exhibiting increased miscalibration, particularly in low-shot settings \cite{ref38}. This issue is exacerbated by advanced techniques like chain-of-thought (CoT) prompting, which, while improving task performance, can produce highly overconfident yet incorrect predictions \cite{ref38}. This creates a challenging trade-off between performance and calibration, where gains in accuracy often come at the cost of reliable confidence estimates \cite{ref39}. The development of recalibration techniques is an active area, but effectively diagnosing and correcting for this overconfidence without degrading the performance benefits of ICL remains an open and pressing problem.

A third major hurdle, closely related to the ossification of priors observed in subjective tasks like emotion recognition discussed in Section 1.4, is the strong and often distorting pull of biases from pre-trained priors. Large language models (LLMs) are not blank slates; they come with powerful pre-existing knowledge that can overshadow the new, specific information provided in the prompt. This leads to a phenomenon where the model relies on the semantic prior of a label rather than the input-label relationship demonstrated in the context, a behavior termed the ``Demonstration Shortcut'' \cite{ref23}. This effect is especially pronounced in tasks where the model's strong and inconsistent priors can ossify its predictions, preventing it from adapting to the nuanced mapping presented in the demonstrations \cite{ref29}. A related manifestation is majority label bias, where the model's predictions are skewed by the imbalanced class distribution in the in-context examples, although robustness to this bias appears to vary significantly across model families and tasks \cite{ref40}. Overcoming these deep-seated priors to enable genuine task learning from context, rather than mere cue-matching triggered by priors, is thus a fundamental challenge for ICL's viability as an adaptation strategy \cite{ref7}.

Practical limitations, particularly concerning context length, present a fourth significant barrier that directly distinguishes the operational constraints of ICL from the weight-updating paradigms of fine-tuning and instruction tuning. The standard transformer architecture imposes a finite context window, which directly limits the number of demonstrations that can be provided. This is a critical bottleneck for tasks requiring extensive exemplars or those involving long documents. The problem is compounded in the many-shot regime and across multimodal settings where inputs like images and video are inherently more data-heavy. Even within the allowed context length, a phenomenon known as the ``lost-in-the-middle'' problem degrades performance when crucial information is placed in the middle of long prompts, a challenge that simple scaling of context windows does not solve. Consequently, developing efficient paradigms for many-shot ICL and methods for context compression without sacrificing critical in-context information are crucial areas of ongoing research.

Finally, a lack of a robust and unified theoretical understanding of ICL remains arguably the most profound open problem. While the previous section described ICL as a form of ``implicit meta-optimization conditioned on pretrained priors,'' the precise mechanisms by which it operates are still actively debated. Several competing and complementary theoretical frameworks have been proposed, including interpreting ICL as implicit gradient descent, Bayesian model averaging over latent concepts, or the operation of specific mechanistic circuits like induction heads. However, a unified theory that reconciles these views and offers predictive power is missing. Key questions persist: How do the properties of pre-training data, such as long-range coherence and token burstiness, govern the emergence of ICL \cite{ref41}? What are the learning dynamics and phase transitions that characterize the ICL capability, and why is it often transient, giving way to in-weights memorization? Without a solid theoretical foundation, the design of more robust, generalizable, and efficient ICL methods remains largely an empirical science, guided by trial-and-error rather than first principles.

In summary, the unique nature of ICL---its reliance on a frozen forward pass for adaptation---directly gives rise to a tightly interconnected web of challenges involving prompt sensitivity, calibration, pre-existing biases, context limitations, and a lack of foundational theory. This survey is structured to provide a comprehensive exploration of these issues. Section 2 delves into the theoretical mechanisms and interpretability efforts aimed at demystifying ICL. Section 3 analyzes how architectural and data distributional choices determine ICL efficacy. Section 4 provides a detailed look at strategies to mitigate instability through optimized demonstration selection and prompt engineering. Sections 5 through 9 then expand this investigation into training paradigms, multimodal frontiers, decision-making applications, and the critical dimensions of robustness and safety, collectively charting the landscape of current solutions and the open frontiers that define the future of in-context learning.

\section{Theoretical Mechanisms and Interpretability of In-Context Learning}

\subsection{Implicit Gradient Descent and Dual Formulations of Attention}

While the Bayesian perspective frames in-context learning (ICL) as an implicit inference process over latent tasks, a complementary and influential theoretical framework posits that the attention mechanism within a Transformer implicitly performs gradient descent on an internal objective function, effectively acting as a meta-optimizer. This perspective provides a compelling mathematical bridge between the forward pass of a static model and the dynamic parameter updates characteristic of explicit learning algorithms, offering a mechanistic counterpart to the Bayesian account's focus on posterior inference. The foundational work in this area draws a formal dual equivalence between multi-head linear attention and gradient descent. This connection is established by showing that the computations involved in linear attention---namely, the formation of value-weighted sums of keys based on query-key similarities---mathematically mirror the updates produced by gradient descent on a linear loss function \cite{ref25}. Under this framework, the forward pass of a Transformer given a prompt with demonstrations can be decomposed into two conceptual stages: first, the model computes ``meta-gradients'' from the demonstration examples, and second, these meta-gradients are applied to the model's foundational weights to build an ICL-tailored model for the query. This duality is not merely a mathematical curiosity; it has been empirically supported by observations that the behaviors of ICL and explicit fine-tuning correlate significantly across multiple real-world tasks, suggesting a deep operational symmetry \cite{ref25}.

Extending this line of inquiry, a contrasting perspective reframes the inference process of ICL not simply as vanilla gradient descent, but specifically as a form of contrastive learning. By leveraging kernel methods, a relationship is established between the softmax-based self-attention mechanism and gradient descent operating within a contrastive learning objective \cite{ref42}. This interpretation moves beyond the linear attention case, accommodating the nonlinearity of the softmax function that is standard in practice. The key insight is that the attention operation inherently discriminates between input tokens, assigning higher weights to ``positive'' examples (demonstrations relevant to the query) against an implicit distribution of ``negatives'' formed by the other tokens in the context window. This framing suggests that ICL is a learned process of identifying an input query's similarity to demonstration patterns in a high-dimensional representational space, a dynamic that can be further improved by referencing advances in the design of contrastive learning algorithms \cite{ref42}.

The model of attention as implicit optimization has also been refined by incorporating more sophisticated optimizer analogues. A basic linear attention model can be interpreted as implementing a form of preconditioned gradient descent, but when learning from data with varying noise levels, linear Transformers can discover and execute far more intricate optimization strategies. These self-discovered algorithms have been reverse-engineered and found to incorporate momentum and adaptive rescaling, mirroring the behavior of advanced optimizers like Adam \cite{ref43}. This finding is theoretically significant because it demonstrates that even simple architectures, when meta-trained on appropriate data distributions, can autonomously learn to implement complex learning rules as part of their forward computation. This concept of momentum-based optimization is not only descriptive but has been leveraged prescriptively. Inspired by the dual form between attention and gradient descent, a momentum-based attention mechanism was designed by analogy with gradient descent with momentum. This modified attention led to improved performance over standard attention mechanisms, providing causal evidence that the algorithmic correspondence can inform effective model design \cite{ref25}.

However, the elegant theoretical equivalence between ICL and gradient descent faces significant challenges when scrutinized in the context of realistic, large-scale pre-trained models. A critical re-evaluation of this hypothesis reveals that many prior validations were conducted under limiting assumptions that diverge from practical settings, often using hand-constructed weights or models explicitly trained for ICL rather than those exhibiting emergent ICL \cite{ref26}. When probing pre-trained models like LLaMA-7B, the empirical alignment between ICL and gradient descent breaks down in several key ways. A fundamental discrepancy is observed in their sensitivity to the order of demonstrations; ICL and gradient descent processes modify the model's output distribution in distinctly different manners, and their performance metrics diverge as a function of various factors such as dataset composition and the number of in-context examples \cite{ref26}. This casts doubt on a literal, unitary interpretation of the equivalence and suggests that while ICL may share algorithmic properties with optimization, the specific objective and optimization path being implemented in a massive, pre-trained model is far more complex.

Further complicating the picture is the concept of ``Layer Causality,'' which highlights a major information flow discrepancy between ICL and direct gradient descent. In explicit gradient descent, error signals propagate from the output backward to earlier layers. In a Transformer's forward pass during ICL, the flow of information is strictly feed-forward. To align the analogy, a modified gradient-based optimization procedure was proposed that respects this layer causality, and it was shown to significantly improve similarity scores with ICL, indicating that if a gradient descent correspondence exists, it must occur in a layer-localized or sequential manner, not as a single end-to-end backward pass \cite{ref10}. The architectural components beyond attention also play a crucial role in this optimized process. The multi-layer perceptron (MLP) block within a Transformer is central to the model's capacity for ICL, as a linear Transformer block without it incurs an irreducible approximation error on tasks like linear regression. The combined linear attention and MLP construct can be directly mapped to a one-step gradient descent estimator with a learnable initialization, demonstrating that the MLP acts as a vital component that reduces error and enables the implementation of a full, effective optimization step \cite{ref44}.

These theoretical insights have spurred the development of new methods that actively exploit the ICL-as-optimization analogy to boost performance. A prominent example is the ``Iterative Forward Tuning'' framework, which explicitly divides the ICL process into a ``Deep-Thinking'' stage followed by an inference stage. During the ``Deep-Thinking'' stage, the model performs iterative forward optimization on the demonstrations, treating them as a training set to be thought through multiple times. This process accumulates meta-gradients by manipulating the Key-Value matrices in the self-attention modules, effectively distilling the demonstrations into a learned state. The subsequent inference stage then applies these learned meta-gradients to the test query without needing the original demonstrations in the prompt, leading to substantial gains in both accuracy and efficiency \cite{ref45}. This represents a practical fusion of theory and application, leveraging the dual form to achieve a form of test-time model adaptation that is more powerful than standard single-pass ICL. Collectively, these theoretical frameworks, despite their acknowledged limitations in fully explaining emergent phenomena in large pre-trained models, provide the most rigorous language for dissecting the mechanics of ICL. Together with the Bayesian perspective, they point toward a unified view where attention orchestrates a family of implicit, feed-forward meta-optimization procedures honed by pretraining data---procedures that can manifest either as Bayesian model averaging over latent tasks or as algorithmic strategies like gradient descent, depending on the diversity and structure of the pretraining distribution.

\subsection{Bayesian Model Averaging and Task Identification}

While the Bayesian perspective frames in-context learning (ICL) as an implicit inference process over latent tasks, a prominent theoretical lens formalizes this by casting the transformer as a Bayesian inference engine that implicitly performs model averaging or task identification. Under this perspective, the pretrained model encodes a posterior distribution over latent tasks or concepts, conditioning on the provided in-context demonstrations to form a predictive distribution for the query input. Several studies have developed this notion, showing that the attention mechanism can naturally parameterize Bayesian model averaging (BMA). Specifically, \cite{ref46} proves that without updating its parameters, a transformer implicitly implements the BMA algorithm by attending to the in-context examples in a manner that approximates the posterior predictive distribution of a latent variable model. This framework also yields an online learning regret bound of \(\mathcal{O}(1/T)\), where \(T\) is the number of in-context examples, characterizing the efficiency with which ICL aggregates information from the prompt.

The Bayesian prism provides a powerful tool for dissecting the inductive biases learned during pretraining. In a stylized meta-learning setup, a transformer trained on sequences of \((x, f(x))\) pairs from a function class can generalize to unseen functions from that same class. As explored in \cite{ref47}, the model's predictions can be compared to the Bayes-optimal predictor computed from the true pretraining distribution. In many cases involving linear and nonlinear function families, high-capacity transformers mimic the Bayesian predictor remarkably well, suggesting that ICL can be understood as a form of hierarchical Bayesian inference over tasks. This aligns with the formulation in \cite{ref46}, where the ICL estimator is shown to correspond to Bayesian model averaging over a set of latent concepts inferred from the prompt.

However, a key tension emerges between recognizing a task from the pretraining distribution and genuinely learning a new task that lies outside it---a tension that directly mirrors the dual algorithmic behaviors discussed in the gradient descent framework. When the pretraining data covers a limited set of tasks, the model's behavior is dominated by its Bayesian prior, acting as an estimator that relies on the previously seen task distribution. \cite{ref48} empirically demonstrates this phenomenon through a task diversity threshold. Below this threshold, a pretrained transformer behaves as a Bayesian estimator with the non-diverse pretraining distribution as its prior, failing to solve fundamentally new regression tasks. Strikingly, once the diversity of tasks in the pretraining data exceeds this threshold, the transformer's ICL behavior shifts: it deviates from the Bayes-optimal estimator with respect to the pretraining distribution and instead aligns with ridge regression, an optimal predictor with a Gaussian prior over all tasks. This transition into a ``non-Bayesian'' regime---where the model outperforms its implicit Bayesian prior---reveals that ICL is not simply retrieving a memorized task but can involve the emergence of an algorithm-like, general-purpose learning rule, complementing the view of attention as a meta-optimizer that can implement strategies like gradient descent.

This trade-off between task recognition and task learning is central to the Bayesian account of ICL. When the model encounters demonstrations, it must simultaneously infer the latent task identity and generate the correct mapping for novel inputs. The Bayesian framework neatly separates these steps: identifying which previously seen concept best explains the provided examples, then applying the corresponding conditional distribution. Nevertheless, the finding that transformers can surpass this identification-based strategy by implementing optimal procedures like ridge regression indicates a more complex internal mechanism that resonates with the explicit optimization procedures described in the previous section. The theoretical result in \cite{ref49} formalizes this by establishing that for a single-layer linear attention model pretrained on linear regression tasks, only a small number of independent tasks is necessary for the model to closely match the Bayes optimal algorithm---specifically, optimally tuned ridge regression---achieving near-Bayes risk on unseen tasks. This demonstrates that the necessary condition for moving beyond simple task retrieval toward learning an optimal algorithm is a sufficient diversity in the pretraining task distribution, bridging the Bayesian viewpoint with the mechanistic understanding of how attention orchestrates learned optimization strategies.

Further nuance is provided by analyses of how the model handles cases where the true task lies entirely outside the span of pretraining experience. \cite{ref47} observes that transformers can generalize to new function classes unseen during pretraining, a capability that inherently requires deviation from the Bayesian predictor tied to the original pretraining distribution. The authors hypothesize that with enough data and scale, the model learns to implement a more universal inference procedure---one that aligns with ridge regression for linear problems---rather than merely averaging over a static set of memorized task posteriors. This aligns with the mechanistic picture of ICL as a form of implicit meta-optimization or algorithm distillation, where the forward pass of the transformer simulates the steps of an iterative learning algorithm, a topic deeply connected to implicit gradient descent and Bayesian model averaging.

The data-centric determinants behind these Bayesian and non-Bayesian behaviors further cement the role of task heterogeneity. The emergence of ridge regression as the learned ICL strategy is driven not only by diversity but also by the statistical structure of the pretraining tasks. \cite{ref48} underscores that when task diversity is high, the model learns to infer a global structure that generalizes beyond the empirical prior. This behavior is replicated in practical settings: \cite{ref50} demonstrates a threshold behavior where, as the number of pretraining tasks grows, the ICL performance switches from that of a minimum mean squared error (MMSE) equalizer with the pretraining prior to an MMSE equalizer with the true data-generating prior, mirroring the transition from a Bayesian prior-dominated solution to a universally optimal one.

On the other hand, when the model is anchored by strong semantic priors from natural language pretraining, the Bayesian prior can exert an overwhelming pull that prevents full adaptation to the information in the prompt. \cite{ref29} finds that for subjective tasks like emotion recognition, LLMs have strong and inconsistent priors that ossify their predictions, leading to suboptimal performance even when contrasting demonstrations are provided. This observation highlights the dual nature of the Bayesian perspective: while the prior encapsulates useful world knowledge that enables zero-shot generalization, it can also constitute a source of bias that ICL struggles to override. The model's predictions remain a compromise between the in-context evidence and the pre-encoded task priors, echoing earlier findings that ICL does not treat all in-context information equally \cite{ref51}.

In summary, the Bayesian perspective offers a unifying framework for understanding ICL as an implicit inference process over latent tasks, carried out through the attention mechanism. It clarifies why selective retrieval of demonstrations and prompt engineering are so impactful: they shift the implicit posterior toward the desired task. Crucially, the interplay between task diversity and optimal predictor emergence---where models transcend their Bayesian priors to implement procedures like ridge regression---connects the Bayesian viewpoint to the broader theoretical narrative of ICL as an emergent meta-learning capability. This synthesis underscores that ICL is neither purely a memory-based retrieval of pretraining concepts nor a blank-slate learner, but a hybrid process where the model leverages its prior over tasks while possessing the capacity to alchemize diverse pretraining experience into general-purpose algorithmic strategies. This hybrid nature provides the theoretical foundation for understanding how such strategies are physically realized in the transformer's circuitry, as explored in the subsequent analysis of induction heads and their role as the mechanistic building blocks of ICL.

\subsection{Induction Heads and Mechanistic Circuit Analysis}

The mechanistic interpretability of in-context learning (ICL) has been significantly advanced by the identification and analysis of \textbf{induction heads}, a specialized type of attention head that implements a simple match-and-copy algorithm. While the Bayesian perspective frames ICL as implicit inference over latent tasks and the task vector perspective frames it as the compression of demonstrations into a reusable representation, the induction head hypothesis provides a concrete, circuit-level account of how these capabilities are physically realized within the transformer's architecture. Initially characterized as heads that complete sequences like \cite{ref52}\cite{ref53} \ldots{} \cite{ref52} \(\rightarrow\) \cite{ref53}, induction heads are hypothesized to constitute the core mechanistic building block driving the majority of ICL behaviors in large transformer models \cite{ref54}. This hypothesis is supported by the striking temporal coincidence between the abrupt development of induction heads during training and a sudden, sharp increase in the model's in-context learning capability, which manifests as a phase change in the training loss. The discovery of induction heads has shifted the discourse from viewing ICL as an opaque emergent property to understanding it as a product of concrete, analyzable circuits, opening a new frontier in decoding the internal algorithms that underpin the Bayesian and algorithmic behaviors described previously.

The fundamental operation of an induction head rests on a two-step process: matching and copying. In the match step, the head leverages the key-query attention mechanism to attend to previous occurrences of the current token or a similar pattern in the context. The copy step then amplifies the output logits of the token that followed the matched occurrence, effectively `predicting' that this historical sequence will repeat. This mechanism allows the model to leverage token co-occurrence statistics from the prompt itself, performing a form of rudimentary pattern completion without any parameter updates. In this sense, the induction head can be viewed as the most elementary distilled algorithm---a hard-coded learning rule that emerges from pretraining and executes during the forward pass, directly instantiating the algorithm distillation framework discussed in the context of task vector formation. The crucial role of induction heads as a mechanistic source of ICL is most strongly demonstrated in small attention-only models, where causal interventions provide compelling evidence, while for larger models with MLPs, the evidence remains more correlational but equally suggestive \cite{ref54}.

A core question in the field is not just \emph{what} induction heads do, but \emph{how} and \emph{why} they emerge so abruptly. This question mirrors the broader inquiry into how transformers transition from Bayesian prior-dominated behavior to general-purpose algorithmic learning, as the emergence of induction heads may represent the mechanistic substrate of that transition. Detailed investigations into training dynamics reveal that this emergence is not a monolithic event but the result of a sequential learning of nested nonlinearities within a subcircuit composed of interacting components. By studying a minimal attention-only network on a simplified in-context classification task, researchers have identified a specific chain of multi-layer operations necessary to achieve ICL. This process can be modeled as the formation of a two-parameter induction head, where the abrupt phase change is traced to the sequential learning of three nested logits, enabled by an intrinsic curriculum that guides the model's development \cite{ref55}. This intrinsic curriculum provides a natural explanation for the sudden, rather than gradual, acquisition of ICL ability, paralleling the threshold effects observed in task diversity experiments where models abruptly transition to non-Bayesian, algorithmic behavior.

To causally dissect this formation process, a novel ``optogenetics-inspired'' framework has been developed, allowing for the targeted clamping or modification of specific activations throughout training. Using this framework, researchers have identified three underlying subcircuits that interact to drive the formation of induction heads. These subcircuits---which are necessary for the head's match-and-copy functionality---must all ``go right'' for an induction head to emerge. The diversity observed in induction heads, where multiple such heads can form, is explained by their additive and potentially interacting nature, with different heads potentially specializing in different contextual patterns \cite{ref56}. This subcircuit-level analysis not only clarifies the `why' of the phase change but also illuminates data-dependent properties of formation, such as the precise timing of the loss bump, demonstrating the promise of in-depth mechanistic understanding for predicting and controlling model behavior.

The concept of induction heads has been extended beyond simple token matching to encompass more sophisticated statistical computations, most notably through the discovery of \textbf{statistical induction heads}. These circuit elements are formed when transformers are trained on tasks that require learning sequential dependencies, such as Markov chains. In this setting, each training example is a sequence sampled from a Markov chain drawn from a prior distribution. Transformers trained on this task learn to compute accurate next-token probabilities based on the bigram statistics of the in-context sequence, effectively becoming ``in-context learning of Markov chains'' (ICL-MC) solvers \cite{ref57}. The learning dynamics are themselves phased, progressing from uniform predictions, to sub-optimal unigram-based predictions, and finally, through a rapid phase transition, to the correct bigram solution that relies on statistical induction heads \cite{ref57}. This multi-phase process reveals a fascinating complexity: the presence of the simpler, albeit suboptimal, unigram solution can actually delay the formation of the correct bigram solution, highlighting a form of competition between different learned heuristics within the model's internal dynamics---a mechanistic echo of the tension between Bayesian prior-dominated strategies and more optimal algorithmic solutions.

Furthermore, the study of induction heads has expanded to encompass semantic relationships, not just surface-level token patterns. The identification of \textbf{semantic induction heads} demonstrates that certain attention heads encode relationships between tokens found in natural language, such as syntactic dependencies or relations within knowledge graphs. These heads exhibit a recall pattern: when attending to a `head' entity, they increase the output logits of the corresponding `tail' entity, akin to recalling a fact. Crucially, the formulation and strength of these semantic induction heads exhibit a close correlation with the emergence of the model's overall in-context learning ability, bridging the gap between mechanistic circuits and high-level semantic understanding \cite{ref58}. This finding connects the circuit-level analysis back to the Bayesian framing, suggesting that semantic induction heads may serve as the physical implementation of the model's prior over semantic tasks, encoding structured world knowledge that can be activated through in-context demonstrations.

Complementary to induction heads, other specialized circuit components have been identified that contribute to the breadth of ICL capabilities. In the domain of algorithmic reasoning, \textbf{successor heads} have been found to increment tokens along a natural ordering, such as days of the week or numbers. These heads implement abstract ``mod-10 features'' that are common across different LLM architectures and scales, from 31 million to 12 billion parameters \cite{ref59}. For more complex, structured reasoning, such as learning to generate strings from formal languages, transformers form \textbf{n-gram heads}, which are higher-order variants of induction heads. These n-gram heads compute input-conditional next-token distributions over longer subsequences, and compellingly, hard-wiring these heads into neural models improves performance not just on in-context language learning tasks but also on standard natural language modeling \cite{ref60}. The evidence that circuits like the indirect object identification (IOI) circuit are largely reused across semantically different tasks further indicates that models might rely on a finite set of interpretable, task-general algorithmic building blocks, a finding that encourages the project of fully reverse-engineering model behavior \cite{ref61}. Collectively, these circuit-level discoveries provide the mechanistic foundation for the higher-level phenomena of task vector compression and algorithm distillation explored in the following subsection, grounding abstract representational concepts in concrete, analyzable neural circuitry.

\subsection{Task Vectors, Representations, and Algorithm Distillation}

A pivotal line of investigation into the mechanisms of in-context learning (ICL) seeks to understand how the information contained within a set of demonstrations is internally represented and utilized by the model. Building on the circuit-level mechanisms described in the previous subsection---where induction heads and their variants implement match-and-copy operations---this line of research addresses a complementary question: how do transformers compress multiple demonstrations into reusable representations that guide downstream processing? Rather than processing each demonstration as an independent, atomic unit, evidence suggests that transformers learn to compress the essential characteristics of a task into a compact, reusable signal, often termed a ``task vector,'' and that this process is deeply intertwined with the model's layered architecture and can be framed as a form of algorithm distillation.

The concept of a task vector posits that the primary function of the demonstrations in ICL is to compute a single, high-level representation of the task, which then modulates the model's behavior when processing the query input. The work by \cite{ref22} provides compelling empirical evidence for this hypothesis. It demonstrates that the complex function learned by a large language model (LLM) through ICL can often be reduced to a simple structure: the transformer operates on the query \(x\) and a single ``task vector'' \(\boldsymbol{\theta}(S)\) derived from the training set of demonstrations \(S\). This finding suggests that ICL is not a process of performing complex reasoning over the entire set of demonstrations for each query but rather one of compressing \(S\) into a compact representation that captures the input-output mapping, which is then used to generate the output. This perspective is further refined by the concept of a ``state vector'' in \cite{ref62}, which draws explicit parallels between these compressed vectors and the parameters trained by gradient descent. The authors propose that these state vectors can be progressively refined through test-time adaptation using inner and momentum optimization methods, directly mirroring the optimization dynamics of explicit parameter updates. This work operationalizes the task vector idea by showing that these internal representations can be optimized in place to improve performance, and they even introduce a divide-and-conquer aggregation method to handle scenarios with many demonstrations, which would otherwise be too lengthy for standard context windows.

The process of forming these task-relevant representations is not monolithic but appears to be distributed across the layers of a transformer in a functionally differentiated manner. A crucial insight comes from studying compositional ICL problems, where the label is a linear function of a complex, non-linear transformation of the input. The study \cite{ref63} provides both theoretical and empirical evidence for a ``dissection'' of labor. The authors show that optimally, a transformer should first learn to transform the inputs using the fixed representation function and then perform linear ICL on the transformed dataset. Their experiments with trained transformers confirm this hypothesis: lower layers are responsible for transforming the raw inputs into a useful representational space, while upper layers are the locus of the linear ICL operation. This layered decomposition is a powerful mechanistic insight, indicating that ICL can involve a form of implicit, compositional representation learning. The model learns to extract and re-represent the data before applying a simpler learning rule, a finding that resonates with the idea that ICL can be understood as a form of structure induction and schema learning, as explored in \cite{ref64}.

A compelling overarching framework for understanding ICL is that of algorithm distillation, where the model does not just learn a specific task from the prompt but implicitly learns and executes a learning algorithm itself. This perspective formalizes the conjecture introduced in the context of induction heads---that these circuits physically instantiate distilled learning rules---by elevating it to a general principle of ICL. The framework is formalized in \cite{ref2}, which frames ICL as an algorithm learning problem where the transformer implicitly constructs a hypothesis function at inference time. The model's ability to generalize to unseen tasks is then analyzed through the lens of multitask learning, with the stability of the implemented algorithm playing a crucial role in its generalization guarantees. This view is directly supported by the evidence of task vector formation and the layered processing of representations. The transformer, through its pretraining on diverse data distributions, has learned to distill a general-purpose learning algorithm---such as gradient descent or its approximations---into its forward pass. The findings in \cite{ref55} further support this by showing that ICL is driven by the abrupt emergence of an ``induction head,'' a specialized circuit that performs a match-and-copy operation. The formation of this head, through a sequential learning of nested nonlinearities, is the physical instantiation of a distilled algorithm that can rapidly adapt to new in-context data. The role of softmax attention in this process is critical, as highlighted by \cite{ref65}, which shows that the attention mechanism learns to implement a nearest-neighbors predictor whose window adapts to the properties of the pretraining tasks, such as Lipschitzness and label noise. This adaptivity is a key feature of a sophisticated algorithm that has been distilled into the network's weights.

The refinement of this internal representation is an active area of research, with methods seeking to make the implicit learning process more explicit. The work in \cite{ref45} proposes a two-stage framework that exploits the dual form between transformer attention and gradient descent. In the ``Deep-Thinking'' stage, the model performs iterative forward passes over the demonstrations, accumulating ``meta-gradients'' by manipulating the Key-Value matrices. The inference stage then applies these learned meta-gradients directly to the query, without needing the demonstrations in the context. This effectively bridges the gap between implicit in-context meta-optimization and explicit gradient-based fine-tuning, showing that the ICL process can be deepened and made more robust by explicitly simulating the iterative optimization of a task vector. Similarly, the concept of anchoring information in label words, as explored in \cite{ref66}, reveals a specific mechanism for representation consolidation. The study shows that semantic information flows into the representations of label words during the processing of shallow layers, and these enriched representations then serve as a reference point for the final prediction. This anchoring mechanism provides a concrete, interpretable pathway for how a task vector is formed and then utilized, turning the abstract notion of compression into a traceable information flow process within the model's layers. Collectively, these findings paint a picture of ICL as a sophisticated, multi-stage process of algorithm distillation where demonstrations are compressed into powerful, updatable task vectors through a layer-wise, compositional learning of representations. This representational perspective directly motivates the subsequent investigation into how specific properties of pretraining data---such as burstiness, parallel structures, and task diversity---provide the training signal necessary for models to develop these compression and distillation capabilities in the first place.

\subsection{Pretraining Data Properties Driving ICL Emergence}

Understanding why in-context learning (ICL) emerges requires extending the mechanistic perspective of algorithm distillation and task vector formation into the pretraining data itself. If ICL represents a distilled learning algorithm implicitly executed within a transformer's forward pass, then the properties of the data on which this distillation occurs are of paramount importance. Rather than viewing ICL as a purely architectural phenomenon or an inevitable consequence of scale, a growing body of research identifies specific properties of pretraining corpora that are critical determinants of whether and how effectively models develop the ability to form compressed task representations and implement learning algorithms in-context. These data-centric properties include burstiness, long-range coherence, token skewness, concept-awareness, and compositional structure, each contributing to the emergence of ICL through distinct mechanisms that shape how models learn to utilize contextual information for the kind of representational compression and algorithm execution discussed in previous sections.

The foundational role of parallel structures in pretraining data has been demonstrated through systematic ablation studies. When researchers remove parallel structures---pairs of phrases following similar templates within the same context window---from pretraining data, language models' ICL accuracy drops by approximately 51\%, compared to only a 2\% reduction from random ablation \cite{ref67}. This dramatic differential effect persists even when excluding common patterns such as n-gram repetitions and long-range dependencies, indicating that parallel structures constitute a diverse and general class of pretraining phenomena that fundamentally enable ICL. These structures span diverse linguistic tasks and long distances within documents, effectively providing the model with implicit training signals that resemble the demonstration-query format of ICL prompts. From the perspective of task vector formation, such parallel structures naturally encourage the model to compress recurring patterns into reusable representations, directly scaffolding the development of the compression mechanisms central to ICL.

The distributional properties of pretraining data, particularly burstiness and token skewness, exert strong control over the emergence and relative strength of in-context versus in-weights learning. In a minimal attention-only network trained on simplified datasets, researchers have demonstrated that burstiness---where tokens appear in concentrated clusters rather than uniformly---along with large dictionaries and skewed rank-frequency distributions, governs the trade-off between these two learning modes \cite{ref55}. This work reveals that ICL is driven by the abrupt emergence of an induction head, which subsequently competes with in-weights learning circuits. The abrupt emergence itself traces to the sequential learning of nested logits through an intrinsic curriculum, suggesting that the statistical structure of pretraining data directly shapes the developmental trajectory of ICL capabilities through multi-layer operations implemented by nested nonlinearities. This finding directly connects data properties to the mechanistic formation of the induction heads that serve as the physical substrate of distilled learning algorithms, and foreshadows the competitive dynamics between ICL and in-weights learning that characterize training.

The challenge of learning from long-range context represents another crucial data property that shapes how models learn to aggregate information across demonstrations. Through iterative, gradient-based approaches to identify subsets of pretraining data that support ICL, researchers have discovered that supportive pretraining data possess distinctive characteristics \cite{ref41}. These supportive subsets contain a higher mass of rarely occurring, long-tail tokens and represent challenging examples where the information gain from long-range context falls below average. This finding suggests that learning to incorporate difficult long-range context actively encourages the development of ICL abilities, as the model must learn to leverage distant dependencies to make accurate predictions. Continued pretraining on these relatively small supportive subsets can improve model ICL ability by up to 18\%, demonstrating that strategic data selection can substantially enhance in-context learning without architectural modifications. This aligns naturally with the layered processing of representations described earlier, where lower layers transform inputs and upper layers perform linear ICL; long-range dependencies likely force the model to develop robust mechanisms for propagating and integrating information across sequence positions.

The importance of concept-awareness and analogical reasoning in pretraining data has led to the development of targeted training strategies. Concept-aware Training (CoAT) constructs training scenarios where capturing analogical reasoning concepts proves beneficial for the language model, resulting in consistent improvements in reasoning ability \cite{ref68}. Notably, in-context learners trained with CoAT on only two datasets of a single question-answering task perform comparably to larger models trained on over 1600 tasks, underscoring that the quality and conceptual structure of training data can substitute for raw task diversity when specifically designed to foster ICL capabilities. This aligns with broader findings about concept-centric pretraining, where models that explicitly learn relational commonsense knowledge about everyday concepts through generative and contrastive objectives demonstrate improved performance on downstream tasks requiring common sense reasoning \cite{ref69}. Such conceptual structure in pretraining data likely facilitates the schema-learning and rebinding mechanisms that enable models to recognize and adapt to novel task structures in-context.

The role of compositional structure and task diversity in pretraining data has been elucidated through studies on linear regression tasks, providing a direct link to the algorithmic perspective on ICL. A critical finding reveals the existence of a task diversity threshold for the emergence of ICL: below this threshold, pretrained transformers cannot solve unseen regression tasks and instead behave like Bayesian estimators constrained by the non-diverse pretraining task distribution as a prior \cite{ref48}. Beyond this threshold, transformers significantly outperform this estimator and align with ridge regression, corresponding to a Gaussian prior over all tasks, including those not seen during pretraining. This transition from Bayesian to non-Bayesian behavior demonstrates that pretraining on data with sufficient task diversity enables transformers to optimally solve fundamentally new tasks in-context, with the critical capability of deviating from the Bayes optimal estimator with the pretraining distribution as the prior. This finding directly supports the algorithm distillation framework: sufficient task diversity enables the model to distill a general-purpose learning algorithm that can adapt to novel tasks, rather than simply memorizing a fixed set of task-specific solutions.

The phenomenon of data curation stabilizing ICL performance further illuminates the relationship between pretraining data properties and in-context capabilities. Careful curation of training subsets can greatly stabilize ICL performance without any modifications to the ICL algorithm itself \cite{ref34}. Methods that score individual training examples by their average development set ICL accuracy when combined with random examples, or that learn linear regressors estimating how each training example influences LLM outputs, yield subsets that improve average accuracy over random sampling by 6-8\%. Surprisingly, these stable subset examples are not characterized by exceptional diversity or low perplexity, challenging assumptions about what constitutes high-quality ICL demonstrations and suggesting that the determinants of stable ICL performance are more nuanced than simple heuristics suggest. This stabilization is particularly relevant given the transient nature of ICL during training, as it suggests that appropriate data curation may help maintain ICL capabilities against the competitive pressure from in-weights learning.

The formal diversity of pretraining data, as measured by the Task2Vec diversity coefficient, has been shown to be high for publicly available LLM datasets, approaching theoretical upper bounds \cite{ref70}. This diversity coefficient increases as the number of latent concepts increases, aligning with intuitive properties of diversity. Such formal diversity likely contributes to the emergence of ICL by exposing models to a wide range of linguistic patterns and task structures during pretraining, facilitating the development of general-purpose in-context learning mechanisms rather than narrow, task-specific heuristics. When viewed through the lens of the layered processing framework, this diversity provides the varied training signal necessary for lower layers to develop flexible representation learning capabilities that can adapt to diverse input transformations.

The interplay between data properties and model capabilities extends to the skills framework, where ordered skill sets exist in training data and enable more advanced skills to be learned with less data when training on prerequisite skills \cite{ref71}. This data-driven skills perspective suggests that ICL itself may emerge as a higher-order skill built upon more fundamental capabilities acquired through exposure to appropriately structured pretraining data, with online data sampling algorithms over mixtures of skills achieving substantial improvements over random sampling in both continual pretraining and fine-tuning regimes. This hierarchical learning of skills mirrors the sequential, intrinsic curriculum observed in the formation of induction heads, where nested nonlinearities are learned in a specific order.

The implications of these findings extend to practical considerations for pretraining data design. The explicit dependencies among tokens in sequences have been identified as critical for downstream performance, with models achieving better results when pretrained on datasets possessing longer ranges of implicit dependencies \cite{ref72}. This reinforces the centrality of long-range coherence and dependency structure in fostering the transferable capabilities, including ICL, that make pretrained models valuable. Together, these investigations establish that ICL emergence is not merely a function of scale but is fundamentally shaped by the statistical and structural properties of pretraining data, with burstiness, token skewness, parallel structures, task diversity, concept-awareness, and long-range dependencies each playing distinct and complementary roles in enabling models to distill learning algorithms, form task vectors, and execute multi-stage representational transformations at inference time. The data properties discussed here directly set the stage for understanding the complex training dynamics, including the abrupt phase transitions and competition between learning strategies, that characterize how these capabilities actually emerge during the training process.

\subsection{Training Dynamics, Phase Transitions, and Transience}

The statistical and structural properties of pretraining data discussed in the previous section directly shape the complex, non-linear dynamics through which in-context learning (ICL) emerges during training. Rather than developing gradually, ICL is characterized by abrupt phase transitions, transient phenomena, and a continual competition with alternative learning strategies---dynamics that are fundamentally rooted in the burstiness, token skewness, task diversity, and long-range dependencies of the training corpus. Understanding these dynamics is crucial for optimizing training, controlling model capabilities, and bridging the mechanistic insights of previous sections with the theoretical guarantees that follow.

A central finding in this area is the \textbf{transient nature} of ICL, which directly manifests the competition between in-context and in-weights learning modes foreshadowed by the data-dependence results. Contrary to the assumption that once ICL emerges, it persists as a stable capability, research shows it can be a temporary stage in a model's developmental trajectory. A study on transformers trained on synthetic data---where both ICL and in-weights learning (IWL) could solve the task---revealed that ICL often first emerges, achieves a peak, and then completely disappears, giving way to IWL, all while the overall training loss continues its asymptotic decrease \cite{ref73}. This indicates that, in the long run, models may develop a preference for memorization-based IWL strategies over context-driven ICL solutions, raising important questions about the ``overtraining'' of models when compact, fast-running solutions are desired. The study suggests that L2 regularization can stabilize ICL and mitigate its transient decay, offering a practical lever to preserve this capability without relying on early stopping---a finding that anticipates the stability analysis and regularization perspectives developed in the theoretical formalisms of the following section.

This transience is a consequence of a deeper \textbf{competition between ICL and IWL circuits} within the model, a dynamic that directly extends the mechanistic basis of data dependence established earlier. The abrupt emergence and subsequent decline of ICL are posited to stem from a competitive dynamic where ICL circuits, which appear rapidly, are later overshadowed and functionally replaced by more slowly developing IWL circuits \cite{ref55}. This competition is not merely a linear trade-off but can involve specific architectural subcomponents, explaining why the pretraining data properties of burstiness and token skewness govern the relative strength of these two learning modes. For instance, the formation of \textbf{induction heads (IHs)} is widely considered a key mechanistic driver for ICL, as these heads match and copy patterns from a sequence's past, serving as the physical substrate of the distilled learning algorithms discussed in the algorithm distillation framework \cite{ref56}. The development of IHs is itself a multi-stage process, where nested nonlinearities within the attention mechanism are learned sequentially via an intrinsic curriculum, directly building on the sequential learning of nested logits identified in minimal attention-only networks. The abrupt emergence of an IH, and thus ICL, is traced to the specific chain of multi-layer operations required to activate it, which is only possible after certain foundational subcircuits are in place \cite{ref55}. Once formed, these ICL circuits must then compete with the parameters dedicated to IWL, foreshadowing the approximation-estimation trade-off formally characterized in the subsequent section's theoretical decomposition.

These dynamics are often accompanied by striking \textbf{abrupt phase changes} during training, which provide the empirical foundation for the phase transitions analyzed in theoretical frameworks. When plotting ICL performance on held-out tasks, one often observes a sharp, discontinuous jump from near-chance to near-perfect accuracy, rather than a smooth curve---behavior that connects directly to the task diversity threshold identified for linear regression tasks, where transformers transition from Bayesian estimation to non-Bayesian ridge regression. Experiments on controlled, synthetic data have linked this phase change directly to the moment the necessary subcircuits for an induction head fully ``click'' into a functional unit \cite{ref56}. This research, using a causal framework akin to optogenetics, identified three specific subcircuits whose interaction drives this sudden formation, demonstrating that a precise configuration of attention patterns is required for ICL to manifest. The timing of this phase change is sensitive to properties of the training data and the model's initial conditions, reinforcing the central role of data properties like parallel structures and long-range coherence in scaffolding the development of these circuits.

However, the training landscape is not only defined by major phase changes but also by prolonged \textbf{learning plateaus} where ICL capability stagnates before improving, echoing the skill hierarchies where higher-order capabilities depend on prerequisite foundations. These plateaus can be understood by decomposing a model's internal representation into a ``context component,'' directly modulated by input examples, and a ``weights component,'' fixed by the model's trained parameters \cite{ref74}. This decomposition mirrors the layered processing framework where lower layers transform inputs into representations and upper layers perform linear ICL operations, and connects naturally to the concept-awareness and analogical reasoning structures that enable models to recognize novel task formats. Research indicates that these plateaus occur precisely when the development of the weights component is compromised. When the weights component fails to effectively interpret or integrate the information extracted by the context component, the model's ICL performance hits a ceiling. By developing strategies to boost the functionality of the weights component---akin to bridging a gap between context processing and stored knowledge---it is possible to break through these plateaus, leading to more efficient and environmentally friendly training of strong in-context learners \cite{ref74}.

Further enriching this picture, the training process itself seems to follow an \textbf{intrinsic curriculum}, where models learn to implement the nested operations required for ICL in a specific, sequential order---a finding that directly extends the hierarchical skill acquisition observed in data-driven skills frameworks. This intrinsic curriculum also manifests in behavior that mirrors human learning, providing a cognitive connection to the meta-learning perspective that will be formalized in the subsequent section's generalization guarantees. For example, the opposing effects of ``blocked'' versus ``interleaved'' learning paradigms, a classic result in human cognitive science, spontaneously emerge. ICL, as an inner-loop learning algorithm operating in activations without weight changes, demonstrates the ``blocking advantage'' seen in humans for rule-like tasks, whereas concurrent IWL exhibits the ``interleaving advantage'' for tasks lacking such structure \cite{ref75}. This not only bridges neural network training with cognitive science but also highlights that the dynamics of ICL and IWL are deeply intertwined, progressing along distinct developmental timelines that collectively define the model's overall learning trajectory---a trajectory whose theoretical characterization in terms of regret bounds, stability guarantees, and generalization error decomposition forms the focus of the following section.

\subsection{Generalization, Stability, and Theoretical Bounds}

Building on the dynamic and often transient nature of ICL's emergence during training, a parallel line of inquiry seeks to formally characterize its behavior at inference time. Understanding the theoretical foundations of in-context learning (ICL) requires a rigorous examination of its generalization guarantees, stability properties, and the conditions under which transformers can implement learning algorithms without parameter updates. This section synthesizes formal frameworks that characterize ICL's sample efficiency, error decomposition, and adaptive behavior, drawing on PAC-based learnability, regret analysis, and stability theory.

A central theoretical perspective treats ICL as a meta-learning problem where pretraining over diverse task distributions yields a model capable of implicit learning at inference time. Within this view, \cite{ref46} provides a comprehensive statistical framework, formulating ICL as Bayesian model averaging over latent concepts. Their analysis decomposes the pretrained model's error into an approximation error and a generalization error, proving that the approximation error decays exponentially with transformer depth, while the generalization error decays sublinearly with the number of pretraining tokens. Importantly, they establish an online learning perspective on ICL performance, deriving an \(\mathcal{O}(1/T)\) regret bound for perfectly pretrained ICL, where \(T\) denotes the number of in-context examples. This regret characterization demonstrates that, under ideal pretraining conditions, ICL's predictive accuracy improves at a rate inversely proportional to the demonstration count, offering a theoretical basis for few-shot learning's empirical success.

The PAC-Bayes framework has emerged as a powerful tool for analyzing meta-learning and ICL generalization. \cite{ref76} derives PAC bounds for gradient-based meta-learning by combining uniform stability at the base learner level with PAC-Bayesian analysis at the meta level. This dual-framework approach yields tighter guarantees when base learners adapt quickly---precisely the goal of ICL. Their bound reveals that the generalization gap shrinks as the meta-learner's adaptation speed increases, providing formal justification for why large language models with strong pretraining priors can generalize from few demonstrations. However, the limitations of PAC-Bayes analysis must be acknowledged: \cite{ref77} demonstrates that even simple learnable tasks, such as 1D linear classification, may resist PAC-Bayesian analysis, suggesting that alternative theoretical frameworks may be necessary for a complete understanding of ICL.

Generalization error decomposition for ICL has been further refined through the lens of pretraining data properties. The PAC-based learnability frameworks adapted from classical learning theory provide bounds that separate the error into components attributable to insufficient demonstration coverage and model capacity constraints. Drawing parallels to interactive estimation problems, \cite{ref78} introduces the dissimilarity dimension, a combinatorial measure that captures learnability in interactive settings. While not directly addressing ICL, this framework suggests that in-context learning's sample complexity may be characterized by analogous structural properties of the prompt space.

Stability conditions constitute another pillar of theoretical ICL analysis, directly relating to the competition between in-context and in-weights circuits discussed in the training dynamics. The correspondence between ICL and gradient descent offers insights into algorithmic stability. When transformers implement optimization procedures through attention computations, the stability of these implicit algorithms governs generalization. \cite{ref79} analyzes the stability of stochastic gradient methods, establishing bounds that depend on Lipschitz constants and aggregated step sizes. These results are relevant to understanding how ICL's implicit gradient updates maintain stability across demonstrations. The uniform stability framework ensures that replacing a single demonstration does not dramatically alter predictions, a property empirically observed in robust ICL systems and conceptually linked to the competition dynamics where stable ICL circuits must resist displacement by in-weights mechanisms.

The adaptivity of softmax attention to task structure has been precisely characterized in theoretical terms. \cite{ref65} demonstrates that attention units learn windows implementing nearest-neighbor predictors adapted to the pretraining task landscape. Critically, they show that this window widens with decreasing Lipschitzness and increasing label noise in pretraining tasks, revealing that softmax attention automatically adjusts its effective receptive field based on task smoothness. On low-rank linear problems, attention learns to project onto appropriate subspaces before inference. This adaptivity relies fundamentally on the softmax activation and cannot be replicated by linear attention, establishing a theoretical distinction between commonly studied linear attention models and practical softmax implementations, much like how the nested nonlinearities of induction heads proved critical for ICL's abrupt emergence.

Lipschitz continuity plays a multifaceted role in ICL theory. \cite{ref80} illuminates how Lipschitz constants govern generalization error bounds in ranking problems, with the choice of norm being crucial---\(\ell_\infty\)-based definitions yielding tighter guarantees. This insight transfers to ICL, where the representation space's geometry affects bound tightness. \cite{ref81} proposes calibration of Lipschitz-margin trade-offs, showing that larger Lipschitz constants need not sacrifice accuracy when properly calibrated. Applied to ICL, this suggests that attention mechanisms can maintain both expressivity and stability through appropriate Lipschitz control, echoing the balancing act observed during training where L2 regularization can stabilize emergent ICL capabilities.

Theoretical bounds for ICL also benefit from connections to reinforcement learning theory. \cite{ref82} and \cite{ref83} introduce uniform-PAC frameworks that simultaneously provide PAC and regret guarantees. These unified analyses offer templates for deriving similar combined guarantees for ICL, where both the probability of large errors and the cumulative prediction quality matter. The multi-level partition scheme developed in the latter work for sample selection mirrors the demonstration selection strategies employed in practical ICL systems.

The implicit optimization perspective on ICL finds theoretical support through analyses of convergence and approximation. \cite{ref84} establishes Lipschitz continuity and positive semi-definiteness conditions for softmax-ReLU regression problems, providing convergence guarantees for greedy Newton-based algorithms. These results underpin the understanding of how attention layers perform implicit regression during ICL. Similarly, \cite{ref85} develops large deviation inequalities for Lipschitz functions of dependent sequences, offering tools for analyzing ICL's generalization when demonstrations are not independent---a realistic scenario in practical prompting.

The interplay between in-context and in-weights learning, a central theme in the training dynamics, has been formalized through competition dynamics and transient behavior analysis. The theoretical bounds emerging from this literature suggest that ICL's transient nature---where capabilities emerge then potentially decay---can be characterized by the relative rates at which implicit (in-context) and explicit (in-weights) components absorb task information. The approximation-estimation decomposition provides a precise language for this trade-off: in-context mechanisms contribute low approximation error but potentially high estimation variance, while in-weights mechanisms exhibit the opposite profile. Understanding this decomposition enables the design of hybrid systems that optimally balance both learning modes, directly addressing the transient dynamics and plateaus observed in the training process.

Finally, the theoretical understanding of how distribution shifts affect ICL generalization continues to develop. Bounds derived under covariate shift assumptions reveal that ICL's robustness depends on the overlap between pretraining task distributions and the target task's structure. When pretraining tasks exhibit certain smoothness or low-rank properties, the resulting attention mechanisms transfer these structural biases to novel contexts, yielding bounded generalization gaps even under moderate distribution shifts. These theoretical insights collectively provide a foundation for understanding when and why ICL succeeds, while highlighting open questions regarding compositionality, long-context scaling, and the precise characterization of pretraining data properties necessary and sufficient for robust in-context generalization.

\section{Architectural and Distributional Determinants of ICL}

\subsection{Comparative Analysis of Model Architectures for ICL}

The transformer architecture, built upon the self-attention mechanism, has long been considered the central engine driving in-context learning (ICL) in large language models. However, a growing body of work has begun to dissect this assumption by systematically comparing the ICL capabilities of transformers against a diverse spectrum of alternative sequence models. This comparative analysis is crucial for understanding whether attention is a necessary condition for ICL, and for identifying the specific computational primitives that enable a model to adapt to novel information presented in its context window. As will be explored in the next subsection, the inductive biases of softmax attention are indeed foundational, but here we first examine whether they are strictly exclusive.

Recent empirical studies have challenged the exclusivity of the transformer advantage. A comprehensive evaluation of thirteen model architectures on synthetic ICL tasks revealed that a wide range of paradigms, including recurrent networks, convolutional models, and state-space model (SSM) inspired designs, are capable of performing in-context learning under broader conditions than previously documented \cite{ref86}. This finding fundamentally reframes ICL not as a magical property exclusive to attention, but as a more general emergent capability that can be realized through different architectural inductive biases. However, the study also uncovered stark differences in statistical efficiency and consistency. While several attention alternatives, such as certain recurrent and convolutional architectures, proved competitive with or even superior to transformers on specific tasks with a limited number of in-context examples, no single architecture demonstrated universal dominance. A common failure mode was observed when the number of in-context examples significantly exceeded the context length seen during training; many non-transformer architectures suffered from performance plateaus or outright degradation, whereas transformers showed greater consistency. This points to a crucial interplay between architectural design and the ability to robustly scale to longer in-context prompts, a property that the following subsection links directly to the adaptive windowing behavior inherent in softmax attention.

State-space models, and the Mamba architecture in particular, have emerged as a leading alternative to transformers, promising linear-time scaling with sequence length \cite{ref87}. A direct comparison of Mamba and transformers on ICL tasks presents a nuanced picture of their relative strengths. On standard regression tasks, which approximate simple function learning, Mamba performs comparably to transformers. Intriguingly, on specific tasks like sparse parity learning, Mamba has been shown to outperform its attention-based counterparts \cite{ref88}. This suggests that Mamba's input-dependent gating and selective state-space mechanisms provide a distinct inductive bias that is particularly well-suited for certain types of rule-based inference. The ability of Mamba to solve ICL problems appears to be grounded in its capacity to incrementally optimize its internal representations, a dynamic analogous to gradient-based optimization \cite{ref89}. This finding builds a bridge between the mechanistic interpretations of ICL in transformers and the operational principles of SSMs, indicating that implicit meta-optimization might be a shared algorithmic signature of successful in-context learners regardless of their architecture---a theme that resonates with the connection between attention and gradient descent discussed later.

Despite these successes, Mamba-like models exhibit clear failure modes that highlight the critical role of attention in specific ICL subroutines. A key weakness is observed in tasks requiring non-standard retrieval functionality \cite{ref88}. This deficiency is further elucidated when the problem domain shifts to in-context language learning (ICLL), where a model must infer a formal language's rules from examples. In this regime, transformers significantly outperform recurrent and convolutional baselines \cite{ref60}. The success of transformers in ICLL is attributed to the emergence of specialized computational heads, such as ``n-gram heads'' that act as higher-order variants of induction heads, computing input-conditional next-token distributions in a manner that RNNs and SSMs struggle to replicate. This failure is rooted in the architectural constraints of these alternatives; the gating and recurrent dynamics, while efficient, may not provide the same facility for the precise, content-based match-and-copy operations that are effortlessly executed by a dedicated attention head, a capability that Section 3.2 will formalize through the concept of induction heads and token selection mechanisms.

The recognition of this performance gap has naturally led to the investigation of hybrid architectures designed to combine the best of both worlds. The MambaFormer architecture, which augments a Mamba backbone with a small number of attention blocks, exemplifies this approach \cite{ref88}. This hybrid model surpasses both its pure Mamba and pure transformer counterparts on tasks where each struggles independently, effectively using attention to patch the specific retrieval and comparison-based shortcomings of the SSM while retaining the SSM's efficiency for other operations. This principle extends beyond SSMs; the finding that hard-wiring n-gram heads directly into existing models improves both ICLL capabilities and natural language perplexity suggests a more general design strategy \cite{ref60}. By embedding the known algorithmic prerequisites for ICL directly into the architecture, one can create hybrid systems with enhanced statistical efficiency and consistency. This pragmatic synthesis of architectural strengths is precisely the direction explored by the sparse modular activation and hybrid gating mechanisms discussed at the end of Section 3.2.

Further expanding the architectural landscape, alternative sequence prediction methods like clone-structured causal graphs (CSCGs) have demonstrated comparable ICL capabilities while offering a crucial advantage: interpretability \cite{ref64}. The CSCG framework explicitly implements a process of schema retrieval and token rebinding, providing a transparent computational model for how context is utilized. The convergence of evidence from CSCGs and transformers on schema-based mechanisms suggests that the core logical operations underlying ICL---pattern completion and context-sensitive variable binding---might be architecturally agnostic. The analysis of inductive biases reinforces this view. For instance, it has been shown that softmax attention adapts its implicit function approximation window based on the Lipschitzness of the pretraining task landscape, a form of sophisticated adaptivity not easily replicated by linear attention variants \cite{ref65}.

In conclusion, the comparative analysis of model architectures reveals that ICL is not a monolithic phenomenon requiring a single architectural solution. While transformers remain the most consistent and general-purpose in-context learners, particularly for tasks requiring complex retrieval, compositional reasoning, and robust scaling to larger context sizes, the assumption that attention is strictly required is false. Alternative architectures like Mamba excel in specific niches and offer dramatic computational advantages, while CSCGs provide interpretable mechanistic insights. The frontier of ICL-capable architectures is therefore moving towards a unified understanding of the underlying computational primitives---such as induction heads, schematic templates, and iterative state updates---which can be implemented using diverse neural substrates. The future of ICL architecture relies on a deeper integration of theoretical understanding with practical design, leading to hybrid models that combine the efficiency of linear-time sequence models with the powerful retrieval and adaptation capabilities inherent in targeted attention mechanisms. As the next subsection details, it is precisely these adaptive, retrieval-oriented inductive biases of softmax attention that make it an indispensable component, even if not an exclusive one, in the expanding ecosystem of in-context learners.

\subsection{The Role of Attention Mechanisms and Architectural Inductive Biases}

Building upon the comparative architectural analysis, which revealed that while ICL is not exclusive to transformers, softmax attention provides unique and indispensable capabilities, we now delve into the specific mechanisms and inductive biases that underpin this advantage. The softmax attention mechanism is the cornerstone of the Transformer architecture and the primary engine driving its in-context learning (ICL) capabilities. A detailed analysis reveals that softmax attention possesses inherent inductive biases that allow it to adaptively adjust its behavior to the underlying structure of the pretraining task distribution. Crucially, this adaptivity is not merely a product of scale but is mechanistically linked to the softmax function itself. Research demonstrates that in an ICL setting where each context encodes a regression task, a single softmax attention unit learns a window of influence to implement a nearest-neighbors predictor. This window widens when pretraining tasks exhibit lower Lipschitzness or higher label noise, showing an explicit adaptation to function properties \cite{ref65}. Furthermore, on low-rank linear problems, the same attention unit learns to project inputs onto the appropriate low-rank subspace before performing inference. This adaptive behavior, which is critical for generalizing across diverse task families and directly explains the robust scaling to longer prompts observed in the previous section, relies fundamentally on the softmax activation and cannot be replicated by linear attention variants frequently used in prior theoretical analyses \cite{ref65,ref90}.

The expressivity of the attention mechanism is further enriched by the multi-head design, which provides the architectural substrate for the hierarchical function learning explored in the following subsection. A theoretical comparison of single-head versus multi-head attention for in-context linear regression reveals that multi-head attention achieves a superior performance guarantee, with a smaller multiplicative constant in its prediction loss bound as the number of in-context examples increases \cite{ref91}. This finding validates the architectural choice of multi-head attention, as it allows the model to simultaneously learn and apply different adaptive strategies, such as focusing on different feature subspaces or implementing distinct windowing functions. However, this powerful mechanism is not without its inherent biases. Without the stabilization provided by skip connections and multi-layer perceptrons (MLPs), pure self-attention is shown to converge doubly exponentially to a rank-1 matrix, a phenomenon described as ``token uniformity'' \cite{ref92}. This rank collapse is countered by the interplay of attention's localization tendencies. When the query-key (QK) eigenspectrum has small variance, attention becomes highly localized, a state that prevents both rank and entropy collapse, thereby fostering better model expressivity and trainability \cite{ref93}.

These architectural inductive biases give rise to specialized computational circuits, most notably the ``induction heads,'' which serve as the mechanistic bridge between the attention mechanism and the algorithmic capabilities that the next subsection formalizes into a hierarchy of learnable functions. Induction heads are widely hypothesized to be the mechanistic source of ICL \cite{ref54}. They implement a simple pattern-completion algorithm, such as copying a token that followed a similar context in the past, and their abrupt development during training coincides precisely with a phase transition marking a sharp increase in ICL ability \cite{ref54}. Extending this concept, ``semantic induction heads'' have been identified, which not only recall tokens but specifically encode relational knowledge from training data. For instance, when attending to a head entity token, these heads can recall the corresponding tail entity and increase its output logits, with their formation being closely correlated with the emergence of ICL \cite{ref58}. The development of these heads is further explained by the implicit bias of gradient descent, which causes parameter norms to grow during training, eventually saturating the network into a discretized state where attention heads specialize into distinct roles---some focusing locally on a narrow context window while others compute global averages for counting \cite{ref94}. This functional specialization is supported by findings that attention heads learn diverse, interpretable roles based on relative position, part-of-speech, and named entities, and that enforcing sparsity via alternatives like \(\alpha\)-entmax can enhance this diversity and interpretability \cite{ref95,ref96}.

Contrasting this with the alternative mechanisms found in State Space Models (SSMs), as introduced in the previous subsection, highlights the crucial role of input-dependent token selection. While standard SSMs rely on gating, convolutions, and fixed recurrent dynamics to process sequences with linear complexity, they fundamentally lack the fine-grained, data-dependent interaction selection that softmax attention provides. This deficiency directly explains the failure modes of Mamba-like models on retrieval tasks and in-context language learning discussed earlier. Hybrid architectures attempt to bridge this gap by sparsely activating attention modules, a design strategy that yields models which excel precisely where pure SSMs fail. A prime example is the Sparse Modular Activation (SMA) mechanism, which uses an SSM's state representations to dynamically decide which sequence elements require attention, activating a Gated Attention Unit (GAU) only for those specific inputs. This approach, implemented in models like SeqBoat, achieves a superior quality-efficiency trade-off by constraining attention to only the most salient parts of the sequence \cite{ref97}. This demonstrates that the inductive bias of attention can be effectively combined with the efficiency of SSMs, but its role is refined from a global operator to a targeted, on-demand processor. Similarly, other works propose combining sigmoid gating mechanisms with softmax attention to boost performance without increasing architectural size, suggesting that a synthesis of gating and attention can yield more powerful and parameter-efficient models \cite{ref98}.

The relationship between these architectural components and ICL can be further understood through the lens of implicit optimization, a perspective that foreshadows the gradient descent and higher-order optimization algorithms that Section 3.3 will show transformers learn to implement. The self-attention mechanism, when trained with gradient descent, has been proven to converge in direction to a max-margin solution that optimally separates locally optimal tokens from non-optimal ones, formalizing attention as a token selection mechanism \cite{ref99}. This creates a direct parallel between the forward pass of a trained Transformer and an implicit optimization process. In fact, a close mathematical similarity has been shown between the data transformations performed by a single self-attention layer and the updates of gradient descent on a softmax regression loss, suggesting that Transformers learn to simulate gradient-based learners within their forward pass \cite{ref100}. This implicit optimization perspective is further nuanced by the ``word sensitivity'' property of softmax attention, which shows that a single word can have an outsized influence on the attention feature map, a capability not shared by standard random features whose sensitivity decays with sequence length \cite{ref101}. This high sensitivity allows attention to precisely capture context-dependent meaning, a core component of the inductive bias that makes it ideal for ICL. In summary, the architectural inductive biases of softmax attention---its adaptivity to task geometry, its tendency to form specialized induction heads, and its implicit connection to optimization---are the principal drivers of ICL. These biases stand in stark contrast to the fixed, non-adaptive sequence mixing of early SSMs, though modern hybrid designs are increasingly exploiting the strengths of both. As the next subsection details, it is precisely these inductive biases that enable transformers to implement a rich hierarchy of in-context learnable functions, from simple linear regression to complex compositional reasoning and adaptive algorithm selection.

\subsection{What Models Learn In-Context: From Simple Functions to Compositional Representations}

Building on the architectural inductive biases of softmax attention, which endow transformers with the capacity to adaptively implement nearest-neighbor predictors, ridge regression, and even second-order optimization, we can now systematically map the landscape of computations that transformers can learn to implement in their forward pass. A productive line of research has investigated this capability by training transformers from scratch on synthetic data, ranging from simple linear function classes to more complex, structured representations, revealing a hierarchy of in-context learnable functions that mirrors the algorithmic expressivity enabled by the attention mechanisms discussed previously.

At the foundational level, transformers can learn to implement standard statistical algorithms for elementary function classes. The seminal work of \cite{ref102} established that standard transformers can be trained to perform ICL of linear functions, achieving performance comparable to the optimal least squares estimator, even under distribution shifts between training prompts and inference-time prompts. This capability extends beyond linear models to more complex function classes, including sparse linear functions, two-layer neural networks, and decision trees, where trained transformers match or exceed task-specific learning algorithms. The algorithmic basis for this learning has been further elucidated, with \cite{ref103} providing evidence that transformer-based in-context learners internally encode standard estimation algorithms. By proving that transformers can construct implementations of gradient descent and closed-form ridge regression, and then showing that trained models closely match the predictors computed by these algorithms, this work demonstrated that in-context learning is understandable in algorithmic terms, with late-layer activations non-linearly encoding weight vectors and moment matrices. This convergence to standard algorithms deepens with model size, with learned predictors transitioning between gradient descent, ridge regression, and exact least-squares solutions as depth and dataset noise vary, directly echoing the adaptive behavior of softmax attention to function properties such as Lipschitzness and noise levels.

This algorithmic perspective deepens when considering the optimization order that transformers implement. While early hypotheses centered on first-order gradient descent, subsequent investigations revealed that transformers learn higher-order optimization methods, aligning with the theoretical expressivity of multi-head attention and its connection to implicit optimization processes. \cite{ref104} demonstrated that for in-context linear regression, transformers learn to implement an algorithm very similar to Iterative Newton's Method, a second-order optimization method, rather than gradient descent. Empirically, predictions from successive transformer layers closely match different iterations of Newton's method, with each middle layer approximating three iterations, conferring an exponentially faster convergence rate than gradient descent. This capability extends to ill-conditioned data where gradient descent struggles. The capacity for higher-order optimization is not limited to linear regression; \cite{ref105} established theoretically that linear attention transformers with ReLU layers can approximate second-order optimization algorithms for logistic regression and can even implement a single step of Newton's iteration for matrix inversion with merely two layers. Collectively, these findings suggest that the transformer architecture possesses a profound capacity to implement complex numerical algorithms beyond simple gradient-based updates, a capacity rooted in the same computational primitives that give rise to induction heads and specialized attention circuits.

Moving beyond the learning of static functions, a critical question is whether models can learn structured, compositional representations in-context, a capability that would naturally interact with the ICL-IWL trade-off dynamics explored in the following section. A key step in this direction is the study of learning with representations, where the true label depends on a complex but fixed representation function composed with a task-specific linear function. \cite{ref63} constructed synthetic problems with this compositional structure and showed theoretically that transformers with mild depth and size can approximately implement the optimal ICL algorithm: first transforming inputs by the representation function in lower layers, then performing linear ICL on the transformed dataset in upper layers. Trained transformers consistently achieved near-optimal performance and exhibited this clear functional dissection, with concrete copying behaviors, linear ICL in upper layers alone, and a post-ICL representation selection mechanism in mixture settings. This work highlights a critical computational subroutine: the compression of a demonstration set into a ``task vector'' that modulates subsequent processing. \cite{ref22} provides strong evidence for this hypothesis, showing that the functions learned by ICL correspond to the transformer LLM whose inputs are the query and a single task vector computed from the training set, effectively compressing the demonstrations into a compact representation that modulates the transformer's behavior to produce the output. This task vector mechanism provides a crucial bridge to understanding how ICL can coexist with IWL, as the compressed in-context signal must compete with or complement the static knowledge encoded in model weights.

The complexity of in-context learnable functions extends to sequential and compositional domains. \cite{ref106} demonstrated that transformers can learn from non-textual sequential function class data distributions, retaining ICL capabilities even when label associations are partially obfuscated. This suggests an inherent robustness in how transformers perceive and process sequentiality, reminiscent of the adaptive windowing behavior of softmax attention. Regarding compositionality, \cite{ref107} showed that autoregressive transformers can learn compositional structures from small amounts of data and generalize to exponentially many functions. They found that generating intermediate outputs when composing functions is more effective for generalization, and that attention layers select which capability to apply while feed-forward layers execute the selected capability, reflecting the functional specialization of attention heads described in the architectural analysis. The theoretical underpinnings for this compositional learning are explored in \cite{ref108}, which argues that ICL relies on the recombination of compositional operations present in pretraining data, deriving information-theoretic bounds to show how ICL abilities arise from generic next-token prediction when the pretraining distribution has sufficient compositional structure. This view is supported by empirical findings in \cite{ref109}, which showed that transformers become more tree-like over training and that this tree-likeness is predictive of better generalization on tests of compositional generalization, suggesting that in practice, transformers learn simpler, hierarchically structured computations, a property that directly influences the balance between flexible ICL and rigid IWL.

Finally, the breadth of in-context learnable functions is underscored by the ability of a single transformer to perform in-context algorithm selection. \cite{ref110} provided a comprehensive statistical theory showing that transformers can implement a broad class of standard ML algorithms in context, including least squares, ridge regression, Lasso, and gradient descent on two-layer neural networks. Building on these base algorithms, they proved that a single transformer can adaptively select between different ICL algorithms based on the input sequence, akin to a statistician choosing an appropriate method, using mechanisms like pre-ICL testing and post-ICL validation. This capability for adaptive algorithm selection points to a highly flexible and structured internal learning process that moves far beyond simple pattern matching toward a form of learned algorithm distillation and structure induction. As the subsequent section will explore, the extent to which such sophisticated ICL algorithms dominate over simpler in-weight memorization strategies is governed by a delicate competition shaped by training data properties, optimization dynamics, and architectural choices.

\subsection{The Interplay and Competition Between In-Context and In-Weights Learning}

Extending the hierarchy of in-context learnable functions discussed in the previous section, a central dynamic in understanding ICL is its competition and interplay with in-weights learning (IWL), the standard process of encoding knowledge directly into model parameters through gradient-based optimization. Rather than being independent capabilities, ICL and IWL represent distinct and often competing modes of generalization within a transformer, governed by a complex trade-off shaped by training data properties, optimization dynamics, and model architecture. This interplay directly determines the robustness of ICL under distribution shifts, as the balance between flexible in-context reasoning and rigid in-weights memorization dictates how models respond to data that diverges from their pretraining distribution, a challenge we will examine in the next section.

A pivotal finding in this interplay is the often transient nature of ICL. While ICL is frequently viewed as a persistent, emergent property, research demonstrates that it can be a temporary phase during training. In a study on the transient nature of ICL, transformers trained on synthetic data that allowed for both ICL and IWL strategies initially developed strong ICL capabilities, only to have them subsequently decay and be overtaken by IWL, even as the overall training loss continued to decrease \cite{ref73}. This indicates an asymptotic preference for memorization and IWL, raising fundamental questions about the training dynamics of large language models (LLMs) and whether standard pretraining protocols inadvertently drive models away from flexible, on-the-fly adaptation toward heavier reliance on stored knowledge. The transient ICL phenomenon is not a peculiarity of small-scale settings; it was observed consistently across a range of model sizes and datasets, underscoring its potential relevance to the development of frontier models \cite{ref73}.

This competition between ICL and IWL is fundamentally a battle of circuits and optimization landscapes, echoing the functional specialization and circuit-level analyses introduced in our architectural and algorithmic discussions. The decay of ICL is posited to be caused by direct competition between ICL and IWL circuits within the network, where weight updates that improve the IWL pathway can dismantle or overshadow the more fragile mechanisms that underpin ICL, such as the induction heads and task vector compression mechanisms described in prior sections \cite{ref73}. This circuit-level conflict is a manifestation of the broader plasticity-stability dilemma common in neural networks, but with a specific in-context character. The strong pull of prior knowledge learned during pretraining, a defining feature of IWL, can ossify model predictions and prevent the full integration of novel information presented in a prompt. This effect is particularly pronounced in tasks where the in-context examples contradict strong model priors, leading to performance saturation at suboptimal levels and a failure to fully learn from contrastive demonstrations \cite{ref29}. The model's existing ``in-weights'' priors exert such a powerful force that they can override the in-context signals, limiting the model's ability to adapt to new tasks or perspectives.

The training objective and data composition serve as the primary governors of the ICL-IWL trade-off. The properties of the pretraining data directly dictate which learning mode is favored, building on the principle introduced earlier that transformers can learn to implement algorithms like ridge regression and Newton's method when the data distribution incentivizes such structured inference. Data distributions with high burstiness, long-range coherence, and a prevalence of challenging examples where contextual information is essential for prediction are believed to create an environment where ICL is a more effective strategy, thereby encouraging its emergence \cite{ref41}. Conversely, data that is more repetitive or allows for simple statistical memorization may favor IWL. The concept of task diversity is critical here; pretraining on a sufficiently diverse distribution of tasks can force a model to develop ICL as a meta-learning strategy to handle new, unseen tasks, whereas a narrow task distribution favors IWL-like specialization \cite{ref48}. This principle extends to concept-aware training strategies, which explicitly construct training scenarios that reward analogical reasoning, thereby directly steering the model away from pure IWL and fostering more robust and persistent ICL abilities, much like the structured, compositional pretraining distributions that give rise to hierarchical task vector representations \cite{ref68}.

Given the competitive and often transient nature of ICL, deliberate training interventions are crucial for preserving and stabilizing it. Early stopping, based on the performance of ICL-style validation tasks rather than solely on next-token prediction loss, has been identified as a straightforward method to capture the model at its peak ICL capability before it is overtrained into IWL \cite{ref73}. Even more effectively, L2 regularization appears to fundamentally alter the learning dynamics in a way that favors ICL. Empirical work has shown that applying weight decay or L2 penalties can lead to the emergence of more persistent ICL, effectively removing the need for early stopping and preventing the later decay into memorization, suggesting that controlling the representational capacity and stability of the network through regularization is a key lever for promoting flexible in-context adaptation \cite{ref73}. This aligns with the broader understanding that ICL, as a form of implicit meta-learning, may thrive under conditions that prevent overfitting and encourage the learning of adaptable, general-purpose representations rather than task-specific shortcuts.

Interestingly, the balance can also be nudged through data curation alone, without any architectural or algorithmic changes. The extreme sensitivity of ICL to the choice of in-context demonstrations can be reframed as an instability in the ICL mechanism itself, where the model struggles to consistently apply the correct in-context strategy. By carefully curating a stable subset of training examples---identified through metrics like their average dev-set ICL accuracy when combined with random training examples---the ICL process can be significantly stabilized \cite{ref34}. These carefully chosen subsets do not necessarily need to be more diverse or have lower perplexity; they simply form a more reliable basis for the model to latch onto and consistently apply its in-context reasoning, strengthening the ICL signal against the background noise of competing IWL priors.

The relationship between ICL and IWL is not always purely competitive; they can coexist and even be synergistic under certain training regimes, directly informing how models will later respond to distribution shifts. The scaling of language data and models reveals a complementary dynamic: while scaled-down data allows ICL-like abilities to emerge in dramatically smaller models, demonstrating that the core computational capability is not strictly tied to massive scale, the largest models exhibit the most dominant, yet sometimes brittle, ICL behaviors \cite{ref16}. Multi-task learning and finetuning frameworks provide further insight, showing that sequential pretraining and finetuning (PT+FT) induces a hybrid operational regime that combines feature learning (characteristic of ICL's flexibility) with a ``nested feature selection'' bias from IWL \cite{ref111}. This suggests that IWL can be guided to extract and reuse a sparse subset of pretrained features, creating a division of labor where IWL provides a robust base of compressed knowledge, and ICL provides a dynamic, non-parametric interface for rapid, example-driven adaptation on top of that base. Ultimately, designing optimal learning systems requires moving beyond viewing ICL and IWL as mutually exclusive forces and instead engineering their dynamic balance through informed choices in data composition, training curricula, and regularization. This interplay is a defining characteristic of modern foundation models, dictating their ability to shift seamlessly from recalling memorized facts to performing novel, compositional reasoning from a few prompts, and as we will see in the following section, fundamentally governing their robustness when confronted with distribution shifts that challenge their pretrained priors.

\subsection{ICL Under Distribution Shifts and Pretraining Data Mixtures}

Building on the complex interplay between in-context learning (ICL) and in-weights learning (IWL) explored in the previous section, we now examine how this dynamic determines the robustness of ICL under distribution shifts. The generalization capabilities of ICL when confronted with data that diverges from the pretraining distribution is a critical determinant of its practical utility, as real-world deployments inevitably involve covariate shifts, out-of-domain task families, and noisy label associations. This section assesses how the composition of pretraining data mixtures fundamentally enables or limits in-context model selection and adaptation, directly connecting the data-centric governance of the ICL-IWL trade-off to out-of-distribution robustness.

A central finding is that ICL is not uniformly robust to distribution shifts. While models can often adapt to mild covariate shifts---where the input distribution changes but the conditional label distribution remains the same---their performance degrades under more severe or complex shifts. For instance, evaluations of large language models on syntactic tasks reveal that ICL generalization is inconsistent and heavily dependent on the specific model and pretraining corpus, with models often relying on superficial heuristics that fail under out-of-distribution (OOD) syntactic structures \cite{ref112}. This fragility extends to multimodal settings, where foundational models like GPT-4V(ision) exhibit significant performance variance across diverse datasets spanning natural, medical, and molecular domains. While these models show some zero-shot generalization capabilities, their limitations become apparent under controlled data perturbations, and ICL is shown to be vulnerable to domain shifts, label shifts, and spurious correlation shifts between the in-context examples and the test data \cite{ref113,ref114}.

This sensitivity to distribution shift directly extends the circuit-level competition between ICL and IWL discussed earlier. The strong prior knowledge encoded via IWL can exert a ``strong pull'' on predictions, causing the model to ossify its predictions around its pre-training domain and fail to fully integrate information from demonstrations that contrast with these priors. This effect is particularly pronounced in emotion recognition tasks, where it becomes more severe with larger models, suggesting that ICL in larger models may actually be less adaptive when the task distribution deviates significantly from the pretraining data \cite{ref29}. Similarly, in the context of majority label bias, where the distribution of labels in the in-context examples is skewed, the robustness of different LLMs varies widely. While some models exhibit high robustness (up to 90\%), others are significantly more susceptible, and this robustness is influenced by model size and the richness of instructional prompts, highlighting a complex interaction between model scale, prompt design, and distributional robustness \cite{ref40}.

The concept of hidden covariate shift, where the goal is to learn a representation such that the label distribution conditioned on this representation is domain invariant, provides a theoretical grounding for understanding the challenges of ICL under distribution change. While domain-invariant representation learning is a common approach, it relies on the assumption that such representations are well-suited for the target task, an assumption that can be violated in practice \cite{ref115}. In the tabular domain, robust bias-aware prediction under covariate shift has been extended through kernel methods, which reweight the source distribution to minimize the regularized expected target risk, offering consistency guarantees and improved performance on datasets with natural covariate shift \cite{ref116}. These theoretical frameworks underscore the difficulty of learning a predictor that is min-max optimal under the worst-case covariate shift, a challenge that is directly relevant to the implicit optimization performed by ICL.

A critical factor governing the robustness of ICL to unseen tasks---and one that ties back to the governing role of pretraining data discussed in the previous section---is the diversity and composition of the pretraining data mixture. The ability of transformers to perform in-context model selection, that is, identifying the correct task family from a prompt and then learning within it, is closely tied to the coverage of these task families in the pretraining data. Empirical studies demonstrate that transformers exhibit near-optimal model selection capabilities when the task families are well-represented in their pretraining data mixtures. However, when presented with tasks or functions that fall outside this pretraining distribution, even for simple extrapolation tasks, transformers exhibit various failure modes and a degradation of generalization. This suggests that the impressive ICL abilities of high-capacity sequence models may be more dependent on the coverage of their pretraining data mixtures than on fundamental, architecture-inherent inductive biases for generalization \cite{ref117}.

This finding is further reinforced by research on the emergence of ICL through task diversity, extending the principle introduced earlier that pretraining task diversity acts as a governor of the ICL-IWL trade-off. A task diversity threshold has been identified for the emergence of ICL on linear regression; below this threshold, a pretrained transformer behaves like a Bayesian estimator with the non-diverse pretraining distribution as its prior and cannot solve unseen tasks. Above the threshold, the transformer significantly outperforms this Bayesian estimator and aligns its behavior with ridge regression, corresponding to a Gaussian prior over all tasks, including those unseen during pretraining. This demonstrates that transformers can optimally solve fundamentally new tasks in-context, but only when the pretraining task diversity is sufficiently high, causing the model to deviate from the Bayes optimal estimator of the pretraining distribution \cite{ref48}. The theoretical underpinnings of this are echoed in statistical task complexity bounds for linear attention models, which show that effective pretraining for ICL requires only a small number of independent tasks to achieve near-Bayes-optimal risk on unseen tasks, corroborating the experimental findings on the importance of task diversity \cite{ref49}.

The transient nature of ICL and its mediation by pretraining data properties further determines robustness under distribution shifts. As discussed in the previous section, supportive pretraining data for ICL---characterized by high mass of rarely occurring, long-tail tokens and challenging examples where information gain from long-range context is low---encourages the model to rely on genuine in-context reasoning rather than memorizing spurious patterns \cite{ref41}. This suggests that curating pretraining data to emphasize long-range coherence and burstiness, the very properties that favor ICL over IWL, can simultaneously harden the model against distribution shifts that would otherwise break brittle in-weights associations. The presence of diverse and mutually supportive data sources in the pretraining mixture can also enable a form of domain generalization, where the model learns to exploit the diverse knowledge contained in various pretrained models to achieve enhanced OOD generalization performance \cite{ref118}.

In summary, the robustness of ICL under distribution shifts is a direct function of the interplay between the pretraining data mixture, the degree of task diversity, and the nature of the shift itself---all factors that govern the ICL-IWL competition. While ICL can exhibit remarkable generalization to new tasks when pretrained on sufficiently diverse data, it remains brittle to distribution shifts that are not well-represented in the pretraining corpus, and its performance is highly sensitive to the specific model, the prompt design, and the alignment between the strong pretraining priors and the target task. A key open challenge is to move beyond the coverage of pretraining data mixtures and develop architectures and training strategies that enable systematic, compositional generalization to fundamentally novel tasks. As we will see in the next section, hybrid architectures that combine complementary computational primitives offer one promising direction for addressing these limitations by expanding the functional repertoire available for in-context adaptation.

\subsection{Hybrid Architectures and Design Principles for Enhanced ICL}

Building on the critical role of pretraining data diversity and task coverage in determining ICL robustness under distribution shifts, we now turn to architectural innovations that can overcome the inherent limitations exposed in the previous section. If the brittleness of ICL under novel distributions is exacerbated by the rigid priors encoded via in-weights learning, then expanding the model's functional repertoire through hybrid architectures offers a direct path to more flexible, generalizable in-context adaptation. The limitations of individual architectures---transformers' quadratic complexity and state-space models' (SSMs) struggles with specific in-context retrieval tasks---have motivated a surge of research into hybrid models that synergistically combine complementary mechanisms, thereby addressing both the scalability and generalization bottlenecks that the ICL-IWL competition reveals.

A central finding driving this direction is that SSMs like Mamba exhibit strong performance on standard regression ICL tasks but fall short in tasks requiring non-standard retrieval, a gap that hybrid models can effectively bridge \cite{ref88}. The MambaFormer architecture, which interleaves Mamba blocks with attention blocks, exemplifies this principle, surpassing the ICL performance of either pure model on tasks where they struggle independently \cite{ref88}. This demonstrates that the distinct computational primitives---gated convolution and input-dependent selection in SSMs versus query-key matching in attention---are not mutually exclusive but can be composed to expand the overall functional repertoire for in-context task solving, directly mitigating the narrow model selection capabilities that arise from limited pretraining data coverage.

A deeper analysis of this synergy is provided by the development of models like H3 (Hungry Hungry Hippos), which was explicitly designed to address the expressivity gap between SSMs and attention in language modeling. The gap was diagnosed through synthetic tasks revealing that SSMs struggle with recalling earlier tokens and comparing tokens across the sequence---both fundamental operations for many ICL applications and precisely the kind of long-range retrieval that enables generalization beyond pretraining task families \cite{ref119}. A hybrid 125M-parameter H3-attention model, which retains just two attention layers, surprisingly outperformed pure Transformers by 1.0 PPL on OpenWebText, while scaling up to 2.7B parameters yielded better perplexity and zero- and few-shot learning performance on a majority of SuperGLUE tasks \cite{ref119}. This work provides a crucial design principle: attention is not required uniformly across all layers; rather, strategic placement of a small number of attention layers can amplify the ICL capabilities of an otherwise efficient SSM backbone, effectively marrying training efficiency with the expressive power needed for robust generalization across diverse task distributions.

The transition from static to dynamic module activation marks a further refinement in hybrid design, mirroring the adaptive, context-dependent nature of ICL itself. The Sparse Modular Activation (SMA) mechanism enables neural networks to sparsely and dynamically activate sub-modules for each sequence element in a differentiable manner \cite{ref97}. The SeqBoat architecture implements this by coupling an SSM with a Gated Attention Unit (GAU), where the SMA constrains the GAU to perform local attention only on the inputs activated based on the SSM's state representations. This achieves linear inference complexity with a theoretically infinite attention span, while the learned sparse activation patterns reveal the precise amount of attention needed for each task \cite{ref97}. This principle of dynamic, state-dependent gating of attention generalizes the hard-coded hybrid approach, suggesting that models can learn to allocate expensive attention computations only when the simpler SSM pathway is insufficient, thereby optimizing the quality-efficiency trade-off for ICL tasks and providing a mechanism for context-dependent adaptation that parallels the ICL-IWL competition at the architectural level.

Beyond the SSM-attention dichotomy, the compositional nature of attention itself has been targeted for improvement through architectural innovation. The Compositional Attention mechanism disentangles the standard head structure into its two fundamental computations: search (query-key interaction) and retrieval (value matrix extraction) \cite{ref120}. By introducing a soft competition stage between the query-key combination and value pairing, the model can dynamically and flexibly compose search and retrieval in a context-dependent manner. This not only reduces redundancy in learned parameters but also enhances generalization, including in out-of-distribution settings, by allowing specialized heads to emerge based on the type of retrieval needed \cite{ref120}. This principle of disentangling and dynamically recomposing sub-operations within attention offers a pathway to more efficient and generalizable ICL by moving away from the rigid, monolithic head structure that can lock in brittle in-weights associations and limit adaptation to novel tasks.

The integration of recurrent inductive biases with self-attention has also proven effective for enhancing ICL, particularly through better modeling of hierarchical structure---a capacity that is essential for compositional generalization beyond flat pattern matching. Hybrid models combining self-attention networks with Ordered Neurons LSTM (ON-LSTM), which introduces a syntax-oriented inductive bias for tree-like composition, have outperformed both individual architectures and standard hybrid models on tasks requiring hierarchical reasoning \cite{ref121}. This suggests that the complementarity between attention and recurrence lies in part in their differing capacities for chunking and composing information---attention excels at flat, global context aggregation, while structured recurrence can capture the nested, compositional structures that underpin many reasoning tasks in ICL and enable systematic generalization to unseen task structures.

The design principles for enhancing ICL are not limited to internal architectural components but also extend to training strategies that align the behavior of heterogeneous models, addressing the sensitivity to demonstration selection that contributes to ICL's fragility under distribution shifts. Bidirectional Alignment (BiAlign) is a distillation method that improves the ICL capabilities of smaller models by fully leveraging the preferences of larger models for demonstration examples, in addition to aligning token-level output distributions \cite{ref122}. By incorporating a novel ranking loss that aligns input preferences, BiAlign ensures that the smaller model not only mimics the output of the larger model but also learns to value the same kinds of in-context demonstrations, leading to consistent improvements on tasks including language understanding, reasoning, and coding \cite{ref122}. This highlights a design principle that extends beyond the model's forward pass: the input selection and preference mechanisms, which are crucial for ICL's sensitivity to demonstrations, can be transferred and aligned across architectures, potentially hardening models against the distribution-dependent brittleness discussed in the previous section.

Efficiency-oriented design principles also emerge from the study of hybrid architectures. The concept of ``hybridizing'' convolution and attention through alternate layouts, rather than simply grafting them, allows all stages of a network to benefit from both spatial convolutions and content-dependent interactions, leading to state-of-the-art accuracy on parameter-limited benchmarks \cite{ref123}. Similarly, the use of key-only attention, which excludes query-key pairwise interactions in favor of a compute-efficient saliency-gate, models local-global interactions with linear complexity, demonstrating that the quadratic bottleneck of attention can be circumvented by redesigning the mechanism while preserving its ICL benefits \cite{ref123}. The Cross-Architecture Transfer Learning (XATL) framework further addresses the practical barrier to adopting new architectures by transferring weights from shared components of pre-trained transformers to linear-cost inference architectures, reducing training time by up to 2.5x and yielding stronger models on language modeling benchmarks \cite{ref124}. This principle of weight transfer and architectural modularity is critical for democratizing the development of efficient ICL systems that can maintain robust performance across diverse distributional scenarios without prohibitive computational costs.

In summary, the design principles for enhanced ICL emerging from hybrid architectures collectively address the central challenge raised in the previous section: how to move beyond the coverage of pretraining data mixtures and develop systems capable of systematic generalization. These principles include: (1) the strategic, non-uniform interleaving of attention and state-space or recurrent layers to combine complementary computational strengths; (2) the dynamic, state-dependent gating of expensive operations to optimize quality-efficiency trade-offs; (3) the disentanglement of monolithic attention operations into composable, context-dependent sub-functions; (4) the incorporation of hierarchical inductive biases from structured recurrence to improve compositional reasoning; (5) the alignment of input preferences and output distributions across architectures during training; and (6) the modular transfer of pre-trained weights to new, more efficient architectures. Collectively, these principles point toward a future of flexibly composed, highly efficient models that maintain the robust ICL capabilities of large transformers while overcoming their fundamental limitations in scalability and specialized reasoning under distribution shifts.

\section{Demonstration Selection, Ordering, and Prompt Engineering}

\subsection{Retrieval-Based Demonstration Selection}

Building upon these foundational strategies, retrieval-based demonstration selection forms the cornerstone of modern in-context learning (ICL), moving beyond the early paradigm of randomly or heuristically chosen few-shot examples. The fundamental insight driving this approach is that the performance of large language models (LLMs) as in-context learners is highly sensitive to the choice of demonstrations, and that selecting examples semantically or lexically similar to a test query can substantially improve task accuracy and robustness. Empirical evidence from retrieval-based ICL research demonstrates that even simple word-overlap similarity measures, such as BM25, can outperform randomly selected demonstrations, underscoring the latent benefits of aligning prompt context with the input distribution \cite{ref19}. This finding has catalyzed a broad spectrum of methods that operationalize similarity along multiple dimensions, ranging from lexical overlap and dense semantic embeddings to syntax-aware metrics, each tailored to the specific demands of the target task.

The most straightforward instantiation of retrieval-based demonstration selection relies on lexical similarity, where the surface-form overlap between the test input and candidate demonstrations is scored using algorithms like BM25 or TF-IDF. Despite their simplicity, lexical methods offer a strong baseline, particularly for tasks where shared vocabulary signals a direct topical or structural alignment with the query. The success of BM25 in retrieval-augmented ICL, even when applied to instruction-finetuned LLMs that have already been exposed to the training data during fine-tuning, highlights that test-time retrieval from the original training corpus can provide additional contextual cues that the model does not fully internalize in its parameters \cite{ref19}. Moreover, in specialized domains such as grammatical error correction (GEC), where the task is fundamentally syntax-oriented, standard word-matching or semantic similarity measures can be suboptimal. Instead, ungrammatical-syntax-based selection strategies that compute similarity based on the syntactic structures of sentences---targeting the most similar ill-formed syntax to the test input---have been shown to outperform generic retrieval methods, demonstrating the value of tailoring the similarity function to the linguistic properties of the task \cite{ref125}.

Dense semantic similarity, computed through neural embedding models, represents a more sophisticated and widely adopted family of retrieval strategies. By encoding both the test input and candidate demonstrations into a shared vector space using pre-trained sentence encoders, these methods capture deeper semantic relationships that lexical overlap can miss. This approach aligns naturally with the embedding-based retrieval architectures that have become standard in information retrieval pipelines. The effectiveness of dense retrieval for ICL is evidenced by the fact that semantically similar demonstrations, when coupled with chain-of-thought (CoT) prompting, can significantly boost the reasoning capabilities of LLMs on complex multi-modal and text-based tasks \cite{ref126}. The generalizability of dense retrieval is further enhanced by unified embedding models, such as the LLM-Embedder, which are specifically optimized for the diverse retrieval augmentation needs of LLMs through multi-task fine-tuning, reward formulation based on LLM feedback, and homogeneous in-batch negative sampling, thereby surpassing both general-purpose and task-specific retrievers across various scenarios \cite{ref127}.

The integration of retrieval-based demonstration selection with instruction-finetuned LLMs and CoT prompting presents both opportunities and notable refinements. For instruction-finetuned models, the retrieval of demonstrations from the very data on which the model was instruction-tuned yields better results than using no demonstrations or random examples, indicating that the model benefits from the explicit, localized context even when the knowledge is already implicitly represented in its weights \cite{ref19}. In the context of CoT reasoning, retrieval mechanisms can be extended to dynamically select demonstration examples based on cross-modal and intra-modal similarities, a strategy that is particularly powerful for multi-modal tasks. By employing stratified sampling to categorize demonstration examples into groups based on their reasoning types and then retrieving examples from different groups, the diversity of reasoning patterns in the prompt is promoted, leading to substantial gains on benchmarks like ScienceQA and MathVista for models such as GPT-4 and GPT-4V \cite{ref126}. This highlights a synergy between retrieval and prompt engineering: retrieval provides relevance, while strategic sampling ensures diversity and coverage of reasoning paradigms.

However, despite their empirical success, retrieval-based methods that score and select demonstrations independently based on similarity to the test input suffer from several critical limitations that motivate the more principled approaches discussed in the following subsection. The most prominent drawback is that independent similarity scoring ignores the crucial influence of the LLM's existing knowledge and decision boundaries. A demonstration that is semantically similar to the test input may not necessarily be informative for the model's learning process; in fact, the model may already handle such examples correctly, offering no marginal gain. This observation has motivated ambiguity-aware selection strategies, where the goal is to choose demonstrations that not only share semantic similarity with the test example but also help resolve the inherent label ambiguity surrounding it. Experiments across text classification tasks have shown that including demonstrations that the LLM previously misclassified---and that fall on the test example's decision boundary---yields the most significant performance improvements, as these examples are more effective at disambiguating the model's predictions \cite{ref36}. Similarly, a purely similarity-driven approach can be outperformed by selection strategies based on example difficulty, where the most challenging or informative demonstrations for the model are prioritized over those that are merely topically related, a principle that has proven effective for both passage ranking and question generation tasks \cite{ref128}.

Another fundamental limitation is that independent similarity scoring treats the demonstration set as a collection of isolated examples, neglecting the compositional and order-dependent nature of ICL. This can lead to a lack of diversity in the selected demonstrations, redundancy in the information provided to the model, and a failure to capture the full spectrum of reasoning strategies required for complex tasks. Active learning principles applied to ICL have further revealed that while similarity-based selection consistently outperforms random sampling and uncertainty-based methods, there remains a nuanced interplay between similarity, uncertainty, and diversity that independent scoring fails to optimize \cite{ref129}. The bias introduced by the model's strong priors and the label distributions in demonstrations further compounds this issue, as models can learn spurious correlations from the retrieved context rather than the intended input-label mapping, necessitating additional calibration techniques to mitigate demonstration shortcuts. These limitations naturally give rise to the more advanced strategies---including influence-based, information-gain, and ambiguity-aware selection methods---that move beyond independent similarity to consider the joint utility of demonstrations and their alignment with the model's internal learning dynamics.

\subsection{Influence, Information Gain, and Ambiguity-Aware Selection}

Building upon retrieval-based similarity measures, a more sophisticated class of demonstration selection strategies has emerged, moving beyond superficial lexical or semantic resemblance to quantify the true \emph{informativeness} of candidate examples. The critical limitations of independently scored similarity methods---their neglect of the model's existing knowledge, decision boundaries, and the joint utility of demonstrations---directly motivate these advanced methods. Grounded in principles from active learning, information theory, and decision theory, they aim to identify demonstrations that most effectively reduce model uncertainty, refine decision boundaries, or maximize the knowledge gained about a task. Central to this endeavor are influence-based metrics, maximum information gain, and ambiguity-aware selection, which provide a more principled lens for constructing few-shot prompts.

A foundational approach in this domain is the application of influence functions, borrowed from robust statistics, to estimate how a candidate demonstration would alter the model's predictions on a test set if it were included in the prompt. This method, which treats the in-context learning process as an implicit form of model training, allows for the selection of examples that have the highest potential positive influence on downstream performance. Similarly, ambiguity-aware selection strategies directly consider the model's decision boundary and its hesitant regions, directly addressing the observation from the previous subsection that demonstrations the LLM previously misclassified yield the most significant performance improvements \cite{ref36}. Instead of merely seeking examples similar to a test query, these methods probe for demonstrations that the model finds confusing or that lie near its class boundaries. The rationale is that such ambiguous examples are highly informative for resolving the task's structure. A model's prior misclassifications on a held-out validation set can serve as powerful indicators of epistemic uncertainty, guiding the selection of demonstrations that actively clarify the correct input-label mapping. This shifts the selection paradigm from a passive, query-centric view to an active, model-centric one, directly targeting and rectifying the model's knowledge gaps.

A particularly compelling framework quantifies this informativeness through the lens of information theory, specifically Information Gain (IG). In this context, IG measures the reduction in uncertainty about a test prediction after observing a candidate demonstration. The work by \cite{ref130} formally addresses the high variance of in-context learning by proposing to sample demonstrations that yield the maximum IG. This approach directly tackles a critical failure mode of random or simple similarity-based selection: even when all other prompt engineering factors are held constant, random example selection induces significant performance instability. By selecting examples based on their expected information reward, the IG framework provides a task-agnostic and theoretically grounded metric for gauging a demonstration's utility. However, its direct application can be confounded by a subtle but pervasive issue: the evaluation of a demonstration's informativeness can be unfairly skewed by the biases inherent in the specific prompt template used to present it.

This challenge, where a demonstration's apparent utility is an artifact of template phrasing rather than its true task-relevance, is known as template bias. The authors of \cite{ref130} identify this problem and introduce a ``Calibration Before Sampling'' strategy to mitigate it. The core idea is to first calibrate the raw predictive probabilities generated by the model for each template, effectively neutralizing the optimistic or pessimistic biases introduced by surface-level formatting. This calibration step ensures that the subsequent IG calculation reflects the genuine information carried by a demonstration's content, disentangling it from spurious correlations with prompt style. This integration of calibration with uncertainty-aware selection represents a critical maturation in prompt engineering, acknowledging that the measurement of informativeness is itself an inference procedure susceptible to model priors and presentation noise. The strategy ensures a fairer comparison among candidate demonstrations, leading to more robust and consistently informative few-shot prompts.

The principles of active learning provide a unifying framework for these advanced techniques. A systematic study connecting active learning heuristics to in-context learning is presented in \cite{ref129}. This work frames the single-iteration demonstration selection process as a pool-based active learning problem and rigorously evaluates classic strategies. A central and counter-intuitive finding is the stark performance disparity between uncertainty-based methods, a cornerstone of traditional active learning, and similarity-based methods in the in-context learning regime. While uncertainty sampling is designed to select examples the model is least confident about, \cite{ref129} consistently demonstrates that it performs poorly for in-context demonstration selection, often underperforming even random selection. In contrast, similarity-based selection, which chooses demonstrations that are most representative of the test query, proves remarkably effective. The analysis suggests that in-context learning, operating without parameter updates, benefits more from high-quality, prototypical examples that exhibit low uncertainty and high similarity to the test instance. This finding fundamentally reframes the notion of ``informative'' for frozen models: an informative demonstration is not necessarily a confusing one, as in traditional active learning, but one that provides a clear, canonical, and representative snapshot of the task mapping.

This insight aligns with the function of prototypical calibration, which aims to build a more robust internal model of task categories. The method proposed in \cite{ref131} moves beyond greedy decoding by estimating a Gaussian mixture distribution over prototypical clusters for each label. This can be viewed as a form of ambiguity-aware inference, where an example's prediction is calibrated by its likelihood relative to these learned prototypes, making the decision boundary adaptive and more robust to the specific permutations and templates that can otherwise cause fragile behavior. While this operates at the decoding stage rather than the selection stage, the underlying principle---leveraging a probabilistic model of category structure to handle ambiguity---is deeply connected to the ambiguity-aware selection strategies discussed above. It suggests a symbiotic relationship where ambiguity-aware selection chooses examples that help the model build a better internal task representation, and ambiguity-aware decoding techniques like prototypical calibration then leverage that representation for robust prediction.

In summary, the shift from similarity-based to influence-, information-gain-, and ambiguity-aware selection strategies represents a profound deepening of our understanding of in-context learning. These methods reframe the prompt from a passive collection of independently scored examples to an actively curated set of informative interventions that account for the model's internal learning dynamics. By borrowing from information theory, they provide a principled metric of a demonstration's utility; by incorporating calibration techniques, they address the measurement biases that can distort this metric; and by connecting to active learning, they uncover the unique behavioral properties of in-context learning, revealing that canonical, low-uncertainty examples are often more powerful than boundary cases. However, as with the similarity-based methods that preceded them, these approaches still treat demonstrations as individually informative units, overlooking the complex interdependencies, compositional effects, and order-dependent dynamics that arise when examples are assembled into a complete prompt---limitations that directly motivate the sequential and reinforcement learning-based selection strategies explored in the following subsection.

\subsection{Sequential and Reinforcement Learning for Demonstration Retrieval}

While conventional similarity-based methods score candidates independently and the advanced informativeness metrics from the previous subsection quantify a single example's potential utility, a fundamental limitation remains: both paradigms overlook the complex interdependencies between demonstrations within a prompt. The composition of a prompt set---and critically, the order in which examples are presented---creates a context that profoundly shapes a model's interpretation and reasoning trajectory. This has motivated the development of sequential and reinforcement learning (RL) approaches that frame demonstration retrieval not as a set of independent selections, but as a sequential decision-making process, where the choice of each subsequent example is conditioned on the preceding ones.

The core idea is to model the construction of a prompt as a trajectory in a Markov Decision Process (MDP). In this framework, the state represents the current set of selected demonstrations and the query input, the action is the selection of the next demonstration from a candidate pool, and the reward is a signal reflecting the performance of the final prompt, such as the likelihood of generating the correct answer or a task-specific evaluation metric. This formulation naturally captures the order-dependent effects that are critical in ICL: the information provided in early examples establishes the interpretive frame for later ones, and a poorly ordered sequence can undermine even an otherwise informative set of demonstrations. By maximizing cumulative reward, an RL agent can learn a policy for sequentially assembling a highly effective prompt, directly optimizing for the collective, ordered impact that set-level diversity and coverage metrics seek to measure.

A prominent example of this approach is the Retrieval for In-Context Learning (RetICL) framework \cite{ref132}. RetICL explicitly models the sequential selection of examples, training a retriever using reinforcement learning to optimize for the performance of the final downstream task. Unlike heuristic or independent scoring methods, the RetICL agent learns to select an example based on the current state of the partially built prompt. This allows it to choose examples that complement each other, creating a coherent chain of reasoning rather than a redundant collection of individually similar instances. The learned policy can, for instance, select an initial example that establishes a problem-solving framework, followed by a second example that addresses a specific nuance---a strategy that case studies have shown RetICL to implicitly learn \cite{ref132}. This stands in sharp contrast to similarity-based methods that might retrieve several examples all highly similar to the query but redundant with each other, failing to provide the diverse perspectives or structural cues that a sequential model can learn to assemble, thereby addressing the very redundancy problem that motivates set-level optimization.

Building upon this sequential paradigm, the \(Se^2\) (Sequential Example Selection) method further refines the process by leveraging the large language model's own feedback to guide the search \cite{ref133}. \(Se^2\) formulates the problem as a sequential selection task and utilizes beam search to explore the space of high-quality example sequences. At each step of sequence construction, the LLM's evaluation of the partially formed prompt serves as a dynamic signal, helping the beam search to capture both the inter-relationships between examples and the sequential information that is crucial for effective ICL. This approach not only enhances the quality and diversity of the selected sequences compared to random or greedy selection but also ensures a consistency between the training objective and inference-time usage that is often missing in ``select then organize'' paradigms. The beam search effectively acts as a planning mechanism, looking ahead to find sequences of demonstrations whose combined effect on the LLM is maximally beneficial, significantly enriching the contextuality of the prompt in a manner that implicitly satisfies the dual goals of diversity and coverage that the following subsection explores in depth \cite{ref133}.

The use of reinforcement learning for discrete prompt optimization extends beyond retrieval from a fixed dataset to the generation of synthetic prompts. Policy gradient methods, a class of RL algorithms, can be used to optimize a discrete policy that constructs prompts token by token. Although not directly covered in the provided papers on demonstration retrieval, this concept is closely related to the broader theme of using sequential decision-making for prompt engineering. The idea is to bridge the gap between the implicit gradients computed by the LLM during in-context learning and explicit parameter updates, a principle seen in methods that treat ICL as a form of meta-optimization \cite{ref134}. This connection strengthens the theoretical grounding for using RL in prompt optimization, as it positions the learned retriever as a meta-learner that is learning to optimize the context for a frozen ``inner-loop'' optimization process occurring within the transformer's forward pass.

Furthermore, the sequential reasoning central to these retrieval methods mirrors the step-by-step nature of many complex tasks and aligns naturally with the compositional task structures that coverage-based methods aim to capture. In domains like mathematical problem-solving, a model must follow a logical chain of deductions, and the selection of in-context examples can be seen as providing a template for this reasoning chain. A sequential retriever can learn to select examples that build a scaffolding for the required reasoning strategy, progressing from simpler to more complex sub-problems or from known concepts to novel applications. This aligns with findings in imitation and reinforcement learning that show the power of hierarchical guidance and skill acquisition for solving long-horizon problems \cite{ref135,ref136}. The retrieval process itself becomes a form of high-level planning, where the agent decomposes the task of ``providing a good prompt'' into a sequence of sub-decisions about which example to add next---a strategy that echoes the divide-and-conquer approach to planning in complex state spaces \cite{ref137}.

In summary, sequential and reinforcement learning-based approaches for demonstration retrieval represent a significant advance over both independent selection methods and the informativeness metrics discussed previously. By formalizing the selection problem as an MDP and leveraging RL algorithms like policy gradients, or by using beam search guided by LLM feedback, these methods explicitly model the crucial interdependencies and order effects within a prompt. They transform the problem from simple similarity matching or uncertainty quantification into a form of compositional planning, where the goal is to construct a coherent, sequentially-ordered context that guides the LLM through a desired reasoning trajectory to the correct answer. This sequential perspective provides a natural bridge to the set-level optimization strategies that follow, which formalize the collective objectives---diversity and coverage---that sequential methods implicitly learn to satisfy through their trajectory-based optimization.

\subsection{Diversity, Coverage, and Set-Level Optimization}

Building upon the sequential and reinforcement learning frameworks that model the interdependencies and ordering of demonstrations, a paradigm shift toward set-level optimization has gained prominence, focusing on the collective properties of the demonstration set---namely diversity, coverage, and their interplay---to maximize in-context performance. While sequential methods implicitly learn to satisfy these properties through trajectory-based optimization, this subsection explicitly analyzes the metrics and algorithms that ensure a set of demonstrations is both diverse and has high coverage of the salient aspects of the test input, avoiding redundancy and better representing complex reasoning patterns. These set-level strategies formalize the very objectives that sequential retrievers approximate heuristically.

A foundational concept in this area is the move from instance-level to set-level metrics. Standard methods rank and select examples independently, which is efficient but fails to account for inter-example relationships. The work on coverage-based example selection for ICL directly confronts this limitation by introducing a set-level objective. It demonstrates that metrics like BERTScore-Recall (BSR) serve as superior selectors for individual examples because they prioritize demonstrations that cover more of the salient content, such as reasoning patterns, of the test input. Critically, the authors extend BSR to a set-level metric, Set-BSR, which explicitly optimizes the collective coverage of the chosen subset. This approach ensures that the selected demonstrations, as a group, encapsulate a broader spectrum of the input's features, leading to substantial performance gains, particularly on compositional tasks where the method outperforms independent ranking by a significant margin without any task-specific training \cite{ref20}.

This principle of coverage is conceptually aligned with techniques developed in coreset selection for data-efficient training. While coreset selection traditionally aims to choose a subset of training data that preserves model performance, its underlying theories of diversity and coverage are directly applicable to demonstration selection. The critical challenge, as identified in coreset literature, is that at high pruning rates, selection methods prioritizing only sample importance suffer catastrophic accuracy drops due to poor data coverage. The proposed Coverage-centric Coreset Selection (CCS) method addresses this by jointly considering overall data coverage and sample importance, an insight that resonates strongly with the goals of ICL demonstration set optimization. The concept of maximizing coverage to ensure the selected set is representative of a broader distribution is a core component of building informative prompts \cite{ref138}.

Beyond coverage, maintaining explicit diversity among demonstrations is crucial to prevent the model from overfitting to a single, narrow pattern. Determinantal Point Processes (DPPs) offer a mathematically elegant framework for this, as they model the probability of selecting a subset where the marginal gain of an item decreases if similar items are already present, thus naturally promoting diversity. The connection between DPPs and the very notion of task-oriented diversity is further unpacked through the lens of Rate-Distortion (RD) theory. The RD-DPP framework establishes a fundamental link, observing a phase transition in the diversity of DPP-selected data, which suggests that diversity-focused selection is most beneficial at the initial stages of sample accumulation before transitioning to an uncertainty-based criterion. This bi-modal strategy provides a principled way to balance diversity with task-specific informativeness, a concept that can be generalized to the ordering and composition of ICL prompts \cite{ref139}.

The need for diversity is also about robustly representing the underlying problem space. In the context of reinforcement learning and unsupervised skill discovery, algorithms based on DPPs, such as ODPP, have been designed to unify the dual goals of diversity and coverage. This work explicitly quantifies both properties through a DPP and optimizes them to discover a set of skills that are both distinct and wide-ranging. This dual focus on diversity and coverage in skill spaces mirrors the objective in ICL to select demonstrations that showcase a wide variety of valid reasoning paths (coverage) while doing so through structurally distinct mechanisms (diversity), thereby enhancing the model's robustness to different query formulations \cite{ref140}.

A specific technique for integrating these concepts into language model prompting is Contrastive Demonstration Tuning. This method moves beyond sampling to introduce a pluggable, extensible approach for tuning demonstrations. By employing a contrastive framework, the model is implicitly trained to distinguish between more and less effective demonstration sets, a process that inherently captures the relational properties of the set as a whole. This approach directly targets the variance introduced by random sampling and aims to produce a more stable and optimized set of demonstrations, improving performance when combined with standard prompt-tuning methods \cite{ref141}.

The logic of set-level diversity is further validated by research showing that decomposing complex problems into a larger set of sub-objectives can implicitly maintain diversity without explicit mechanisms. In the domain of deceptive problem-solving, optimizing a potentially large set of extracted objectives using lexicase selection was found to outperform explicit quality-diversity algorithms like MAP-Elites. This suggests that ensuring the prompt covers multiple, decomposed aspects of a task can be more effective than a monolithic, single-objective approach to demonstration selection, as the model can leverage diverse in-context signals to navigate complex reasoning paths \cite{ref142}.

The practical challenge of measuring diversity in high-dimensional feature spaces is addressed by introducing metrics that go beyond simple aggregation. The Contributing Dimension Structure (CDS) metric identifies that standard similarity measures obscure the contribution of individual feature dimensions, causing traditional methods to fail in adequately capturing true diversity. By enforcing a diversity constraint based on CDS, one can select a coreset with a richer variety of feature-dimension structures, a principle that can be extended to selecting demonstrations that differ not just in superficial wording but in their underlying feature-level representations, thus correcting a hidden failure mode of naive diversity metrics \cite{ref143}.

In summary, the optimization of demonstration sets for ICL is evolving from independent similarity-based retrieval to holistic, set-level strategies that jointly model diversity and coverage. Techniques ranging from set-level metrics like Set-BSR, to mathematical frameworks like DPPs, to training-based contrastive tuning, and refined metrics like CDS, all converge on a central insight: a demonstration set is more than the sum of its parts. Its effectiveness is determined by the collective information it covers and the absence of redundant signals, ensuring the model is exposed to a rich and non-overlapping set of reasoning patterns that robustly characterize the target task, a principle that aligns with findings on the importance of data diversity for robust instruction tuning \cite{ref144}. This set-level perspective directly sets the stage for the following subsection, which extends the logic of diversity and coverage from the selection of whole demonstrations to the internal compositional quality of their reasoning chains through chain-of-thought prompting strategies.

\subsection{Chain-of-Thought, Pattern-Aware, and Synthetic Prompting}

While the previous subsection examined the holistic optimization of demonstration sets through diversity and coverage, the structure of the demonstrations themselves---their internal reasoning patterns---introduces an orthogonal and equally critical dimension for in-context learning. Chain-of-thought (CoT) prompting has emerged as a cornerstone technique for eliciting complex reasoning from large language models (LLMs) by structuring prompts to include intermediate reasoning steps. Although the manual crafting of CoT demonstrations has proven highly effective, it is labor-intensive and brittle. Consequently, significant research has focused on automating and enhancing CoT prompting through pattern-aware strategies, latent skill discovery, and synthetic demonstration generation, alongside ensemble optimization methods to improve reasoning stability. These efforts extend the set-level logic of the previous subsection from the selection of whole demonstrations to the internal compositional quality of the reasoning chains themselves.

A central challenge in CoT prompting is the selection of demonstrations that best guide the model toward correct reasoning, directly paralleling the shift from instance-level to set-level metrics discussed earlier. The effectiveness of these demonstrations is not solely dependent on their semantic accuracy; the underlying reasoning patterns they embody are equally crucial. The Pattern-Aware CoT (PA-CoT) method addresses this by explicitly considering the diversity of reasoning patterns in demonstrations, such as step length and the structure of intermediate reasoning processes \cite{ref145}. By incorporating a diverse set of patterns, PA-CoT mitigates the biases induced by a narrow set of demonstrations and enhances generalization to diverse problem scenarios. This approach recognizes that the \emph{how} of reasoning---the procedural structure---is as important as the \emph{what}---the factual content---in guiding the model's own reasoning generation, complementing the coverage and diversity principles that govern optimal demonstration set composition.

Moving beyond predefined patterns, another line of work focuses on discovering the latent reasoning skills required for a given task, extending the concept of coverage to the space of cognitive operations. The Reasoning Skill Discovery (RSD) method introduces a technique for creating a latent space representation of rationales, effectively decomposing the reasoning process into a set of intrinsic skills \cite{ref146}. RSD uses unsupervised learning to map rationales to this skill space and then learns a reasoning policy to determine the requisite skills for a novel question. This policy then guides the selection of demonstrations that exemplify those specific skills, ensuring that the chosen examples are pedagogically aligned with the cognitive demands of the target problem. This approach is theoretically grounded, sample-efficient, and model-agnostic, offering a sophisticated alternative to heuristic-based selection by focusing on the compositional skill requirements of complex reasoning tasks.

Synthetic prompting represents a paradigm shift by automating the generation of CoT demonstrations entirely, reducing the reliance on human-annotated exemplars and providing a principled way to generate the diverse and coverage-rich sets advocated in the previous subsection. The Synthetic Prompting method leverages a small set of handcrafted examples to bootstrap a model into generating a larger pool of diverse and effective demonstrations through a backward-forward generation process \cite{ref147}. In the backward process, the model generates a question that plausibly matches a sampled reasoning chain, ensuring the question is solvable and well-defined. The forward process then refines the reasoning chain for the newly generated question, improving its quality and detail. This iterative cycle produces a set of high-quality, synthetic exemplars that can outperform manually crafted prompts on numerical, symbolic, and algorithmic reasoning tasks. Similarly, the Automate-CoT strategy automates the augmentation and selection of rational chains from a small labeled dataset, using a variance-reduced policy gradient to select the optimal combination of machine-generated rationales for prompting \cite{ref148}. This approach bypasses the need for human engineering of CoT examples, enabling quick adaptation to new tasks. Moreover, the Self-Contemplation Prompting strategy (SEC) entirely eliminates the need for external demonstrations by having the LLM first create its own demonstrations based on the query before generating the final output \cite{ref5}. This self-referential approach challenges the necessity of human-crafted data, demonstrating that contemporary LLMs possess a sufficient internal competence to guide their own reasoning in many scenarios.

A critical vulnerability of naive greedy decoding in CoT prompting is its tendency toward repetitiveness and local optimality, which can lead to incorrect answers even when the model possesses the requisite knowledge. Ensemble-optimization methods address this by sampling multiple diverse reasoning paths and then aggregating them to find the most consistent answer, effectively applying the diversity principle at the output level. The foundational technique in this area is Self-Consistency, which replaces greedy decoding by sampling a diverse set of reasoning paths and selecting the answer that is the most consistent across all paths, based on the intuition that a complex problem admits multiple valid ways of reasoning to a single correct answer \cite{ref149}. This simple yet powerful marginalization strategy leads to significant performance gains.

While highly effective for tasks with structured answer formats, Self-Consistency's reliance on exact answer extraction limits its applicability to free-form generation tasks. To overcome this, Universal Self-Consistency (USC) leverages the LLM itself as the arbiter of consistency \cite{ref150}. USC prompts the model to select the most consistent response from a pool of candidate outputs, thereby extending the benefits of ensemble optimization to open-ended tasks like code generation, long-context summarization, and open-ended question answering. The Self-Agreement method further generalizes this approach by prompting the model to determine the optimal answer by selecting the most agreed-upon answer among a diverse set of sampled reasoning paths, proving effective across scenarios where the answer format is either known or unknown \cite{ref151}. These methods collectively refine the raw output of CoT through a process of consensus-building, significantly improving the stability and accuracy of model reasoning.

The iterative refinement of reasoning is another powerful axis of improvement that builds directly on the output of ensemble methods. Progressive-Hint Prompting (PHP) enables automatic multi-turn interactions where the model's previous answers serve as progressive hints to guide the model toward the correct answer in subsequent turns \cite{ref152}. This approach is orthogonal to CoT and Self-Consistency, and when combined, it yields state-of-the-art results by leveraging the model's own output as a dynamic scaffold for correction. The Iterative Bootstrapping in Chain-of-Thoughts (Iter-CoT) approach similarly focuses on autonomously rectifying errors in reasoning chains and selecting exemplars of moderate difficulty to enhance generalizability \cite{ref153}. This iterative self-correction mechanism is fundamental to closing the gap between a model's initial reasoning attempt and a verifiably correct solution.

In summary, the frontier of CoT prompting is defined by a transition from static, hand-crafted demonstrations to dynamic, automated, and self-improving reasoning systems. By integrating pattern-aware selection, latent skill discovery, and synthetic generation, these methods create more robust and generalizable prompts that embody the set-level principles of diversity and coverage within the internal structure of reasoning chains. Coupled with ensemble-optimization and iterative refinement strategies that build consensus and correct errors, these advanced techniques form a powerful toolkit for stabilizing and enhancing the complex reasoning capabilities of large language models. Together with the structured and soft-prompt engineering methods discussed in the following subsection, they represent complementary approaches to moving beyond discrete, hand-crafted demonstrations toward more sophisticated and adaptive prompt formulations.

\subsection{Structured, Soft-Token, and Hierarchical Prompt Engineering}

Whereas the previous subsection explored the automated construction and internal structuring of reasoning chains within demonstrations, a parallel and highly active research thrust focuses on the design and optimization of the prompt's fundamental building blocks themselves, moving from hand-crafted natural language instructions to learned, structured, and hierarchical formulations. This line of work addresses the brittleness of manual prompt engineering and complements chain-of-thought methods by introducing parameter-efficient, data-driven techniques that can discover optimal prompt representations in both continuous and discrete spaces, further extending the theme of moving beyond static, hand-crafted elements.

The core innovation in this area is the shift from discrete ``hard'' prompts to continuous ``soft'' prompts, which are learnable vectors prepended to the model's input sequence. These soft prompts are optimized via gradient descent on a small amount of task-specific data while keeping the pre-trained language model entirely frozen. This paradigm, known as prompt tuning, was established as a parameter-efficient alternative to full fine-tuning, capable of matching its performance while modifying only a tiny fraction of the total parameters \cite{ref154}. However, a key challenge identified in this framework is that a single, static soft prompt may not generalize well to the diverse inputs within a task. To address this, the Vector-quantized Input-contextualized Prompts (VIP) framework extends soft prompt tuning by conditioning the prompt on the specific input instance \cite{ref155}. VIP employs a small sentence encoder to generate contextualized prompt tokens, which are then mapped to a learned codebook of vectors through a vector quantization network. This process not only captures instance-specific information but also leads to more robust prompt representations, demonstrating significant improvements in out-of-domain generalization and multi-task settings.

The sensitivity of soft prompt tuning to initialization and its tendency to overfit in few-shot scenarios have motivated the integration of meta-learning into the prompt engineering pipeline. The central idea is to learn a universal prompt initialization that can rapidly adapt to new, unseen tasks. The MetaPrompting approach directly applies a model-agnostic meta-learning algorithm to find a better initial set of soft prompt parameters, enabling fast adaptation with superior performance over randomly initialized or hand-crafted prompts \cite{ref156}. Building on this, the SUPMER framework introduces a self-supervised meta-learning approach that not only learns a universal prompt initialization from unlabeled data but also jointly meta-learns a gradient regularization function \cite{ref157}. This regularizer transforms raw task-specific gradients into a domain-generalizable direction, explicitly combating overfitting and enhancing cross-domain performance. A similar concept is extended to the vision-language domain by the GRAM framework, which meta-learns both a soft prompt initialization and a lightweight gradient-regulating function, demonstrating that such regularization is a model-agnostic plugin that can consistently improve various prompt tuning methods \cite{ref158}.

To further enhance the capacity of prompt tuning without a corresponding increase in computational overhead, structured and decomposed approaches to parameterization have been proposed. Residual Prompt Tuning reparameterizes soft prompt embeddings using a shallow network with a residual connection, which significantly improves performance and stability, allowing for a 10x reduction in prompt length without sacrificing accuracy \cite{ref159}. Similarly, Decomposed Prompt Tuning (DePT) decomposes a standard soft prompt into a shorter soft prompt and a pair of low-rank matrices optimized with different learning rates \cite{ref160}. This decomposition addresses the sequence length and computational cost issues introduced by long soft prompts, achieving superior performance with substantial memory and time savings, particularly as the underlying model size scales. The concept of a ``prompt pool'' offers a different form of structure, where a set of key-value prompts is learned and combined via instance-aware attention. This approach, as seen in MetaPrompter, allows different instances to draw upon distinct task knowledge encoded in the pool, constructing instance-dependent prompts that outperform a single shared initialization \cite{ref161}.

Hierarchical and multi-level prompt structures represent a further advancement, designed to capture task, domain, and context information at different granularities. The MPrompt method for machine reading comprehension exemplifies this by utilizing task-specific, domain-specific, and context-specific prompts \cite{ref162}. By enforcing an independence constraint among domain-specific prompts and using a prompt generator that incorporates contextual knowledge, this hierarchical approach achieves a richer and more nuanced conditioning of the pre-trained model, leading to state-of-the-art results. Discourse-Aware Soft Prompting similarly applies a structural inductive bias to the prompt parameters themselves, using hierarchical blocking and attention sparsity to simulate the discourse structure of text, resulting in more coherent and relevant generated text \cite{ref163}.

While soft prompts achieve high performance, they often lack interpretability. This has led to a resurgent interest in automatic methods for discovering discrete ``hard'' prompts, which can be interpreted, shared, and directly used with any black-box API---a crucial capability for the self-adaptive optimization techniques discussed in the following subsection. A Bayesian optimization framework for hard prompt tuning treats the search over sequences of tokens as a combinatorial optimization problem executed in a continuous embedding space, enabling sample-efficient discovery of effective prompts without gradient access to the language model \cite{ref164}. Alternatively, gradient-based discrete optimization can automate the creation of hard prompts for both text-to-image and text-to-text applications, bridging the gap between the performance of soft prompts and the interpretability of hard prompts \cite{ref165}. This is closely related to automatic label sequence generation, where the goal is to find the optimal verbalizer for a given task. The AutoSeq method tackles this by searching over sequences of indefinite length on sequence-to-sequence models, using beam search to generate candidate label sequences and contrastive re-ranking to select the best one, often outperforming manually curated verbalizers \cite{ref166}. Finally, when the model is a complete black-box, the BBT and BSL frameworks have shown that effective soft prompts can be learned by performing derivative-free optimization in a low-dimensional, randomly generated subspace, sidestepping the need for backpropagation entirely and achieving competitive performance with gradient-based methods \cite{ref167,ref168}. These diverse strands of research collectively demonstrate that prompt engineering has evolved into a sophisticated discipline, encompassing continuous optimization, hierarchical structuring, and discrete search to unlock the full potential of frozen language models. This progression sets the stage for the self-adaptive and evolutionary strategies in the next subsection, which automate the very process of discovering such optimal prompt configurations.

\subsection{Self-Adaptive and Evolutionary Prompt Optimization}

While the previous section detailed methods for structured, soft-token, and hierarchical prompt engineering, a complementary and burgeoning area of research moves beyond the static design of prompt representations to automate the prompt engineering process itself. These self-adaptive and evolutionary methods treat prompt optimization as a dynamic learning problem, where the prompt is iteratively refined through feedback loops, evolutionary algorithms, or reinforcement learning, enabling the discovery of highly effective, often non-intuitive, prompt formulations without manual intervention.

A prominent direction in this space is self-referential evolution, exemplified by frameworks like Promptbreeder \cite{ref169}. This approach conceptualizes prompt optimization not merely as improving a task-prompt, but as evolving an entire self-improving ecosystem. Promptbreeder maintains a population of task-prompts and, crucially, a set of mutation-prompts that govern how task-prompts are modified. An LLM acts as the evolutionary engine, applying the mutation-prompts to generate new variants of task-prompts. These are then evaluated for fitness on a training set. In a self-referential twist, the mutation-prompts themselves are subject to evolution, allowing the system to get better at generating better prompts over time. This recursive self-improvement has been shown to outperform established hand-crafted strategies like Chain-of-Thought on reasoning benchmarks.

Complementing evolutionary strategies, ensemble learning methods with feedback loops offer another pathway for automated prompt refinement. The PREFER framework \cite{ref170} addresses the instability of LLMs by automatically constructing an ensemble of prompts. It operates via a feedback-reflect-refine cycle: a feedback mechanism identifies the inadequacies of existing ``weak learner'' prompts, particularly on hard examples. Based on this reflection, the LLM synthesizes new prompts to address these specific weaknesses, iteratively refining the ensemble. A novel prompt bagging method, incorporating forward and backward thinking, enhances the stability of evaluating a prompt's effect, which is critical for both the feedback signal and the final weight calculation in the boosting-style ensemble. This approach achieved state-of-the-art results in multiple tasks, demonstrating the power of systematic, feedback-driven prompt synthesis.

Reinforcement learning (RL) provides a principled framework for discrete prompt optimization when the underlying ``reward'' is a downstream task metric. By framing the selection or generation of prompt tokens as a sequential decision-making problem, an RL policy can be trained to maximize task performance. The DP\(_2\)O method \cite{ref171} utilizes a policy gradient approach where a relatively small policy network is trained to match prompts to inputs optimally. It first generates an initial set of candidate prompts via a multi-round dialogue with a strong LLM like GPT-4, screens them for quality, and then uses RL to learn the best prompt-input pairings, outperforming prior methods in few-shot settings with high parameter efficiency. This paradigm can also be extended to multi-step tasks, where a score function evaluating the final outcome guides an RL agent to propose prompt edits. Integrating human-designed feedback rules can further direct the optimization, as seen in the PROMST framework \cite{ref172}, which combines heuristic performance prediction with human-aligned reward signals to efficiently search the prompt space.

The pursuit of optimal prompts has also led to investigations into the counter-intuitive properties of what constitutes an effective prompt. A remarkably powerful and simple baseline is the use of random separators \cite{ref173}. This work demonstrated that randomly sampled tokens from the model's vocabulary, used as separators between the instruction and the input, can perform as well as or better than human-curated and even some LLM-optimized prompts across a wide array of text classification tasks. With an average relative improvement of 12\% over strong human baselines and a greater than 40\% chance that a random separator outperforms a human one, this finding challenges the assumption that an effective prompt must be human-readable or semantically relevant. It suggests that for some models and tasks, the language space is dense with effective ``token-level'' triggers that the model is sensitive to, which may be exploited without semantic coherence.

Further challenging conventional wisdom, studies on the robustness of automatically learned discrete prompts have revealed their brittle nature \cite{ref174}. While they can achieve high performance on in-domain data, these prompts often contain noisy and counter-intuitive lexical constructs and are highly sensitive to perturbations like token shuffling or deletion, and they generalize poorly across different datasets. This underscores a critical trade-off between peak performance and robustness in discrete prompt learning. Such brittleness directly motivates the calibration and debiasing techniques discussed in the next section, as both automatic prompt optimization and calibration ultimately aim to produce more reliable and generalizable in-context learning.

To navigate this trade-off and improve the reliability of prompt selection, the concept of prompt flatness has been introduced as a key robustness metric \cite{ref175}. Drawing inspiration from flatness regularization in statistical learning, where a model is considered robust if its loss remains stable under small parameter perturbations, prompt flatness quantifies the expected utility of a prompt by assessing the stability of the model's output when small, meaningless perturbations (like synonym substitutions or token reordering) are applied to the prompt itself. A flatter prompt is one whose performance does not degrade sharply with such perturbations, indicating greater reliability. Theoretically, this metric connects to existing selection methods, and empirically, incorporating flatness significantly improves both the accuracy and sample efficiency of prompt selection, leading to an average increase of 5\% in accuracy and 10\% in Pearson correlation compared to metrics that only consider raw task performance. This points toward a future where prompt optimization is not just about peak accuracy but about finding robust, generalizable prompts that perform consistently across diverse inputs and conditions.

\subsection{Calibration of Label Biases and Mitigating Demonstration Shortcuts}

A central challenge in harnessing in-context learning (ICL) is the propensity of large language models (LLMs) to rely on spurious priors embedded in demonstrations rather than genuinely learning the input-label relationships. This phenomenon, termed the ``Demonstration Shortcut'' \cite{ref23}, describes a failure mode where the model's predictions are unduly influenced by its pre-trained semantic knowledge of the label words or their distribution, instead of the novel mapping presented in the context. Such brittleness is a direct counterpart to the robustness challenges discussed in the previous section, where automatically optimized prompts were shown to be highly sensitive to perturbations and prone to exploiting non-semantic token-level triggers. This section addresses how label distributions and prior semantic knowledge bias model predictions and surveys calibration techniques designed to rectify these biases, ultimately aligning model behavior with the intended task and laying the groundwork for the position and order sensitivities explored in the following section.

The root of the demonstration shortcut lies in the powerful semantic priors that LLMs acquire during pretraining. When presented with few-shot demonstrations, a model may default to predictions consistent with the pre-existing meaning of label words, ignoring contradictory evidence in the prompt. For instance, if tasked with a sentiment analysis task where the label ``positive'' is flipped to mean a negative review, smaller models tend to persist with the original semantic association, while the capacity to override these priors and adhere to the flipped mapping emerges with model scale \cite{ref176}. This reliance on priors is not merely about word semantics but extends to the label distribution. Models can develop a ``vanilla-label bias'' favoring certain labels independent of the input, a ``context-label bias'' where the relative positions of examples in a prompt skew predictions, and a ``domain-label bias'' where the mere presence of words from a specific task domain triggers a pre-learned label preference, even when the input is non-informative \cite{ref177}. This complex interplay of biases reveals that ICL predictions are often a mixture of learning from the given demonstrations and a fallback on internal, potentially spurious, knowledge, echoing the finding from the previous section that even random token separators can inadvertently activate such pre-existing model sensitivities.

A foundational approach to mitigating these biases is in-context calibration, which aims to adjust the model's output distribution to nullify the influence of unwanted priors. The principle is to estimate the bias directly from the model's behavior and then subtract or correct for it during inference. One straightforward strategy is ``contextual calibration,'' where the model's prediction for a test input is adjusted by its propensity to predict certain labels on a content-free input, such as an empty string or random in-domain words \cite{ref177}. By estimating the model's prior label preference in the absence of meaningful input, this baseline bias can be removed from the predictions on genuine examples, significantly recovering performance on tasks with strong domain-label bias. A more unified perspective frames calibration as controlling for the ``contextual bias from the batched input'' \cite{ref178}. Batch calibration (BC) computes a bias vector from a batch of test inputs, representing the model's average prediction behavior across the batch, and uses it to calibrate individual predictions. This approach is zero-shot and inference-only, elegantly unifying prior methods and demonstrating state-of-the-art performance across various natural language understanding and image classification tasks by controlling for biases like formatting and verbalizer choice.

Moving beyond simple bias subtraction, calibration by marginalization offers a principled probabilistic framework to disentangle the model's understanding of input-label relationships from label marginals. A key theoretical and empirical finding identifies that a significant part of the performance instability in ICL is due to ``label shift,'' where the LLM maintains a well-calibrated class-conditional distribution \(p(x|y)\) but distorts the label marginal \(p(y)\) \cite{ref179}. Generative calibration exploits this insight by using the LLM itself to estimate the correct label marginal via Monte-Carlo sampling. By generating sample texts for each label and measuring the model's likelihood of assigning them to their respective labels, an unbiased estimate of the label prior is obtained. This prior then replaces the distorted marginal in a Bayes-optimal prediction rule, leading to dramatically more consistent and accurate performance across different prompt configurations and demonstration orders---directly prefiguring the order sensitivity issues that the next section will address. Similarly, prototypical calibration constructs a Gaussian mixture model over the representations of each label class from the demonstrations and then calibrates the test prediction by its probability under these prototypical clusters, creating a robust decision boundary that is less sensitive to example permutations and class imbalance \cite{ref131}.

A more diagnosis-driven approach to mitigating shortcuts involves the explicit design of demonstrations that prevent the model from exploiting trivial correlations. The concept of ``Comparable Demonstrations'' (CDs) \cite{ref8} directly tackles the shortcut by constructing minimally contrasting example pairs. By editing the input text to flip the label, CDs force the model to attend to the features that genuinely distinguish the classes, promoting a causal rather than correlational understanding of the task. This method has proven effective, especially in out-of-distribution scenarios where spurious correlations are broken. This strategy is conceptually aligned with the broader goal of context-consistent ranking, where the objective is to calibrate prediction confidence not just based on a single demonstration set but through a process of mutual consistency checks among different ways of presenting the same information, effectively filtering out shortcuts that do not hold across contexts---a principle that resonates with the permutation-based aggregation methods detailed in the following section.

Ultimately, the effectiveness of all calibration techniques hinges on a precise diagnosis of the bias. Research has shown that models do learn label relationships from context, but this learning is often incomplete and intermixed with pre-existing knowledge \cite{ref51}. The challenge is that ICL ``does not consider all in-context information equally,'' making it susceptible to shortcuts. An effective calibration pipeline therefore requires a holistic view, recognizing that the choice of demonstrations, their labels, and the model's own output statistics are all potential sources of bias. A unified strategy involves selecting demonstrations based on information gain while calibrating the selection criterion itself against template bias \cite{ref130}, and then applying prediction-time calibration methods like batch or generative calibration \cite{ref178,ref179}. By systematically addressing the semantic and distributional biases present in both the prompt and the model's prior, these calibration frameworks are essential for transforming ICL from a brittle, prior-driven process into a robust mechanism for genuine task learning from context. This pursuit of debiased, reliable predictions directly motivates the subsequent section's focus on order and position sensitivities, as both lines of inquiry converge on the shared goal of making in-context learning invariant to arbitrary and spurious features of the prompt format.

\subsection{Mitigating Sensitivity to Demonstration Order and Position}

A critical and well-documented vulnerability of in-context learning (ICL) is its sensitivity to the order and position of demonstrations within the prompt. While the previous section detailed how semantic priors and label distributions can bias model predictions away from the intended task, this section turns to an equally pervasive source of instability: the arbitrary sequence in which in-context examples are arranged. This sensitivity manifests as significant performance variance when the sequence of examples is permuted, undermining the reliability and fairness of ICL-based systems. Research has identified a key architectural source of this brittleness: the auto-regressive attention mask in causal language models (CausalLMs). Unlike prefix language models, which allow bidirectional attention over the context, CausalLMs restrict each token's access to only its preceding tokens. This constraint creates varying receptive fields for demonstrations at different positions, leading to representation disparities. As shown in \cite{ref180}, CausalLMs exhibit substantially greater sensitivity to demonstration order than their prefix counterparts, directly attributing this to the causal attention mechanism. The effect is not limited to language models; it extends to reasoning tasks, where the logical ordering of premises drastically impacts accuracy. For instance, permuting the premise order in deductive reasoning tasks can cause a performance drop of over 30\% \cite{ref181}.

The positional bias is pervasive across various task formats. In multiple-choice question answering, reordering answer options can lead to a performance gap ranging from 13\% to 75\% \cite{ref182}. This is particularly pronounced when the model is uncertain among the top-2 or -3 choices, with specific option placements (e.g., locating the top two choices as the first and last options) amplifying the bias. Similarly, in listwise ranking tasks, LLMs exhibit a strong positional preference that skews ranking outputs. The ``lost-in-the-middle'' phenomenon further illustrates this problem in long-context scenarios, where information positioned centrally within a lengthy prompt is processed less effectively than information at the beginning or end. These findings point to an inherent positional preference within pre-trained models, a challenge that simple prompt-based solutions often fail to overcome \cite{ref183}. An analysis of instruction fine-tuning reveals that the conventional placement of task instructions before the input can lead to ``instruction forgetting'' for long sequences. Shifting the instruction's position to follow the input has been proposed as a straightforward yet effective way to re-weight the model's learning focus and improve instruction following in sequence generation tasks \cite{ref184}.

To combat these pervasive order and position sensitivities, a variety of methodological interventions have been proposed, spanning inference-time strategies, fine-tuning paradigms, and causal intervention frameworks. At the inference level, permutation-based aggregation seeks to marginalize out order biases by evaluating the model on multiple shuffled versions of the same prompt and then aggregating the results. A prominent example is permutation self-consistency for listwise ranking, which repeatedly shuffles the list in the prompt, passes each variant through the LLM, and computes a central ranking that is order-independent. This method is theoretically robust, proving convergence to the true ranking under random perturbations, and has demonstrated substantial empirical gains of up to 7-18\% on ranking benchmarks \cite{ref185}. This approach essentially trades computational cost for predictive stability and reliability, echoing the principle seen in calibration methods from the previous section where multiple inference passes are leveraged to cancel out spurious signals.

A more fundamental line of attack involves fine-tuning models to be inherently less sensitive to order. One such method is an unsupervised fine-tuning approach termed Information-Augmented and Consistency-Enhanced (InfoAC) learning. This technique employs contrastive objectives to align the representations of identical in-context examples across different positions within the prompt, while a consistency loss encourages the model to produce similar output distributions for different permutations of the same input. The result is a model with robust generalizability, maintaining performance even when the number of demonstrations or their source pool changes from the training phase \cite{ref180}. Another parameter-efficient fine-tuning (PEFT) strategy, Position-Aware Parameter Efficient Fine-Tuning (PAPEFT), combines a data augmentation technique with a specialized adapter module. This approach is designed to foster a more uniform attention distribution across the input context, directly counteracting the model's inherent positional preferences and improving its effectiveness on tasks requiring the retrieval and use of externally sourced knowledge from long contexts \cite{ref183}.

Beyond direct architectural and representational alignment, other works frame the problem through a causality lens, drawing a conceptual link to the earlier discussion of spurious correlations and the demonstration shortcut. For example, the ``reversal curse,'' a phenomenon where an LLM trained on ``A precedes B'' contexts fails to infer the reverse relationship, is identified as a product of the unidirectional nature of next-token prediction. A proposed solution, Bidirectional Causal Language Modeling Optimization (BICO), modifies the causal attention mask to function bidirectionally during fine-tuning, combined with a mask denoising objective. This approach has been shown to dramatically improve a model's resistance to this specific type of order-induced failure, lifting accuracy from near zero to around 70\% on targeted assessments \cite{ref186}. Furthermore, the framework of causal prompting uses structural causal models to diagnose biases in prompting and proposes a front-door adjustment method. By employing the model's own chain-of-thought as a mediator variable, this technique aims to calculate and mitigate the direct causal effect of spurious, order-based correlations without accessing model parameters or logits \cite{ref187}. Collectively, these methods move beyond treating order sensitivity as a nuisance to be averaged out and instead target its root causes in model architecture, training objectives, and inference-time decision processes, striving for in-context learners that are both accurate and predictively consistent. The goal is to decouple the model's task understanding from the arbitrary sequence of its inputs, a foundational requirement for trustworthy deployment that complements the calibration efforts aimed at neutralizing semantic and distributional biases.

\section{Training, Distillation, and Efficiency for Enhanced In-Context Capabilities}

\subsection{Meta-Training Paradigms for In-Context Learning}

Building upon the foundational objective of transferring in-context learning (ICL) capabilities to smaller models, as discussed in the distillation literature, meta-training for ICL represents an even more fundamental shift in how we prepare models: moving from optimizing static task performance to explicitly cultivating the skill of learning from context itself. Unlike standard fine-tuning or instruction tuning, which often target a model's zero-shot or single-task adherence, meta-training paradigms are designed to optimize the model's few-shot ICL ability across a wide distribution of tasks. The core principle is that by training a model on a diverse collection of tasks where the solution must be inferred from provided demonstrations, the model internalizes the general algorithm of ``learning-to-learn in context,'' leading to more robust and generalizable adaptation at test time. While distillation focuses on transferring this capability between models, meta-training constitutes the upstream process of instilling it in the first place.

A foundational framework in this space is MetaICL (Meta-training for In-Context Learning) \cite{ref188}, which directly addresses the objective of improving ICL. The methodology involves fine-tuning a pre-trained language model on a vast and diverse set of training tasks, each formatted as an ICL episode with a few demonstrations followed by a test input. Crucially, no parameter updates occur at test time; the meta-training simply conditions the model to perform this specific form of learning more effectively. The power of this approach lies in its task diversity. Studies show that MetaICL significantly outperforms standard ICL (without meta-training) and multi-task zero-shot learning, with the gains being particularly pronounced when the target tasks exhibit domain shifts from the meta-training distribution. This underscores a critical finding that resonates with the importance of input-level alignment in distillation: meta-training with a highly diverse set of tasks is key to fostering generalizable ICL capabilities, allowing the model to approach or even surpass the performance of models fully fine-tuned on the target task, and to outperform substantially larger models with nearly 8x the parameters \cite{ref188}. This explicitly optimizes the ``promptability'' of the model, making it a more effective and data-efficient in-context learner, much like how bidirectional alignment sharpens a student model's sensitivity to informative demonstrations.

The effectiveness of meta-training is deeply connected to the properties of the training data distribution, with task diversity acting as a critical catalyst. Empirical investigations into the emergence of ICL have demonstrated a ``task diversity threshold'' \cite{ref48}. When pre-trained on data below this threshold, transformer models fail to solve fundamentally new, unseen tasks in-context, instead behaving as Bayesian estimators with a prior overfitted to the narrow pre-training distribution. However, once task diversity surpasses this threshold, a qualitative shift occurs: the model's behavior aligns with ridge regression, adopting a Gaussian prior over all tasks, including those never seen during pre-training. This transition reveals that sufficient diversity allows the model to deviate from a memorized Bayesian prior and instead learn a more generalizable algorithmic strategy. This insight is reinforced by the study of supportive pre-training data, which found that the data subsets most responsible for ICL are not necessarily topically similar to downstream tasks but rather consist of challenging examples where leveraging long-range context is unusually difficult \cite{ref41}. A continued pretraining on this small, difficult subset can drastically improve ICL performance, highlighting that meta-training's success is linked to mastering the integration of complex, long-range contextual information---a capability that bidirectional alignment and similar distillation techniques aim to replicate in compact models.

Complementing these findings, curriculum learning strategies within meta-training offer a structured pathway to mastering this complexity. In the multi-task learning (MTL) setting for ICL, the order and schedule of task presentation can significantly affect learning efficiency and stability. A ``mixed curriculum,'' which progressively introduces more difficult tasks while interleaving previously learned ones, has been shown to allow ICL models to achieve higher data efficiency and more stable convergence than training on a static, random mixture \cite{ref189}. This mirrors a developmental learning process, where the sequential mastery of nested operations prevents the model from getting stuck in a learning plateau. This perspective is also supported by theories on the emergence of induction heads, where a curriculum-like progression through subcircuits is necessary for their abrupt formation. Such findings suggest that meta-training is not just about the data's composition but also its dynamics, with a carefully orchestrated curriculum enabling the model to build increasingly complex reasoning capabilities in a compositional manner. The concept of ``concept-awareness'' further refines this data-centric view. By specifically constructing training scenarios where analogical reasoning on latent concepts is beneficial, methods like Concept-aware Training (CoAT) improve a model's intrinsic ability to utilize new concepts from demonstrations \cite{ref68,ref190}. Remarkably, models trained with such strategies on just a single task can rival the ICL performance of models trained on over 1600 distinct tasks, proving that the structure and intent behind the data are as crucial as its volume.

These data-centric and curriculum-based approaches converge on a unified theme that directly complements the distillation paradigm: meta-training for ICL functions as a form of meta-regularization, shaping the model's inductive biases to favor the extraction of algorithmic patterns over statistical memorization of the pre-training distribution. Strong theoretical and empirical bridges exist between MTL and gradient-based meta-learning (GBML), with proofs that for sufficiently deep over-parameterized neural networks, the learned predictive functions of MTL and GBML are highly similar \cite{ref191}. This connection implies that large-scale multi-task meta-training implicitly performs a kind of first-order meta-optimization, an order of magnitude faster than explicit second-order GBML methods, to instill a robust ``in-context gradient'' mechanism---the very mechanism that distillation methods like bidirectional alignment seek to transfer to student models by aligning both input preferences and output distributions. This perspective is further supported by research demonstrating that meta-learning success relies on navigating the interplay between task diversity, entropy, and difficulty, rather than simply maximizing diversity, which can lead to under- or over-fitting \cite{ref192}. An adaptive task sampling strategy that balances these factors can significantly outperform uniform sampling. In essence, the modern meta-training paradigm for ICL synthesizes massive task diversity, principled curricula, and concept-aware data construction to bridge statistical pre-training and the desired meta-learning objective. This synergy effectively distills the general ability to learn, rather than a set of memorized task-specific solutions, into the model's forward pass, creating an in-context learner that is not only more capable but also more robust, efficient, and adaptable to entirely new challenges. It is precisely this generalized learning capability that subsequent distillation efforts aim to compress and transfer to smaller, more deployment-friendly models, forming a coherent pipeline from the genesis of ICL abilities in large models to their efficient deployment in resource-constrained settings.

\subsection{In-Context Learning Distillation and Knowledge Transfer}

While the meta-training paradigm is essential for instilling in-context learning (ICL) abilities from the ground up, its application is typically limited to large, computationally expensive models. This limitation has motivated a complementary and synergistic line of research: transferring the already-emergent ICL capabilities of large language models (LLMs) to smaller, more efficient counterparts through distillation and knowledge transfer. This subsection examines the primary methodologies designed to imbue compact models with competitive few-shot ICL performance, drawing on paradigms such as multitask instruction tuning, bidirectional alignment of input preferences and output distributions, and other techniques that bridge the gap between large teacher models and their smaller student models. This forms a natural downstream pipeline from meta-training, where the ``learning-to-learn'' skill that is cultivated in large models is subsequently compressed for efficient deployment.

A foundational approach to transferring ICL abilities involves framing the process as a meta-training or multitask distillation problem, directly echoing the core principles discussed previously but now applied in a teacher-student context. The work on \textbf{in-context learning distillation} \cite{ref193} explicitly proposes combining in-context learning objectives with standard language modeling objectives to distill both the ability to read and interpret in-context examples and the underlying task knowledge. This study introduces two key paradigms: Meta In-context Tuning (Meta-ICT) and Multitask In-context Tuning (Multitask-ICT). Meta-ICT optimizes the student model to learn from few-shot examples in a meta-learning loop, while Multitask-ICT jointly trains on a mixture of diverse tasks presented in an ICL format. The findings reveal that these two objective types are complementary under the Multitask-ICT paradigm, with the combination of in-context learning and language modeling losses yielding the best performance on benchmarks like LAMA and CrossFit. This underscores the necessity of not only teaching the student model \emph{what} the task is but also \emph{how} to learn from the contextual demonstrations themselves, mirroring the distinction between task memorization and meta-learning that underpins effective meta-training.

While output-level alignment through distillation effectively encodes task-specific answers, a critical limitation of early methods was the neglect of the input side of the ICL process---a dimension that the earlier discussion of task diversity and concept-aware training identified as fundamental. The sensitivity of ICL to the choice of demonstrations is a well-documented phenomenon, indicating that a model's preference for and processing of demonstration inputs is a core component of its ICL competence. To address this gap, the \textbf{Bidirectional Alignment (BiAlign)} framework \cite{ref122} was proposed to fully leverage the ICL example preferences of larger teacher LLMs. BiAlign goes beyond standard token-level output distribution alignment by introducing a novel ranking loss that aligns the input preferences between the smaller student and larger teacher models. This forces the student model to not only produce correct outputs but also to internally recognize and prioritize informative demonstrations in a manner that mirrors the teacher. Extensive experiments across language understanding, reasoning, and coding tasks demonstrate that this bidirectional approach consistently outperforms baselines that focus solely on output distillation. This method highlights that a crucial aspect of ICL ability is a kind of ``meta-evaluation'' skill---the capacity to assess the utility of provided examples---which can be successfully transferred, much like the way concept-aware training cultivates sensitivity to latent task structures.

The relationship between ICL and parameter-efficient fine-tuning (PEFT) further complicates the distillation landscape. While distillation aims to create a single model capable of ICL at inference time, an alternative is to forgo ICL entirely and use the teacher's demonstrations to directly fine-tune a smaller model using PEFT methods. Research rigorously comparing these paradigms \cite{ref30} has shown that few-shot PEFT with methods like (IA)\(^3\) and the T-Few recipe can not only achieve higher accuracy than ICL but do so with dramatically lower computational costs. This creates a productive tension: distillation seeks to maintain the flexibility of on-the-fly ICL, while direct PEFT offers a more efficient static alternative. However, insights from this comparison inform distillation by demonstrating that the implicit knowledge in a set of demonstrations can be effectively compressed into a small number of model parameters---just as gradient-based meta-learning implicitly encodes task-solving procedures within the model's weights. The choice between distilling a general ICL ability and distilling a specific set of demonstrations into weights becomes a matter of deployment requirements, with the former retaining superior adaptability to distribution shifts, a hallmark of well-meta-trained models.

The line between ICL and instruction tuning (IT) is also increasingly blurred in the context of knowledge transfer. The observation that ICL functions as a form of ``implicit instruction tuning'' \cite{ref9} provides a theoretical underpinning for many distillation strategies. If providing demonstrations in the context during inference alters a model's hidden states in a way that resembles the effect of instructionally tuning it, then a distillation process that aligns the hidden states, output logits, and input preferences of a student model with a teacher performing ICL is essentially learning this implicit gradient update in a compressed, transferable form. Data-centric approaches like LESS \cite{ref194} further enable this transfer by identifying the most influential subsets of instruction-tuning data. By using a small set of exemplar demonstrations to query a gradient datastore, LESS can select the most relevant training data for a specific downstream capability, effectively acting as a targeted distillation process that produces a highly specialized student model without needing to distill the full breadth of the teacher's knowledge. This data-aware selection process parallels the supportive pretraining data identification strategies that proved essential for instilling ICL during meta-training.

In the multimodal domain, the distillation of ICL capabilities must also navigate the complexities of cross-modal interactions. The challenge is not only in compressing model size but in aligning complex input modalities. Techniques involving lightweight tuning modules, such as Conditional Mixture of LoRA \cite{ref195} or prompt tuning-based adapters \cite{ref196}, provide a pathway for transferring multimodal ICL. These methods can be interpreted as distilling the ability to dynamically adapt to in-context multimodal inputs into a small set of parameters that modulate a frozen backbone. Instead of the student model learning ICL entirely from scratch, it learns a highly efficient set of coefficients that enable it to interpret and act upon multimodal prompts, effectively distilling the teacher's flexible in-context adaptation process into a parameter-efficient mechanism. This principle is also evident in the transfer of ICL capabilities across modalities and model scales, where a foundational principle is the alignment of input spaces and the preservation of preferences for salient cross-modal demonstrations, ensuring that the smaller model's ``attention'' to different parts of a multimodal prompt mirrors the teacher's more sophisticated processing.

Overall, the distillation of in-context learning is evolving from simple output mimicry to a holistic transfer of the ICL process, encompassing input preference alignment, meta-training curricula, and parameter-efficient adaptation. The complementary nature of these techniques---aligning output distributions, input rankings, and learned adapters---forms a robust toolkit for creating smaller, deployment-friendly models that retain the powerful emergent ability to learn from context. This distillation pipeline directly complements the data-centric strategies explored next, where instead of compressing an existing large model's ICL abilities, the focus shifts to engineering pretraining data that encourages the emergence of these capabilities natively, even in smaller models trained from scratch. Together, model compression and data-centric cultivation position efficient ICL as an attainable goal across the entire spectrum of model scales and resource constraints.

\subsection{Concept-Aware and Data-Centric Training Strategies}

While much of the progress in in-context learning (ICL) has been driven by scaling model and task sizes, a growing body of evidence has shifted the understanding toward the primacy of pretraining data properties. Data-centric strategies focus on curating and constructing pretraining corpora that inherently encourage models to learn analogical reasoning and context utilization, forming a crucial complement to the model-centric distillation and bridging techniques discussed in the surrounding sections. This section analyzes training approaches that enhance ICL by manipulating data composition and pretraining curricula, demonstrating how structural properties like parallel structures, burstiness, and concept-awareness drive the emergence of robust in-context learners---effectively distilling the principles of ICL directly into the pretraining objective.

A foundational insight in data-centric ICL is that specific patterns in the pretraining data are prerequisites for the capability to emerge. The work by \cite{ref67} provides strong causal evidence for this, demonstrating that language models' ICL ability depends critically on the presence of parallel structures in the pretraining data. These structures are defined as pairs of phrases following similar templates within the same context window. The study found that selectively removing these parallel structures reduced ICL accuracy by 51\%, a drop far greater than the 2\% reduction from random data ablation. This effect persisted even when common patterns like n-gram repetitions and long-range dependencies were excluded, underscoring the diversity and generality of parallel structures. These structures essentially provide a natural, self-supervised analogue to the demonstration-based prompts used during ICL inference, teaching the model to map between analogous patterns by conditioning on shared context. Such findings pivot the focus from architectural tweaks to a fundamental question: what must the data look like for ICL to arise?

Moving beyond identifying patterns in natural data, another line of research actively constructs training data to instill ICL capabilities even in smaller models or from limited data regimes. The Concept-aware Training (CoAT) framework, introduced in \cite{ref68} and elaborated upon in \cite{ref190}, proposes a method for creating training scenarios that make it beneficial, and indeed necessary, for a language model to capture analogical reasoning concepts from demonstrations. Instead of merely training on more tasks, CoAT explicitly constructs examples where the model must derive a latent concept from provided in-context demonstrations to successfully predict the output. This approach consistently improves reasoning ability, and remarkably, in-context learners trained with CoAT on just a couple of datasets for a single task can perform comparably to much larger models trained on over 1600+ tasks. This demonstrates that the quality and pedagogical structure of training data can be a more powerful lever than sheer data quantity or model scale, validating the hypothesis that ICL is a skill that can be deliberately taught.

The effectiveness of concept-aware training is closely tied to findings regarding the specific distributional properties of supportive pretraining data. Research by \cite{ref41} adopted an iterative, gradient-based approach to identify small subsets of pretraining data that are particularly supportive of ICL. Continued pretraining on this small, identified subset improved ICL ability by up to 18\%. A contrastive analysis of this supportive subset against random data revealed three key characteristics. First, domain relevance to downstream tasks was not the primary driver; the supportive data was not necessarily from a similar domain. Instead, the data was characterized by a higher mass of rarely occurring, long-tail tokens, suggesting that exposure to lexically diverse and sparse contexts is crucial. Second and most importantly, the supportive data consisted of challenging examples where the information gain from long-range context was below average. This indicates that ICL is fostered by data that forces the model to learn to incorporate difficult long-range context to make predictions, directly training the very skill required for using multiple demonstrations spread across a prompt. This provides a precise, instance-level view of the pretraining curriculum that naturally selects for ICL.

The mechanistic basis for how such data properties drive learning dynamics is further clarified by \cite{ref55}. This study shows that distributional properties inherent in language, such as burstiness, large dictionaries, and skewed rank-frequency distributions, control the trade-off and competition between ICL and in-weights learning (IWL). In a minimal attention-only network, ICL was driven by the abrupt emergence of an induction head, a circuit that competes with IWL. The study constructed a phenomenological model showing that this abrupt emergence is a consequence of a sequential learning of nested nonlinearities, forming an intrinsic curriculum enabled by the data distribution. This work synthesizes the data-centric perspective with mechanistic interpretability: concept-aware or bursty data provides the necessary curriculum that guides the model through a sequence of learning phases, ultimately leading to the phase transition that creates an induction head and thereby unlocks ICL. This mechanistic understanding directly complements the theoretical insights presented in the following subsection, which formalize this implicit curriculum through the lens of gradient-based optimization.

From a theoretical standpoint, these findings converge into a data-centric learning theory. The work by \cite{ref49} provides a statistical task complexity bound, showing that effective pretraining for ICL, even for simple linear regression, requires only a small number of independent tasks, provided they are structured correctly. This aligns with the empirical finding that a small, well-curated concept-aware dataset can be as effective as massive multi-task learning. Furthermore, \cite{ref6} offers a unifying conceptual framework by reinterpreting ICL through the lens of skill learning and skill recognition. In this view, concept-aware training enables ``skill learning,'' where the model genuinely learns new data generation functions from in-context data, as opposed to merely recognizing and applying a pre-learned function (``skill recognition''). Data that emphasizes burstiness and long-range coherence forces true skill learning, leading to more generalizable and robust ICL.

Finally, the efficacy of data curation as a standalone stabilizer is demonstrated by \cite{ref34}. This work shows that carefully selecting a stable training subset, using methods like CondAcc and Datamodels, dramatically reduces the high variance that ICL typically suffers from when random examples are used. Surprisingly, the examples in these stable subsets were not especially diverse in content or low in perplexity, challenging conventional wisdom and suggesting that data curation methods targeting stability capture a different, more fundamental signal. This reinforces the overarching theme: that the pathway to powerful, reliable ICL lies in understanding and engineering the data that facilitates its emergence. By constructing training scenarios that emphasize analogical reasoning \cite{ref68}, leveraging supportive properties like long-tail tokens and challenging long-range dependencies \cite{ref41}, and aligning with distributional characteristics like burstiness that guide the emergence of ICL circuits \cite{ref55}, researchers are developing a principled framework for data-centric in-context learning that stands as a compelling alternative to pure scale, and which ultimately feeds into the theoretical unification of implicit learning dynamics and explicit optimization that follows.

\subsection{Bridging Implicit In-Context Gradients with Explicit Parameter Updates}

While the preceding section established that careful data engineering can distill the principles of in-context learning directly into a model's parameters, a parallel theoretical thread reveals a deeper unity: in-context learning (ICL) can be mathematically understood as performing implicit gradient-based optimization. This perspective not only demystifies why data-centric strategies like concept-aware training work---by structuring the pretraining curriculum to guide meta-optimization---but also paves the way for novel algorithms that bridge the implicit meta-optimization occurring within a transformer's forward pass and the explicit, gradient-based parameter updates used in fine-tuning. By exploiting this duality, researchers have developed momentum-based attention mechanisms, reweighted ICL algorithms that approximate unbiased estimators, and hybrid frameworks like instruction prompt tuning that synergistically combine learned parameters with in-context demonstrations, directly addressing the same goal of stabilizing and enhancing ICL that motivates the curriculum strategies in the following section.

The foundational insight driving this bridge is the dual form between linear attention and gradient descent. A landmark study \cite{ref25} established that transformer attention can be interpreted as performing implicit fine-tuning. The authors theoretically demonstrated that a linear attention layer has a dual form identical to a gradient descent step on a linear regression loss. In this framework, the in-context demonstrations (input-label pairs) are used to compute ``meta-gradients,'' which are then applied to the model's initial parameters to construct an ICL model tailored to the task. This understanding was empirically validated by showing that the behavior of ICL on real tasks resembles that of explicit fine-tuning from multiple perspectives. Building on this dual form, the work introduced a momentum-based attention mechanism, directly inspired by the momentum in gradient descent optimizers. This simple yet effective architectural tweak improved performance over standard attention, providing strong evidence that optimizing the ICL process through the lens of gradient descent can yield tangible benefits, effectively designing an architecture that accelerates the meta-learning dynamics that data-centric methods aim to induce.

This equivalence, however, is not without its nuances. A subsequent investigation \cite{ref10} revisited the ICL-GD correspondence on realistic NLP tasks and models, uncovering gaps in prior evaluations. They found that even untrained models could achieve comparable ICL-GD similarity scores, prompting a deeper analysis of ``Layer Causality''---the flow of information through the model's layers. Their work revealed a major discrepancy: while ICL processes information in a causally layered manner, explicit GD applied to the full model does not respect this sequential order. By proposing a simple GD-based optimization procedure that respects layer causality, they significantly improved similarity scores, demonstrating that a faithful bridge must account for the architectural constraints of the transformer. This refined understanding indicates that bridging ICL and explicit updates is not merely about global loss optimization but requires alignment with the forward computation graph---a mechanistic constraint that mirrors how supportive pretraining data was shown to force models to learn long-range dependencies in a specific sequential fashion.

The exploitation of the ICL-GD duality has led to practical inference algorithms that overcome the brittleness of standard ICL, a challenge that data curation and curriculum design both seek to mitigate. The \cite{ref197} framework treats ICL as a meta-optimization process, explaining the known sensitivity of LLMs to the order of demonstrations. By deconstructing an N-shot prompt into N separate 1-shot forward computations and aggregating their resulting meta-gradients, Batch-ICL produces a final prediction that is agnostic to example order. This batched meta-optimization approach, which includes a variant with multiple ``epochs'' of meta-gradient aggregation, consistently outperforms standard ICL permutations and approaches the performance of the optimal order, all while reducing computational overhead. This method effectively bridges the gap by performing a more principled form of implicit optimization that mimics the effect of training on multiple examples simultaneously, directly connecting to the subsequent section's focus on overcoming plateaus through structured learning sequences.

Another direct line of attack for bridging the paradigms involves fine-tuning language models to better approximate the implicit learning process. The \cite{ref198} work introduced RICL (Reweighted In-Context Learning), an algorithm that directly addresses the bias and imbalance in input prompts that can degrade ICL performance. RICL fine-tunes the model using an unbiased validation set to learn an optimal weight for each input-output example, thereby approximating an unbiased in-context learner. This effectively translates the explicit knowledge of an unbiased data distribution into the implicit weighting mechanism of ICL. A low-cost linear approximation of this reweighting, LARICL, further bridges the gap by providing a computationally efficient path to debiased ICL, demonstrating a substantial improvement over standard prompting and classic fine-tuning. This approach mirrors the data-centric insight that carefully selected, stable training subsets can dramatically reduce variance, but accomplishes this through a parametric intervention rather than a data curation one.

The convergence between implicit ICL and explicit instruction tuning (IT) is further illuminated at the representation level. An empirical study \cite{ref9} found that ICL acts as implicit IT, changing an LLM's hidden states as if the demonstrations were used to instructionally tune the model. This convergence underscores the potential for hybrid approaches that explicitly leverage the strengths of both. A prime example is Instruction Prompt Tuning (IPT), which concatenates a natural language demonstration with learned, tunable prompt embeddings \cite{ref31}. This method directly combines an explicit parameter update (the learned soft prompt) with an implicit one (the in-context demonstration). The study revealed that IPT not only consistently reduces the high variance associated with prompt tuning alone but also enables positive transfer of learned prompts between tasks when paired with appropriate demonstrations. This hybrid framework represents a practical fusion of meta-optimization in context with explicit parameter-efficient fine-tuning, validating the bidirectional bridge between the two learning paradigms and setting the stage for curriculum strategies that explicitly guide this fusion process through careful sequencing of tasks.

In summary, the exploration of the ICL-GD duality has moved from a theoretical observation to a practical toolkit for enhancing model adaptation. The development of momentum-based attention, layer-causal optimization, meta-gradient aggregation, and reweighted fine-tuning all represent distinct yet converging pathways for bridging implicit in-context gradients with explicit parameter updates. These methods collectively demonstrate that treating ICL as a form of implicit optimization not only explains its mechanisms but also empowers the design of more robust, efficient, and unbiased adaptation algorithms that stand at the interface of prompting and fine-tuning. This bridging perspective provides the theoretical and algorithmic foundation upon which the subsequent section's curriculum learning strategies build---by viewing curriculum design as a way to explicitly orchestrate the implicit learning dynamics, guiding models through the sequential learning of nested operations needed to reliably form the ICL circuits that these hybrid algorithms exploit.

\subsection{Training Curricula and Overcoming Learning Plateaus}

While the previous section explored methods for explicitly bridging implicit in-context gradients with fine-tuning, a parallel and equally critical challenge lies in cultivating the capacity for in-context learning (ICL) itself during model training. The training dynamics of ICL often exhibit a distinctive non-monotonic trajectory, characterized by abrupt phase transitions and developmental plateaus. A central challenge is that these abilities do not emerge smoothly with scale or training time; instead, they can appear suddenly after a period of stagnation, or even decay and become transient, giving way to in-weights learning (IWL) \cite{ref55}. This stands in stark contrast to the efficient, batched, and compressed inference techniques discussed in the following section, which assume a model already possesses robust ICL capabilities. Curriculum learning (CL), a training strategy that presents examples in a meaningful order---typically from easy to hard---has emerged as a powerful tool to navigate these training dynamics, improving data efficiency, accelerating convergence, and steering models toward more generalizable in-context solutions \cite{ref199,ref200}.

The mechanistic basis for the abrupt emergence of ICL is intimately linked to the formation of specialized circuits, such as induction heads, which perform a ``match-and-copy'' operation to bridge past context with the current query. Research analyzing a minimal attention-only network revealed that this sharp transition is not a continuous process but arises from the sequential learning of three nested logits, enabled by an intrinsic curriculum of operations \cite{ref55}. This phenomenological model suggests that the network must learn a specific chain of multi-layer operations, and the abruptness of the ICL phase change is a result of nested nonlinearities being mastered in a necessary order. A plateau in learning occurs when the model is stuck between these stages, unable to fully assemble the induction head circuit until the prerequisite sub-operations are solidified. Explicit curriculum strategies that guide the model through this sequence can therefore directly overcome these developmental bottlenecks, ensuring that the bridge from parameter updates to implicit gradients, as characterized in the previous section, can reliably form.

Empirical studies on multi-task training for in-context capabilities have corroborated the efficacy of curriculum design, particularly when utilizing function classes. A key finding is that ICL models can effectively learn difficult tasks by training on progressively harder tasks while mixing in prior tasks, a strategy termed ``mixed curriculum'' \cite{ref189}. This approach prevents catastrophic forgetting of earlier learned functions while the model is challenged with increased complexity, leading to higher data efficiency and more stable convergence. The mixed curriculum directly addresses the learning plateau by ensuring that the foundational representations required for simpler tasks are not discarded as the model adapts to harder ones, thus maintaining the nested set of abilities necessary for the final induction head assembly. This structured progression toward robust ICL is precisely what enables the scaling to many-shot regimes and the application of efficient, order-agnostic processing methods discussed subsequently.

The link between human and machine curriculum effects provides a deeper theoretical grounding for these strategies. Strikingly, the classic trade-off between blocking and interleaving observed in human learning has been shown to spontaneously emerge with ICL in both meta-learned neural networks and large language models. In tasks governed by succinct, rule-like structures, a blocked curriculum (where related examples are grouped) is more effective for in-context learning, whereas interleaving is superior for in-weights learning on tasks lacking such structure \cite{ref75}. This finding implies that the ``inner-loop'' algorithm of ICL operates under similar principles to human cognition. When designing curricula to overcome plateaus in ICL, a blocked progression that emphasizes the rule-like regularities of the target task may be crucial for forcing the model to abstract the mapping function rather than memorize individual instances, thereby catalyzing the formation of induction heads and, consequently, enabling the efficient many-shot inference techniques explored in the following section.

The broader machine learning literature on curriculum learning offers a wealth of strategies that have been successfully adapted to address the learning plateaus and instabilities observed in ICL training. The core challenge of designing a curriculum---defining a ``difficulty measurer'' and a ``training scheduler''---is critical \cite{ref199}. For ICL, difficulty can be measured not just by task complexity but by the model's own developmental stage. Methods like Self-Paced Learning (SPL) allow the model to dictate its own curriculum by weighting easy samples more heavily in the early stages of training, a principle that has been formalized probabilistically and applied to reinforcement learning to avoid poor local optima \cite{ref201}. In the context of ICL, this translates to starting with demonstrations that are highly salient and have clear input-label mappings, and only gradually introducing more ambiguous or noisy examples as the model's loss decreases, smoothing the optimization landscape and preventing the model from collapsing into a pure memorization strategy before ICL circuits can solidify.

When dealing with the multi-task nature of pretraining, selecting the right sequence of tasks is paramount. Model-based approaches to curriculum design, where the model itself evaluates the difficulty of datasets or individual instances, have been shown to result in strong representations and significant performance improvements, particularly on the most difficult instances \cite{ref202}. This is directly applicable to ICL, where the ultimate goal is to learn from a few demonstrations of a novel, hard task. By dynamically assessing the model's current competence on a held-out validation set of in-context tasks, a scheduler can sequence from tasks with high structural similarity to those requiring more complex analogical reasoning. This approach directly counters the ``transient nature'' of ICL, where the ability emerges and then decays in favor of in-weights learning, by constantly reinforcing the meta-learning objective through a carefully orchestrated sequence of tasks that are just beyond the model's current comfort zone.

Further, the mechanics of optimization itself reveal why poorly designed curricula fail. When curricula are employed in combination with adaptive optimizers like Adam, they can inadvertently learn to adapt to suboptimally chosen optimization parameters, leading to brittle results that do not outperform a well-tuned baseline \cite{ref203}. This insight is crucial for ICL training, as it suggests that the benefit of a curriculum is not automatic; it must be designed to shape the learning dynamics in a way that is synergistic with the optimizer. A curriculum that prematurely increases the difficulty of in-context tasks can lead to a sudden spike in loss, causing Adam's preconditioner to adapt to a noisy gradient regime, which can destabilize the fragile circuit formation. A well-paced curriculum, by contrast, provides a smoother gradient signal, allowing the optimizer to gradually navigate the loss landscape towards the basin of attraction that corresponds to a robust induction head circuit, ultimately equipping the model with the stable ICL capability that the subsequent section's efficiency and scalability techniques can then leverage.

Finally, the concept of an ``in-sample curriculum'' offers a granular approach to bridging the gap between simple pattern recognition and complex in-context generation. Inspired by the ``easy-to-hard'' intuition, a model can be trained to first complete the last few tokens of a target sequence, with the generation window gradually extending to cover the entire output \cite{ref204}. This strategy is particularly effective for overcoming the developmental plateau in ICL because it directly scaffolds the core mechanism of an induction head: attending to and copying a token or pattern from the past. By initially training the model on a trivial in-context copy task before gradually increasing the contextual distance and semantic complexity of the mapping, the curriculum forces the sequential learning of nested operations that phenomenologically models the abrupt emergence of ICL. This targeted, bottom-up approach to curriculum design, informed by the mechanistic understanding of induction head formation, provides a principled path to cultivating more stable and generalizable in-context learners, ready to be deployed in the efficient, many-shot, and long-context regimes that are the focus of the next section.

\subsection{Efficiency and Scalability for Many-Shot and Long-Context Regimes}

The structured curricula and model-guided training schedules described in the previous section are essential for cultivating stable and generalizable in-context learning (ICL) circuits. Once these abilities are reliably established, the next frontier is deploying them efficiently at scale. However, the remarkable success of ICL is tempered by significant computational and memory bottlenecks that emerge when scaling to large demonstration sets or processing long context windows. The standard ICL paradigm, which concatenates all demonstrations with the query into a single, monolithic prompt, incurs quadratic complexity in the self-attention mechanism with respect to the total sequence length. This fundamental limitation constrains both the number of in-context examples that can be practically utilized and the efficiency of inference, creating a pressing need for methods that decouple ICL performance from raw computational cost. This section addresses these challenges by surveying techniques for efficient inference, context-free alternatives, many-shot scaling, and context compression---all of which assume the model already possesses the robust ICL capabilities that the preceding curriculum strategies aim to establish.

\textbf{Batch Inference and Order-Agnostic Processing.} A primary direction for improving ICL efficiency involves restructuring the inference process itself. Standard N-shot ICL processes all demonstrations and the query in a single forward pass, which not only scales poorly but also introduces a notorious sensitivity to the order of demonstrations. The Batch-ICL framework \cite{ref197} reformulates this by treating ICL as a meta-optimization process, directly leveraging the principles of implicit in-context gradients explored in earlier sections. Instead of one N-shot forward computation, Batch-ICL performs N separate 1-shot forward passes and aggregates the resulting meta-gradients. These aggregated meta-gradients are then applied to the zero-shot query computation to produce the final prediction. This batch processing approach renders the model agnostic to demonstration order, consistently outperforming most random permutations, and in some cases exceeding the best possible ordering, all while reducing computational overhead. The multi-epoch variant of Batch-ICL further explores implicit permutations of examples, enhancing performance without exhaustive enumeration.

Extending the batch processing paradigm, Batch Prompting \cite{ref205} enables LLMs to process multiple inference samples simultaneously within a single prompt, rather than sequentially. By grouping several queries into one batched input, the number of inference calls is drastically reduced, achieving near-linear inverse scaling of cost with batch size. Theoretically and empirically, this approach reduces both token consumption and time costs by up to a factor of five while maintaining or improving downstream performance on tasks ranging from commonsense QA to arithmetic reasoning. This method is particularly impactful for API-based LLM deployments where per-call latency and financial cost are primary constraints. These efficiency gains are complementary to the data efficiency improvements achieved through the curriculum strategies discussed previously, together forming a holistic pipeline from training to deployment.

\textbf{Calibration-Free and Context-Free Inference.} Moving beyond the constraints of finite context windows entirely, kNN Prompting \cite{ref206} proposes a paradigm shift away from the standard prompt-completion formulation. The method first queries the LLM with all available training data to obtain distributed representations for each example. During inference, a test instance is predicted by simply referencing its nearest neighbors in this representation space, leveraging the LLM's output distribution to align test and training instances rather than directly mapping to a label space. This approach is inherently calibration-free, as it sidesteps the biases introduced by prompt formatting and verbalizer selection that plague conventional ICL. Crucially, kNN Prompting is a ``beyond-context'' solution, scaling effectively with up to 1024 shots and across model sizes from 0.8B to 30B parameters, thereby bridging data scaling with model scaling in a gradient-free deployment paradigm. This represents a natural evolution from the in-weights versus in-context learning dynamics discussed in the curriculum section, offering a hybrid that leverages distributed representations learned during training without being constrained by context window limitations.

\textbf{Scaling to Many-Shot Regimes.} The expansion of context windows in modern LLMs has opened the door to many-shot ICL, where hundreds or even thousands of demonstrations are provided. The investigation into Many-Shot In-Context Learning \cite{ref15} reveals significant performance gains when transitioning from few-shot to many-shot settings across diverse generative and discriminative tasks. However, this regime is often bottlenecked by the scarcity of human-generated examples. To address this, the study introduces Reinforced ICL, which uses model-generated chain-of-thought rationales in place of human examples, and Unsupervised ICL, which removes rationales entirely and prompts the model solely with domain-specific questions. Both strategies prove effective in the many-shot regime, particularly for complex reasoning tasks, and demonstrate a capacity to override pretraining biases and learn high-dimensional functions with numerical inputs---capabilities not observed in few-shot learning. This ability to override biases and learn complex functions is precisely the outcome that well-designed curricula aim to achieve by preventing premature collapse into in-weights memorization strategies.

\textbf{Context Compression and Efficient Caching.} As context lengths grow, the memory footprint of key-value (KV) caches in transformer models becomes a dominant bottleneck. Several compression techniques have emerged to address this. Compressed Context Memory \cite{ref207} introduces a system that continually compresses accumulating attention KV pairs into a compact memory space using a lightweight conditional LoRA module integrated into the forward pass, without full model fine-tuning. This approach achieves performance comparable to a full context model with a five-fold reduction in memory size and supports streaming settings with unlimited context lengths.

From a structural perspective, Structured Prompting \cite{ref208} scales ICL to thousands of examples by encoding demonstrations separately with well-designed position embeddings and attending to them jointly with the test example using a rescaled attention mechanism. This achieves linear complexity scaling with the number of exemplars, rather than quadratic, enabling substantial end-task performance improvements and reduced evaluation variance as the demonstration count increases. Similarly, context compression via sentinel tokens \cite{ref209} provides a plug-and-play method to incrementally compress intermediate activations of specified token spans into compact representations, reducing both memory and computation for subsequent context processing without sacrificing fluency or semantic fidelity.

Adaptive approaches to KV cache management, such as FastGen \cite{ref210}, profile the intrinsic structure of attention heads to build the cache adaptively. By evicting long-range contexts on heads that emphasize local information, discarding non-special tokens on special-token-focused heads, and retaining the standard cache only for broadly attending heads, substantial GPU memory reductions are achieved with negligible generation quality loss, all without resource-intensive fine-tuning. Cross-attention-based architectures, such as XC-Cache \cite{ref211}, further circumvent the quadratic self-attention cost by conditioning generation on reference text through cross-attention, drastically reducing the space footprint of KV caching by two orders of magnitude while outperforming standard ICL. These compression and caching innovations ensure that the robust induction head circuits and generalizable in-context abilities painstakingly cultivated through curriculum learning can be deployed in resource-constrained environments and scaled to the demands of real-world applications.

\section{Expanding Modalities: Vision, Speech, and Cross-Modal In-Context Learning}

\subsection{Foundations of Multimodal In-Context Learning}

Multimodal In-Context Learning (M-ICL) represents a paradigm shift in machine learning, extending the remarkable emergent abilities of Large Language Models (LLMs) to the multi-sensory world. While the following subsection will explore pure vision ICL, where tasks and prompts are expressed entirely in pixels, this section establishes the foundational principles of learning from demonstrations that span multiple modalities. Formally, Multimodal ICL can be defined as the capability of a model to adapt to new, unseen tasks by conditioning on a prompt that contains interleaved modalities---such as text, images, audio, and video---without any gradient updates or parameter fine-tuning. This definition positions M-ICL as a direct descendant of text-only ICL, which was popularized by the observation that large models could perform few-shot learning by simply prepending a few input-output examples to a query \cite{ref212}. Unlike its unimodal predecessor, M-ICL must not only infer the task's underlying mapping from input to output but also resolve the complex, often heterogeneous, interactions between different sensory streams within the prompt context.

The emergence of M-ICL is intrinsically linked to the evolution of foundational models that natively process multiple modalities. Early multimodal models were predominantly task-specific and relied on modality-specific encoders fused through cross-attention, a paradigm that required extensive fine-tuning on paired data \cite{ref213,ref214}. The transition to M-ICL began with the scaling of vision-language models and the adoption of transformer architectures that could treat multimodal inputs as a unified sequence. Models like Kosmos-1 demonstrated that a Multimodal Large Language Model (MLLM) trained from scratch on web-scale text and image-caption pairs could learn in context, performing tasks like visual question answering and image captioning with few-shot examples \cite{ref11}. This marked a critical departure from earlier cascade systems, such as those that separately processed speech and text, by enabling intrinsic cross-modal conversational abilities, as seen in SpeechGPT, which leverages discrete speech representations to perform cross-modal instruction following \cite{ref215}.

A crucial distinction must be drawn between M-ICL, multimodal fine-tuning, and multimodal instruction tuning. Multimodal fine-tuning involves updating the model's parameters using a downstream task-specific dataset, typically through supervised learning. This is a static adaptation process that requires computational resources and data for every new task. Multimodal instruction tuning, on the other hand, is a specialized form of fine-tuning where the model is trained on a diverse set of tasks formatted as instructions, with the goal of improving zero-shot task generalization without explicit demonstrations at inference time \cite{ref216}. M-ICL fundamentally differs from both. It operates entirely at inference time, leveraging the model's pretrained knowledge and its ability to find patterns in the provided context. Methods like MMICT (Multi-Modal In-Context Tuning) and M\(^2\)IXT (MultiModal In-conteXt Tuning) blend the boundaries of these paradigms by introducing lightweight modules or tuning strategies that enhance a frozen model's ability to process in-context multimodal demonstrations, thereby boosting the ICL capability itself without full fine-tuning \cite{ref217,ref218}. This hybrid approach aims to combine the efficiency of ICL with the performance boost of targeted adaptation.

The core problem formulation for M-ICL involves a prompt structure that can be decomposed into three parts: a set of multimodal demonstrations \(D = \{ (x_1, y_1), ..., (x_k, y_k) \}\), a query input \(x_{query}\), and a task specification, which may be implicit in the demonstrations or stated as an instruction. Each input \(x_i\) can be a composite of multiple modalities, such as an image-text pair, an audio-video clip, or a 3D point cloud with a textual description. The model's objective is to generate the output \(y_{query}\) by modeling the conditional probability \(P(y_{query} | D, x_{query})\). This formulation introduces challenges that are nonexistent in text-only ICL. The model must learn the cross-modal correspondence and alignment within each demonstration, a task complicated by the fact that the semantic information of different modalities at the same layer may not be uniform, leading to a collapse of cross-modal interaction into limited semantic information if not properly managed \cite{ref219}. The concept of ``task vectors'' in this context must therefore capture not just the input-output mapping but also the intricate cross-modal relationships that define the task.

Key challenges in M-ICL are dominated by the complexity of cross-modal interaction. A primary hurdle is understanding how different modalities contribute to the learning signal. Counterintuitive findings from recent studies reveal that in many MLLMs, M-ICL performance is predominantly driven by textual information, with the visual modality in demonstrations having little to no influence on the final outcome \cite{ref220,ref221}. This suggests that models may struggle to effectively extrapolate from visual contextual examples, a challenge that M\(^2\)IXT directly addresses by designing a module to enhance the perception of labeled examples from multiple modalities \cite{ref218}. This phenomenon also highlights the issue of \textbf{modality imbalance}, where one modality dominates the learning process, potentially leading to suboptimal performance and a failure to leverage the full richness of multimodal data \cite{ref222}. As will be seen in the subsequent discussion of visual prompting, overcoming this imbalance is a central theme that has driven the development of image-centric formulations where the visual modality is placed on equal footing with text.

Further challenges include the \textbf{modality gap} and the need for robust, fine-grained interaction. The modality gap refers to the inherent representational differences between modalities, making it difficult for a single model to perform joint reasoning \cite{ref223}. To bridge this gap, Deeply Coupled Cross-Modal Prompt Learning (DCP) methods propose mechanisms like Cross-Modal Prompt Attention (CMPA) to enable progressive and strong mutual exchange of representations \cite{ref224}. This is essential for tasks requiring fine-grained understanding, such as multimodal sentiment analysis, where a token-level contrastive learning approach with modality-aware prompting can align features from text, video, and audio effectively \cite{ref225}. The challenge of \textbf{sensitivity to demonstration selection} is also amplified in M-ICL. While text-only ICL is sensitive to example order and choice, multimodal demonstrations add a dimension of complexity, requiring the selection of examples that are not only similar in task but also cohesive in their cross-modal interactions. Strategies like Mixed Modality In-Context Example Selection (MMICES) which consider both visual and language modalities, have been proposed to select more effective demonstrations \cite{ref221}.

The scope of M-ICL research spans a broad range of domains, forming the overarching framework within which the specialized techniques of visual prompting and other modality-specific ICL paradigms are developed. In vision, this includes visual prompting for pure vision tasks and text-to-image ICL, where models learn to generate images from contextual examples, a task that remains in its infancy and presents significant difficulties for current MLLMs \cite{ref226}. In speech and audio, the challenge lies in integrating these continuous, time-series modalities into a framework originally designed for discrete text, with models like SpeechGPT pioneering the use of discrete representations to enable cross-modal conversational ICL \cite{ref215}. Cross-modal domains, such as 3D perception and video understanding, introduce further complexity with temporal and spatial dynamics. The ultimate goal is the development of unified, general-purpose M-ICL systems that can achieve high-modality representation learning, scaling to a large set of diverse modalities and demonstrating the crucial behavior that performance continues to improve with each new modality added \cite{ref227}. This foundation sets the stage for a deep dive into the specific mechanisms, methodologies, and frontiers that define the rapidly evolving landscape of multimodal in-context learning, beginning with the specialized domain of vision in-context learning and visual prompting.

\subsection{Vision In-Context Learning and Visual Prompting}

Building directly upon the foundational principles of Multimodal In-Context Learning (M-ICL), this subsection examines a radical formulation of the paradigm: pure vision ICL, where tasks, examples, and queries are expressed entirely in pixels, eliminating the textual intermediaries that often dominate multimodal prompts, as analyzed in Section 6.3. This approach addresses the modality imbalance and cross-modal alignment challenges introduced earlier by placing the visual modality on equal footing, demanding novel mechanisms for learning visual functions and transformations directly from pixel-space demonstrations. We survey the key paradigms of visual prompting, including generalist painter models, inpainting-based approaches, and diffusion-based generation, tracing their evolution toward unified, general-purpose visual ICL systems.

A foundational challenge in visual ICL, inherited from the broader M-ICL formulation of heterogeneous multimodal inputs, is the vast diversity of output representations across vision tasks, ranging from low-level image processing to high-level semantic segmentation. The Painter model {[}Images Speak in Images A Generalist Painter for In-Context Visual
Learning{]} addressed this by redefining the output of all core vision tasks as images and specifying task prompts as also images. This ``image-centric'' solution allows for an extremely simple training process: standard masked image modeling is performed on the stitched input-output image pairs. During inference, a pair of input and output images from a given task serves as the visual prompt, conditioning the model to perform the same task on a query image. This paradigm directly mirrors the M-ICL problem formulation of conditioning on demonstration pairs \(D = \{ (x_1, y_1), ..., (x_k, y_k) \}\), but with every \(x_i\) and \(y_i\) existing in a unified pixel space, thereby elegantly sidestepping the cross-modal interaction friction discussed in the previous subsection. Painter achieved competitive performance against task-specific models across seven representative vision tasks, demonstrating that generalist visual ICL is feasible without any task-specific architectural modifications.

Building on this generalist philosophy, inpainting-based approaches frame visual ICL as a conditional generation problem, offering a unified pixel-level objective that resonates with the unified sequence modeling of transformer-based MLLMs from Section 6.1. IMProv \cite{ref228} is a generative model capable of in-context learning from multimodal prompts, including textual descriptions and visual examples. By training a masked generative transformer on a dataset of figures from computer vision papers, IMProv learns to inpaint the corresponding output given a task description and input-output visual examples. This work demonstrated that scaling the dataset size and incorporating text conditioning improves in-context learning for tasks like foreground segmentation and single object detection, with vision and language prompts proving complementary---a finding that prefigures the multimodal prompt engineering synergies discussed in the following subsection. The inpainting formulation elegantly unifies diverse vision tasks under a common pixel-level generation objective, circumventing the need for task-specific output heads and directly addressing the output heterogeneity challenge that complicates general M-ICL systems.

A key insight from M-ICL, that demonstration selection and quality dramatically impact performance, is profoundly amplified in the visual domain. The efficacy of visual ICL is highly sensitive to the quality of prompts, motivating methods that optimize the in-context pairs themselves. Instruct Me More (InMeMo) \cite{ref229} introduces learnable perturbations to augment in-context pairs, exploring the potential of random prompting in visual ICL. By adding lightweight trainable parameters to the prompt, InMeMo boosts performance significantly over baselines, achieving gains of 7.35 and 15.13 in mIoU for foreground segmentation and single object detection, respectively. This indicates that even simple, learned adjustments to the visual prompts can unlock substantial performance improvements, mirroring the lightweight tuning philosophies that will be fully explored in Section 6.4. Complementary work on prompt selection and fusion \cite{ref230} identified prompt selection and prompt fusion as two major factors directly impacting inference performance. Their prompt-SelF framework uses pixel-level retrieval to select suitable prompts and then employs different prompt fusion methods to activate knowledge stored across the large-scale model, ensembling predictions for final results. Notably, prompt-SelF was the first to outperform OSLSM-based meta-learning in 1-shot segmentation, underscoring the potential of visual ICL and directly addressing the sensitivity to demonstration selection challenge that plagues both text-only and multimodal ICL.

Diffusion-based generative models have emerged as a powerful substrate for visual ICL, capable of conditioning generation on visual examples to perform a wide array of tasks. This represents a natural extension of the generative underpinnings that enable M-ICL's emergence from scaled multimodal models. Prompt Diffusion \cite{ref231} is a framework that enables in-context learning in diffusion-based generative models. By proposing a vision-language prompt that models a range of vision-language tasks, and training a diffusion model jointly over six different tasks, Prompt Diffusion becomes the first diffusion-based vision-language foundation model capable of in-context learning. The model performs high-quality in-context generation on trained tasks and generalizes effectively to new, unseen vision tasks with their respective prompts. Extending this paradigm, improved Prompt Diffusion (iPromptDiff) {[}Improving In-Context Learning in Diffusion Models with Visual
Context-Modulated Prompts{]} integrates an end-to-end trained vision encoder that converts visual context into an embedding vector used to modulate the token embeddings of text prompts, thereby fostering the deep cross-modal alignment highlighted as critical in Section 6.1. When equipped with a ControlNet structure, iPromptDiff exhibits versatility across diverse training tasks and excels in in-context learning for novel vision tasks like normal-to-image or image-to-line transformations. Context Diffusion \cite{ref232} further advances this line by proposing a framework that separates the encoding of visual context from preserving the structure of query images, enabling the model to learn from visual context in few-shot settings and enhancing image quality and fidelity over counterpart models.

The evolution toward unified visual ICL frameworks, driven by the same scaling principles that gave rise to generalist MLLMs, has also seen the development of multi-task generalist interfaces. InstructDiffusion \cite{ref233} casts diverse vision tasks into a human-intuitive image-manipulating process whose output space is a flexible pixel space. Built upon the diffusion process, it predicts pixels according to user instructions, handling both understanding tasks like segmentation and generative tasks like editing and enhancement. UniControl {[}UniControl A Unified Diffusion Model for Controllable Visual Generation
In the Wild{]} consolidates a wide array of controllable condition-to-image tasks within a singular framework, using a task-aware HyperNet to modulate the diffusion models for adaptation to different tasks simultaneously, demonstrating impressive zero-shot generation abilities with unseen visual conditions. These unified frameworks embody the ultimate goal of M-ICL from Section 6.1: developing systems capable of handling a broad range of tasks without task-specific modifications.

A critical emerging question, connecting back to the core M-ICL challenge of how models represent and leverage cross-modal task vectors, is how models internally represent and generalize transformations learned in context. The Im-Promptu framework \cite{ref234} investigates whether analogical reasoning can enable in-context composition over composable elements of visual stimuli. By training agents with different levels of compositionality---vector representations, patch representations, and object slots---it reveals trade-offs between extrapolation abilities and the degree of compositionality, with object-centric tokenizers coupled with cross-attention modules generating consistent and high-fidelity solutions. This suggests that the inductive bias of compositional object representations is crucial for generalization in visual ICL, providing a mechanistic insight into how vision models might construct internal ``task vectors'' analogous to those hypothesized in text-only ICL but operating over visual primitives.

Finally, a technique that directly bridges pure visual ICL with the multimodal prompt engineering of Section 6.3 is the consolidation of all prompt information into a single aggregated image. In-Image Learning (I\(^2\)L) \cite{ref235} combines demonstration examples, visual cues, and chain-of-thought reasoning into a single aggregated image. This consolidation reduces inaccurate textual descriptions of complex images, provides flexibility in positioning demonstrations, and avoids multiple input images and lengthy prompts---directly addressing the cross-modal interaction friction and textual bias that Section 6.3 identifies as central hurdles in multimodal ICL. The method proved effective on complex multimodal reasoning tasks, with further improvements from a hybrid strategy that automatically selects between I\(^2\)L and other ICL methods per task instance. This conceptually connects to the structured prompting techniques like Scaffold prompting discussed subsequently, where visual overlays promote explicit vision-language coordination, but here the entire prompt is treated as a visual scene to be reasoned over, effectively demonstrating a limit case where the visual modality fully subsumes the textual.

Overall, visual in-context learning has rapidly progressed from task-specific adaptations to generalist, unified frameworks that learn visual functions and transformations directly from pixel-space examples. The key insight across these works is that redefining tasks in terms of image-to-image mappings, combined with scalable generative architectures and optimized prompt selection and fusion strategies, enables robust and flexible in-context visual reasoning without textual intermediaries. This visual-centric paradigm not only addresses the modality imbalance highlighted in Section 6.1 by placing vision on equal footing but also sets the stage for the multimodal prompt engineering techniques of Section 6.3, which must ultimately reconcile these powerful pixel-space learning strategies with the textual biases and cross-modal alignment challenges inherent in prevailing Vision-Language Models.

\subsection{Vision-Language Models and Multimodal Prompt Engineering}

Building upon the visual prompting paradigms that learn pixel-space transformations and the critical enablers of lightweight tuning and cross-modal alignment, the extension of In-Context Learning (ICL) to Vision-Language Models (VLMs) introduces a new dimension of complexity: the prompt is no longer a purely textual sequence or a unified visual scene but an interleaved structure of images and text. This subsection examines the strategies for engineering these multimodal prompts, examining how demonstrations are selected, the often-surprising dominance of textual information over visual information, and the advanced prompting techniques that have emerged to better leverage the synergy between modalities, ultimately bridging the gap between visual ICL formulations and the parameter-efficient adaptation mechanisms discussed in the previous and following sections.

A central concern in multimodal ICL is the selection of effective demonstrations. Early approaches often adapted text-based retrieval methods, but the presence of images necessitates a joint consideration of both modalities. The standard paradigm involves selecting demonstration pairs---each consisting of an image and a corresponding question-answer text---to condition the model for a new query. A foundational finding in this area is the counter-intuitive dominance of the textual modality. Several studies have systematically demonstrated that the textual information in demonstrations is the primary driver of ICL performance in many VLMs, while the visual information in the provided examples often has a negligible impact \cite{ref221,ref220}. For instance, one study found that ICL in VLMs is ``predominantly driven by the textual information in the demonstrations,'' with model information flow analysis confirming this bias \cite{ref221}. This suggests that large VLMs, often built upon powerful LLM backbones, are highly adept at leveraging the linguistic task specifications and label spaces provided in text, sometimes to the point of ignoring the accompanying visual examples. This phenomenon is closely related to the established ability of text-only ICL to regulate the label space and output format, which can independently account for a significant portion of performance gains, highlighting the need for explicit mechanisms to foster deeper cross-modal interaction as pursued by the lightweight tuning methods in Section 6.4 \cite{ref7}.

In response to this textual bias, advanced multimodal demonstration selection strategies have been proposed to more effectively incorporate visual signals. The Mixed Modality In-Context Example Selection (MMICES) approach explicitly considers both visual and language modalities when selecting demonstrations, thereby counteracting the dominance of text and yielding better ICL performance than text-only selection methods \cite{ref221}. Furthermore, sophisticated retrieval pipelines have been developed to align the visual content of demonstrations with the task intent. The Visual In-Context Learning (VICL) method, for example, employs a ``Retrieval \& Rerank'' paradigm for images, coupled with intent-oriented image summarization and demonstration composition. This process not only selects relevant images but also parses them with task-specific visual understanding and converts them into concise language descriptions to reduce token count and alleviate cross-modal interaction friction \cite{ref236}. This approach represents a key trend in multimodal prompt engineering: using language as a mediating bridge to efficiently convey visual task information gathered from demonstrations, echoing the image-to-text formulations seen in some visual ICL methods.

Moving beyond the selection of discrete examples, prompt engineering in VLMs involves designing the structure and content of the prompt itself. A significant line of work focuses on learning continuous, soft prompts in a parameter-efficient manner, which can be seen as a form of model-level prompt engineering and directly connects to the lightweight adapter and prompt tuning techniques detailed in the subsequent subsection. Early methods like MaPLe proposed learning prompts jointly for both the vision and language branches of a model like CLIP, promoting strong coupling to ensure mutual synergy and prevent independent uni-modal solutions \cite{ref237}. This multi-modal prompting strategy was further refined by APLe, which tunes tokens for both modalities in a sequential, layer-wise manner \cite{ref238}, and by ProMPT, which uses a progressive, iterative evolution of multi-modal prompts to align vision-language features more deeply \cite{ref239}. A unified approach like APoLLo combines adapters and prompts, using cross-attention to strengthen cross-modal alignment in few-shot settings, foreshadowing the unified adapter frameworks discussed next \cite{ref240}. These methods, while often not framing themselves strictly as ``in-context learning'' with multiple demonstration examples, share the core ICL principle of conditioning model behavior on external input without full fine-tuning, using learned prompt vectors as a highly compressed form of task demonstration.

For tasks requiring complex reasoning over dense visual information, structured prompting techniques have proven highly effective. A notable method is ``Scaffold prompting,'' which overlays a dot matrix onto the image as visual anchors and provides multi-dimensional coordinates as textual positional references. This simple yet powerful visual prompt scheme promotes explicit vision-language coordination, significantly outperforming standard textual Chain-of-Thought (CoT) prompting on challenging tasks that require spatial reasoning \cite{ref241}. This approach directly addresses the cross-modal interaction problem by providing a shared, grounded coordinate system that links the visual and textual reasoning streams, akin to how the inpainting-based visual ICL methods provide a unified pixel-level objective.

Beyond coordinate grounding, other advanced techniques focus on decomposing and explicitly aligning the reasoning processes of different modalities. The Bi-Modal Behavioral Alignment (BBA) prompting method tackles complex multimodal reasoning in professional domains by first instructing the VLM to generate separate reasoning chains for visual and domain-specific language (DSL) representations. It then explicitly aligns these chains by resolving inconsistencies, leading to substantial performance improvements on tasks like geometry problem solving and molecular property prediction \cite{ref242}. Similarly, the concept of an ``inner monologue'' has been adapted for VLMs, where a model is trained to engage in a self-questioning and answering dialogue to simulate a cognitive reasoning process that synthesizes visual and textual information more effectively \cite{ref243}. These alignment-centric techniques resonate with the cross-modal alignment objectives that are central to the lightweight tuning methods in Section 6.4.

Another innovative paradigm, directly extending the visual ICL concept from Section 6.2, is the consolidation of all prompt information into a single, aggregated image, a technique introduced as In-Image Learning (I\(^2\)L) \cite{ref235}. This method combines demonstration examples, visual cues, and even chain-of-thought reasoning steps into one composite image. By doing so, it circumvents the limitations of lengthy, interleaved text-image prompts, avoids potentially inaccurate textual descriptions of complex visual scenes, and leverages the strong holistic image processing capabilities of Large Multimodal Models. This represents a radical departure from standard prompt formats, effectively treating the entire prompt as a visual scene to be ``read'' and reasoned over by the model, and serves as a conceptual bridge between the pure visual ICL approaches and the multimodal prompt engineering discussed here.

Finally, a critical aspect of multimodal prompt engineering is understanding and mitigating the fragility and biases of these systems, a consideration that likewise motivates the robust, parameter-efficient adaptation methods that follow. Evaluations have revealed that even the most advanced models like GPT-4V can exhibit a significant drop in trustworthiness when tasks become more challenging, performing inconsistently across modalities \cite{ref244}. A simple yet effective method to address this is ``Vision Description Prompting,'' which explicitly injects a descriptive, language-based summary of the visual content into the prompt, effectively guiding the model's visual reasoning and improving its reliability \cite{ref244}. Furthermore, manipulating the label space of demonstrations to increase ``knowledge density''---through techniques like label distribution enhancement or by adding visual descriptions to labels---can improve classification performance with fewer shots, demonstrating that the format and content of textual labels are powerful levers in multimodal prompt design \cite{ref245}. As the field matures, comprehensive benchmarks like VL-ICL Bench are essential for dissecting these complex interactions, revealing that multimodal ICL involves a spectrum of challenges from perception and reasoning to managing long context lengths, and that even the most advanced models find many of these tasks genuinely challenging \cite{ref13}. These insights into model behavior and fragility directly inform the design of the lightweight yet powerful tuning and alignment strategies that are the focus of the next subsection.

\subsection{Lightweight Tuning and Cross-Modal Alignment for ICL}

As discussed in the previous subsection, the effectiveness of multimodal in-context learning (ICL) is frequently hampered by a predominant reliance on textual cues and shallow cross-modal interaction, which limits the model's ability to fully exploit interleaved visual and linguistic demonstrations. Lightweight tuning and explicit cross-modal alignment have therefore emerged as critical enablers for bridging this gap, addressing the prohibitive computational costs of full fine-tuning while fostering the deep integration between modalities that robust multimodal ICL demands. As multimodal large language models (MLLMs) scale, parameter-efficient fine-tuning (PEFT) methods become essential to adapt these foundation models to downstream tasks without updating all parameters. However, the direct application of ICL in multimodal settings often suffers from two intertwined challenges: the computational burden of processing extensive interleaved multimodal demonstrations in the context window, and the shallow cross-modal alignment that limits the model's ability to effectively reason across visual and textual information. A growing body of research demonstrates that purpose-built lightweight modules and alignment strategies can significantly enhance multimodal ICL capabilities while maintaining a minimal parameter footprint, directly tackling the textual dominance and fragile visual grounding issues identified previously.

A foundational approach to this problem is the design of dedicated in-context tuning modules that can be prepended to existing multimodal unified models. The MultiModal In-conteXt Tuning (M\(^2\)IXT) module exemplifies this paradigm, acting as a lightweight, expandable component that perceives an extended context window to incorporate labeled examples from multiple modalities, including text, image, and coordinate data \cite{ref218}. By prepending M\(^2\)IXT to diverse architectures like OFA, Unival, and LLaVA, and training it with a mixed-tasks strategy on as few as 50K multimodal data points, researchers have observed significant boosts in few-shot ICL performance, achieving an 18\% relative increase for OFA and state-of-the-art results on tasks such as visual question answering, image captioning, and visual grounding, all while being approximately 20 times smaller than comparable models like Flamingo. Similarly, the Context-Aware MultiModal Learner (CaMML) provides a hierarchical, lightweight module that seamlessly integrates multimodal contextual samples into large models, enabling them to derive knowledge from analogous, domain-specific information and make grounded inferences, with CaMML-based models surpassing strong baselines like LLaVA-1.5 on numerous benchmarks without external resources \cite{ref246}. These modules underscore a design principle that mirrors the structured prompting techniques from Section 6.3: instead of retraining the entire model to handle multimodal ICL, strategically inserting a small number of trainable parameters can effectively condition the frozen backbone for in-context adaptation, much like how scaffold prompts condition spatial reasoning without model modification.

Beyond standalone modules, the integration of adapter-based methods and prompt tuning into multimodal contexts provides a systematic framework for lightweight cross-modal knowledge transfer that extends the continuous prompt learning paradigms touched upon earlier. Methods like UniAdapter unify unimodal and multimodal adapters for parameter-efficient cross-modal adaptation on pre-trained vision-language models, distributing adapters across different modalities and their interactions with partial weight sharing to create powerful cross-modal representations while tuning only 1.0\%-2.0\% of the pre-trained parameters \cite{ref247}. This unified design has proven to outperform full fine-tuning on several cross-modal benchmarks, including video-text retrieval and visual question answering. The Cross-Modal Adapter (XMAdapter) further advances this direction by establishing cache models for both text and image modalities and using retrieval through visual-language bimodal information to gather clues for inference, dynamically adjusting the affinity ratio to achieve cross-modal fusion and decoupling modal similarities to assess their respective contributions \cite{ref248}. This retrieval-based approach allows for sample-efficient adaptation, outperforming previous adapter methods in accuracy, generalization, and efficiency, and connects back to the retrieval and reranking strategies for demonstration selection discussed in multimodal prompt engineering. In the realm of prompt learning, deep cross-modal fusion mechanisms have proven particularly impactful for combating the shallow interaction that characterizes naive ICL. The Deeply Coupled Cross-Modal Prompt Learning (DCP) method introduces a Cross-Modal Prompt Attention (CMPA) mechanism that enables the mutual exchange of representations through a well-connected multi-head attention module, achieving progressive and strong interaction between vision and language branches \cite{ref224}. This deep coupling contrasts with earlier methods that learned prompts in isolation for each modality or engaged in shallow interaction, leading to superior few-shot generalization and robust domain adaptation. Similarly, Multi-modal Deep-symphysis Prompt Tuning (MuDPT) extends independent multi-modal prompt tuning by learning a model-agnostic transformative network for deep hierarchical bi-directional prompt fusion, achieving synergistic alignment that improves both recognition and generalization performance \cite{ref249}, directly addressing the cross-modal coordination challenges that scaffold prompting tackles through coordinate grounding.

The concept of meta-alignment and the explicit transfer of ICL capabilities across modalities represent a frontier in this area, pushing toward the generalizable few-shot behavior required for the broader multimodal ICL vision. The \(\pi\)-Tuning method introduces a universal parameter-efficient transfer learning approach that aggregates the parameters of lightweight task-specific experts learned from similar tasks, based on similarities predicted in a unified modality-independent space \cite{ref250}. This framework flexibly explores both intra- and inter-modal transferability, demonstrating that knowledge can be effectively transferred not just from unimodal to multimodal tasks, but across different multimodal tasks, which is a foundational requirement for general-purpose multimodal ICL systems. Gradient-Regulated Meta-prompt learning (GRAM) addresses the sensitivity and overfitting issues of prompt tuning in few-shot scenarios by jointly meta-learning an efficient soft prompt initialization and a lightweight gradient regulating function using only unlabeled image-text pre-training data, leading to consistent improvements in cross-domain and cross-dataset generalization \cite{ref158}. A critical insight from the study of instruction prompt tuning (IPT) is that while combining PT with ICL can reduce performance variance and enable positive transfer from source tasks, the in-context demonstration must be semantically similar to the test input for improvements to materialize, highlighting that the effectiveness of such combined methods is highly conditional on the specific alignment between the demonstration and the query \cite{ref31}, an insight that resonates with the nuanced demonstration selection criteria identified in multimodal prompt engineering.

The pursuit of more expressive and adaptive cross-modal fusion has led to innovations in conditional computation and mixture of experts, strategies that dynamically tailor the model's adaptation to each input instance. The Mixture of Prompt Experts (MoPE) technique disentangles vanilla prompts to adaptively capture dataset-level and instance-level features, and then routes the most effective prompt on a per-instance basis using multimodal pairing priors \cite{ref251}. This conditional prompting exhibits greater expressiveness, scales better with training data, and, through an emergent expert specialization where different experts focus on different concepts, achieves state-of-the-art results matching or surpassing fine-tuning with only 0.7-0.8\% of trainable parameters. The Conditional Mixture of LoRA (MixLoRA) applies a similar philosophy to low-rank adaptation, dynamically constructing low-rank matrices tailored to each input instance to mitigate task interference during multimodal instruction tuning, outperforming conventional LoRA \cite{ref195}. Even within the transformer architecture itself, the tuning of normalization layers has been shown to be a deceptively powerful strategy; by conceptualizing the transformation of an LLM into an MLLM as a domain adaptation process, tuning only the LayerNorm parameters within attention blocks yields strong performance while reducing trainable parameters and GPU memory usage dramatically compared to LoRA, suggesting that much of the cross-modal adaptation can be achieved by recalibrating internal feature normalization rather than altering weight matrices extensively \cite{ref252}.

The synergy between lightweight tuning and explicit cross-modal alignment is crucial for bridging the gap between the implicit learning of ICL and the robust, parameter-efficient adaptation required for real-world multimodal applications, and it directly addresses the fragility and modality inconsistency problems highlighted in analyses of VLMs. Methods like the Multi-way Adapter (MWA) explicitly incorporate an ``Alignment Enhancer'' to deepen inter-modal alignment, ensuring high transferability with minimal tuning effort and maintaining model effectiveness while reducing training time by up to 57\% \cite{ref253}. The Multi-modal Alignment Prompt (MmAP) for CLIP aligns text and visual modalities during the fine-tuning process itself, forming the basis for a multi-task prompt learning framework that leverages gradient-driven task grouping and task-specific prompts to achieve significant performance gains over full fine-tuning with only 0.09\% of trainable parameters \cite{ref254}. These efforts collectively demonstrate that the future of multimodal ICL lies not in a single technique, but in the principled combination of lightweight architectural add-ons, deep cross-modal fusion mechanisms, and meta-learned initialization strategies. By enabling frozen foundation models to dynamically leverage multimodal contextual examples with high fidelity and low overhead, these methods lay the essential groundwork for extending ICL to more complex data domains, including the 3D perception and point cloud tasks explored in the next subsection, where the challenges of cross-modal alignment and data-efficient adaptation are even more pronounced.

\subsection{3D Perception and Point Cloud In-Context Learning}

Building upon the lightweight tuning and cross-modal alignment strategies that have proven transformative for image and video domains, extending in-context learning (ICL) to 3D perception represents a critical frontier. The physical world is inherently three-dimensional and demands spatial reasoning capabilities beyond 2D image understanding, yet unlike 2D vision tasks that can leverage abundant image-text pairs from the internet, 3D perception faces a fundamental data scarcity challenge. This makes the ability to rapidly adapt to new tasks from a few 3D demonstrations particularly valuable and directly motivates the development of parameter-efficient and alignment-focused methods analogous to those discussed in previous sections. Research in this domain has progressed along three primary axes: pre-trained transformers for 3D vision-text alignment, joint multi-modal cue integration for robust 3D representations, and instruction tuning paradigms that encode spatial and geometric context for comprehensive scene understanding. These axes collectively mirror the broader trends in multimodal ICL, where deep cross-modal fusion and lightweight adaptation are paramount.

A foundational challenge for 3D ICL lies in establishing effective alignment between 3D geometric data and natural language, a problem that parallels the cross-modal alignment challenges addressed by methods like UniAdapter and MmAP in the 2D domain. Models like 3D-VisTA \cite{ref255} represent significant strides in this direction by simplifying the architecture to use self-attention layers for both single-modal modeling and multi-modal fusion, avoiding complex task-specific modules. This unified design enables easier adaptation to downstream tasks, which is a hallmark of ICL systems. To address the critical bottleneck of limited 3D-text paired data, synthetic data generation has emerged as a powerful strategy, much like how data augmentation bolsters few-shot learning in other modalities. For instance, the ScanScribe dataset, containing over 278K scene descriptions generated from existing 3D-VL tasks and language models, was constructed to pre-train 3D-VisTA, demonstrating how data augmentation can bootstrap the ICL capabilities of 3D models. Similarly, DR-Point \cite{ref256} leverages differentiable rendering to generate depth images from point clouds, creating tri-modal (RGB, depth, point cloud) training triplets that enrich the representational space and improve alignment between visual and 3D modalities.

Joint multi-modal cues are proving essential for overcoming the information degradation that occurs when aligning 3D data with only single-view 2D images or coarse-grained text categories, a challenge that resonates with the deep cross-modal fusion techniques discussed in the previous subsection, such as the deeply coupled prompt learning of DCP and MuDPT. The JM3D framework \cite{ref257} directly addresses this by introducing a Structured Multimodal Organizer (SMO) that enriches vision-language representations with contiguous multi-view images and hierarchical subcategory text. Their Joint Multi-modal Alignment (JMA) technique models the joint modality by incorporating language knowledge into the visual modality, rather than aligning each modality independently. This approach, which has demonstrated state-of-the-art zero-shot 3D classification results, is closely aligned with the principles of ICL, where models must infer task structure from multimodal contextual cues. The extension JM3D-LLM further marries these enriched 3D representations with large language models through efficient fine-tuning, pointing toward a future where 3D ICL can be performed by directly conditioning LLMs on point cloud demonstrations, analogous to how M\(^2\)IXT conditions frozen multimodal models.

The vision of a unified, omni-modal ICL system that seamlessly incorporates 3D data is being advanced by frameworks like ViT-Lens \cite{ref258}. This approach perceives novel modalities by projecting their signals into a shared embedding space through a modality-specific lens, which is then processed by a powerful, pre-trained Vision Transformer (ViT). By aligning these encoded multimodal representations with a modal-independent space defined by foundation models, ViT-Lens exhibits emergent zero-shot 3D classification and even question-answering capabilities without task-specific fine-tuning. This is a core attribute of ICL: the ability to perform new tasks by leveraging pre-existing knowledge when prompted with a few, or even zero, examples from a new data domain, much like the meta-learned generalization capabilities fostered by methods like GRAM and \(\pi\)-Tuning. The concept of an autoencoder as a cross-modal teacher is also pertinent here. The ACT framework \cite{ref259} shows that a 2D Transformer can act as a teacher for a 3D student model by distilling knowledge through discrete variational autoencoding, effectively transferring the rich structural priors learned from 2D data to the 3D domain. This transfer learning paradigm is a crucial enabler for 3D ICL, as it mitigates the scarcity of native 3D training data.

Instruction tuning for 3D scene understanding is another critical development that brings ICL closer to practical applications in embodied AI and robotics, extending the instruction-following paradigms that have proven powerful in video understanding as will be discussed in the next subsection. The 3DMIT paradigm \cite{ref260} eliminates the need for a separate alignment stage between 3D scenes and language, instead directly extending instruction prompts with 3D modality information, including entire scene context and segmented objects. This promotes an ICL-like behavior where the model can follow textual instructions to perform tasks like 3D question answering, grounding, and conversation by conditioning on the provided spatial prompt, similar to how the Four-tiered Prompts framework enriches video representations with semantic knowledge. The importance of learning rich, transferable 3D representations for robotics is highlighted by frameworks like SUGAR \cite{ref261}, which jointly pre-train on semantic, geometric, and affordance properties of objects from 3D point clouds. SUGAR's use of a transformer to address five distinct tasks (knowledge distillation, masked point modeling, grasping pose synthesis, instance segmentation, and referring expression grounding) fosters a multi-task representation that is ideally suited for ICL, where a single model must be able to adapt to a variety of downstream queries based on contextual demonstrations.

The challenges of achieving robust ICL in 3D are further compounded by the need for viewpoint-agnostic representations, a spatial reasoning requirement that distinguishes 3D ICL from its 2D and video counterparts. The 3D Token Representation Layer (3DTRL) \cite{ref262} addresses this by estimating 3D positional information from 2D patches and applying geometric transformations, enabling transformers to learn representations that are invariant to camera viewpoint. This type of spatial reasoning is fundamental for a true 3D ICL, as a user's demonstration might come from any arbitrary viewpoint. Furthermore, the interactive nature of 3D environments calls for ICL systems that can ground symbols in multi-modal instructions, as demonstrated in table-top manipulation scenarios where models learn to associate linguistic concepts like color and shape with physical objects from a few demonstrations \cite{ref263}. By integrating point clouds, multi-view images, hierarchical text, and instruction-based tuning, the 3D vision community is progressively constructing the foundational architectures and pre-training recipes necessary for powerful, few-shot learning in the physical world. This progression naturally sets the stage for the temporal dynamics of video understanding, where the spatial reasoning capabilities developed for static 3D scenes must be extended to reason over sequences of frames and evolving events.

\subsection{Video Understanding and Temporal In-Context Learning}

While the previous subsection detailed the development of static 3D perception and the spatial reasoning foundations necessary for in-context learning (ICL), extending these capabilities to the video domain introduces the critical dimension of time. Unlike the static point clouds or multi-view images central to 3D ICL, videos require models to reason over sequences of frames, capturing spatio-temporal dependencies, object interactions, and evolving events. This demand for temporal reasoning directly parallels the sequential processing needed for audio signals, which will be discussed in the next subsection, establishing a unifying thread of dynamic, time-series understanding across multimodal ICL. This subsection analyzes the application of ICL to video understanding, focusing on the adaptation of multimodal video transformers for vision-language integration over time and the use of in-context demonstrations for tasks like temporal reasoning, tracking, and action recognition.

The core of video ICL lies in the adaptation of Vision-Language Models (VLMs) and transformer architectures to process temporal information, a progression from the viewpoint-agnostic spatial representations and instruction-tuned 3D scene understanding models like 3DMIT and ViT-Lens. Early work in video understanding often relied on 3D convolutional networks or two-stream architectures, but the paradigm has decisively shifted towards transformer-based models due to their capacity for long-range dependency modeling, a capacity equally critical for capturing the spectral and temporal complexities of audio tokens. The ``all-in-one Transformer'' concept, which unifies video and text encoding into a single backbone, has been a significant step forward, mirroring the unified multi-modal fusion techniques seen in the JM3D framework for 3D data \cite{ref264}. By embedding raw video and textual signals into joint representations, these models can leverage multimodal ICL more seamlessly. A key technical challenge addressed by such frameworks is the encoding of temporal dynamics from video clips in a non-parametric manner, such as through token rolling operations, which allows a modality-agnostic transformer to process spatio-temporal data effectively, a design principle that resonates with the unified tokenization approaches foundational to audio-text ICL models like VATLM \cite{ref264}.

A central theme in video ICL is the integration of vision and language to provide a richer context for temporal reasoning. Models pre-trained on large-scale video corpora with transcribed speech, such as MERLOT, demonstrate that learning multimodal script knowledge in a self-supervised manner enables strong out-of-the-box representations of temporal commonsense \cite{ref265}. These models learn to contextualize events globally over time, matching images to temporally corresponding words and reasoning about the past, present, and future. This pre-training strategy is directly conducive to ICL, as the model develops a strong prior for how sequences of frames map to narrative structures, which can be activated by a few in-context demonstrations, much like how SUGAR learns transferable priors for 3D affordances and semantics. The importance of temporal context is further underscored by methods that explicitly model it to improve action recognition. For instance, a transformer-based multimodal model can ingest video and audio, using an explicit language model to provide action sequence context, thereby attending to surrounding actions and enhancing predictive performance \cite{ref266}.

The use of language as a mediating medium for temporal context is a powerful technique in video ICL, creating a direct conceptual bridge to the instruction-following paradigms common to both 3D and forthcoming audio ICL systems. Instead of relying solely on visual features, models can leverage natural language descriptions of the past to predict the future. The TransFusion model exemplifies this by summarizing past action context into natural language using pre-trained captioning models and then fusing this linguistic summary with the next video frame to forecast object interactions \cite{ref267}. This approach demonstrates that language-based context summaries can be more effective than vision-only approaches for anticipation tasks, highlighting the common-sense and generalization capabilities added by large pre-trained language models. In a similar vein, zero-shot multimodal video classification can be achieved entirely through textual proxies, using LLMs to reason about multimodal textual descriptions of a video's visual and aural aspects without any additional fine-tuning, a testament to the cross-modal transfer power that ICL seeks to harness \cite{ref268}.

In-context demonstrations for action recognition and detection are a critical area of investigation. The challenge lies in adapting short-term pre-trained vision transformers (ViTs), which are excellent feature extractors for individual clips, into long-form video transformers capable of capturing inter-snippet relationships for tasks like temporal action detection in untrimmed videos, directly extending the few-shot adaptation goals of 3D ICL to the temporal axis \cite{ref269}. This involves designing cross-snippet propagation modules that allow information exchange between different temporal segments, enabling the model to perform ICL with entire video sequences as context. For zero-shot spatio-temporal action detection, pre-trained visual-language models can be prompted to align interaction features from video clips with label-based text features, effectively performing ICL without any task-specific training on the target dataset \cite{ref270}.

The design of effective prompts is as crucial in video domains as it is in text and image ICL, and it sets the stage for the structured demonstration formats needed in the nascent field of audio ICL. The Four-tiered Prompts (FTP) framework provides a structured approach by aligning ViT video encodings with VLM outputs that focus on specific aspects of human action: the action category, its components, a description, and contextual information \cite{ref271}. This enriches the visual representations with semantic knowledge from VLMs, making them more generalizable and comprehensive. Although the VLM is used during training to generate these prompts, the learned richer representations can then be used in an ICL setting during inference. Similarly, language-based action concept spaces can be used to adapt image CLIP models to the video domain through self-supervised learning, where objectives operating in an action concept space derived from textual prompts enforce relations between actions and their attributes, improving zero-shot and linear probing performance \cite{ref272}.

Self-supervised learning (SSL) is a foundational technique for training video models that can later be used for ICL, a theme that resonates with the pre-training strategies used to overcome data scarcity in 3D perception and to build robust representations for audio. The goal is to learn robust spatio-temporal representations without explicit labels, which can then be prompted with a few examples. Temporal contrastive learning is a dominant paradigm, with methods like Long-Short Temporal Contrastive Learning (LSTCL) forcing a model to match a short-term view of a video to a long-term view, capturing temporal context effectively \cite{ref273}. This principle is extended by Temporal-aware Contrastive self-supervised learning (TaCo), which uses temporal transformations as both data augmentation and a source of self-supervision, enhancing the model's understanding of temporal dynamics \cite{ref274}. The Temporal Contrastive Graph (TCGL) framework further elaborates on this by jointly modeling inter-snippet and intra-snippet temporal dependencies through a hybrid graph contrastive learning strategy, explicitly capturing multi-scale temporal relationships \cite{ref275}. Representations learned from such SSL methods, including those focused on visual tempo consistency, have been shown to transfer well to downstream tasks like action detection and anticipation, forming a strong basis for in-context adaptation, much like how differentiable rendering and synthetic data in DR-Point bootstrap robust 3D representations \cite{ref276}.

End-to-end multi-modal video temporal grounding, which involves identifying a time interval based on a text query, also benefits from multi-modal fusion models that can be prompted with few-shot demonstrations, a capability that mirrors the omni-modal fusion ambitions of unified systems like i-Code and VALOR \cite{ref277}. These models integrate appearance, motion, and structural information from RGB, optical flow, and depth maps, using dynamic fusion schemes with transformers to model inter-modal interactions. The ability to learn coherent video paragraph captioning, which generates multi-sentence descriptions of untrimmed videos, is another task where ICL principles apply. The Visual-Linguistic Transformer-in-Transformer (VLTinT) models the scene through global visual environments, local visual agents, and linguistic scene elements, capturing the semantic coherence of intra- and inter-event contents \cite{ref278}. This decomposition into visual and linguistic components allows for a structured approach to prompt engineering, where demonstrations can highlight different aspects of narrative coherence.

Finally, the evaluation of video ICL is intrinsically linked to the quality of multimodal fusion. The degree to which multimodal video transformers integrate information from vision and text is partially aligned with how the human brain processes multimodal information, with vision enhancing language processing and vice versa \cite{ref279}. Fine-tuning these models on tasks requiring vision-language inferences can improve their brain alignment, suggesting that the same objectives that enhance ICL performance also lead to more cognitively plausible representations. The tracking domain also illustrates the potential of vision-language ICL, where a unified-adaptive VL representation learned through asymmetric architecture search and a modality mixer can significantly improve multiple tracking baselines \cite{ref280}. This body of work collectively demonstrates that video ICL is not merely a straightforward extension of image ICL but requires a deep integration of temporal modeling, structured multimodal prompting, and robust self-supervised pre-training. This foundational work on processing dynamic visual sequences directly informs the next frontier of multimodal ICL: the assimilation of the purely auditory temporal dynamics of speech and sound, where models must learn from acoustic demonstrations in a truly in-context regime.

\subsection{Speech and Audio In-Context Learning}

While the previous subsection explored video understanding, which inherently incorporates an audio track in many real-world scenarios, this section isolates the specific challenges and advancements in applying in-context learning (ICL) to the speech and audio domains. This represents a critical frontier in building truly general-purpose multimodal systems, moving beyond the predominantly visual and textual focus to encompass the full spectrum of human sensory experience. Unlike vision and text, which have a longer history of joint modeling, the integration of audio into ICL frameworks is a nascent but rapidly evolving field. This section surveys the incorporation of audio modalities into ICL, focusing on the challenges and methodologies for enabling models to learn from acoustic demonstrations without parameter updates. The core challenge lies in effectively representing and fusing the temporal, spectral, and semantic complexities of audio signals with other modalities within a unified prompt structure, a challenge that echoes the temporal modeling concerns of the video domain but is specific to one-dimensional, non-visual signals.

Early efforts in this direction were not framed as in-context learning but laid the essential groundwork by learning joint audio-text or audio-visual representations. Pioneering models like CLAP \cite{ref281} and audio-language pretraining frameworks \cite{ref282} extended the paradigm of contrastive learning, successful in vision-language models like CLIP, to the audio domain. CLAP, for instance, uses two encoders and a contrastive learning objective to bring audio and text descriptions into a shared multimodal space, enabling zero-shot audio classification. However, these models are typically used in a zero-shot transfer manner, evaluating on pre-defined categories without the capacity to adapt to new tasks by conditioning on a sequence of audio-text demonstration pairs in the same way a large language model (LLM) uses text demonstrations. The transition to true in-context learning for audio involves enabling a model to infer a task from a prompt containing several \textless audio, query, answer\textgreater{} tuples, a more complex endeavor that bridges the gap between static representation learning and dynamic, demonstration-driven adaptation.

A significant body of research has focused on developing unified multimodal models that can accommodate audio alongside vision and text, which serves as the architectural foundation for audio ICL. These models aim to create a shared representational space where in-context demonstrations can bridge modalities, directly supporting the cross-modal transfer paradigms discussed in the following subsection. The i-Code framework \cite{ref283} is a prominent example, integrating vision, speech, and language into a single representation space. By employing pretrained single-modality encoders and a novel multimodal fusion network with cross-modality contrastive learning, i-Code can dynamically process single, dual, and triple-modality data. This flexibility is a prerequisite for multimodal ICL, as prompts may contain a mix of speech instructions, visual demonstrations, and textual answers. Similarly, the VATLM model \cite{ref284} employs a unified backbone network to model modality-independent information, using a masked prediction task on unified tokens to integrate visual, speech, and text inputs. This unified tokenization approach is crucial for ICL, as it allows a single decoder to process interleaved sequences of different modalities, a key mechanism demonstrated in unified visual ICL frameworks \cite{ref285} and a foundational principle for the omni-modal representation spaces explored next.

The specific challenge of building vision-language models (VLMs) that can effectively handle audio has been directly addressed by extending architectures like CLIP. The CLIP4VLA model \cite{ref286} adapts CLIP to the vision-language-audio triad by applying inter-modal and intra-modal contrastive learning. A key innovation is the design of an audio type token to dynamically learn different audio information types, distinguishing between verbal (speech) and nonverbal (environmental sounds, music) cues. This distinction is vital for ICL, as a model must recognize whether a speech demonstration is providing a verbal instruction or whether a non-verbal audio sample is serving as an example of a sound event to be classified. The VALOR model \cite{ref287} further advances this omni-perception capability by jointly modeling the relationships between vision, audio, and language using separate encoders and a multimodal conditional text generation decoder. Pretrained on a large-scale dataset of audiable videos, VALOR can answer questions conditioned on vision, audio, or both, a capability directly applicable to prompting with multimodal demonstrations. The cognition-inspired CoAVT model \cite{ref288} employs a joint audio-visual encoder for non-verbal information and a separate text encoder for verbal information, bridged by a query encoder that extracts the most informative audiovisual features for a given text. This separation of verbal and non-verbal processing streams is a crucial inductive bias for managing the heterogeneity of audio signals in an ICL context, a design principle that resonates with the mixture-of-experts strategies for decoupling complex modal interactions in unified architectures.

A particularly promising direction for speech ICL is the development of unified mixed-modal pretraining frameworks that can perform zero-shot transfer to unlabeled modalities, an ability that is essentially a form of in-context modality adaptation. The u-HuBERT model \cite{ref289} demonstrates this by using a unified masked cluster prediction objective and modality dropout during pretraining. A single fine-tuned model can then achieve competitive performance on audio-visual, audio-only, and even visual-only speech recognition, effectively learning to leverage the available context, whether auditory or visual, at inference time. This points toward a future where a speech model prompted with a few visual lip-reading demonstrations could perform silent speech recognition in-context. The Fine-grained Audio-Visual Joint Representations (FAVOR) framework \cite{ref290} extends a text-based LLM to perceive speech and audio events at the frame level, using a causal Q-Former to fuse the streams and align them with the LLM's input space. This architecture enables a form of ICL where the LLM can reason over fine-grained temporal interactions within multimodal prompts, directly extending the temporal reasoning capabilities discussed in the video ICL section to the audio domain.

Despite this progress, the application of explicit in-context learning---where a prompt consists of multiple independent demonstration examples---remains underexplored for audio. Many current models are evaluated on their ability to answer questions about an audio clip (a single-context, zero-shot task) rather than learning a mapping from several audio-to-text pairs. This gap is analogous to the early stages of vision-language ICL before the development of structured multi-modal demonstration formats. Key challenges include the design of effective audio-text demonstrations and the scarcity of suitable data. Studies on vision-language ICL have revealed that textual information in demonstrations often dominates over visual information \cite{ref221}, and a similar analysis for audio is needed to understand whether acoustic or linguistic cues drive in-context task inference. Furthermore, the robustness of audio-text models to complex natural language concepts is still limited, with state-of-the-art models struggling to fully leverage sequential or concurrent ordering of sound events from text descriptions \cite{ref291}. This suggests that constructing demonstrations that precisely convey an acoustic task through natural language remains a significant obstacle.

The path forward involves developing more sophisticated multimodal ICL benchmarks that include diverse speech and audio tasks, such as speaker identification, emotion recognition, and sound event detection, formatted as few-shot prompts. Just as the video domain has advanced through structured prompt designs like the Four-tiered Prompts framework, the audio domain requires analogous innovations for acoustic concept specification. Efforts like the Multi-modal In-Context Learning (MIC) dataset, which enhances a VLM's ability to understand complex multi-modal prompts \cite{ref14}, provide a blueprint that can be extended to the audio domain. Ultimately, achieving a paradigm of ``in-context listening'' will require models that can be lightly adapted with parameter-efficient tuning modules to handle audio streams and then conduct rapid, demonstration-driven adaptation for entirely new acoustic concepts. This vision aligns with the broader unification paradigms discussed in the next subsection, where a single model with a shared core and modality-specific adapters seamlessly processes text, images, video, audio, and other modalities presented in its context, moving beyond zero-shot classification to a true learning-from-demonstration regime in the audio modality.

\subsection{Cross-Modal Transfer and Unification Paradigms}

Building on the audio-specific challenges and architectural foundations discussed in the previous subsection, this section broadens the perspective to the overarching paradigms that enable in-context learning (ICL) to transfer across the full spectrum of modalities. While text-based ICL benefits from a relatively unified token space, vision, speech, and 3D data reside in fundamentally different geometric and statistical manifolds, making cross-modal generalization a fundamental obstacle for the unified architectures evaluated in the following subsection. This subsection investigates three key strategies for transferring ICL capabilities across these diverse modalities: omni-modal representation learning, unified architectural encoding, and multi-modal mixture-of-experts for decoupling complex interactions---each providing a pathway toward the modality-agnostic reasoning that next-generation benchmarks demand.

One promising direction for cross-modal transfer is the development of omni-modal representation spaces that can encode any modality into a shared, language-aligned embedding. Foundations for this are laid by frameworks like i-Code, which introduces an integrative and composable multimodal learning framework where vision, speech, and language are processed by pretrained single-modality encoders and then fused into a unified vector representation \cite{ref283}. This approach enables the flexible combination of modalities during both training and inference, a crucial property for ICL where demonstrations may span an arbitrary subset of available senses---directly supporting the flexible audio-visual-text prompt formats explored in the previous subsection. Building on this, the ViT-Lens series proposes a more radical unification: rather than designing separate encoders, it repurposes a powerful pretrained Vision Transformer (ViT) as a universal backbone. A modality-specific ``lens'' is tuned to project any-modal signals, such as 3D point clouds, depth, audio, tactile, and EEG data, into an intermediate embedding space digestible by the frozen ViT \cite{ref292}. The ViT, carrying rich pre-trained visual knowledge, then processes these signals, and the resulting representations are optimized to align with a pre-defined, modal-independent space from off-the-shelf foundation models. This paradigm is particularly powerful for cross-modal ICL transfer because it leverages the emergent reasoning capabilities of a single, shared backbone. A well-trained lens for a novel modality immediately inherits the in-context behavior of the ViT model, enabling zero-shot understanding and generation tasks for that modality without any further adaptation, as demonstrated by integrating a 3D lens into InstructBLIP for zero-shot 3D question answering \cite{ref258}. This zero-shot transfer capability is precisely the kind of cross-modal generalization that diagnostic benchmarks are designed to measure.

The construction of a unified representation space is further advanced by methods that connect pre-existing, independently trained, multi-modal contrastive representations (MCRs). The C-MCR framework proposes to project two MCRs, pre-trained on (A, B) and (B, C) modality pairs, into a new space and align them using data from the overlapping modality B \cite{ref293}. This alignment is then transferred to the non-overlapping pair (A, C), effectively creating a unified space for modalities A, B, and C without any paired (A, C) data. Ex-MCR extends this principle to flexibly learn a unified contrastive space for more than three modalities by integrating multiple existing MCRs into a single ``based'' MCR, such as aligning CLAP (audio-text) and ULIP (3D-vision) into CLIP (vision-text) \cite{ref294}. The resulting unified representations enable emergent semantic alignment between modalities that were never paired during training, like audio and 3D, and achieve state-of-the-art performance on cross-modal retrieval tasks. This line of work is critical for multimodal ICL because it provides a pathway to scale the number of modalities a model can reason over in-context without the prohibitive cost of collecting all pairwise combinations of supervised data, directly addressing the data scarcity challenges that benchmarking efforts like VL-ICL Bench seek to quantify.

A second major paradigm involves designing unified architectures that can natively handle diverse inputs, moving beyond the post-hoc alignment of modality-specific encoders---a theme that echoes the unified mixed-modal pretraining frameworks for speech discussed previously. A key insight from the Modality-Shared Contrastive Language-Image Pre-training (MS-CLIP) framework is that a mostly unified transformer encoder for vision and language can outperform separate encoders while reducing parameters \cite{ref295}. This suggests that sharing parameters encourages the encoding of semantic concepts from different modalities closer together in the embedding space, facilitating the transfer of common semantic structures like attention patterns. Such architectural weight sharing is foundational for building a single model that can perform ICL across modalities by applying the same learned transformation rules to differently-sourced tokens---a prerequisite for the omni-modal reasoning that benchmarks like MULTI and Blink evaluate. The challenge of efficiently scaling to a high number of modalities (a ``high-modality'' regime) is addressed by HighMMT, which introduces information-theoretic metrics for modality and interaction heterogeneity \cite{ref227}. By quantifying how similar two modalities are and how similarly pairs of modalities interact, HighMMT can automatically prioritize parameter sharing and fusion decisions, resulting in a single transformer model that scales to 10 modalities and 15 tasks. This principled approach to architecture design ensures that performance continues to improve with each added modality and that the model can transfer to entirely new modalities and tasks during fine-tuning, a crucial step toward general-purpose multimodal ICL systems that can be rigorously evaluated across the diverse task spectrum outlined in the following subsection.

The third paradigm focuses on decoupling the complex interactions between modalities, a necessity for robust ICL when the model must selectively attend to and reason over different modalities in context---a capability whose absence is precisely what vision-language-consistency analyses and diagnostic benchmarks expose. Mixture-of-Experts (MoE) architectures provide a natural framework for this, as they allow for specialized sub-networks to handle different input types or interaction patterns. In the context of multimodal generative models, the Variational Mixture-of-Experts Autoencoder (MMVAE) demonstrates how a mixture-of-experts prior can successfully learn a generative model that decomposes the latent space into shared and private subspaces across modalities, enabling coherent joint and cross-generation \cite{ref296}. For large multi-modal models, the prohibitive cost of replicating experts is mitigated by Omni-SMoLA, which uses a Soft Mixture of Low-rank Experts to residually learn specialized knowledge, either per-modality or multimodally, on top of a foundational backbone \cite{ref297}. This architecture allows the model to efficiently adapt to a vast collection of tasks without performance degradation, a critical capability for a system that must learn to follow new in-context instructions across modalities while maintaining the robust cross-modal consistency that evaluation frameworks demand. The principle of alternating or decoupled optimization is also explored in Multimodal Learning with Alternating Unimodal Adaptation (MLA), which reframes joint multimodal learning as an alternating unimodal learning process to minimize interference between modalities, while cross-modal interactions are gradually captured through a shared head \cite{ref222}. This approach resonates with the needs of ICL, where the model must be able to switch its focus between different modalities demonstrated in the context without one dominating the learning signal, directly addressing the modality dominance issues identified in the benchmarking literature.

Collectively, these paradigms are converging toward a unified, modality-agnostic ICL framework---the very system that the rigorous evaluation protocols and diagnostic benchmarks described next are designed to validate. The vision is a system where a single model parameterized by a shared core and lightweight, modality-specific adapters can seamlessly process text, images, audio, 3D, and structured data presented in its context. The model would leverage a pre-aligned omni-modal representation space to understand cross-modal correspondences, use a unified architecture to apply consistent reasoning patterns across data types, and deploy a mixture-of-experts mechanism to decouple and specialize its processing for complex multi-modal interactions. This unification is a prerequisite for the next generation of agents that can be prompted with a few demonstrations involving any combination of senses to perform novel tasks, from understanding a 3D scene through verbal description and visual sketches to generating a video from acoustic cues and a text script---capabilities that the comprehensive benchmarking suites discussed in the following subsection are explicitly designed to measure and stress-test.

\subsection{Benchmarks, Evaluation, and Analysis of Multimodal ICL}

Building on the unified architectures and representational paradigms discussed in the previous subsection, the rigorous evaluation of multimodal in-context learning (ICL) capabilities is critical for measuring true progress toward modality-agnostic reasoning. However, evaluation presents unique challenges that extend beyond those found in text-only or traditional vision-language benchmarks. Foundational evaluations of vision-language models (VLMs) have historically centered on static, single-task benchmarks such as VQAv2, COCO Caption, and NLVR2. While these datasets have been instrumental, their standard evaluation protocols often fail to capture the dynamic, example-driven nature of ICL that the architectures above are designed to enable. Recent research has highlighted that conventional few-shot visual question answering (VQA) and image captioning evaluations, which have been the predominant focus for multimodal ICL investigations, neither fully exploit the strengths of ICL nor rigorously test its limitations \cite{ref13}. This has motivated the community to develop more comprehensive and diagnostic benchmark suites that systematically probe the multifaceted abilities of models when learning from multimodal context.

A primary limitation of using standard VQA and captioning benchmarks for ICL evaluation is their insufficient scope. They often fail to disentangle a model's ability to leverage textual versus visual information within the provided demonstrations---a capability that directly depends on the quality of the cross-modal alignment and unified representations discussed previously. For instance, research has shown that in many multimodal ICL scenarios, performance is predominantly driven by the textual content of the demonstrations, with the visual information having a negligible effect \cite{ref221}. Therefore, benchmarks are needed that can isolate and quantify the contribution of each modality, revealing potential spurious successes where a model performs well by ignoring the visual context altogether, thus failing to exercise the omni-modal reasoning capabilities that unified architectures aim to provide.

To address these shortcomings, several comprehensive benchmark suites have been introduced that span a wide spectrum of tasks, from low-level perception to high-level reasoning. \textbf{VL-ICL Bench} is a significant contribution that encompasses a broad range of tasks involving both images and text as inputs and outputs, explicitly designed to test different types of challenges, including perception, reasoning, and handling long context lengths \cite{ref13}. This benchmark reveals that even the most advanced models, such as GPT-4V, find these structured ICL tasks challenging, highlighting significant room for improvement in the unified ICL paradigm. Similarly, \textbf{MULTI} (Multimodal Understanding Leaderboard with Text and Images) serves as a cutting-edge benchmark for evaluating models on understanding complex tables and images and reasoning with long context, mimicking realistic examination styles with over 18,000 questions \cite{ref298}. The inclusion of a hard subset, MULTI-Elite, and an external knowledge component, MULTI-Extend, allows for fine-grained analysis of model capabilities in knowledge-intensive and complex multimodal ICL settings, directly testing the generalization capacity of cross-modal transfer mechanisms.

Another crucial direction in evaluation is the assessment of a model's foundational visual perception abilities that are often taken for granted in higher-level VLM benchmarks. The \textbf{Blink} benchmark reformulates 14 classic computer vision tasks into 3,807 multiple-choice questions paired with visual prompting, focusing on core visual perception skills like relative depth estimation, visual correspondence, and multi-view reasoning that can be solved by humans ``within a blink'' \cite{ref299}. The finding that even top models like GPT-4V and Gemini achieve only slightly above random guessing on many Blink tasks underscores that robust visual perception has not yet fully ``emerged'' in current multimodal ICL frameworks, despite the advances in unifying backbones like ViT-Lens. This is consistent with findings from low-level vision benchmarks. \textbf{Q-Bench} and its subsequent extension for image pairs systematically evaluate MLLMs on low-level visual perception, description, and quality assessment, revealing that while MLLMs possess preliminary low-level visual skills, these are often unstable and imprecise \cite{ref300,ref301}. These benchmarks indicate that ICL's effectiveness is bottlenecked by the model's intrinsic visual processing capabilities, an area often overshadowed by linguistic reasoning prowess and one that must be addressed for the unified ICL vision to be realized.

A critical aspect of multimodal ICL analysis involves diagnosing the internal consistency and cross-modal interaction gaps within models---failures that mixture-of-experts and unified architectures explicitly aim to resolve. The framework introduced for vision-language-consistency (VLC) analysis systematically quantifies the capability disparities between a model's vision and language modules in multimodal settings \cite{ref244}. This work reveals that while models like GPT-4V perform consistently across modalities on simple tasks, the trustworthiness of the vision modality significantly degrades as task complexity increases. This ``lost in translation'' phenomenon highlights a fundamental failure mode in multimodal ICL where the fusion of information is not robust. Building on these findings, methods like ``Vision Description Prompting'' have been proposed to mitigate this gap by forcing explicit alignment \cite{ref244}.

Further specialized evaluations probe other dimensions of multimodal ICL relevant to the decoupled reasoning capabilities that MoE architectures target. \textbf{MultipanelVQA} challenges models to comprehend multi-panel images, a common real-world format, using synthetically generated data to isolate the effects of layout and subfigure composition on reasoning \cite{ref302}. \textbf{IllusionVQA} tests models on inherently unreasonable or ambiguous visual stimuli like optical illusions, discovering a potential weakness in ICL: models can fail to locate optical illusions even when the correct answer is provided as a few-shot example in the context window \cite{ref303}. This challenges the assumption that ICL always benefits from more examples and points to a brittleness in visual reasoning under adversarial or non-standard perceptual conditions that unified models must overcome.

The evaluation protocol itself must also be robust to align with the flexible input-output paradigms of unified architectures. Many existing benchmarks rely on task-oriented input-output formats that are ill-suited for the free-form text generation of modern LVLMs, making automatic assessment challenging. \textbf{ReForm-Eval} addresses this by re-formulating existing benchmarks into unified LVLM-compatible formats, enabling comprehensive and quantitative evaluation without requiring manual evaluation of free-form text for every new benchmark \cite{ref304}. Furthermore, evaluation must move beyond accuracy alone. The integration of uncertainty quantification into benchmarking, using methods like conformal prediction, reveals that the models with the highest accuracy may also exhibit the highest uncertainty, and this uncertainty is often correlated with the language model component rather than the visual one \cite{ref305}. Finally, the critical issue of out-of-distribution (OOD) generalization is paramount. Cross-dataset evaluations for VQA have shown that pretrained V\&L models often learn to solve the specific benchmark rather than acquiring the high-level ICL skills required for the task, with automatic metrics often penalizing correct responses due to rigid n-gram matching \cite{ref306}. These findings collectively call for a multi-faceted evaluation approach that considers accuracy, modality-specific contribution, internal consistency, perceptual robustness, and calibration under distribution shift to truly gauge the progress and reliability of multimodal ICL as it moves toward the unified, modality-agnostic paradigm envisioned by the community.

\section{In-Context Learning for Sequential Decision-Making and Reinforcement Learning}

\subsection{Foundations of Sequential Decision-Making with In-Context Learning}

The intersection of in-context learning (ICL) and sequential decision-making represents a paradigm shift in how we conceptualize adaptive behavior in reinforcement learning (RL). At its core, in-context learning for sequential decision-making frames adaptation not as an iterative process of gradient-based optimization, but as a direct inference capability emerging from the model's ability to condition on prior experience presented within its context window. This subsection establishes the formal foundations of this approach, contrasting it with classical meta-reinforcement learning (meta-RL), and sets the stage for the subsequent exploration of contextual bandits as a concrete instantiation of these principles.

\subsubsection{Formal Problem Setting}

In the classical meta-RL formulation, an agent is trained over a distribution of tasks \(p(\mathcal{T})\), where each task \(\mathcal{T}_i\) is a Markov Decision Process (MDP) defined by the tuple \((\mathcal{S}, \mathcal{A}, P_i, R_i, \gamma)\) with state space \(\mathcal{S}\), action space \(\mathcal{A}\), task-specific transition dynamics \(P_i(s'|s,a)\), reward function \(R_i(s,a)\), and discount factor \(\gamma\). Traditional meta-RL algorithms, such as MAML or RL\(^2\), require gradient updates---either at meta-training time or during adaptation---to modify the agent's policy for a new task. In stark contrast, in-context decision-making reformulates adaptation as a sequence modeling problem where the policy \(\pi_\theta(a_t | \tau_{:t-1}, \mathcal{D}_{\text{ctx}})\) conditions on both the current trajectory prefix \(\tau_{:t-1}\) and a set of demonstration trajectories \(\mathcal{D}_{\text{ctx}}\) that implicitly specify the task \cite{ref307}.

The input prompt structure in this setting comprises a concatenation of demonstration episodes: multiple sequences of states, actions, and rewards from one or more tasks, followed by the current episode's history. The model must infer the underlying task structure---its dynamics, reward function, or desired behavior---directly from this contextual information, generating appropriate actions for the query state without any parameter updates. This formulation fundamentally treats RL as a conditional generative process, where the context serves as a non-parametric specification of the task \cite{ref308}. This perspective directly prefigures the contextual bandit formulation discussed in the next subsection, where the demonstration history is simplified to a sequence of context-action-reward tuples, and the model must infer the reward function to guide exploration.

\subsubsection{In-Context Learning as Meta-Learning Without Gradients}

The emergence of in-context decision-making capabilities represents a form of meta-learning where the ``learning algorithm'' is implemented entirely through the forward pass of a sequence model. This stands in marked contrast to classical meta-RL approaches that explicitly optimize for rapid gradient-based adaptation. Works such as \cite{ref309} highlight the limitations of purely sequence-based meta-RL methods like RL\(^2\), which struggle with asymptotic performance and out-of-distribution generalization compared to algorithms that incorporate explicit value function components. However, when scaled to sufficiently large models and trained on diverse task distributions, in-context adaptation can exhibit properties that rival or exceed gradient-based methods, as demonstrated by the AdA agent which achieved human-timescale adaptation across a vast space of embodied 3D problems \cite{ref308}.

The key insight is that transformers and other attention-based architectures can internally implement optimization algorithms through their forward pass. Rather than executing explicit gradient steps on a loss function, the model's attention mechanism dynamically re-weights prior experiences to form predictions---effectively performing a form of implicit meta-gradient computation. This perspective connects to theoretical frameworks viewing ICL as implicit gradient descent or Bayesian inference, which we explore in depth in Section 2. For decision-making, this means that the model can distill patterns from demonstration trajectories to compose novel behaviors, effectively learning from the context without changing its weights \cite{ref307}. As we will see in the contextual bandit setting, this capacity for implicit inference allows a model to approximate sophisticated exploration strategies like Thompson sampling or Upper Confidence Bound purely through in-context conditioning, without explicit algorithmic specification.

\subsubsection{Contrasts with Classical Meta-RL}

The divergence between in-context decision-making and classical meta-RL manifests across several critical dimensions. First, in terms of the \textbf{adaptation mechanism}, classical meta-RL algorithms explicitly separate meta-parameters (which encode prior knowledge shared across tasks) from task-specific parameters (which are updated during adaptation). Context-based meta-RL methods, such as those using learned context encoders to infer task representations from collected trajectories, occupy a middle ground by compressing experience into a latent context variable that conditions policy behavior \cite{ref310,ref311}. However, these methods still require training specialized context encoders and often rely on recurrent aggregation of transitions. In-context decision-making eliminates even this specialized machinery, relying entirely on the model's pretrained attention mechanisms to extract task-relevant patterns from raw sequences of interactions \cite{ref312}. This distinction is particularly salient in bandit settings, where the ``context encoder'' is simply the transformer's self-attention over the history of interactions, and exploration strategies like those learned by MELEE \cite{ref313} emerge naturally from the model architecture.

Second, the \textbf{nature of the training distribution} differs fundamentally. Classical meta-RL typically assumes parametric task distributions where tasks vary along well-defined axes (e.g., goal positions, friction coefficients). In-context methods, particularly those trained through algorithm distillation or meta-training over large, diverse environments, can handle broader and non-parametric task distributions. The TIGR algorithm \cite{ref314} and MoSS framework \cite{ref311} demonstrated that Gaussian mixture models and self-supervised task representations could capture multi-modal, non-parametric task variations---but these still rely on specialized inference modules. In-context approaches go further by letting the sequence model itself serve as the inference mechanism, potentially encompassing an even wider range of task specifications. This is precisely the capability exploited in contextual bandit meta-learning, where training on a distribution of synthetic bandit tasks enables zero-shot adaptation to novel bandit instances at test time.

Third, \textbf{sample efficiency and computational cost} present a complex trade-off. Classical meta-RL incurs computational cost during adaptation through gradient updates, but can be efficient in terms of interaction samples if meta-training is well-designed. In-context methods shift this burden to the pretraining phase and the forward-pass cost at inference, which can be substantial for long contexts but requires no gradient computation during adaptation. The Guided Meta-Policy Search approach \cite{ref315} attempted to decouple meta-training from on-policy data collection, and Concurrent Meta-RL \cite{ref316} introduced parallel execution to address temporal credit assignment---both highlighting the sample complexity challenges that in-context methods must also confront. In the bandit domain, this trade-off is crystallized: the quality of in-context exploration depends critically on the informativeness of the trajectories presented in the context, motivating techniques like contrastive learning for context construction that we will discuss.

\subsubsection{The Role of Demonstration Trajectories}

A defining characteristic of in-context decision-making is the use of whole demonstration trajectories as the medium for task specification. This draws directly from imitation learning and learning-from-demonstration paradigms, but with a crucial distinction: the model is not trained to mimic the demonstrations via behavior cloning, but rather to \emph{extract the underlying intention, constraints, or optimality principles} that generated them. The EMRLD framework \cite{ref317} showed how even sub-optimal demonstration data could be leveraged to improve meta-training, providing guidance in sparse reward settings. In-context approaches generalize this insight by allowing the model to condition on arbitrary quality demonstrations at inference time, inferring task structure without requiring the demonstrations to be expert-level. This principle directly underpins the success of in-context exploration in bandits, where the model conditions on a history of interactions to infer the reward function and adapt its action-selection strategy---performing exploration via inductive inference from the provided exemplars, as elaborated in the following subsection.

This connects to the broader theme of algorithm distillation, where models learn to implement entire learning algorithms through exposure to the training histories of RL agents. The In-Context Policy Iteration (ICPI) algorithm \cite{ref307} explicitly demonstrated that a frozen language model, without any gradient updates or domain-specific pretraining, could perform policy iteration entirely through prompt engineering---updating its policy via trial-and-error interactions stored in the context. This blurs the boundary between ``learning'' and ``inference,'' suggesting that with sufficiently capable sequence models, the distinction collapses into a single conditional generation process. In the contextual bandit setting, this distillation manifests as the model's ability to implicitly perform model selection over a prior distribution of bandit tasks, effectively meta-learning an adaptive posterior inference procedure through in-context conditioning alone.

\subsubsection{Scope and Organization}

This section surveys the burgeoning landscape of in-context decision-making, organized to reflect the spectrum from theoretical underpinnings to practical applications. We begin with contextual bandits (Section 7.2), the simplest sequential decision problem, where in-context exploration strategies condition on past contexts and rewards to make decisions---a framework that crystallizes the core ICL principles of task inference and adaptive exploration without gradient updates. We then examine offline meta-RL and context-based adaptation (Section 7.3), where context encoders learn to infer task representations from pre-collected datasets---addressing the critical challenge of context distribution shift between training and testing. Representation learning for context (Section 7.4) explores how contrastive, mutual information, and self-supervised objectives can produce compact and sufficient task embeddings \cite{ref318,ref312}.

Subsequent sections broaden the scope to imitation learning with in-context demonstrations (Section 7.5), algorithm distillation and sequence modeling for meta-RL (Section 7.6), curriculum learning strategies that structure the presentation of tasks (Section 7.7) \cite{ref319,ref320}, skill acquisition and hierarchical decomposition (Section 7.8) \cite{ref321}, sim-to-real transfer and domain adaptation (Section 7.9) \cite{ref322}, and finally generalization challenges and open problems (Section 7.10) \cite{ref323,ref324}.

Throughout these subsections, we maintain a focus on the interplay between in-context mechanisms and the unique demands of sequential decision-making: the need for exploration, credit assignment over long horizons, robustness to distribution shift, and the composition of skills across temporally extended tasks. The central thesis is that in-context learning represents not merely an alternative to gradient-based meta-RL, but a fundamentally different approach to building adaptive agents---one where learning is indistinguishable from inference, and where the boundaries between pretraining, fine-tuning, and deployment dissolve into a unified framework of context-conditioned sequential generation. The forthcoming subsection on contextual bandits grounds this thesis in the simplest possible sequential setting, demonstrating how the core principles of in-context exploration, task inference, and algorithm distillation operate in a tractable yet richly expressive domain.

\subsection{Contextual Bandits and In-Context Exploration}

Building directly on the formal foundations established in Section 7.1---where we characterized in-context decision-making as meta-learning without gradients through forward-pass conditioning on demonstration trajectories---contextual bandits provide the simplest yet most illuminating instantiation of these principles. In a standard contextual bandit, an agent observes a context, selects an action, and receives a reward, all while seeking to minimize cumulative regret. In-context learning, in this framing, allows a model---typically a large language model or a transformer trained via meta-learning---to condition on a history of past interactions (contexts, actions, and rewards) presented as a prompt. This enables the model to infer the underlying reward structure and adapt its action-selection strategy on the fly, essentially performing exploration via inductive inference from the provided exemplars. This directly mirrors the core mechanism described previously: adaptation through conditioning on a context window of prior experience, without any parameter updates.

A core insight is that in-context exploration can be understood as a form of meta-learned exploration strategy, crystallizing the meta-learning-without-gradients thesis from the foundations subsection. The work of \cite{ref313} introduced MELEE, a meta-learning algorithm explicitly designed to learn an exploration policy for interactive contextual bandits. Rather than relying on hand-crafted exploration bonuses like UCB or Thompson sampling, MELEE learns to explore by imitating optimal behavior on synthetically generated bandit tasks. This imitation learning-based approach to crafting an exploration policy is inherently compatible with in-context learning: the model learns a mapping from a history of contexts and rewards to an action-selection strategy, which it can then apply to novel bandit problems at test time without parameter updates---precisely the in-context adaptation paradigm. The effectiveness of this approach is underscored by its superior performance on hundreds of real-world datasets, particularly when the reward gap between actions is large, highlighting how a meta-learned, context-conditioned strategy can outperform traditional analytic exploration heuristics.

The ability to generalize in-context to new bandit tasks is significantly amplified by the quality of the representations used, directly connecting to the demonstration quality principles discussed in Section 7.1. \cite{ref313} further demonstrated that a rich feature representation is crucial for learning how to explore. This aligns with the broader ICL paradigm, where the model's capacity to identify the task from demonstrations hinges on a representational space that captures the latent structure of the problem domain---the same principle that underlies the use of demonstration trajectories for task specification in the general sequential setting. In the bandit setting, this means that the encoded contexts and their historical associations with rewards must form a representation from which the model can distinguish between different reward functions. More formally, this connects to the Bayesian perspective on ICL, where a model implicitly performs model selection over a prior distribution of bandit tasks. The meta-learning process, as in \cite{ref325} and \cite{ref326}, shapes this prior by exposing the learner to a distribution of bandit instances. \cite{ref325} learns a prior over bandit models and uses Thompson sampling to explore efficiently, effectively meta-learning an adaptive posterior inference procedure. This posterior inference is precisely what ICL replicates: by conditioning on a sequence of demonstrated interactions, the model updates its belief about which bandit instance it is facing and acts accordingly, directly instantiating the implicit Bayesian inference framework introduced earlier.

Several works have formalized how a pre-trained model can be used to approximate an optimal Q-function or exploration policy in a zero-shot manner through in-context conditioning, extending the algorithm distillation concept from Section 7.1 to the bandit domain. \cite{ref327} proposed a meta-learning pipeline to approximate an optimal Q-function for contextual bandits, using a synthetic training environment and a scheduled exploration policy. While this work focused on meta-learning an explicit parameterized policy, the natural extension is to cast this Q-function as an in-context learner that infers action values directly from the prompt. Similarly, \cite{ref326} developed a framework for optimizing differentiable bandit policies using policy gradients, where the policy can be designed to condition on the entire history of interactions, a direct antecedent to the prompting format used in ICL. The resulting policies can capture complex, history-dependent exploration strategies that are far more flexible than standard myopic heuristics, embodying the transition from gradient-based adaptation to inference-time conditioning that defines the in-context approach.

The interplay between exploration during meta-training and in-context adaptation at test time raises a crucial challenge---the distribution of exploration trajectories---which directly anticipates the context distribution shift problem that will dominate our discussion of offline meta-RL in Section 7.3. In-context learning for bandits relies on the model having seen informative trajectories during its pre-training or meta-training phase. \cite{ref310} addresses this by proposing a contrastive learning objective to train a compact and sufficient context encoder, and, critically, a separate exploration policy optimized for maximizing information gain. This can be seen as a method for generating the most informative contexts to condition on during in-context adaptation. The theoretical derivation of an information-gain-based objective for trajectory collection ensures that the in-context prompt contains the maximal possible signal about the task identity and reward structure, enabling more sample-efficient zero-shot adaptation. This is essential, as the quality of ICL is directly bounded by the informativeness of the demonstrations---providing a model with uninformative or redundant history is a known failure mode that foreshadows the representation robustness challenges in the offline setting.

Moreover, the success of in-context exploration in bandits is tied to the model's ability to implicitly perform algorithm distillation, a phenomenon highlighted in Section 7.1 through the ICPI algorithm's demonstration that frozen language models can implement policy iteration via prompt engineering. The work \cite{ref328} showed that meta-trained LSTMs learn in ways that are congruent with Bayes-optimal inference, uncovering task structure concurrently rather than in the staggered manner of standard deep learning algorithms. When such a meta-trained model is prompted in-context, it effectively distills the exploration and inference algorithm it was exposed to during meta-training. This means an ICL system can mimic sophisticated strategies like Thompson sampling or Upper Confidence Bound (UCB), or even more complex adaptive exploration tactics that balance immediate reward with long-term information gain, all without explicitly programming these algorithms. This is powerfully demonstrated in \cite{ref329}, where a language model recursively improves its own in-context learning strategy simply through in-context conditioning on a history of tasks. In a two-armed bandit domain, this meta-in-context learning allows the model to adaptively reshape its priors over expected tasks and modify its exploration-exploitation strategy, effectively learning to learn in context---a recursive form of the algorithm distillation process.

Practical applications of contextual bandits in recommendation and advertising often involve resource constraints, where every exploration action has a non-zero cost, adding a layer of complexity that tests the limits of in-context adaptation. \cite{ref330} introduced HATCH, which adaptively allocates exploration resources based on remaining budget and estimated reward distributions. While this is a learned, hierarchical policy, the principle of using contextual information to guide adaptive, budget-aware exploration is directly transferable to an ICL paradigm. A model prompted with a history that includes budget and cost signals could learn in-context to modulate its exploratory actions, becoming more exploitative as the budget is exhausted. Similarly, the challenge of learning from offline, logged bandit data with deficient support \cite{ref331} can be approached through ICL by providing the model with a prompt that includes not just successful trajectories but also counterfactual estimates or indicators of data deficiency, allowing it to infer a more robust policy---a theme that directly transitions to the offline meta-RL methods discussed next. The integration of causal inference with bandit learning, as explored in \cite{ref332}, where balancing methods are used to reduce estimation bias, could also be implicitly prompted to an ICL model by structuring the historical context to highlight distributional shift or bias, guiding it toward less biased in-context estimates. The promise of ICL in this domain lies in unifying the exploration strategy, task identification, and reward estimation into a single forward pass of a transformer, conditioned on a prompt that seamlessly integrates raw context data with the meta-structure of exploration trade-offs---a unification that prefigures the robust context-based adaptation frameworks that form the backbone of offline meta-RL, which we now examine.

\subsection{Offline Meta-Reinforcement Learning and Context-Based Adaptation}

Building directly on the principles of exploration and adaptation introduced in the contextual bandit setting, offline meta-reinforcement learning (OMRL) constitutes a critical frontier for applying in-context learning principles to real-world sequential decision-making, where online environment interaction is expensive, risky, or infeasible. The central objective of OMRL is to meta-train an agent on a static, pre-collected dataset of trajectories from multiple tasks, such that it can rapidly adapt its policy to unseen target tasks using only a small amount of context data, without any further gradient updates. This paradigm inherently mirrors the in-context learning process found in large language models, where the context---comprising trajectories of state-action-reward tuples---serves as the prompt from which the agent infers the latent task identity and desired behavior. A dominant approach for achieving this is context-based OMRL, which employs a context encoder to compress collected experience into a compact task representation, to which a task-conditioned meta-policy is then conditioned for adaptation \cite{ref333}.

Mirroring the challenge of distribution shift in context representations discussed in the previous section, a foundational problem that plagues context-based OMRL is the context distribution shift between the meta-training and meta-testing phases. During meta-training, the context encoder is trained on data generated by a mixture of behavior policies present in the offline dataset. At meta-test time, however, the agent must infer the task representation from trajectories collected by its own exploration policy, which may be substantially different from any behavior policy in the training corpus. This distributional mismatch causes the context encoder to produce unreliable or uninformative task representations, severely degrading the policy's adaptive performance. Several works have directly diagnosed and addressed this problem. The work on context distribution shift in task representation learning \cite{ref334} demonstrated empirically that a context encoder trained solely on offline data suffers a significant performance drop when deployed with a novel exploration policy, and proposed a hard-sampling strategy to train a more robust encoder by explicitly focusing on difficult-to-predict contexts during training. Similarly, the CSRO framework \cite{ref335} tackled this issue by designing a max-min mutual information objective to learn representations that are minimally dependent on the specific behavior policy, combined with a non-prior context collection strategy during meta-testing to reduce the effect of the exploration policy on the inferred context.

Extending the representation learning themes of invariant and sufficient encoders explored previously, robust task representation learning is paramount for effective context-based adaptation. A key desideratum is that the task representation should capture the causal dynamics and reward structure of the task, rather than spurious correlations with the behavior policy used to collect the offline data. The CORRO framework \cite{ref336} achieved this via a bi-level contrastive learning objective that maximizes mutual information between representations from the same task under different behavior policies while pushing apart representations from different tasks. This regularizes the encoder to be invariant to policy-induced distribution shifts. Relatedly, the issue of behavior policy confounding was addressed through adversarial data augmentation \cite{ref337}, which generates adversarial examples designed to fool the task encoder, forcing it to learn policy-invariant features that more accurately reflect the task itself. The FOCAL algorithm \cite{ref338} further contributed a stable distance metric learning approach on a bounded context embedding space, combined with standard offline RL behavior regularization to combat the bootstrapping errors that arise from out-of-distribution actions. To provably improve the robustness of these representations, subsequent work incorporated an intra-task attention mechanism and inter-task contrastive learning objectives, demonstrating theoretical and empirical gains in the face of sparse rewards and distribution shift \cite{ref339}.

More structured approaches to representation learning have emerged to handle limited data regimes and the need for generalizable prior experience, aligning with the broader trajectory-level and disentanglement objectives discussed in the context representation learning literature. The GENTLE algorithm \cite{ref340} was designed for scenarios with both limited training tasks and limited behavioral diversity. It uses a task auto-encoder (TAE) that reconstructs state transitions and rewards, capturing the generative structure of the task model itself and producing representations that generalize well even when training tasks are scarce, without relying on contrastive learning which can fail under limited coverage. For decoupling complex task structures, the DCMRL framework \cite{ref321} learns separate, clustered latent spaces for task contexts and reusable skills using Gaussian quantization variational autoencoders, enabling the agent to compose generalizable prior experience for rapid adaptation. The MoSS algorithm \cite{ref311} further extended meta-RL to broad non-parametric task distributions by modeling task representations with a Gaussian mixture model, enabling strong generalization and online adaptation in non-stationary environments.

A notable trend, directly building on the self-supervised and contrastive representation learning techniques surveyed in the prior section, is the use of self-supervised objectives to enhance the quality and sample efficiency of context representation learning specifically in the offline setting. The Trajectory Contrastive Learning (TCL) method \cite{ref318} proposed a simple yet effective auxiliary task where the encoder learns to predict whether two transition windows originate from the same trajectory, leveraging the natural temporal structure of experience to accelerate meta-training. The CCM framework \cite{ref310} coupled a contrastive context encoder with an information-gain-based exploration policy, explicitly collecting informative trajectories that reduce task uncertainty. The efficiency gains from such methods are substantial; the PEARL algorithm's formulation using probabilistic latent context variables already demonstrated that disentangling task inference and control through off-policy RL can achieve 20-100x improvements in sample efficiency over prior on-policy meta-RL methods \cite{ref333}, and these self-supervised extensions push this boundary further.

Addressing the broader challenge of offline data scarcity and coverage, hybrid methods have been proposed to bridge the offline-online gap. The hybrid approach \cite{ref341} uses the offline dataset for meta-training a base adaptive policy, but then collects additional, reward-free online data to bridge the distribution shift caused by the change in exploration strategy. Crucially, this online data requires no task-specific reward labels, drastically reducing the cost of data collection compared to full online meta-RL, while matching its performance on challenging locomotion and manipulation domains. To handle the fundamental tension between exploring via the meta-policy and exploiting the offline data, the MerPO algorithm \cite{ref342} learns a meta-model and a regularized meta-policy, using a novel model-based actor-critic that balances staying close to both the behavior policies and the meta-policy, theoretically guaranteeing policy improvement on new tasks. The broader theoretical community has also investigated the integration of distributionally robust learning principles with offline RL \cite{ref343}, providing a principled minimax framework for mitigating the shift and ensuring sample-efficient learning under partial data coverage, a formalization that closely aligns with the robustness goals of OMRL. These advances collectively underscore that robust context-based adaptation, achieved through careful representation learning and explicit distribution shift mitigation, is central to unlocking the practical potential of in-context learning for sequential decision-making, bridging the foundational bandit exploration strategies with the sophisticated representation learning techniques required for complex, dynamic environments.

\subsection{Representation Learning for Context in Meta-RL}

Extending the representation learning principles introduced in the contextual bandit and offline meta-RL discussions, this subsection delves into the core methodologies for acquiring robust and generalizable task context representations in Meta-Reinforcement Learning. In context-based Meta-RL, the effectiveness of adaptation critically hinges on the quality of the learned task context---a latent representation that encodes task-specific information from a history of interactions. A compact, sufficient, and robust context representation enables a policy to rapidly generalize to new tasks with minimal additional experience, directly addressing the distribution shift and robustness challenges highlighted in the previous subsection on offline meta-RL. This subsection surveys key advances in representation learning for context in Meta-RL, focusing on methodologies that leverage contrastive learning, mutual information optimization, and trajectory-level self-supervised objectives to improve sample efficiency and generalization.

A fundamental challenge in learning useful context representations, one that lies at the heart of the offline meta-RL distribution shift problem discussed earlier, is ensuring that the encoder captures all information relevant to the task while discarding spurious details from the behavior policy or environment stochasticity. Early approaches often relied on optimizing the context encoder solely through the RL objective, but this can be sample-inefficient and lead to representations that are heavily coupled to the specific policy used for data collection, causing the very distribution shift problems during meta-testing that methods like CSRO and CORRO aim to mitigate. To address this, a line of work has formalized context learning through the lens of information theory, providing a unifying perspective that connects robust representation learning with the generalization guarantees sought in offline meta-RL. The unified information-theoretic framework presented in \cite{ref344} demonstrates that many existing Context-based Offline Meta-RL (COMRL) algorithms are essentially optimizing a mutual information objective between the task variable and its latent representation via various approximate bounds. This insight is pivotal as it provides a principled foundation for designing representation learning objectives that naturally combat the policy-induced distribution shift identified in the previous subsection. Building on this, the algorithm UNICORN is proposed, which leverages the information bottleneck principle to learn task representations that are maximally informative about the task while being compressive about the observed trajectory, achieving state-of-the-art performance across diverse benchmarks.

The application of contrastive learning has emerged as a particularly powerful technique for shaping the geometry of the task representation space, building on the same principles that underpin the CORRO framework's invariance to behavior policies and the self-supervised methods in offline meta-RL. The core idea is to bring representations of trajectories or transition windows from the same task closer together while pushing apart those from different tasks. \cite{ref310} proposes the CCM framework, which directly addresses the dual questions of how to train a compact and sufficient context encoder and how to collect informative trajectories. CCM leverages the contrastive nature between different tasks to train the encoder, ensuring that similar task dynamics are mapped to nearby points in the latent space, while an information-gain-based objective guides an exploration policy to gather trajectories that maximize the distinctiveness of the context. Similarly, \cite{ref339} introduces a framework that combines an intra-task attention mechanism with an inter-task contrastive learning objective. The attention mechanism helps the context encoder focus on the most task-relevant segments within a trajectory, while the contrastive loss robustifies the learned representations against distribution shift and sparse rewards, with theoretical guarantees provided for the improved performance---directly extending the theoretical robustness results sought in the offline meta-RL literature.

Trajectory-level self-supervised learning objectives have been specifically designed to exploit the hierarchical structure inherent in meta-RL data, resonating with the trajectory-centric approaches used in offline meta-RL frameworks like GENTLE and MoSS. \cite{ref318} proposes Trajectory Contrastive Learning (TCL), a self-supervised auxiliary task that trains the context encoder to predict whether two transition windows originate from the same trajectory. TCL makes minimal assumptions about the task structure and can be seamlessly integrated into existing context-based meta-RL algorithms. By operating at the level of trajectory windows, TCL encourages the encoder to learn temporally coherent and task-consistent features, which accelerates the training of the context encoder and leads to substantial performance gains on both MuJoCo and Meta-World benchmarks. This trajectory-centric view is also explored in other self-supervised paradigms that aim to learn compositional and functionally-aware state representations, as seen in \cite{ref345}, which reinterprets mutual information maximization as a representation learning method for goal reaching, offering a principled way to acquire latent goals or skills that can serve as building blocks for context---a concept that directly feeds into the skill-based imitation learning architectures discussed in the following subsection.

The challenge of learning disentangled representations is particularly acute when task dynamics are influenced by multiple confounding factors, mirroring the decoupling objectives of DCMRL in the offline meta-RL setting. In real-world applications, changes in transition dynamics may stem from a combination of independent causal mechanisms. \cite{ref312} addresses this by introducing DOMINO, a framework that explicitly learns a disentangled context representation. DOMINO maximizes the mutual information between the context and historical trajectories while simultaneously minimizing a state transition prediction error. The theoretical analysis shows that learning disentangled representations overcomes the underestimation of mutual information caused by multi-confounded dynamics, leading to more sample-efficient and robust generalization. This pursuit of decoupling factors of variation is further echoed in \cite{ref321}, which proposes DCMRL. This framework uses contrastive learning to pull together representations from the same task and push apart those from different tasks, and then employs a Gaussian quantization variational autoencoder to cluster and decouple the continuous latent spaces of task contexts and skills. The resulting discrete, representative codebooks of task context and skills act as strong, generalizable priors that facilitate effective adaptation to unseen target tasks, directly paralleling the hierarchical and compositional representations that will be leveraged for imitation learning in the next subsection.

A persistent problem that connects the offline meta-RL discussion to the core of context representation learning is the context shift phenomenon, where the representation learned from a fixed offline dataset, collected by a behavior policy, fails to generalize to the data collected by the exploration policy during meta-testing. \cite{ref335} proposes CSRO to tackle this by minimizing the influence of the policy on the context representation. During meta-training, a max-min mutual information learning mechanism is designed to make the task representation invariant to the behavior policy's idiosyncrasies, while a non-prior context collection strategy is introduced for the meta-test phase to reduce the distributional mismatch. Complementary to this, a hard-sampling strategy is presented in \cite{ref334} to directly train a more robust context encoder by exposing it to challenging, under-represented parts of the offline data distribution. These approaches embody the same principles of invariance and robustness that are essential for the context-conditioned imitation learning methods discussed subsequently, where policies must generalize across diverse demonstration qualities without retraining.

Beyond these core methods, the notion of bootstrapping self-supervised learning for representation improvement has been explored through a meta-learning lens, offering a paradigm for continuous refinement that anticipates the algorithm distillation capabilities covered later. \cite{ref346} proposes a framework that formulates the learning of general visual representations from small data as a bi-level optimization problem, mimicking a human-like ability to learn from experience. By reconstructing tasks to combine the strengths of meta-learning and self-supervised learning, and employing a bootstrapped target based on meta-gradients, the model becomes its own teacher, iteratively refining its representations. This concept of self-distillation and using the model's own knowledge to guide representation learning is a powerful paradigm for improving the compactness and generalizability of context encoders, bridging the representation learning focus of this subsection with the broader theme of learning-to-learn that underpins both offline meta-RL and in-context imitation.

In summary, the evolution of representation learning for context in Meta-RL has moved from simple encoder optimization to a suite of sophisticated techniques grounded in information theory and self-supervision. Contrastive learning and mutual information maximization have become standard tools for ensuring that context representations are sufficient and invariant to irrelevant factors, while trajectory-level and disentanglement objectives address the specific temporal and compositional challenges of sequential decision-making. These methods collectively enable a more sample-efficient meta-training process and equip agents with the robust, compact task representations necessary for rapid and effective zero-shot or few-shot generalization in dynamic environments. The principles and techniques surveyed here form the representational backbone for the in-context imitation learning approaches discussed in the following subsection, where similar ideas of contrastive learning, trajectory-level reasoning, and invariant representations are repurposed for learning from demonstrations without explicit reward signals.

\subsection{Imitation Learning with In-Context Demonstrations}

Building on the representation learning paradigms discussed in the previous subsection---where context encoders are trained to extract compact, sufficient task representations from interaction histories---the field of imitation learning has undergone a parallel transformation through in-context learning. Imitation learning (IL) has traditionally focused on acquiring policies from expert demonstrations via behavioral cloning (BC) or inverse reinforcement learning (IRL), where the learned policy is parameterized by explicit weights. The emergence of in-context learning (ICL) in large models offers a radically different paradigm for imitation: rather than updating model parameters through gradient descent, an agent can infer a task and generate appropriate behaviors simply by conditioning on a sequence of in-context demonstrations provided in its prompt or context window. This subsection explores how in-context demonstrations are used for behavior cloning and goal-conditioned imitation, including methods that leverage offline prior data, skill-based frameworks, and transformer architectures to stitch and generalize from sub-optimal trajectories without explicit reward functions, forming a direct bridge to the algorithm distillation and sequence modeling approaches that follow.

At its core, ICL for imitation treats the problem of policy learning as a sequence modeling problem, akin to how context-based meta-RL encodes task information from trajectories into a latent representation. Given a set of demonstration trajectories, which are sequences of states, actions, and optionally rewards, a causal transformer or similar architecture is trained to autoregressively predict the next token---which in this case corresponds to an action or a state-action pair. Crucially, during inference, the model can be prompted with one or more full or partial trajectories from a new task, and it will generate behaviors that match the demonstrated patterns without any weight updates. This is akin to a form of meta-imitation learning, where the model has learned to imitate from the structure of its pretraining data, much as meta-RL systems learn to extract task context from prior experience. Architectures such as the Behavior Transformer and Decision Transformer exemplify this shift, encoding entire trajectories as sequences and leveraging the context to adapt to new task specifications.

A central advantage of in-context imitation is its capacity to leverage sub-optimal or heterogeneous data, extending the robustness principles seen in representation learning for context. Traditional BC methods, as discussed in much of the offline RL and IL literature, suffer from compounding errors and poor generalization when the expert data is limited or of mixed quality \cite{ref347,ref348}. In-context approaches, however, can be meta-trained on vast offline datasets containing trajectories of varying optimality, learning to identify and replicate the underlying intent or goal from a few examples---much as contrastive and mutual information objectives enable context encoders to extract task-relevant signals from noisy trajectories. For instance, the TRAIL framework learns a latent action space from suboptimal data to facilitate downstream imitation, suggesting that such structured representations can be absorbed into an ICL paradigm for few-shot policy inference \cite{ref349}. Similarly, SubIQ learns from suboptimal demonstrations by regularizing the learned reward landscape, a process that can be implicitly mirrored when a transformer's attention mechanism conditions on successful versus unsuccessful trajectory snippets \cite{ref350}.

Goal-conditioned imitation learning is a particularly natural fit for in-context methods, extending the trajectory-level reasoning that underpins self-supervised context learning. By providing a prompt that includes the desired goal state or a trajectory leading to a goal, the model can generate the sequence of actions required to reach that state. The Generalization Through Imitation (GTI) framework demonstrated that policies can be trained to generalize to novel start-goal pairs by intercepting intersecting trajectories in the state space \cite{ref351}. In an ICL setting, this idea is extended: a transformer can be shown a few examples of goal-reaching trajectories and then conditioned on a new goal to produce a novel path by composing previously seen segments. GO-DICE, a hierarchical offline IL method, learns goal-conditioned policies by decomposing long-horizon tasks into sub-tasks through stationary distribution matching \cite{ref352}; an in-context learner could internalize this hierarchical decomposition to stitch sub-skills without explicit sub-task labels, much as disentangled context representations decouple independent factors of variation in meta-RL.

Stitching, the capability to recombine fragments of sub-optimal trajectories into an optimal or near-optimal policy, is a hallmark of advanced offline RL and a key rationale for in-context imitation, directly foreshadowing the algorithm distillation capability discussed in the next subsection. The Trajectory Stitching (TS) algorithm and its model-based variants demonstrated that augmenting a dataset with synthetic transitions that connect disconnected but promising states can significantly boost the performance of a BC policy \cite{ref353,ref354}. In-context learners, particularly those based on transformer sequence models, can perform a form of implicit planning and stitching by attending over all states and actions in their context. Context-Former explicitly combines contextual imitation learning with sequence modeling to stitch sub-optimal trajectory fragments by emulating the representations of expert trajectories, achieving competitive performance on D4RL benchmarks without explicit dynamic programming \cite{ref355}. This demonstrates how the attention mechanism can serve as a soft retrieval-and-composition engine, blending segments from different demonstrations to construct novel, improved trajectories---a precursor to the full algorithm distillation where transformers learn to replicate entire RL learning procedures.

The integration of diffusion models into behavior cloning also presents a powerful synergy with in-context learning. Diffusion model-augmented behavioral cloning (DBC) models both the conditional and joint distribution of expert data, showing that generative modeling of trajectories can improve generalization over standard BC \cite{ref356}. In an ICL context, a diffusion model can be conditioned on in-context examples to iteratively denoise a trajectory into one that satisfies the demonstrated constraints. Cold Diffusion on the Replay Buffer demonstrates that routing a diffusion model's denoising process through a buffer of known good states increases the feasibility of generated plans, effectively using in-context states as waypoints \cite{ref357}. This is a direct implementation of the ICL philosophy: the model's behavior is entirely dictated by the provided context (the replay buffer and demonstration goals) at inference time, echoing how context encoders in meta-RL condition on historical trajectories without weight updates.

Handling heterogeneous or suboptimal demonstrations without online interaction is a persistent challenge, directly paralleling the context shift reduction problem in offline meta-RL. Methods like Learning to Discern (L2D) learn latent quality representations to weigh the contribution of different demonstrations \cite{ref358}. A transformer trained for ICL can similarly learn to implicitly weight more successful trajectory segments higher by correlating their successor states with higher returns, as seen in the context of offline reinforced imitation learning (ORIL) \cite{ref359}. ORIL learns a reward function by contrasting expert and unlabeled trajectories and then uses this for offline RL; an in-context model could bypass the explicit reward learning step and directly absorb the contrastive signal during meta-training, thereby producing actions that shift the distribution towards expert-like behavior when similar contexts are provided.

Robustness to initial state distribution shift is another area where ICL-based imitation shows promise, mirroring the robustness guarantees sought in context representation learning through information-theoretic objectives. The Return-to-Distribution Planning (POIR) algorithm combines BC with a planner that nudges the agent back to familiar states, demonstrating improved robustness \cite{ref360}. An in-context learner, having seen a wide variety of recovery behaviors during its training on offline data, can exhibit similar robustness by pattern-matching the current off-distribution state to the nearest trajectories in its prompt and generating corrective actions. The backward models used for robust imitation in \cite{ref347} can be thought of as a specialized form of in-context retrieval, where the model learns to imagine trajectories that lead back to the demonstration manifold.

Practical applications, such as autonomous driving and robotic manipulation, have begun to adopt context-conditioned imitation paradigms. The context-conditioned imitation learning approach for urban driving maps a context state directly to a future trajectory, sidestepping covariate shift by redefining the prediction target \cite{ref361}. This aligns perfectly with the in-context philosophy of conditioning a model on a rich representation of the environment state to directly output a plan. In robotics, frameworks like SWBT leverage both expert and easy-to-collect imperfect demonstrations without environment interaction, using a similarity-weighted behavior transformer to focus on the most relevant contexts during learning \cite{ref362}.

In summary, in-context demonstrations for imitation learning represent a convergence of sequence modeling, offline RL, and meta-learning. By treating policy learning as a conditional generation problem over trajectories, transformer-based and diffusion-based architectures can stitch together sub-optimal data, generalize to new goals, and adapt to novel tasks without any parameter updates. The principles extracted from robust BC, trajectory stitching, and goal-conditioned hierarchical methods are being unified under the ICL umbrella, pointing toward general-purpose agents that reuse and recombine past experience in a purely context-driven manner. This paradigm directly anticipates the more general algorithm distillation frameworks discussed next, where transformers are trained not merely to imitate optimal policies, but to internalize and execute the full process of reinforcement learning improvement itself.

\subsection{Algorithm Distillation and Sequence Modeling for Meta-RL}

The ability to infer a task and generate appropriate behaviors from in-context demonstrations, as discussed in the previous subsection, represents a powerful form of imitation. A more profound and general capability, however, is the capacity to act as universal in-context learners not merely for single tasks, but for entire reinforcement learning algorithms. This capability shifts the paradigm from learning a single policy to learning the process of policy improvement itself---a direct leap from imitating optimal behavior to internalizing the very algorithms that produce it. This transition forms the core of this subsection and establishes the sequence modeling foundations that the subsequent subsection on curriculum learning will build upon, as those automated training strategies are essential for generating the vast and diverse data required for algorithm distillation to succeed.

Two closely related frameworks embody this idea: sequence modeling for meta-RL, where the transformer's internal state mimics Bayesian inference or memory reinstatement, and algorithm distillation, where the transformer is explicitly trained to replicate the learning and adaptation dynamics of a full RL algorithm. The foundational insight for sequence modeling in meta-RL is that the transformer architecture is uniquely suited to implement the core computational processes of a meta-learner. As argued in the introduction of \cite{ref363}, the self-attention mechanism inherently performs a form of fast adaptation by attending over a context of past experiences to construct a task representation. The resulting model, TrMRL, proposes that transformers ``mimic the memory reinstatement mechanism,'' recursively building an episodic memory through its layers. This process is formalized as the self-attention layers computing a consensus representation that minimizes the Bayes Risk, providing a principled grounding for the model's ability to infer the latent task from a history of interactions \cite{ref363}. This perspective is elegantly generalized by theoretical frameworks that recast memory-based meta-learning as amortized Bayesian filtering \cite{ref364}. In this view, a transformer processing a trajectory \((s_0, a_0, r_0, \ldots)\) is learning to approximate Bayes-optimal inference over the latent task variables or parameters of the environment's dynamics and reward function, directly translating data into optimal actions without an explicit parameter update step---extending the representation learning principles of previous sections into a fully Bayesian framework. Empirical validation of this principle shows that transformers and other memory-based models can indeed learn to ``accurately approximate known Bayes-optimal algorithms'' in piecewise stationary environments with unobserved changepoints, behaving as if performing true probabilistic inference over the underlying latent structure \cite{ref365}.

Moving beyond inference of a static latent task, a more ambitious goal is to model the full RL learning process. This is the core of algorithm distillation, which can be viewed as a supervised meta-learning approach where the training dataset consists of the entire learning histories of an expert RL agent. Instead of the transformer learning to be an optimal policy for a specific task, it is trained to take in a long sequence of interactions from a task and output actions that match an expert's learning progression. The transformer must thus learn to represent and execute the operations of the expert algorithm's own learning rule, such as computing sufficient statistics, exploring, and updating a value function. A rigorous theoretical analysis of this paradigm is provided by a framework that formalizes this as ``supervised pretraining for ICRL'' \cite{ref366}. This work proves that when a transformer is trained on data generated by an expert algorithm, it will learn to ``imitate the conditional expectation of the expert algorithm given the observed trajectory'' under realizability conditions. This provides a theoretical guarantee for algorithm distillation: the transformer becomes a simulator of the expert algorithm's behavioral output, not just the final optimal policy, achieving meta-generalization to unseen tasks with a scaling error related to model capacity and a distribution divergence factor \cite{ref366}.

A powerful emergent property of transformers trained in this way is their ability to perform ``stitching,'' or data composition, where the model intelligently recombines sub-optimal trajectory segments from different behaviors into an optimal policy. This capability directly extends the trajectory stitching seen in in-context imitation to the algorithm learning process itself, as showcased by the Agentic Transformer (AT) \cite{ref367}. AT is trained on sequences of trajectory experience sorted by their total returns, conditioned on target returns and task completion tokens. This relabeling and sequence-level training allows AT to learn the process of improvement: it can ``learn to improve upon itself both at training and test time'' by treating sub-optimal data as stepping stones for better performance \cite{ref367}. This capability is a distinct advantage over standard imitation learning, which would simply replicate the average or best behavior in the data, and it aligns with the theoretical findings that sequence models can implement not just one algorithm, but can adaptively select among different learning strategies based on the context \cite{ref110}.

The architectural capacity to support these complex in-context algorithms is an active area of research. The Transformer in Transformer (TinT) construction demonstrates how a large transformer can simulate the full fine-tuning of an internal, smaller model within a single forward pass, offering a precise architectural mechanism for how an ICL model could internally implement a complex learning subroutine like updating an entire neural network \cite{ref368}. This moves the needle from simulating simple regression to simulating the multi-step adaptation of a pre-trained language model, a primitive highly relevant to replicating sophisticated RL updates. For RL specifically in partially observable environments requiring long-range memory, world models built on transformer state-space architectures, like the TransDreamer agent, show how explicit sequence modeling of dynamics can outperform recurrent network baselines in complex 3D tasks by providing a more powerful and stable memory for policy learning \cite{ref369}. The convergence of these lines of work---from Bayes-optimal inference, to algorithm distillation and stitching, to increasingly powerful internal model simulations---paints a picture where sequence models serve as a unified computational substrate for both representing the knowledge gained across an agent's lifetime and executing the algorithms to acquire and improve upon that knowledge. The result is a class of meta-learners that can achieve rapid adaptation, zero-shot generalization to new tasks, and even outperform the demonstrators on which they were trained, all through a purely in-context forward pass. However, realizing this vision at scale critically depends on the agent's exposure to a sufficiently rich and appropriately structured distribution of training tasks, a challenge that motivates the automated curriculum generation methods explored in the following subsection.

\subsection{Curriculum Learning and Self-Paced Contextual Reinforcement Learning}

Building on the sequence modeling and algorithm distillation frameworks discussed in the previous subsection---which implicitly require vast and diverse training data to induce general-purpose in-context learning---a fundamental challenge in meta-reinforcement learning (meta-RL) is how to design or automatically generate the distribution of training tasks, or contexts, over which an agent should practice. Manually specifying this distribution is often prohibitive, particularly in open-ended or procedurally generated environments where the space of possible tasks is vast and the difficulty of any given task is not known \emph{a priori} \cite{ref199}. This challenge becomes acute when the goal is to train agents capable of the kind of hierarchical skill composition and compositional generalization described in the next subsection, where the agent must ultimately master a combinatorially large space of multi-stage tasks. To address this, curriculum learning (CL) and its automated variants (autocurricula) have emerged as critical components for improving both the sample efficiency and asymptotic generalization of context-based meta-RL agents. These methods frame the training process not as a static sampling from a fixed task distribution, but as an adaptive sequence that tailors task presentation to the agent's current capabilities, thereby accelerating learning and avoiding local optima \cite{ref201}.

The core principle underlying these approaches is to structure the learning experience from ``easier'' to ``harder'' contexts, mimicking the pedagogical strategy of a human teacher. Self-paced learning (SPL) operationalizes this by having the agent itself, or a teacher model, select training instances based on a difficulty metric, often related to the agent's current loss or learning progress \cite{ref370,ref371}. In the contextual meta-RL setting, a task is defined by a context variable that parameterizes the environment's dynamics or reward function. The challenge then becomes how to dynamically sample these contexts to form an effective curriculum---a progression that, ideally, lays the foundation for the kind of complex, multi-step learning processes that algorithm distillation seeks to replicate. The framework proposed by \cite{ref372} (SPaCE) exemplifies this approach by leveraging the agent's value function estimates to automatically generate an online curriculum. SPaCE prioritizes contexts where the agent is making meaningful progress, effectively creating a developmental schedule without computationally expensive optimization, and has been shown to accelerate training by up to 10x and improve generalization to unseen task instances.

Moving beyond simple loss-based heuristics, optimal transport theory provides a principled mathematical lens for curriculum generation. Methods like GRADIENT \cite{ref373} and related work \cite{ref374} frame the curriculum as a geodesic interpolation between a source (easy) and target (hard) task distribution. By treating the curriculum design as a constrained optimal transport problem over a space of contextual distributions, these algorithms generate a smooth sequence of intermediate task distributions. This gradual domain adaptation mitigates the large distributional shift that can catastrophically derail learning, ensuring a more stable and effective transfer from simple to complex contexts. This perspective is particularly powerful when the context space is continuous and high-dimensional, and the definition of task difficulty is non-trivial---precisely the conditions under which a transformer-based meta-learner must operate to achieve the in-context algorithm emulation described in Section 7.6.

Another significant advancement is the concept of meta-automatic curriculum learning (Meta-ACL), which seeks to generalize curriculum generation itself, directly mirroring the meta-learning philosophy of learning to learn. Standard ACL algorithms are typically \emph{tabula rasa} processes that must explore the task space anew for every learning agent, a costly proposition. Meta-ACL, as introduced by AGAIN \cite{ref375,ref376}, distills this exploration into a reusable pedagogical policy. By training a curriculum generator over a distribution of learners, AGAIN learns a prior for effective task sequencing that can rapidly specialize to a new, previously unseen agent. This ``classroom teaching'' paradigm abstracts the principles of good curricula, making the training of new meta-RL agents significantly more efficient and providing a curriculum-level analogue to the algorithm distillation approach where the teaching process itself is learned from data.

In cases where any explicit reward signal design is challenging, fully unsupervised methods for task generation and selection become essential. Unsupervised Environment Design (UED) formulates the interaction between the agent and a task-generating adversary to produce a curriculum of environments at the frontier of the agent's capabilities \cite{ref377,ref378}. CLUTR \cite{ref379} improves upon this by decoupling the process into two stages: first learning a smooth latent manifold of tasks via a variational autoencoder, and then training a teacher to propose sequences of latent codes from this manifold based on minimax regret. This decoupling overcomes the non-stationarity that often plagues jointly trained UED systems, resulting in robust zero-shot generalization. Similarly, unsupervised methods for inducing a meta-training task distribution by modeling the agent's trajectory data distribution can scaffold an autocurriculum, enabling vision-based meta-RL agents to learn exploratory and control policies before ever seeing a task-specific reward function \cite{ref319}.

These curriculum strategies are deeply intertwined with the quality of context representations---a theme that resonates with the context-conditioned skill execution and compositional generalization discussed in the following subsection. The effectiveness of adapting to a new context hinges on the encoder's ability to extract sufficient task information from a history of interactions. Self-supervised objectives, such as trajectory contrastive learning, directly enhance this process by training the context encoder to map trajectories from the same task to similar embeddings, thereby providing a cleaner and more structured task space over which a curriculum can operate \cite{ref318,ref310}. When contexts are well-disentangled, a curriculum in the latent space can translate more directly to a meaningful progression of task complexities---a progression that can ultimately scaffold the acquisition of the reusable skill libraries central to hierarchical RL. Further, the structural prior of a domain can be explicitly encoded to guide the curriculum. For instance, integrating reward machines---finite-state machines representing non-Markovian task structure---into the self-paced RL process provides a strong inductive bias that stabilizes curriculum generation and enables the solution of long-horizon tasks where purely data-driven curricula fail to find any progress niche \cite{ref380}.

Despite these successes, open challenges remain in bridging the gap between an automated curriculum's notion of difficulty and the compositional structure of the target task distribution. A curriculum that progresses by varying a single quantitative parameter (like goal distance) may fail on tasks requiring the sequential mastery of qualitatively different skills---precisely the kind of multi-stage compositional tasks addressed by the hierarchical and skill-based approaches of Section 7.8. Future work must focus on curricula that respect the compositional and causal structure of the environment \cite{ref381}, potentially through hierarchical or modular curriculum strategies that sequence the acquisition of skills, not just tasks \cite{ref321}. The integration of these structured autocurricula with lifelong, non-stationary adaptation remains a grand challenge for developing truly autonomous, self-improving contextual agents capable of the open-ended compositional generalization that lies at the heart of intelligent behavior.

\subsection{Skill Acquisition, Hierarchical RL, and Compositional Contexts}

Building directly on the curriculum and self-paced learning strategies of Section 7.7---which provide the pedagogical sequencing for acquiring increasingly complex behaviors---we now examine how these behaviors are structured once learned. A central challenge in deploying agents for complex, real-world tasks lies in overcoming the prohibitively long horizons and sparse rewards that plague standard reinforcement learning (RL) algorithms. A powerful solution, drawing a direct parallel to how humans decompose novel problems into familiar subroutines, is to structure agent behavior hierarchically. In this paradigm, agents leverage a repertoire of reusable skills, where context, as explored throughout this survey, serves as the critical binding mechanism for composing, adapting, and transferring these skills to new tasks and environments. This subsection analyzes approaches that integrate skill discovery and hierarchical RL (HRL) with the principles of in-context learning and meta-learning, focusing on how context conditions the execution of learned sub-policies---a capability that is foundational for the sim-to-real transfer and domain adaptation methods discussed in the following subsection, where these structured skills must be robustly deployed under varying dynamics.

The foundational premise of skill-based RL is to break down long-horizon tasks into a sequence of temporally extended actions or skills. In a contextual framework, these skills are not monolithic; they are context-conditioned policies. For example, a ``navigation'' skill might exhibit different behaviors depending on whether the context specifies a point or region goal, as demonstrated by multi-skill mobile manipulation methods that train navigation with region goals to provide flexibility for downstream manipulation skills \cite{ref382}. The process of acquiring such skills often involves meta-RL, echoing the meta-learning principles that underpin the algorithm distillation and autocurricula of previous sections. Frameworks like DCMRL explicitly decouple the learning of task contexts from skill representations, using contrastive objectives to pull context representations of the same task together while pushing away those from different tasks. This disentanglement allows for the creation of discrete skill and context codebooks via Gaussian quantization, providing a generalizable prior experience that enables rapid adaptation to unseen target tasks during meta-testing---a structural separation that directly anticipates the dynamics vs.~skill disentanglement central to sim-to-real techniques like DOMINO \cite{ref321}.

The discovery of these skills from data, rather than hand-engineering them, is critical for scalability and mirrors the automatic task generation challenges addressed by curriculum learning. Several approaches leverage the temporal structure inherent in demonstrations. A bottom-up method constructs a hierarchical task structure from unsegmented demonstrations through agglomerative clustering, identifies recurring patterns as skills, and then uses a meta-controller to compose these discovered skills for new tasks \cite{ref383}. Other works integrate large language models (LLMs) to guide this discovery process, with algorithms using LLM-proposed trajectory segmentations as a prior for a hierarchical variational inference framework, discovering reusable skills that accelerate learning on new long-horizon tasks \cite{ref384}. The quality and alignment of these skills can be further refined through human feedback, as seen in algorithms that learn a model over human preferences to extract human-aligned skills from offline data, significantly outperforming methods that rely on biased generative models of demonstrations \cite{ref385}.

Once a skill library is established, meta-learning and context-based adaptation become central to their reuse, forming a direct link to the rapid adaptation mechanisms required for closing the reality gap. Parameterized skill-learning algorithms leverage off-policy Meta-RL to learn transferable skills and synthesize them into a new action space for efficient learning, effectively constructing a three-level hierarchical framework that models a Temporally-extended Parameterized Action Markov Decision Process \cite{ref386}. A key innovation is the use of context to define the sub-goals or parameters of these skills. For instance, compositional imitation learning approaches can perform one-shot generalization to previously unseen tasks by using a single reference trajectory to generate adaptive ``near future'' subgoals through compositional arithmetic in a learned latent representation space \cite{ref387}.

The interplay between hierarchical structures and context is most potent when addressing the compositional nature of real-world tasks. Many tasks are not a single monolithic problem but a sequence of distinct sub-tasks, requiring an agent to not only identify the right skill but also understand the transitions between them. The OCEAN framework tackles this online task inference problem by modeling both global context variables (representing the mixture of sub-tasks) and local context variables (capturing the transitions between sub-tasks) in a joint latent space. This allows for flexible, sequential context adaptation that boosts performance on complex, multi-stage tasks---precisely the kind of compositional challenges that autocurricula must scaffold and that robust sim-to-real policies must ultimately solve \cite{ref388}. This compositional ability naturally aligns with the semantic reasoning offered by LLMs. Hierarchical agents can exploit the planning and in-context reasoning capabilities of LLMs to guide a high-level policy, making learning significantly more sample-efficient. The LLM provides the compositional plan as a form of context for the low-level RL skills, and once trained, the agent no longer requires the LLM at deployment \cite{ref389}.

The mechanisms for skill diversity and specialized composition are also refined through architectural choices. Mixture-of-experts (MoE) models provide a natural framework where each expert represents a distinct skill formalized as a contextual motion primitive. By optimizing a maximum entropy objective over both the experts and their associated context distributions, these methods can learn diverse skills where each expert focuses on its best-performing sub-region of the context space, creating an automatic curriculum that resonates with the adaptive difficulty scaling of Section 7.7 \cite{ref390}. Beyond static skill selection, adaptive skill sequencing learns to chain existing skills for exploration but decouples this from the final policy learning, which operates directly on raw experience. This prevents the learned policy from being constrained by sub-optimal skill priors and significantly outperforms classical skill transfer methods \cite{ref391}.

To overcome the domain-specificity trap, where learned skills fail to transfer to environments with different dynamics or visual appearances---the very challenge that sim-to-real transfer addresses---context again provides the key to generalization. The DuSkill framework employs a guided diffusion model to generate versatile skills from limited offline data. It disentangles the skill embedding space into domain-invariant behaviors and domain-specific variations, allowing it to generate skills that bridge different domains and accelerate high-level policy learning for new environments, embodying the same separation of universal and context-specific factors seen in Contextualized World Models \cite{ref392}. Similarly, a semantic skill translator framework (SemTra) uses multi-modal models and the reasoning of a pre-trained language model to extract skills from interleaved multi-modal snippets and adapt them to new target domains. This zero-shot adaptation occurs at two levels: task adaptation (translating a sequence of skills) and skill adaptation (parameterizing each skill for the new context) \cite{ref393}.

Ultimately, the synergy between hierarchical RL, skill acquisition, and context-driven meta-learning points towards agents capable of lifelong and compositional learning---a vision that motivates the generalization and continual adaptation challenges examined in subsequent sections. Frameworks that leverage diverse suboptimal world models can automatically decompose a complex task into simpler, modular subpolicies in a bottom-up manner, acquiring the necessary skills and a controller concurrently \cite{ref394}. This process is further augmented by using language not just as a command interface, but as a compositional abstraction itself. Learning an instruction-following low-level policy and a high-level policy that reuses abstractions across tasks allows agents to reason using structured language, where the compositional nature of language is critical for systematic generalization to new sub-skills and tasks \cite{ref395}. The integration of task planning with RL-based skill learning, as in the LEAGUE framework, closes the loop by using the planner's symbolic interface to guide skill learning and create abstract state spaces for reuse, all while the agent continuously grows its capability and the set of solvable tasks in-situ---a crucial step toward the open-ended, self-improving contextual agents that in-context learning ultimately promises to deliver \cite{ref396}.

\subsection{Sim-to-Real Transfer and Domain Adaptation with Context-Aware Meta-RL}

Building upon the compositional skill architectures and hierarchical planning frameworks discussed in the previous subsection, we now turn to a critical challenge that determines whether these sophisticated agents can leave the simulator and operate in the physical world: the ``reality gap.'' The deployment of reinforcement learning (RL) agents in real-world robotic systems is fundamentally bottlenecked by this discrepancy between simulated training environments and the complex, variable dynamics of the physical world \cite{ref397}. Context-aware meta-reinforcement learning (meta-RL) has emerged as a powerful paradigm to bridge this gap by enabling agents to rapidly infer and adapt to the unknown dynamics of a target environment at deployment time. In direct continuity with the skill composition and hierarchical adaptation methods previously described, context-based meta-RL treats the mismatch between simulation and reality not as noise to be averaged out, but as a hidden task variable to be inferred, training a latent-conditioned policy that can modulate its behavior based on a small amount of recent interaction data, effectively performing online system identification and adaptation.

A central technique in modern sim-to-real transfer is domain randomization (DR), which bridges reality by massively varying the physical parameters of the simulator during training---such as friction, mass, and visual textures---forcing the policy to learn robust strategies that are invariant to these variations \cite{ref398}. Context-aware meta-RL provides a formal structure for DR by treating each set of randomized parameters as a distinct ``task'' and learning to infer a context variable that captures the relevant dynamics. This directly echoes the skill-codebook and task-context disentanglement principles seen in DCMRL, but applied to dynamics rather than task objectives. However, naïve application of DR with fixed, uniform randomization distributions often leads to overly conservative policies or fails to provide the curriculum needed for effective learning. Several works have proposed automated and adaptive DR schemes to address this. For instance, Automatic Domain Randomization (ADR) dynamically increases the complexity of the environment distribution based on the agent's current performance, creating an emergent curriculum that was instrumental in solving complex manipulation tasks like solving a Rubik's Cube with a dexterous robot hand, where the learned policy exhibited signs of emerging meta-learning at test time \cite{ref399}. Similarly, Bayesian Domain Randomization (BayRn) formulates the search over randomization parameter distributions as a Bayesian optimization problem, using sparse real-world data to adapt the distribution itself and maximize real-world performance \cite{ref400}. Other methods like Continual Domain Randomization (CDR) address the problem of jointly learning over all randomized parameters by enabling sequential training on subsets of randomizations using continual learning, thereby reducing task difficulty while maintaining robustness and outperforming joint or sequential randomization baselines \cite{ref401}. The principle of maximizing internal representation robustness is also captured by DORAEMON, which frames DR as a constrained entropy maximization problem, directly maximizing the entropy of the dynamics parameter distribution while keeping success probability high, resulting in policies that successfully zero-shot transfer to unknown real-world parameters \cite{ref402}.

Beyond explicit randomization of simulator parameters, context-based meta-RL enables sim-to-real transfer through the direct learning of structured latent representations of dynamics---a natural extension of the latent skill and context spaces explored in compositional learning. A key challenge is inferring these dynamics from high-dimensional, pixel-based visual observations, which are the primary input for many robotic systems. Latent state models (LSMs) provide a compelling solution by learning a compact, low-dimensional representation of the underlying state from image sequences, which can then be used for planning and policy learning. The MELD algorithm (meta-RL with latent dynamics) demonstrates that an LSM, when provided with an observation space augmented by previous actions and rewards, can inherently perform task inference, enabling a real robotic arm to learn to insert an Ethernet cable into new locations after only 8 hours of real-world meta-training \cite{ref403}. This approach side-steps the need for explicit, hand-crafted DR by directly learning a representation that captures the essential variability needed for adaptation. This concept of representation learning is further refined by focusing on context variables themselves. For example, the Decomposed Mutual Information Optimization (DOMINO) framework addresses the challenge of multi-confounded dynamics by learning a disentangled context representation that maximizes mutual information with historical trajectories while minimizing state transition error, leading to more sample-efficient generalization in unseen environments \cite{ref312}. Similarly, the ProtoCAD model learns a prototypical context-aware dynamics representation that enforces temporal consistency in the context variable, significantly improving generalization across diverse dynamics in high-dimensional control tasks compared to standard recurrent state-space models \cite{ref404}. Mirroring the decoupling strategies introduced in the previous subsection for task and skill representations, DCMRL uses contrastive learning and Gaussian quantization to create discrete, generalizable codes for task contexts and reusable skills, demonstrating that the separation of what varies (the task) from what is reusable (the skill) is equally powerful for dynamics adaptation as it is for task composition \cite{ref321}.

Knowledge distillation offers another potent avenue for context-aware sim-to-real transfer by leveraging privileged information that is available in simulation but not in the real world---a strategy consistent with the teacher-student dynamics observed in LLM-augmented hierarchical agents. In simulation, a ``teacher'' policy and world model can be trained with full access to ground-truth state information, such as object poses, velocities, and friction coefficients. This teacher's internal representations can then be distilled into a ``student'' model that operates purely on realistic, domain-randomized images, which is the model eventually deployed on the robot. The TWIST framework (Teacher-Student World Model Distillation) exemplifies this by training a teacher world model on state information and then supervising a student world model that only sees domain-randomized images, efficiently transferring the learned latent dynamics and outperforming naïve domain randomization \cite{ref405}. A related approach, the Historical Information Bottleneck (HIB), distills privileged knowledge by learning a representation from historical trajectories that captures underlying changeable dynamic information, which theoretically and empirically helps reduce the value discrepancy between an oracle policy trained with privileged knowledge and the learned policy deployed in the real world \cite{ref406}. Bi-directional Domain Adaptation (BDA) extends this by bridging the gap in both directions: it adapts real-world visual data to look like simulation (real2sim) for the perception model and adapts the simulation's dynamics to match the real world (sim2real) for the control policy, dramatically improving sample efficiency for embodied navigation agents \cite{ref407}.

Ultimately, the success of these methods lies in their ability to leverage context not just as a passive summary, but as an active, structured inference mechanism---a theme that directly anticipates the generalization challenges and compositional reasoning demands explored in the next section. The Contextualized World Models (ContextWM) framework explicitly separates the modeling of universal dynamics from context-specific environmental factors, allowing it to effectively pre-train on abundant but diverse ``in-the-wild'' videos and then rapidly fine-tune to specific robotic manipulation and locomotion tasks \cite{ref408}. This is mirrored by the contextual recurrent state-space model (cRSSM), which incorporates context directly into a Dreamer-style world model, enabling it to disentangle latent state from context and achieve zero-shot generalization to worlds with unseen dynamics \cite{ref409}. While these methods have demonstrated impressive results, open challenges remain---challenges that form the core of the generalization failures and lifelong adaptation limitations discussed in the following subsection. Efficiently integrating real-world interaction data, scaling to the extreme diversity of real-world non-stationarity, and ensuring safety during the online adaptation phase are critical frontiers. Nevertheless, the convergence of context-aware meta-RL, adaptive domain randomization, structured latent representations, and privileged knowledge distillation forms a robust and maturing toolkit for finally making sim-to-real transfer a reliable pipeline for the compositional, hierarchical, and generalist robotic agents that in-context learning promises to deliver.

\subsection{Generalization, Robustness, and Open Challenges in In-Context Decision-Making}

Building directly upon the convergence of context-aware meta-RL, adaptive domain randomization, structured latent representations, and privileged knowledge distillation discussed in the previous subsection, we now confront the critical generalization challenges that determine whether these sim-to-real transfer successes can extend to truly open-world deployment. The capacity for generalization in in-context decision-making agents is a multifaceted challenge, tested severely by the dynamic and often unpredictable nature of real-world environments. A core promise of context-aware meta-reinforcement learning lies in its ability to infer task identity from a support set of experiences and adapt policy execution accordingly, without gradient updates---precisely the same online system identification principle that enables zero-shot sim-to-real transfer. However, the robustness of this inference mechanism degrades under distribution shifts, sparse rewards, and unobserved confounders, revealing critical gaps between empirical success and a principled understanding of in-context learning. This section examines the generalization landscape, analyzes failure modes, and identifies open challenges that define the frontier for reliable in-context sequential decision-making.

A primary test for any learning agent is its performance under distribution shift. For context-aware agents, this occurs when the deployment task distribution diverges from the training environment distribution---a more severe version of the dynamics gap that domain randomization and context inference were designed to bridge. The challenge is profound in settings where reward mechanisms change entirely, as a policy must disentangle consistent dynamics from variable objectives, mirroring the separation of universal dynamics from context-specific factors that ContextWM and cRSSM achieve but extending it to reward structure. Methods like Task Aware Dreamer (TAD) have demonstrated that integrating reward-informed features into world models can help identify latent characteristics consistent across tasks, thereby improving generalization even when reward functions shift significantly \cite{ref410}. However, these models still rely on a fundamental assumption of overlapping structure. A more severe challenge is out-of-variable generalization, where agents encounter environments with entirely novel combinations of causal variables never jointly observed during training. This requires the agent to re-use past marginal information over subsets of variables, a capability that goes beyond standard distributional robustness and touches upon the compositional nature of world knowledge \cite{ref411}.

The presence of unobserved confounders and complex dynamics further complicates in-context adaptation, extending the multi-confounded dynamics challenge that DOMINO addresses in the sim-to-real context. In real-world scenarios, multiple latent factors can simultaneously influence transition dynamics, making it difficult to infer a single, compact context vector for decision-making. Just as DOMINO learns a disentangled context representation by maximizing mutual information with historical trajectories while minimizing state transition prediction error, theoretical analysis shows this approach overcomes the underestimation of mutual information caused by multi-confounded environments, benefiting both model-based and model-free algorithms in unseen dynamics \cite{ref312}. This principled treatment of context disentanglement is crucial for transferring policies learned in simulation to reality, a domain where context-based adaptation must bridge visual and dynamics gaps---precisely the challenge that BDA addresses through bidirectional domain adaptation. Further underscoring this theme, approaches that co-evolve a context recognition module with a skill module have been shown to yield significantly more robust behavior in environments with previously unseen effects, echoing the task-skill decoupling principles of DCMRL \cite{ref412}.

Sparse rewards present another compounding factor that amplifies all previously discussed challenges. In long-horizon tasks, an agent may receive feedback only upon task completion, making it exceptionally hard to credit individual actions or infer the correctness of an inferred context from limited demonstration trajectories---the very setting in which meta-RL with latent dynamics models like MELD must operate. To address this, goal-conditioned and skill-based frameworks have emerged. The GEASD framework, for example, captures environmental structural patterns through an adaptive skill distribution that optimizes the local entropy of achieved goals, enhancing deep exploration in states with familiar structures even under sparse reward signals \cite{ref413}. This resonates with the concept of creating generally capable agents through open-ended learning, where a dynamic, evolving curriculum of tasks prevents the agent from stalling on a single objective and promotes the discovery of robust, reusable skills---a philosophy directly aligned with the emergent curricula generated by adaptive domain randomization schemes like ADR \cite{ref414}.

Despite these advances, several open challenges limit the leap from structured benchmarks to truly general in-context decision-makers. The foremost of these is compositional generalization: the ability to recombine known sub-tasks, skills, or primitive concepts to solve novel tasks unseen in training or context demonstrations. This is a well-documented weakness in deep learning, often stemming from a model's inability to induce systematic, sparse correspondences between input components and output behaviors. For in-context meta-RL, a single context vector may be insufficient to capture the compositional nature of a novel task that requires sequencing several previously learned abilities in a new order---the same limitation that motivates hierarchical skill composition and codebook architectures. The OCEAN framework tackles this by moving beyond a single Gaussian context to a joint latent space of global and local variables, where global variables represent a mixture of required sub-tasks and local variables capture the transitions between them, enabling online task inference for multi-stage problems \cite{ref388}. Yet, achieving this in a scalable, unsupervised manner that mirrors human-like one-shot generalization from a single demonstration of a composite task remains an open problem. Research suggests that agents infused with hierarchical latent language for both task inference and sub-task planning are better equipped to describe and replicate unseen procedural tasks, pointing toward a synergy between linguistic compositionality and behavioral generalization that will be explored further in the next section \cite{ref415}.

Long-horizon reasoning is inextricably linked to this compositional challenge. While an agent may infer a task from a few demonstrations, extrapolating that understanding to plan thousands of steps into the future without compounding model error is a formidable barrier that even the teacher-student world model distillation frameworks like TWIST do not fully resolve. Model-based approaches that learn task-agnostic dynamics models naturally suffer from compounding one-step errors. A transformative alternative is the use of generalized occupancy models (GOMs), which model the distribution of all possible long-term outcomes from a state under the coverage of a stationary dataset, thereby avoiding compounding error while retaining generality across arbitrary reward functions \cite{ref416}. Translating such long-horizon reasoning into an in-context learning paradigm, where a few trajectories must reveal the long-term outcome distribution for a new task, is a critical next step.

Lifelong in-context adaptation represents a further horizon that pushes beyond the single-task transfer scenarios discussed in sim-to-real pipelines. Currently, most agents are evaluated on their ability to adapt to a single new task from a static context. A more realistic scenario is a non-stationary stream of tasks, where the agent must continually adapt without catastrophic forgetting of previous contexts and without allowing the immediate in-context prompt to grow unboundedly---a challenge that continual domain randomization methods like CDR begin to address but do not fully solve. The interplay between generalization and forgetting is not zero-sum; methods that improve generalization through robust representation learning, such as enforcing shape-texture consistency, can also mitigate catastrophic forgetting in a continual learning framework, suggesting a mutually reinforcing relationship that could be harnessed for lifelong in-context learners \cite{ref417}.

Finally, a significant theoretical chasm exists between the empirical heuristics of in-context decision-making and formal guarantees. While methods like self-paced contextual reinforcement learning introduce principled, information-theoretic objectives to structure the curriculum of context distributions---echoing the entropy maximization principles of DORAEMON---translating these into PAC-style generalization bounds for in-context adaptation is nascent \cite{ref320}. Many algorithms demonstrate empirical improvements in sample efficiency and robustness but lack guarantees on how many in-context examples are needed for a policy to generalize safely under a defined distribution shift. Bridging this gap requires a theory that accounts for the emergent nature of in-context learning in transformers, moving beyond linear meta-learning analogies to explain the non-linear, compositional inferences that pre-trained sequence models perform. Solving these interconnected challenges---compositional understanding, long-horizon coherence, lifelong streaming adaptation, and a rigorous theory---is essential to deploy in-context decision-making agents in the open, unpredictable worlds that lie beyond the laboratory, forming the very generalization and adaptation demands that the following section addresses.

\section{Robustness, Calibration, Safety, and Privacy in In-Context Learning}

\subsection{Miscalibration and Overconfidence in In-Context Learning}

Building directly on the understanding that ICL predictions are deeply sensitive to prompt design, the calibration properties of models under In-Context Learning (ICL) further reveal a complex and often precarious relationship between predictive performance and uncertainty quantification, a cornerstone for the safe deployment of machine learning systems. While the previous section discussed how sensitivity manifests as inconsistency, this section delves into how it manifests as miscalibration. Unlike traditional supervised models where calibration is a post-training concern, ICL introduces unique calibration dynamics that evolve with the prompt structure and the number of demonstrations provided. The fundamental issue of miscalibration in deep neural networks, where a model's predicted confidence diverges from its empirical accuracy, is well-documented across various paradigms \cite{ref418}. However, ICL presents a distinct landscape where these issues are both amplified and reshaped by the very context that enables learning.

A central finding in the analysis of ICL calibration is its non-monotonic relationship with the number of in-context examples. Comprehensive studies across a broad spectrum of natural language understanding and reasoning tasks have revealed that models initially exhibit increased miscalibration as the first few examples are added, before eventually achieving better calibration in higher-shot regimes \cite{ref38}. This pattern indicates a critical vulnerability in low-shot settings, where practical ICL applications often operate. The initial rise in miscalibration can be attributed to the model rapidly updating its implicit task beliefs from a small, potentially biased sample, leading to a high degree of overconfidence in predictions derived from this sparse information. This observation challenges the assumption that simply providing more examples linearly improves all aspects of model reliability, revealing a ``valley of miscalibration'' that practitioners must navigate.

This miscalibration problem is further exacerbated by prompting techniques designed to enhance task performance. Notably, chain-of-thought (CoT) prompting, while significantly improving accuracy on complex reasoning tasks, has been observed to produce overconfident yet incorrect predictions \cite{ref38}. The generation of intermediate reasoning steps can create an illusion of reliability; the model's confidence in its final, flawed conclusion may remain high because it is conditioned on a coherent-sounding, yet logically faulty, rationale. This phenomenon is linked to the tendency of models to ``overthink'' when supplied with misleading contextual information, a direct analogue to the sensitivity issues caused by spurious correlations in demonstrations. Research on how language models process false demonstrations has identified ``overthinking'' as a behavior where model accuracy given incorrect demonstrations progressively decreases after a certain critical layer, driven by ``false induction heads'' that copy false information \cite{ref419}. This mechanistic insight suggests that CoT prompting can inadvertently create a self-reinforcing loop of confident error, where the model's own generated context misleads its subsequent generation layers, decoupling confidence from correctness.

The challenge of achieving simultaneous gains in both task performance and calibration is a recurring theme. An in-depth analysis contrasting supervised fine-tuning (SFT) and ICL has confirmed that a state of overconfidence, or miscalibration, is a pervasive problem across all learning methods in low-resource scenarios, and simultaneous improvements in accuracy and calibration are difficult to achieve \cite{ref39}. This exposes a fundamental trade-off that complicates model selection and deployment, forcing a prioritization between pure performance and trustworthy, well-calibrated confidence estimates. This tension is particularly critical in high-stakes domains where acting on a confident but incorrect prediction has severe consequences.

To combat miscalibration, a range of recalibration techniques have been developed and adapted for the ICL setting, paralleling the mitigation strategies for prompt sensitivity discussed earlier. Benchmark studies have shown that a scaling-binning calibrator can consistently reduce calibration errors across various tasks \cite{ref38}. This post-hoc method learns a transformation of the model's output logits or probabilities based on a held-out calibration set, serving as a pragmatic remedy after the ICL process. More advanced methods have been proposed to address the contextual biases that contribute to miscalibration, unifying various prior approaches into a single framework. For instance, Batch Calibration (BC) is a zero-shot, inference-only technique that controls the contextual bias arising from the batched input, effectively mitigating the effects of format and verbalizer biases in prompts \cite{ref178}. This method highlights how the source of miscalibration in ICL is not just a property of the model's internal knowledge but is also significantly induced by the prompt engineering process itself.

The interplay between calibration, accuracy, and adversarial robustness creates a complex web of model reliability. As noted in the context of prompt sensitivity, a critical finding is that predictions which are sensitive to perturbations in the prompt are less likely to be correct, demonstrating a strong negative correlation between ICL sensitivity and accuracy \cite{ref33}. This insight has motivated selective prediction methods that abstain from making predictions when sensitivity is high, thereby improving the overall reliability of the system. The connection between adversarial robustness and calibration runs deeper: inputs for which a model is sensitive to adversarial perturbations are more likely to have poorly calibrated predictions \cite{ref420}. This relationship implies that a model's worst-case reliability under attack is intertwined with its average-case reliability in expressing uncertainty. Methods that improve adversarial robustness, such as adaptive label smoothing based on a sample's adversarial vulnerability, have been shown to also improve model calibration, demonstrating a synergistic rather than adversarial relationship between these two objectives \cite{ref420}.

This synergy, however, is not universal and can be exploited. A new class of threats, termed ``calibration attacks,'' has been introduced, where adversaries specifically aim to miscalibrate a model without altering its accuracy, thereby destroying its trustworthiness while maintaining a facade of performance \cite{ref421}. These attacks can render any decision-making based on confidence scores dangerously unreliable, expanding the surface area of model vulnerabilities beyond simple input manipulations. The interplay is further complicated by model explanations; leveraging model explanations to identify non-inductive attributions can improve calibration, as it allows the model to be less confident when token attributions are not clearly indicative of any class \cite{ref422}. This convergence of interpretability, robustness, and calibration research underscores that a holistically reliable ICL model must be accurate, resistant to adversarial manipulation, and capable of expressing its uncertainty in a manner that is both faithful to its internal state and interpretable to the end-user. The path toward trustworthy ICL, therefore, involves not just recalibrating outputs or desensitizing predictions in isolation, but understanding and correcting the deeply intertwined mechanisms of sensitivity, overconfidence, and adversarial vulnerability that emerge within the context window.

\subsection{Sensitivity to Prompt Components and Label Noise}

Building directly on the understanding that ICL predictions are deeply sensitive to prompt design, a critical barrier to reliable deployment is the pronounced fragility of in-context learning (ICL) to variations in prompt design, a phenomenon that undermines the consistency and trustworthiness of model predictions. While the previous section on calibration examined how this fragility manifests as miscalibration and overconfidence, this section analyzes its roots in the fundamental sensitivity to prompt configuration. This sensitivity manifests across multiple axes of the prompt, including the selection and ordering of demonstrations, the phrasing of templates, and the presence of noise or biases in the provided labels. A systematic investigation into these factors reveals that ICL performance is not solely a function of the model's underlying task capabilities but is deeply entangled with the specific prompt configuration, leading to high variance and unstable generalization \cite{ref32}.

The sensitivity to the choice and ordering of in-context examples is a well-documented source of this instability, directly setting the stage for the adversarial demonstration attacks discussed in the following subsection. Models can exhibit significant performance swings based on which demonstrations are included, even when all examples are correctly labeled. This is partly attributed to the model's reliance on spurious correlations or shallow heuristics within the provided examples, rather than consistently inferring the true task. A key insight from recent work is the strong negative correlation between a model's sensitivity to prompt perturbations and its predictive accuracy: predictions that are less robust to changes in the prompt are more likely to be incorrect \cite{ref33,ref423}. This relationship suggests that sensitivity itself can serve as an unsupervised proxy for model confidence or correctness, echoing the calibration challenges where confidence and accuracy become decoupled. Building on this, an analysis of the auto-regressive attention mask in causal language models shows how the position of information within the prompt creates a recency bias, further contributing to order-dependent instability \cite{ref1}.

The sensitivity of ICL extends profoundly to the label information within demonstrations, a vulnerability that adversarial attacks can exploit by manipulating demonstration labels. A counter-intuitive finding from early ICL research suggested that the correctness of input-label mappings was less critical than expected, but deeper analysis using novel metrics like Label-Correctness Sensitivity and Ground-truth Label Effect Ratio (GLER) clarifies that the impact of ground-truth labels is highly context-dependent \cite{ref424}. The verbosity of the prompt template and the scale of the language model itself act as controlling factors that modulate the model's resilience to noisy labels. This nuanced view is complemented by the identification of different categories of label bias that can skew model predictions. A comprehensive typology defines three such biases: vanilla-label bias, context-label bias, and the previously undetected domain-label bias, which can force models to near-random performance regardless of the demonstrations chosen \cite{ref177}. The presence of majority label bias, where a skewed distribution of labels in the demonstrations leads the model to favor the majority class, is another critical challenge. The robustness of LLMs to this specific bias varies widely depending on the model architecture, its size, and the richness of the instructional prompts provided \cite{ref40}.

The model's prior knowledge, as shaped by pretraining, creates strong inductive biases that interact with the prompt to determine the final sensitivity. Even when demonstrations are underspecified---where two features are equally predictive of the correct label---LLMs exhibit clear and sometimes difficult-to-overcome feature biases, such as preferring sentiment over shallow lexical patterns \cite{ref425}. This interaction between the model's intrinsic priors and the provided prompt context means that the same demonstration set can lead to different behaviors depending on the specific template phrasing. A holistic evaluation of these factors confirms that spurious correlations between input distributions and labels, while a known issue in traditional fine-tuned models, form only a minor problem for prompted models compared to the massive impact of template design and other prompt-level choices \cite{ref37}. These findings underscore that the challenge is not just about what is shown to the model, but how it is shown, a principle that extends directly to the adversarial prompt manipulations and calibration attacks explored subsequently.

To mitigate this pronounced sensitivity and preempt the security vulnerabilities it creates, a range of calibration and selective prediction methods have been developed. The understanding that sensitivity is a negative proxy for accuracy has led to the creation of selective prediction frameworks. The proposed method SenSel, for instance, explicitly abstains from making predictions that are deemed sensitive to perturbations, outperforming standard confidence-based and entropy-based abstention methods \cite{ref33}. This approach is complemented by sensitivity-aware decoding, which incorporates a sensitivity estimation as a penalty term during the standard greedy decoding process, proving particularly beneficial when the input information is scarce \cite{ref423}. A more direct approach to combating label-specific biases is the domain-context calibration method, which estimates a language model's label bias using random in-domain words from the task corpus and then calibrates predictions accordingly, yielding substantial performance gains on tasks with high domain-label bias \cite{ref177}. These methods represent a crucial shift from merely optimizing for average accuracy to building systems that are robustly consistent, acknowledging that a correct prediction is only valuable if it can be reliably reproduced across reasonable prompt variations. The ultimate goal is to dissociate the model's task understanding from the brittle surface forms of the prompt, ensuring that ICL serves as a reliable inference paradigm rather than a fragile prompt-engineering exercise---a foundational requirement for defending against the adversarial exploitation of these same sensitivities that will be examined next.

\subsection{Adversarial Attacks on In-Context Learning}

While the preceding discussion centered on how benign prompt variations and model priors undermine stability, a more deliberate threat emerges when an adversary actively crafts inputs to compromise model behavior. The vulnerability of deep learning models to adversarial examples---inputs intentionally perturbed to cause misclassification---has been extensively documented, and in-context learning (ICL) introduces a unique and critical attack surface that extends the fragility concerns into the security domain \cite{ref426,ref427}. Unlike traditional fine-tuning, where an adversary might need to poison the training data or craft input-specific perturbations, ICL's reliance on externally provided, untrusted prompts creates opportunities for a new class of attacks that target the prompt's structure, its demonstrations, or the model's calibration. This subsection surveys the landscape of adversarial threats specific to ICL, encompassing direct manipulation of demonstrations, prompt-level backdoors, and more subtle calibration attacks that undermine trust without reducing accuracy. These threats represent a deliberate exploitation of the same prompt sensitivity discussed earlier, transforming an instability concern into a concrete security vulnerability.

A direct and potent threat comes from adversarial manipulation of the in-context demonstrations themselves. The work by \cite{ref428} introduced the advICL attack, which manipulates only the demonstrations prepended to a query, without altering the test input. They found that as the number of poisoned demonstrations increases, the robustness of ICL can paradoxically decrease, and that these adversarial demonstrations are transferable, meaning an attacker can craft a generic poisoned prompt that misleads the model on unseen test inputs. This is particularly dangerous because it aligns with a realistic threat model where a service provider or a compromised data source supplies a prompt with malicious examples. Similarly, \cite{ref429} proposes a transferable attack that learns imperceptible, gradient-based adversarial suffixes appended to demonstrations. This attack hijacks the LLM's attention mechanism, distracting it towards adversarial tokens to generate a targeted, unwanted output, effectively allowing an attacker to control the model's response through a poisoned few-shot prompt. These attacks highlight that ICL's performance is heavily dependent on the integrity of the demonstration set, which, if compromised, can become a primary vector for adversarial influence. This extends the earlier observation about sensitivity to demonstration choice and ordering into a security context, where the threat is no longer random variance but targeted manipulation. The problem extends beyond language-only scenarios. In the visual domain, prompt learning, closely analogous to ICL, is also vulnerable. \cite{ref430} introduced BadVisualPrompt, demonstrating that a malicious visual prompt provider can implant a backdoor that is activated by a specific trigger pattern, achieving very high attack success rates with minimal impact on clean accuracy. This work further identified that traditional backdoor defenses from the model, prompt, or input levels are often ineffective or impractical against such attacks, revealing a critical vulnerability in the visual prompt as a service (VPPTaaS) paradigm.

Another significant class of threats targets the calibration of a model's confidence rather than its accuracy, directly weaponizing the miscalibration and sensitivity issues discussed in the previous subsection. These calibration attacks are insidious because they break the trustworthiness of a model's predictions without necessarily changing the predicted label, making them harder to detect. \cite{ref421} formalizes this threat, introducing a framework for generating four types of calibration attacks: underconfidence, overconfidence, maximum miscalibration, and random confidence attacks. These attacks can succeed in both black-box and white-box settings with a relatively low query budget, significantly damaging the reliability of confidence scores, which are crucial for risk assessment, selective prediction, and downstream decision-making. This directly undermines methods like SenSel, which rely on sensitivity as a proxy for correctness, by corrupting the very confidence estimates upon which selective prediction depends. The danger of miscalibration is further amplified by the phenomenon of the ``illusion of adversarial robustness'' \cite{ref431}. This work reveals that a deliberately or accidentally miscalibrated model can mask gradients, causing standard gradient-based adversarial attack search methods to fail and creating a false appearance of high robustness. An adversary aware of this can simply apply test-time temperature scaling to recalibrate the model and nullify this illusion, exposing its true vulnerability. This finding is a stark warning for the evaluation of ICL systems, urging the integration of test-time calibration into robustness benchmarks to ensure that any measured robustness gains are genuine and not artifacts of extreme overconfidence or underconfidence.

The interplay between adversarial robustness and ICL-specific phenomena also introduces novel threat vectors. For instance, the use of chain-of-thought (CoT) prompting, while beneficial for performance, can produce overconfident and miscalibrated predictions \cite{ref38}. This miscalibration can be exploited by an adversary to make the model confidently wrong, a particularly hazardous failure mode in high-stakes reasoning tasks. Moreover, attackers can craft prompts that fool the model into self-attacking. The \cite{ref432} (PromptAttack) framework demonstrates that an LLM can be induced by a carefully designed attack prompt to output an adversarial sample that fools its own prediction mechanism, effectively turning the model against itself through in-context instructions.

Defending against these multifaceted adversarial threats requires a combination of evaluation rigor and novel defensive strategies that build upon the calibration and selective prediction methods introduced earlier. From an evaluation standpoint, the findings on calibration attacks and the illusion of robustness mandate the use of reliable, standardized evaluation protocols. As championed by benchmarks like \cite{ref433} and new attack ensembles \cite{ref434}, robustness evaluations must account for gradient masking and miscalibration, often by using adaptive attacks that incorporate test-time calibration, such as temperature scaling, to pierce the illusion \cite{ref431}. For backdoor and demonstration attacks, defenses might involve certifying the integrity of demonstration sources, developing robust prompt formatting that is less susceptible to trigger patterns, or using test-time detection mechanisms that identify poisoned examples before they influence the model's output. In the realm of visual prompts, research like \cite{ref435} shows that attacks can be designed to survive even fine-tuning-based defenses, necessitating the development of more robust, defense-aware training paradigms. In the context of text-based ICL, methods like confidence-calibrated adversarial training, which biases models towards low confidence on adversarial examples, offer a pathway to generalizing robustness to unseen attacks by enabling a reject option based on low confidence \cite{ref436}. These defensive strategies form a natural counterpart to the privacy protections discussed in the following subsection, as both domains grapple with the challenge of securing a model whose primary interface---the prompt---is inherently untrusted and exposed. Ultimately, securing ICL systems against adversarial attacks involves recognizing the prompt as a critical and exposed attack surface, dynamically verifying its contents, and rigorously evaluating model robustness in a calibrated, attack-aware manner to prevent illusionary security and ensure genuine, trustworthy adaptation.

\subsection{Privacy Leakage and Membership Inference Attacks in ICL}

Extending the security analysis from adversarial manipulation to unauthorized information disclosure, the same reliance on untrusted, ephemeral prompts that creates attack vectors for adversaries also introduces fundamental privacy risks. In-context learning (ICL) fundamentally operates by exposing a model to a prompt that may contain proprietary, personal, or otherwise sensitive information, raising critical questions about what the model memorizes and what can subsequently be extracted or inferred. While the previous subsection detailed how adversaries can actively inject malicious content into prompts, this subsection examines how passive adversaries can exploit that same prompt interface to extract information, weaponizing the model's sensitivity to its input rather than its manipulability. The primary threat vectors for privacy leakage in ICL stem from the prompt itself and the demonstrations it contains. Unlike fine-tuning, where sensitive data is embedded into model weights through gradient updates, ICL presents information ephemerally at inference time, but this does not guarantee immunity from privacy attacks. Indeed, the risks manifest through several distinct mechanisms: the potential for models to memorize and later regenerate prompt content, the susceptibility of models to membership inference attacks that determine whether specific data was part of the prompt or the pretraining corpus, and more subtle forms of property inference that reveal aggregate characteristics of the data.

A central concern, which directly parallels the adversarial focus on prompt integrity in the prior subsection, is the vulnerability of prompts to membership inference attacks (MIAs), a threat extensively studied in the context of machine learning models \cite{ref437,ref438}. These attacks aim to determine whether a particular data record was part of a model's training data, but in the ICL setting, the attack surface expands to include the prompt demonstrations themselves, meaning the very same demonstrations an adversary might poison could also be targeted for extraction. Research on visual prompt learning has demonstrated that learned prompts are vulnerable to both property inference and membership inference, with adversaries capable of mounting successful attacks even with limited resources and relaxed assumptions \cite{ref439}. This finding has direct parallels in the language domain, where prompt-tuning---a process that learns continuous prompt embeddings while keeping the language model frozen---has been shown to leak private user information. A proposed attack framework exploits the prompt module with user-specific signals to infer sensitive attributes, revealing that prompt-tuning does not inherently ensure privacy \cite{ref440}.

The comparative vulnerability of ICL to fine-tuning-based adaptation methods has been systematically investigated, revealing a nuanced landscape where no single adaptation technique is universally more secure \cite{ref441}. This study evaluated Low-Rank Adaptation (LoRA), Soft Prompt Tuning (SPT), and ICL against membership inference, backdoor, and model stealing attacks, finding that each technique exhibits different strengths and weaknesses across the threat landscape. This lack of a uniformly superior approach echoes the adversarial domain, where the choice of defense must also be tailored to the specific threat model. Importantly, the privacy risks of ICL extend beyond the demonstrations provided at inference time to encompass the model's pretraining data, as models can leak information about their training corpus through sentence-level membership inference and reconstruction attacks \cite{ref442}. This creates a dual privacy concern: the prompt data itself may be compromised, and the model may inadvertently reveal pretraining data that surfaces during in-context processing, a vulnerability that compounds the contextual privacy failures discussed later.

Critically, the effectiveness of MIAs against models used in ICL settings has been significantly underestimated by prior evaluations, a problem that extends to the broader privacy literature and mirrors the critique of robustness evaluations suffering from the illusion of security identified in Section 8.3. Empirical privacy defenses have been critiqued for relying on weak attacks, failing to characterize leakage for the most vulnerable samples, and avoiding comparisons with practical differential privacy baselines \cite{ref443}. A systematic evaluation of privacy risks demonstrated that prior work underestimated leakage by an order of magnitude, and that under stronger evaluation protocols, empirical defenses fail to compete with properly tuned DP-SGD baselines \cite{ref444}. This critique is particularly salient for ICL, where the transient nature of in-context data might falsely suggest inherent privacy protection, much like gradient masking created an illusion of adversarial robustness. Label-only membership inference attacks, which require only predicted labels rather than confidence scores, have been shown to perform on par with confidence-based attacks, breaking defenses that rely on confidence masking \cite{ref445}. This has direct implications for ICL systems that expose only final predictions to users, as they remain vulnerable to inference attacks, undermining the output-privatization strategies that will be explored in the following subsection on differential privacy.

The landscape of privacy attacks against language models has evolved to include user inference attacks, where an adversary determines whether a specific user's data was used for fine-tuning, achieving near-perfect success rates in some scenarios \cite{ref446}. While focused on fine-tuning, these findings highlight the vulnerability of models to inferring the presence of individual data sources, a concern that translates to ICL when prompts contain data attributable to specific individuals or organizations. Furthermore, the emergence of privacy backdoors represents a sophisticated threat where pre-trained models are intentionally poisoned to amplify privacy leakage during downstream adaptation \cite{ref447}. In such attacks, an adversary releases a maliciously crafted pre-trained model; when victims use it for ICL or fine-tuning with their private data, the model leaks information at significantly higher rates. This threat directly parallels the adversarial backdoors in visual prompts discussed in the previous subsection, demonstrating a troubling convergence between security and privacy attacks that both exploit the supply chain of models and prompts. This is compounded by frameworks like PreCurious, which demonstrate how a seemingly innocent pre-trained model can act as a privacy trap, escalating both membership inference and data extraction risks even when victims employ parameter-efficient or differentially private techniques \cite{ref448}.

The problem of information leakage detection has been formalized through statistical learning theory and information theory, establishing that mutual information between observable and secret information can be accurately estimated by approximating the log-loss and accuracy of the Bayes predictor \cite{ref449}. This theoretical framework provides grounding for evaluating privacy risks in ICL systems, though practical application remains challenging. The concept of maximal leakage, an information-theoretic measure, has been proposed as an alternative to differential privacy for bounding information disclosure in adaptive data analysis settings \cite{ref450}, offering potential relevance for analyzing the iterative exchange of information in multi-turn ICL interactions, a setting that will demand more rigorous privacy accounting than single-turn scenarios.

A critical and often overlooked dimension of privacy in ICL is contextual privacy, which concerns whether language models can appropriately reason about what information to share, with whom, and for what purpose within a given context. This directly weaponizes the miscalibration and sensitivity issues explored in Section 8.2 by causing models to leak information not through explicit attacks but through failures in social reasoning. The ConfAIde benchmark was designed to identify weaknesses in the privacy reasoning capabilities of instruction-tuned LLMs, revealing that even state-of-the-art models like GPT-4 and ChatGPT disclose private information in contexts where humans would not, at rates of 39\% and 57\% respectively \cite{ref451}. This leakage persists even when employing privacy-inducing prompts or chain-of-thought reasoning, underscoring the inadequacy of current prompt-based privacy safeguards and connecting directly to the overconfidence issues of CoT prompting discussed in the adversarial context. The risks are further compounded by the fact that models can leak personally identifiable information (PII) even when dataset curation techniques like scrubbing are applied, as scrubbing is imperfect and algorithmic defenses like differential privacy, while reducing risk, still allow approximately 3\% of PII sequences to leak \cite{ref442}.

The sensitivity of membership inference has been shown to follow a power-law dependence on the number of examples per class, with large gradients at the end of training strongly correlated with vulnerability \cite{ref452}. While this finding emerged from image classification studies, the underlying principles regarding memorization and gradient-based leakage provide insights into when and why ICL demonstrations might be most vulnerable to extraction or inference, informing the defensive allocation of privacy budgets that follows. Model explanations, often touted as transparency mechanisms, can paradoxically increase privacy risk, as backpropagation-based explanations have been found to leak significant information about individual training datapoints by revealing statistical properties of decision boundaries \cite{ref453}. This creates a tension between explainability and privacy that is particularly acute in ICL, where explanations might be provided alongside demonstrations to improve performance but could simultaneously expose sensitive prompt content, a trade-off that the differential privacy mechanisms in the next subsection must navigate.

In summary, the privacy vulnerabilities of ICL are multifaceted, encompassing direct leakage from prompts, susceptibility to membership and property inference attacks, risks from poisoned pre-trained models, and failures in contextual privacy reasoning. The comparison between ICL and fine-tuning reveals that neither approach is categorically more private, but their vulnerability profiles differ, with ICL's inference-time data exposure creating unique attack surfaces that demand the specialized defenses and rigorous evaluation methodologies discussed subsequently. These vulnerabilities, ranging from the extractive threats that exploit prompt memorization to the contextual failures that mirror the miscalibration issues of earlier sections, collectively motivate the principled privacy-preserving frameworks centered on differential privacy that form the core of the following subsection.

\subsection{Differential Privacy Mechanisms for In-Context Learning}

Building on the multifaceted privacy vulnerabilities outlined above---ranging from membership inference and prompt leakage to contextual privacy failures---differential privacy (DP) has emerged as the gold standard for providing formal, quantifiable privacy guarantees, and its integration into in-context learning (ICL) represents a critical frontier for enabling trustworthy deployment of LLMs on sensitive data. The core challenge is that ICL fundamentally relies on the inclusion of exemplar data directly within the prompt, creating a direct channel for potential privacy leakage that compounds the risks discussed in the previous section. The LLM's responses can memorize and regurgitate sensitive information contained in those exemplars \cite{ref454}, making the ephemeral nature of ICL insufficient as a privacy safeguard. Existing research has thus focused on adapting DP mechanisms to the unique ICL paradigm, where traditional approaches like DP-SGD are inapplicable because there is no gradient-based parameter update. Instead, privacy-preserving frameworks for ICL must operate on the data input or the model output, leading to a novel taxonomy of perturbation points that addresses the specific attack surfaces identified earlier.

A foundational framework for privatizing ICL is the DP-ICL paradigm, which achieves differential privacy through a noisy consensus mechanism among an ensemble of LLM responses \cite{ref454}. The key insight is to partition the available in-context exemplars into disjoint subsets, generate independent predictions or text generations from the LLM conditioned on each subset, and then aggregate these outputs via a differentially private mechanism. For text classification, this aggregation can involve a noisy majority vote, while for language generation, it can be a noisy selection among the generated candidate texts. This approach effectively privatizes the response at the output perturbation level, and it has demonstrated a strong utility-privacy tradeoff on both classification and generation benchmarks, highlighting that ICL can be made private without catastrophic performance degradation. Importantly, this ensemble-based strategy directly mitigates the membership inference risks detailed in the prior section by ensuring that no single exemplar deterministically dictates the model's output.

A parallel line of work focuses on protecting the data itself before it is even placed into the prompt, particularly for structured data like tabular records. Frameworks such as DP-TabICL formalize the application of DP to tabular ICL by privatizing the data prior to its serialization into natural language \cite{ref455}. This work investigates both local differential privacy (LDP) and global differential privacy (GDP) scenarios, directly engaging with the tension between privacy and fairness that will be explored in depth in the following section. In the LDP setting, noise is injected into individual records before they are shared with the LLM, while in the GDP setting, noise is added to aggregate statistics of the demonstration set. The empirical results from these frameworks show that DP-based ICL can protect the privacy of underlying tabular data while achieving comparable performance to non-LLM baselines, especially under stringent privacy budgets. This is crucial for domains like healthcare and finance, where tabular data is prevalent and highly sensitive, and where privacy regulations often mandate formal guarantees that go beyond the empirical defenses critiqued earlier.

The aforementioned approaches highlight the need for a holistic understanding of where to apply perturbation in the ICL pipeline. A comprehensive utility analysis framework, DP-UTIL, provides a systematic way to evaluate the impact of DP across different perturbation points: input, objective, gradient, output, and prediction \cite{ref456}. While not all are directly applicable to ICL's gradient-free nature, the framework's analysis of input perturbation (as in DP-TabICL) and output prediction perturbation (as in DP-ICL's noisy consensus) is directly relevant. One of the key findings from such holistic analyses is that prediction perturbation consistently achieves the lowest utility loss across various models and datasets, which provides a strong theoretical backing for the output-aggregation strategies employed in DP-ICL. Furthermore, the interaction between privacy and utility is not uniform; the choice of perturbation mechanism must be carefully matched to the specific ICL task and data modality, as the dynamics between DP mechanisms, the number of classes, and the privacy budget significantly influence the final outcome.

A particularly challenging sub-problem is differentially private text generation, which is central to ICL when the task involves generating long-form output rather than selecting a label. This challenge is exacerbated by the contextual privacy failures identified in the prior section, where models leak information even when prompted to maintain confidentiality. Private text rewriting mechanisms, such as those built on DP guarantees, often struggle with the high-dimensional, discrete nature of text and the fixed global sensitivity of noise addition, leading to sub-optimal privacy-utility trade-offs \cite{ref457}. The problem is compounded in ICL, where the generated text is conditioned on private exemplars. A simple and practical recipe has been shown to be effective: fine-tuning a pretrained generative language model with DP and then using it for synthetic text generation, producing high-quality text with strong privacy protection \cite{ref458}. This approach, while involving a training step, points to a hybrid future where ICL might be performed by a model that has been pre-privatized for certain tasks. The challenge, however, is that reparametrization tricks and implementation errors can easily break the formal DP guarantees in text representation learning, underscoring the need for rigorous empirical sanity checks and reproducible frameworks like DP-Rewrite to validate privacy claims \cite{ref459,ref460}. This concern for rigorous evaluation echoes the broader critique from the prior section that empirical privacy defenses often rely on weak attacks and misleading assessments.

Another critical dimension of DP in ICL is the relationship between privacy and other trustworthiness pillars, particularly fault tolerance and calibration. The concept of ``budget recycling'' introduces a framework, BR-DP, which allows for soft-bounded noisy outputs, improving utility by probabilistically recycling out-of-bound noisy results without violating the DP guarantee \cite{ref461}. This mechanism is inherently fault-tolerant, as it can gracefully handle the excessive noise that sometimes causes DP mechanisms to produce unusable outputs. In the context of ICL, where a single very noisy response could break the user's trust, a budget recycling mechanism can enhance the reliability and robustness of the private LLM. This aligns with the goal of creating holistically trustworthy ICL systems, where privacy, utility, and robustness are not isolated objectives but are jointly optimized.

The interplay between fairness, bias, and privacy also emerges as a central concern, bridging the privacy vulnerabilities of this section with the fairness challenges that follow. DP mechanisms can sometimes exacerbate bias and unfairness for different groups, an issue analyzed under the lens of differentially private empirical risk minimization \cite{ref462}. When applying DP to ICL, the strong priors in the LLM can interact with the noise added for privacy to produce disparate impacts on prediction accuracy across different demographic groups. This necessitates a ``fair and wise'' allocation of the privacy budget, where the noise injected is not uniform but is distributed across features or exemplars to maximize predictive accuracy while bounding potential disparities \cite{ref463}. Such strategies are essential for ensuring that privacy-preserving ICL systems do not inadvertently sacrifice equity for the sake of confidentiality, a theme that will be explored comprehensively in the following section on fairness and privacy.

In summary, the integration of DP into ICL is a multi-faceted endeavor that spans data privatization, output privatization, and the careful balancing of the utility-privacy-fairness tradeoff. The primary mechanisms involve adapting output perturbation through noisy consensus ensembles \cite{ref454}, input perturbation for structured data serialization \cite{ref455}, and hybrid approaches that leverage DP-trained models for text generation \cite{ref458}. The field is moving towards holistic frameworks that analyze the entire ICL pipeline for privacy vulnerabilities \cite{ref456} and towards mechanisms that are not only private but also robust and fair \cite{ref461,ref463}. The unique challenges of private text generation and the critical need for sound implementation practices remain key areas of active investigation, as the practical deployment of private ICL hinges on moving beyond theoretical guarantees to verifiable, end-to-end privacy-preserving systems. These efforts collectively address the privacy attack vectors identified earlier while setting the stage for the integrated fairness-privacy perspective developed in the subsequent section.

\subsection{Fairness, Bias, and the Interplay with Privacy in ICL}

Building on the multifaceted privacy vulnerabilities and the differential privacy (DP) frameworks explored in the previous section, this subsection turns to the parallel and deeply intertwined challenge of social bias and fairness in In-Context Learning (ICL). ICL inherits the powerful priors encoded within large language models (LLMs) during pretraining, yet these same priors can introduce or amplify social biases related to gender, race, and religion. When sensitive attributes are directly included in demonstrations, models may rely on spurious correlations, leading to discriminatory predictions. Even when sensitive attributes are omitted, strong priors can cause LLMs to infer and act upon protected characteristics from correlated proxy features \cite{ref464}. This presents a fundamental challenge for trustworthy ICL that mirrors the privacy challenges discussed earlier: the need to achieve fairness while simultaneously protecting the privacy of sensitive information. Addressing this dual objective is critical, as privacy regulations and ethical considerations often prevent the collection or use of sensitive attributes, creating a semi-private setting where fairness must be achieved without full access to the data needed for conventional debiasing \cite{ref465,ref466}. This tension directly extends the fairness-privacy nexus introduced in the prior section's discussion of budget allocation and disparate impacts.

The interplay between fairness and privacy in machine learning creates a complex tension, and ICL introduces unique dimensions to this relationship. Work on fair classification has explored how models can be debiased when sensitive attributes are partially or entirely unavailable, often by leveraging domain adaptation to transfer fairness knowledge from a source domain with sensitive labels \cite{ref467} or by inferring a sensitive proxy through causal variational autoencoders \cite{ref466}. These methods, however, do not provide the formal privacy guarantees that DP-based frameworks like DP-ICL and DP-TabICL offer. When differential privacy is introduced to protect sensitive attributes, new trade-offs emerge that directly connect to the perturbation point analysis from DP-UTIL discussed previously. A central finding is that DP does not necessarily worsen fairness and can, under certain conditions, even reduce disparities. Empirical studies on local differential privacy (LDP) for multi-dimensional sensitive attributes demonstrate that LDP can act as an efficient approach to reduce group disparity, with the method of combining independent versus correlated noise mattering primarily at low privacy budgets \cite{ref468,ref469}. This challenges the common belief that privacy mechanisms inevitably exacerbate bias and provides empirical grounding for the theoretical promise of jointly optimizing privacy and fairness.

The amplification of bias by DP has been observed, however, particularly in the fine-tuning of LLMs, echoing the concerns about implementation errors and rigorous evaluation raised in the context of private text generation. It has been shown that DP-SGD can amplify gender, racial, and religious bias, producing models that are more biased than non-private counterparts due to a disparity in gradient convergence across subgroups \cite{ref470}. This finding is directly relevant to ICL, as the pretrained models used for prompting may have been fine-tuned with DP, and the in-context demonstrations themselves can interact with these pre-existing biases, much as they can leak memorized information. Fortunately, established fairness interventions like Counterfactual Data Augmentation (CDA) can simultaneously mitigate this DP-induced bias amplification, suggesting that fairness and privacy objectives can be jointly addressed rather than traded off \cite{ref470}. The integration of privacy and fairness is further demonstrated in representations learned for text, where combining DP with adversarial training for fairness can lead to private text encoders that also induce fairer downstream models, indicating a potential positive reinforcement between the two objectives that parallels the budget recycling mechanisms for robustness and privacy \cite{ref471}.

A primary framework for achieving fairness under privacy constraints, extending the data privatization strategies of DP-TabICL, is the use of local differential privacy (LDP) for fair representation learning. The principle of information obfuscation, grounded in an information bottleneck (IB) framework with LDP randomizers, provides theoretical justification for this synergy. By formalizing the encoding process, it can be shown that the disclosure of sensitive information is bounded by the privacy budget of the LDP randomizer. This constraint forces the optimization process within the IB framework to suppress sensitive information while preserving utility-relevant features, effectively using obfuscation as a tool for fairness \cite{ref472}. A key advantage of this approach is that it can be implemented with a non-adversarial variational encoding method that does not require a predefined prior, making it practically attractive for ICL settings where adversarial training may be unstable or computationally expensive, and aligning with the ensemble-based output privatization of DP-ICL in providing practical deployment pathways.

For classification tasks within ICL, frameworks have been developed to achieve differential privacy and fairness simultaneously, providing the theoretical underpinnings for the holistic perturbation analysis advocated by DP-UTIL. The calibrated functional mechanism adds Laplace or Gaussian noise directly to the polynomial coefficients of the objective function, with the noise calibrated to account for fairness constraints, enabling a jointly private and fair classifier with provable guarantees \cite{ref473}. Stochastic optimization algorithms extend this capability to large-scale settings, providing the first differentially private stochastic algorithm for fair learning that converges with minibatches and supports multiple fairness notions like demographic parity and equalized odds \cite{ref474}. The challenge of designing machine learning pipelines that do not implicitly favor one objective over another has led to the principle of impartiality, where models are trained to recover the full Pareto frontier between fairness, privacy, and utility, allowing practitioners to make informed trade-offs rather than accepting a single, biased solution \cite{ref475}. This impartiality principle directly addresses the compositional challenges identified in the prior section, where privacy mechanisms and fairness interventions can interact in unforeseen ways.

The practical application of these theoretical frameworks to ICL synthesizes the mechanisms introduced in the previous section with fairness-aware design. The DP-TabICL framework demonstrates a concrete approach for protecting tabular data used in ICL by applying DP mechanisms for data privatization prior to serialization and prompting, with instantiations in both local and global DP scenarios that can be audited for fairness impacts \cite{ref455}. For general ICL tasks, the DP-ICL paradigm generates differentially private responses through a noisy consensus mechanism among an ensemble of LLM responses based on disjoint exemplar sets, providing a principled way to privatize both text classification and language generation tasks in ICL, and this ensemble diversity can additionally mitigate bias by preventing over-reliance on any single potentially biased exemplar \cite{ref454}. Similarly, generating synthetic few-shot demonstrations from private datasets with formal DP guarantees enables effective ICL while protecting the underlying data, with the opportunity to proactively control for fairness in the synthetic generation process \cite{ref476}. A holistic utility analysis framework like DP-UTIL helps practitioners navigate the perturbation choices across the machine learning pipeline, revealing, for instance, that prediction perturbation often achieves the lowest utility loss, while gradient or objective perturbation can minimize privacy leakage depending on model convexity, with fairness considerations adding an additional dimension to this optimization \cite{ref456}.

The relationship between fairness and privacy is not always harmonious, underscoring a critical lesson that will be central to the holistic trustworthiness discussion in the following section. Algorithmic fairness interventions can themselves increase privacy risks. Analyzing the privacy cost of group fairness through membership inference attacks reveals that fair models can leak more information about their training data, and this cost is disproportionately borne by unprivileged subgroups---the very groups fairness aims to protect \cite{ref477}. The more biased the training data, the higher the privacy cost of achieving fairness for these subgroups, creating a pernicious feedback loop. This finding has profound implications for ICL, where demonstrations are drawn from potentially biased historical data and where the membership inference vulnerabilities detailed in the prior section can be exacerbated by fairness interventions. Data imbalances, even small ones, can cause disparate impacts under DP, underscoring the need for careful fairness auditing when deploying private ICL systems \cite{ref478}. These findings directly motivate the integrated evaluation frameworks and compositional perspectives that will be explored in the next section.

Emerging research suggests that the negative interaction between DP and fairness is not inevitable, offering a constructive path forward that bridges the privacy mechanisms of the previous section with the holistic approaches to follow. Rigorous empirical analysis has shown that selecting an optimal model architecture for DP-SGD can result in fairness metrics that decrease or are negligibly impacted compared to non-private baselines, challenging the narrative that DP necessarily exacerbates unfairness \cite{ref479}. Furthermore, pre-trained foundation models fine-tuned with DP can achieve similar accuracy and fairness to non-private classifiers, even under significant distribution shifts, marking a milestone toward practical and reliable private training \cite{ref480}. For ICL, these findings indicate a path forward that integrates the lessons from both privacy-preserving and fairness-aware methods: by carefully designing the prompting pipeline, selecting demonstrations with fairness and privacy in mind, and leveraging techniques like flatness-guided optimization to improve the privacy-generalization trade-off \cite{ref481}, it is possible to build in-context learning systems that are simultaneously fair, private, and accurate. The future of trustworthy ICL lies in frameworks that jointly address robustness, calibration, fairness, and privacy as interconnected pillars rather than isolated objectives, setting the stage for the integrated, holistic perspective developed in the subsequent section.

\subsection{Toward Holistically Trustworthy In-Context Learning Systems}

The preceding subsections have delineated critical challenges to the trustworthiness of In-Context Learning (ICL), spanning miscalibration, adversarial vulnerabilities, privacy leakage, and biases. As the discussion of fairness and privacy in the prior section made clear, these pillars do not operate in isolation; privacy mechanisms can either amplify or mitigate bias, and fairness interventions can incur privacy costs. A growing consensus in the broader AI community, reflected across multiple foundational papers, therefore asserts that true trustworthiness is not merely the sum of isolated solutions but an emergent property of a system where robustness, calibration, fairness, privacy, and explainability are deeply interconnected and holistically managed. Synthesizing these dimensions is paramount, as interventions targeting one pillar often introduce unforeseen trade-offs or, conversely, can be strategically aligned for mutual reinforcement. This section surveys the theoretical motivations and emerging methodologies for building holistically trustworthy ICL systems, moving beyond siloed evaluations towards integrated frameworks.

The primary impetus for a holistic approach arises from the empirically observed and theoretically grounded interactions between trustworthiness pillars. For instance, the relationship between calibration and fairness is a focal point of recent research. The widely recognized incompatibility of certain fairness criteria, such as calibration within groups and balance for positive/negative classes, necessitates methods that can navigate these trade-offs. The FAIM algorithm proposed by \cite{ref482} directly addresses this by continuously interpolating between competing fairness metrics, allowing for a weighted, context-dependent optimization that is crucial for real-world ICL deployments where legal and ethical definitions of fairness may conflict. Similarly, the interplay between uncertainty and fairness is highlighted by \cite{ref483}, which extends multi-calibration---a fairness concept---to the domain of conformal prediction, yielding a fairness notion for uncertainty quantification. This unified approach demonstrates that predictive coverage guarantees can and should be distributed equitably across demographic subgroups, directly linking the calibration and fairness pillars within a single framework.

The composition of these pillars becomes even more complex when adversarial robustness, privacy, and explainability are added. An empirical study on the composition of functions for trustworthiness by \cite{ref484} reveals that sequentially applying interventions from different pillars (e.g., a bias mitigation algorithm followed by a post-hoc explainability method) leads to non-trivial outcomes. The final fairness and fidelity of explanations are not a simple sum of the individual components; instead, compositions can produce emergent behaviors where a privacy transformation might inadvertently undermine the effectiveness of an adversarial defense. This necessitates a compositional perspective, where ICL pipelines are evaluated from end-to-end rather than by the performance of individual components in isolation. The relationship between fault tolerance and adversarial robustness, explored in \cite{ref485}, further complicates the landscape, showing that while techniques like differential privacy can enhance fault tolerance, they can be at odds with measures designed for adversarial robustness, highlighting a critical trade-off specific to safety-critical ICL applications.

To navigate this interconnected space, several unified evaluation frameworks have been proposed. The ComplAI framework \cite{ref486} operationalizes this by providing a single ``Trust Factor'' score that aggregates performance across explainability, robustness, fairness, and model drift. This model-agnostic approach enables direct comparison of different ICL models not just on accuracy, but from a holistic responsibility perspective, facilitating actionable recourse. In the data-centric realm, the holistic auditing framework from \cite{ref487} introduces a trustworthiness index to rank synthetic datasets, an approach directly applicable to the in-context demonstrations that underpin ICL. By auditing demonstration sets for bias, fidelity, and privacy risks, one can curate prompts that inherently promote more trustworthy model behavior before any inference-time intervention. This proactive design philosophy is echoed in a proposed secure development lifecycle for AI \cite{ref488}, which advocates for threat modeling and penetration testing to anticipate and mitigate vulnerabilities at the design stage, a practice urgently needed for ICL systems exposed to prompt injection and data poisoning.

The convergence of these efforts points towards a set of emerging methodologies that jointly address multiple trust dimensions. Conformal prediction stands out as a particularly versatile theory for this integrative role, extending the synergy between fairness, privacy, and uncertainty quantification discussed previously. The framework for private prediction sets developed by \cite{ref489} elegantly unifies rigorous uncertainty quantification with differential privacy, ensuring that prediction sets provided by an ICL model have a statistical coverage guarantee without leaking sensitive information about the calibration data. This concept is extended to distributed and non-exchangeable settings critical for federated ICL applications in \cite{ref490} and \cite{ref491}, which propose and solve the challenge of providing valid, personalized uncertainty estimates when data across clients or demonstrations is heterogeneous. These methods collectively bridge robustness to distribution shift and privacy preservation, a dual requirement for trustworthy ICL in sensitive, decentralized domains like healthcare.

Another promising integrative technique is self-supervised probing, introduced in \cite{ref492}. This approach serves as a flexible ``wrapper'' that can be applied post-hoc to enhance calibration, misclassification detection, and out-of-distribution detection simultaneously. By probing a model's internal representations with self-supervised tasks, it can effectively surface overconfidence and correct it, improving multiple trustworthiness metrics without retraining the core ICL model. This aligns with the finding from \cite{ref493} that calibration directly influences the fidelity of post-hoc explanations, creating a positive reinforcement loop where a better-calibrated ICL model yields both more reliable predictions and more interpretable reasoning.

Ultimately, moving toward holistically trustworthy ICL systems requires a shift in perspective from a patchwork of post-hoc fixes to a principled design philosophy where trustworthiness is a primary optimization objective. This involves not only developing new algorithms but also embracing comprehensive evaluation protocols that measure the Pareto frontier of trade-offs between, for example, accuracy, fairness, and privacy, as demonstrated by \cite{ref494}. Foundational theoretical work that provides a unified information-theoretic view of these trade-offs, such as \cite{ref495}, is essential for analytically characterizing what is fundamentally achievable. The path forward lies in building ICL systems that are not just accurate few-shot learners, but are designed from the ground up to be transparent in their uncertainty, provably robust to manipulation, equitably fair in their outcomes, and respectful of the data they are conditioned on, as envisioned by comprehensive roadmaps for the field \cite{ref496,ref497}.

\section{Applications Across Diverse Domains and Specialized Tasks}

\subsection{Information Extraction and Semantic Parsing}

Information Extraction (IE) constitutes a foundational challenge in Natural Language Processing, aiming to automatically distill unstructured text into structured, machine-readable knowledge. Core sub-tasks such as Named Entity Recognition (NER), Relation Extraction (RE), and Event Extraction (EE) are pivotal for constructing knowledge bases and powering downstream applications \cite{ref498}. The advent of In-Context Learning (ICL) has revolutionized IE by offering a paradigm that bypasses the prohibitive cost of task-specific fine-tuning, instead conditioning powerful Large Language Models (LLMs) on instructions and a handful of demonstrations. This shift is particularly transformative for low-resource domains where annotated data is scarce, enabling competitive performance through prompt engineering alone \cite{ref499}.

Early explorations into zero-shot IE demonstrated the latent potential of LLMs by recasting extraction tasks as multi-turn question-answering problems. For instance, frameworks like ChatIE leveraged the powerful conversational abilities of ChatGPT to decompose complex IE tasks into a two-stage dialogue, achieving impressive results that sometimes surpassed fully supervised models on datasets like NYT11-HRL \cite{ref500}. However, the prevailing trend for enhancing precision involves the ICL paradigm, where a small number of input-output examples are prepended to the query. The structure of these demonstrations is critical; a naive approach using only positive examples neglects the rich signal available in negative instances. To address this, contrastive ICL methods have been proposed. The c-ICL technique enriches prompts by incorporating not only correct samples but also hard negative samples and the reasoning behind them. By exposing the model to contrasting cases, this strategy helps LLMs identify and correct potential errors in entity and relation boundaries, leading to substantial performance improvements across a spectrum of NER and RE tasks \cite{ref501}.

A significant challenge in ICL for IE lies in the inherent complexity and underspecification of task instructions within a limited context window. Traditional prompts often fail to capture the intricate schemas, edge cases, and label semantics that a human annotator would naturally internalize. The Guideline Learning (GL) framework tackles this by introducing a reflective mechanism: during a learning phase, the system automatically synthesizes a set of actionable guidelines by analyzing error cases from an initial pass. These guidelines are then retrieved and integrated into the prompt during inference, effectively bridging the task comprehension gap between humans and LLMs and significantly boosting in-context performance on event and relation extraction \cite{ref502}. This schema-driven extraction paradigm is further realized by structuring outputs. The Unified Structure Generation for Universal Information Extraction framework (UIE) encodes different extraction structures via a structured extraction language and a schema-based prompt mechanism. This allows a single model to adaptively generate targeted structures for NER, RE, EE, and sentiment extraction, demonstrating strong transferability and state-of-the-art results in low-resource and few-shot settings \cite{ref503}. Similarly, the code generation perspective offers a universal solution, where extraction is reframed as generating Python classes that instantiate task schemas. The Code4UIE framework leverages this concept, using retrieval-augmented ICL to provide LLMs with semantically similar code examples, which formalizes the structural knowledge and enables universal extraction across diverse tasks \cite{ref504}.

Another major axis of innovation involves data augmentation and generation to support ICL, especially in low-resource scenarios where even a handful of demonstrations might be unavailable or insufficient. Synthetic data generation via LLMs has proven remarkably effective. The STAR framework operates on a structure-to-text principle, where LLMs first generate target structures (Y) and then generate corresponding passages (X), guided by fine-grained step-by-step instructions. Through iterative self-reflection and refinement, STAR produces high-quality, task-aligned data that can even surpass human-curated data in boosting performance for low-resource event and relation extraction \cite{ref505}. For NER specifically, the ProgGen methodology instructs LLMs to self-reflect on a domain to generate attribute-rich training data. Crucially, it preemptively generates entity terms before developing the surrounding context, cleverly bypassing LLMs' struggles with complex structural generation and yielding superior datasets for both general and niche domains \cite{ref506}.

The principles of ICL for IE extend seamlessly into semantic parsing, particularly the complex task of translating natural language questions into executable SQL queries. Text-to-SQL parsing requires models to ground linguistic utterances to a structured schema, a task that benefits immensely from well-chosen in-context demonstrations. Systematic studies reveal that prompt design strategies significantly impact LLM performance. Leveraging the syntactic structure of SQL queries to retrieve demonstrations that balance diversity and similarity leads to enhanced parsing accuracy. Furthermore, augmenting prompts with database-related knowledge enables LLMs to better understand schema constraints, outperforming specialized fine-tuned systems on benchmarks like Spider and demonstrating the adaptability of LLMs to structured knowledge tasks \cite{ref507}. The underlying philosophy connecting IE and semantic parsing is the translation of augmented natural languages. The TANL framework treats all structured prediction tasks, including joint entity-relation extraction and coreference resolution, as a translation problem between a textual input and an augmented output containing task-specific markers. By treating all tasks homogeneously, this approach achieves state-of-the-art results even when a single model is trained to solve multiple tasks simultaneously, with particularly strong gains in low-resource settings due to improved exploitation of label semantics \cite{ref508}.

These foundational advances in general-domain IE directly inform and enable the specialized applications discussed next. The same principles of contrastive demonstration selection, retrieval-augmented generation, and synthetic data creation are acutely relevant in high-stakes fields like biomedicine, where the scarcity of annotated data and the complexity of specialized vocabularies pose even greater challenges. Despite the promise, challenges remain. Many sophisticated IE models require comprehensive understanding of sentence structure and domain knowledge, and well-tuned Language Models still often dominate in specific benchmarks \cite{ref499}. Yet, the extraordinary flexibility of ICL allows for the effortless integration of external knowledge. Retrieval-Augmented Generation (RAG) approaches, such as those proposed for relation extraction, combine the power of LLMs with retrieved examples to overcome hallucinations, demonstrating superior performance on established benchmarks like TACRED and SemEval \cite{ref509}. As LLMs continue to advance, their role is evolving from mere extractors to universal, instruction-following data engines. By synergizing contrastive learning, self-reflexive data generation, and structured schema prompting, the field is rapidly overcoming the annotation bottleneck, confirming that ICL is not just a viable alternative but a leading strategy for information extraction in both general and highly specialized domains such as biomedicine, where these techniques find their most critical and demanding application \cite{ref510}.

\subsection{Biomedical and Clinical Text Processing}

Building directly on the foundational ICL strategies for general information extraction, the biomedical domain presents a uniquely challenging and high-impact application area characterized by a severe scarcity of annotated data, a complex and specialized vocabulary, and high-stakes decision-making where errors can have critical consequences. ICL has emerged as a powerful paradigm to address these challenges, enabling large language models (LLMs) to perform sophisticated biomedical text processing tasks without the need for extensive task-specific fine-tuning. This section analyzes the application of ICL across key biomedical tasks, including biomedical concept linking, clinical named entity recognition, adverse drug event extraction, and automated ICD coding, while examining the synergistic roles of retrieval-augmented generation, knowledge base integration, and model distillation---techniques that directly mirror the retrieval-augmented and contrastive selection methods developed for general-domain IE.

For biomedical concept linking, the task of mapping textual mentions of diseases, drugs, or other medical concepts to entries in a standardized knowledge base, ICL methods have demonstrated remarkable effectiveness, echoing the retrieval-augmented approaches in general relation extraction. A two-stage retrieve-and-rank framework leveraging ICL has been shown to achieve competitive performance against fully supervised methods \cite{ref511}. In this approach, biomedical concepts are first embedded using language models, and the top candidates are retrieved via embedding similarity. These candidates, along with their contextual information, are then incorporated into a prompt and processed by an LLM to re-rank the concepts, achieving an accuracy of 90\% on BC5CDR disease entity normalization and 94.7\% on chemical entity normalization. This represents a significant improvement, with an over 20-point absolute increase in F1 score on an oncology matching dataset, highlighting the power of ICL to disambiguate concepts by leveraging the model's pre-trained knowledge of biomedical contexts. The integration of knowledge graphs (KGs) further amplifies these capabilities, paralleling the schema augmentation techniques seen in Text-to-SQL. A task-agnostic Knowledge Graph-based Retrieval Augmented Generation (KG-RAG) framework, which leverages the massive biomedical KG SPOKE, has consistently enhanced the performance of LLMs across various prompt types, including drug repurposing queries and biomedical true/false questions \cite{ref512}. Notably, KG-RAG provided a 71\% boost in the performance of the Llama-2 model on a challenging multiple-choice question dataset, demonstrating the framework's capacity to empower even open-source models for domain-specific questions. This approach effectively combines the explicit knowledge of the KG with the implicit knowledge of the LLM, enhancing the adaptability of general-purpose models to specialized biomedical queries.

In clinical Named Entity Recognition (NER), the task of identifying mentions of clinical concepts like problems, treatments, and tests in unstructured clinical notes, ICL strategies have proven crucial for overcoming data scarcity, much as they do for low-resource IE tasks. Research has revealed that the strategic selection of in-context examples can yield a notable 15-20\% increase in F1 score across benchmark datasets for few-shot clinical NER, strongly resonating with the contrastive demonstration selection principles like c-ICL \cite{ref513}. This underscores the sensitivity of ICL performance to the quality and composition of demonstrations. Beyond simple few-shot prompting, integrating external resources through prompting strategies inspired by Retrieval-Augmented Generation (RAG) can bridge the gap between general-purpose LLM proficiency and the specialized demands of medical NER, boosting the F1 score of LLMs in zero-shot settings \cite{ref513}. The development of unsupervised tools like the Medical Concept Annotation Tool (MedCAT) has also been pivotal, using unsupervised machine learning to disambiguate entities and achieving state-of-the-art performance for large-scale entity extraction \cite{ref514}. While unsupervised learning provides a strong foundation, the limitations of small datasets can be addressed by the active learning and supervised learning extensions supported by such tools, forming a comprehensive ecosystem where ICL can be applied for rapid prototyping and low-resource adaptation, directly reflecting the data generation philosophies of frameworks like STAR and ProgGen.

The extraction of adverse drug events (ADEs) from biomedical literature and clinical notes is another critical task where ICL and distillation synergize powerfully, extending the self-reflexive data generation concepts from general IE. While LLMs like GPT-4 possess a decent competency in structuring biomedical text, their out-of-the-box performance can be significantly improved and made more efficient through distillation into task-specific student models via self-supervised learning \cite{ref515}. This approach, where a GPT-3.5 distilled PubMedBERT model attained comparable accuracy to supervised state-of-the-art models without using any labeled data, demonstrates that ICL capabilities can be transferred to smaller, more deployable models, a theme that will recur in the model specialization strategies for Text-to-SQL. The distilled model outperformed its teacher GPT-3.5 by over 6 absolute points in F1 score, highlighting the benefits of cost, efficiency, and white-box model access. The integration of explicit medical knowledge further enhances ADE detection. A knowledge-augmented graph neural network with a concept-aware attention mechanism, which incorporates medical knowledge from the Unified Medical Language System, has shown performance competitive with recent advances by learning features differently for different types of nodes in a heterogeneous text graph \cite{ref516}. This illustrates how ICL principles can be complemented by structured knowledge to improve the relational learning required for identifying complex drug-event associations.

Automated ICD coding, the task of assigning International Classification of Diseases codes to patient visits based on clinical notes, represents a large-scale multi-label challenge that ICL and related architectures are beginning to address. The process is traditionally manual, time-consuming, and error-prone \cite{ref517}. Transformer-based architectures have been applied to capture the interdependence among tokens in a document and use code-wise attention mechanisms to learn code-specific representations, outperforming recurrent neural network baselines by a significant margin \cite{ref518}. The use of label distribution aware margin (LDAM) loss functions further handles the extreme imbalance in code frequencies. To make the problem tractable, a hierarchical approach reformulates the extreme multi-label problem into a simpler multi-class problem by first identifying focus sentences containing medically relevant events and then classifying each sentence to an ICD code, achieving large improvements in subset accuracy and micro-F1 \cite{ref519}. This decomposition mirrors the step-by-step reasoning that ICL can facilitate and prefigures the decomposition strategies used in Text-to-SQL for complex query generation. Furthermore, contextualized code relation learning frameworks that consider the context of clinical notes to model intricate code dependencies have been shown to boost performance over state-of-the-art baselines \cite{ref520}. These methods, while often involving fine-tuning, are deeply aligned with the ICL paradigm of leveraging contextual information to make predictions, and future work may see the direct application of many-shot ICL to such complex code assignments.

The effectiveness of ICL in these biomedical tasks is significantly amplified by RAG, which addresses the information overload problem and the long-tail distribution of biomedical knowledge, a challenge that directly parallels the retrieval of diverse and representative demonstrations for Text-to-SQL parsing. A novel graph-based retriever that leverages a knowledge graph to downsample over-represented conceptual clusters has been shown to achieve about twice the retrieval performance of embedding similarity alternatives on both precision and recall, enabling LLMs to surface rare but critical associations between drugs, genes, and diseases \cite{ref521}. This hybrid model, combining embedding similarity and knowledge graph retrieval, points to the future of ICL in biomedicine, where prompts are dynamically augmented with the most relevant and comprehensive knowledge. Finally, the distillation of LLMs into task-specific student models, such as PubMedBERT, emerges as a pivotal strategy for overcoming annotation scarcity and deployment constraints, turning the implicit knowledge and ICL capabilities of massive models into efficient, specialized systems for biomedical knowledge extraction \cite{ref515}. The synergy between LLMs and human expertise also accelerates the creation of labeled datasets, with LLM-assisted annotation significantly reducing the human burden while maintaining high accuracy for medical information extraction tasks \cite{ref522}. Through these multifaceted approaches, ICL is not just a prompting technique but a foundational paradigm for building robust, scalable, and knowledge-rich biomedical NLP systems, forming a direct bridge to the structured output generation and schema-grounding challenges that define the Text-to-SQL domain discussed next.

\subsection{Text-to-SQL and Database Query Generation}

Mirroring the cross-lingual transfer and retrieval-augmented strategies that have proven transformative in machine translation, the task of Text-to-SQL---translating natural language questions into executable SQL queries---has become a prominent testbed for In-Context Learning (ICL). It presents stringent requirements for domain generalization, compositional reasoning, and structured output generation, making it a natural parallel to the challenges of bridging linguistic domains. ICL has fundamentally shifted the paradigm from specialized, fine-tuned models to generalist Large Language Models (LLMs) that can adapt to unseen database schemas with only a few demonstrations. The core challenge, akin to selecting aligned examples for low-resource languages, lies in constructing prompts that effectively communicate the target schema, the mapping between questions and SQL logic, and the syntactic constraints of the query language, all without updating the model's parameters.

A primary research thrust in this domain is the strategic selection and organization of demonstrations, a problem directly analogous to the syntax-aware selection strategies in cross-lingual ICL. Early approaches often relied on retrieving examples based on superficial question similarity. However, research has shown that leveraging the syntactic structure of SQL queries is far more effective, just as syntactic tree similarity improves translation quality. For instance, selecting demonstrations based on the similarity of SQL parse trees or by ensuring coverage of diverse SQL clauses and operators significantly boosts performance. A study on prompt design strategies found that pursuing both diversity and similarity in demonstration selection, using the syntactic structure of an example's SQL query for retrieval, leads to enhanced execution accuracy on benchmarks like Spider \cite{ref507}. This principle of structural coverage is echoed by other works that aim to cover all required output structures collectively. By selecting diverse demonstrations that represent the full range of SQL clauses, functions, and operators (e.g., \texttt{JOIN}, \texttt{GROUP\ BY}, nested queries), models are better encouraged to generalize compositionally to novel query structures absent from the training set \cite{ref523}. This is a direct countermeasure to the brittleness of models on compositional splits, where a model must recombine known components in unseen ways, much like cross-lingual models must recombine linguistic primitives for new language pairs.

Moving beyond simple retrieval, several methods frame demonstration selection as a set-level optimization problem, reflecting the cross-lingual field's move toward holistic prompt engineering. One algorithm devises an offline sampling strategy to synthesize a minimal set of few-shot examples with maximal domain and clause coverage within a token budget, forming a fixed, generic prompt that avoids expensive test-time retrieval and adapts well to new databases \cite{ref524}. Others propose hybrid demonstration pools that fuse out-of-domain examples with synthetically generated in-domain examples, leveraging the strengths of both sources to improve robustness \cite{ref525}. The challenge of generating high-quality synthetic data is addressed by frameworks that use iterative, human-free fusing to build diverse demonstration pools, overcoming the diversity limitations and high costs of human-labeled examples \cite{ref526}. The quality of synthetic data is paramount; models pre-finetuned on data synthesized with strong typing, key relationship awareness, and schema-distance-weighted sampling have achieved new state-of-the-art results, proving that the data's fidelity is more critical than its mere volume \cite{ref527}.

A significant innovation in ICL for semantic parsing, which echoes the distillation of LLM capabilities into specialized student models in the biomedical domain, is the use of intermediate representations or general-purpose programming languages as the target output instead of SQL. This approach capitalizes on the pretraining biases of LLMs, which are heavily exposed to languages like Python. Research has demonstrated that prompting a model to generate Python code, augmented with a structured domain description of available classes and functions, dramatically improves accuracy over generating rare domain-specific languages (DSLs). This strategy nearly closes the gap between easy i.i.d. and hard compositional splits, suggesting that the resemblance of the target language to general-purpose code is a stronger success factor than the language's mere prevalence in training data \cite{ref528}. This technique aligns with the development of modular query plan languages (QPL) that decompose complex SQL into simple, regular sub-queries, making the parsing task more interpretable and compositionally easier for the model \cite{ref529}.

To further enhance generalization, especially in cross-domain and cross-lingual settings, ICL prompts are often augmented with explicit domain knowledge and structural cues. This practice mirrors the integration of knowledge graphs in biomedical ICL. Such augmentations can take the form of table documentation that contextualizes tables and columns, or schema linking discriminators that enforce correct question-schema alignments \cite{ref530,ref531}. For cross-lingual tasks, where a non-English question must be mapped to an English schema, the XRICL framework retrieves relevant English exemplars and includes global translation examples to bridge the language gap, demonstrating effective cross-lingual transfer without target-language annotations \cite{ref532}. This directly parallels the PARC pipeline and other cross-lingual retrieval methods that use high-resource languages as a pivot for low-resource tasks. Complementary to these methods, post-hoc reranking of n-best hypothesis lists using query plan models and schema linking algorithms can correct errors that even well-prompted models make, particularly on complex, multi-turn conversational queries \cite{ref533}.

In summary, the application of ICL to Text-to-SQL parsing has moved from simple few-shot mimicking to a sophisticated orchestration of prompt components. The frontier involves the strategic synthesis and selection of demonstrations that are simultaneously diverse and syntactically representative, the adoption of model-friendly intermediate target languages, the integration of explicit domain schema into the prompt, and the use of post-hoc correction and cross-lingual retrieval. These methodologies collectively tackle the core challenges of compositional generalization and domain adaptation, establishing a shared toolkit with other ICL-intensive fields. By uniting insights from cross-lingual alignment and knowledge-grounded biomedical prompting, ICL is becoming a powerful and increasingly dominant paradigm for building flexible and accurate natural language interfaces to databases.

\subsection{Machine Translation and Cross-Lingual Transfer}

In-context learning (ICL) has emerged as a transformative paradigm for machine translation (MT) and cross-lingual transfer, offering a compelling pathway to bridge the performance gap between high-resource languages (HRLs) and low-resource languages (LRLs) that have been historically underserved by traditional supervised models. Much like the structural and compositional challenges seen in Text-to-SQL parsing, MT demands an acute sensitivity to linguistic structure and semantic alignment, making the prompt's construction the central lever for successful cross-lingual ICL. Unlike parameter-updating methods, ICL allows models to perform translation or cross-lingual tasks by conditioning on a few demonstration examples within the prompt, making it especially valuable for the thousands of languages lacking extensive parallel corpora. The foundational capability of LLMs to perform zero-shot and few-shot translation by leveraging in-context demonstrations has been extensively documented, yet the quality of this translation is profoundly influenced by the prompt's construction, particularly the selection of examples and the alignment of semantics between languages.

A central challenge in cross-lingual ICL for MT, directly paralleling the demonstration selection problem in Text-to-SQL, is the selection of highly informative demonstrations. Early approaches often relied on random selection, but research has shown this is severely limited in cross-lingual settings due to a lack of alignment in both input and output spaces. A pivotal advancement comes from syntax-aware selection strategies. The work by \cite{ref534} demonstrates that moving beyond superficial word-level features to deep syntax-level knowledge is critical for improving MT quality, reinforcing the same principle that syntactic tree similarity improves Text-to-SQL accuracy. By computing syntactic similarity between the dependency trees of source sentences using Polynomial Distance, their method selects examples that share structural properties with the test input, leading to more coherent and linguistically grounded demonstrations. Furthermore, an ensemble strategy that combines both word-level and syntax-level criteria was shown to yield the highest COMET scores across numerous translation directions, underscoring the necessity of integrating structural linguistic insights into prompt engineering for translation tasks.

For the particularly challenging domain of LRLs, the direct application of ICL is often hampered by the very scarcity of high-quality parallel data that makes these languages low-resource in the first place. This has spurred significant innovation in cross-lingual retrieval-augmented in-context learning, a strategy that mirrors the retrieval-based demonstration selection pipelines developed for cross-domain Text-to-SQL. A prominent line of work proposes augmenting the prompt for an LRL task with semantically similar examples retrieved from a HRL, effectively using the HRL as a pivot to provide the structured knowledge that the LRL lacks. The paper \cite{ref535} introduces the PARC pipeline, which retrieves semantically similar sentences from an HRL to construct prompts, leading to substantial improvements in zero-shot performance on tasks like sentiment classification and topic categorization across multiple low-resource language families. This approach is directly analogous to the XRICL framework in Text-to-SQL, which retrieves English exemplars to bridge the language gap for non-English questions. The principle is corroborated by \cite{ref536} and \cite{ref537}, which show that cross-lingual retrieval-augmented prompts, particularly when sourcing from HRLs, can steadily boost the performance of generative MPLMs like BLOOMZ on low-resource languages such as Bangla. These findings highlight that the effectiveness of ICL is not solely a property of the model but is heavily contingent on the quality and semantic alignment of the external knowledge integrated into the context, a lesson that equally applies to enriching Text-to-SQL prompts with domain documentation and schema linking.

However, the efficacy of retrieval-augmented ICL is not uniform across all task types, much like how the choice between SQL and intermediate Python-like languages matters for semantic parsing. A nuanced analysis in \cite{ref536} reveals a critical performance dynamic: while cross-lingual retrieval augmentation yields steady improvements for classification tasks, it faces significant challenges in generation tasks, such as those central to MT. This discrepancy suggests that while classification can benefit from a broad semantic match, generation requires a more precise and structurally faithful alignment that simple semantic similarity from a HRL may not guarantee. This observation motivates more sophisticated alignment techniques and echoes the insight that generating intermediate representations in a model-friendly language can outperform direct SQL generation. A critical insight from the study of ICL for low-resource text classification in \cite{ref538} is the identification of the limitations of random input-label pair selection for cross-lingual prompts, primarily due to a lack of coherence in semantics and task-based alignment. Their proposed method, X-InSTA, explicitly injects coherence into the semantics of input examples and aligns the task across the source and target languages, substantially outperforming random selection. Similarly, for machine translation specifically, \cite{ref539} identifies the shortcomings of standard in-context label alignment and introduces ``query alignment'' as a more effective alternative. This method focuses on closing the language gap by aligning the query with semantically relevant information in both the target low-resource language and the high-resource language the model is proficient in, thereby enhancing the model's understanding quality.

The broader theme of bridging the language gap is central to advancing ICL for MT and cross-lingual transfer, just as bridging the modality gap between text and tables is fundamental to the next frontier of ICL. This involves not just aligning queries or input examples, but also addressing the fundamental representational disparities between languages. The importance of explicit cross-lingual alignment is further underscored by work on cross-lingual sentence embeddings, which form the basis of many retrieval methods. Research such as \cite{ref540} demonstrates that cross-lingual word representations in current models are notably under-aligned between LRLs and HRLs. By explicitly aligning words using off-the-shelf word alignment models and training with aligned word prediction and translation ranking objectives, sentence embeddings and thus downstream retrieval for ICL can be significantly improved. The process of leveraging explicit alignment extends to the broader pretraining and data augmentation landscape. For instance, \cite{ref541} demonstrates that mining and training iteratively on sentence pairs mined using a model's own encoder outputs can continuously improve cross-lingual alignment and translation ability. This approach achieved state-of-the-art unsupervised machine translation results and massive improvements on sentence retrieval tasks, showcasing a virtuous cycle where better alignment fuels better ICL retrieval, which in turn enables more effective cross-lingual task performance---a cycle reminiscent of how synthetic data generation iteratively improves Text-to-SQL parsing.

Furthermore, the transfer of linguistic knowledge through ICL is not limited to semantic content but also involves structural concepts. The potential for explicitly aligning conceptual correspondences between languages is explored in \cite{ref542}. By analyzing the spaces of structural concepts (e.g., syntactic properties) across 43 languages, the study reveals a high degree of alignability, and a meta-learning-based method is proposed to learn to align these conceptual spaces. This facilitates zero-shot and few-shot generalization in concept classification and offers profound insights into the cross-lingual ICL phenomenon itself, suggesting that models can be guided to recombine known linguistic primitives for novel cross-lingual tasks, much like how diverse demonstrations encourage compositional generalization in Text-to-SQL.

To overcome the extreme data scarcity in many LRLs, researchers have also explored methods to democratize the English-dominant abilities of LLMs. The approach in \cite{ref543} assembles synthetic exemplars from a diverse set of HRLs to prompt an LLM to translate from any language into English. This process creates high-quality intra-lingual exemplars for the target LRL, allowing for unsupervised prompting that performs on par with, or even outperforms, supervised few-shot learning on translation tasks for numerous African and Indic languages. This effectively leverages the strong alignment between multiple HRLs and English to bootstrap performance on LRLs where direct data is absent, and serves as a natural precursor to the tabular domain, where universal predictors are similarly trained to generalize across diverse structured datasets without task-specific data. The ultimate goal of these diverse approaches---from syntax-based selection and retrieval-augmented context to explicit semantic and query alignment---is to transform ICL from a brittle few-shot trick into a robust, generalizable mechanism that can transfer knowledge across language boundaries. By aligning the semantics and structures between the resource-rich and resource-scarce linguistic worlds, ICL acts as a dynamic bridge, enabling models to utilize their vast pretrained knowledge to serve as universal translators and cross-lingual learners, a bridging capability that the following subsections will extend from linguistic domains into the structured worlds of databases and tabular data.

\subsection{Tabular Data Understanding and Prediction}

Mirroring the cross-modal alignment challenges in machine translation and scientific discovery, the application of in-context learning (ICL) to tabular data represents a significant paradigm shift, bridging the gap between the structured world of databases and spreadsheets and the flexible, natural language reasoning of large language models (LLMs). Tabular data, characterized by rows of samples and columns of features, is the lingua franca of countless industries, from finance and healthcare to logistics and e-commerce \cite{ref544}. The primary challenge in adapting ICL to this domain lies in an inherent impedance mismatch akin to the language gap in cross-lingual transfer: LLMs are trained on sequential, unstructured text, while tabular data is structured and heterogeneous. Consequently, a core area of innovation is the development of serialization techniques that convert tabular data into natural language representations that LLMs can effectively process, enabling zero-shot and few-shot prediction without any parameter updates.

A foundational approach to this challenge involves converting individual data records into textual prompts. Frameworks like TabText have demonstrated that by extracting contextual information from tabular structures, including column header descriptions, and converting the content into language, pre-trained LLMs can be leveraged to generate high-performing feature representations for downstream models \cite{ref545}. This approach has shown significant improvements in healthcare prediction tasks, such as patient discharge and mortality, where simply augmenting pre-processed tabular data with LLM-generated text representations improved the AUC of standard machine learning models by up to 6\% \cite{ref545}. Similarly, the P-Transformer architecture encodes diverse modalities from both structured and unstructured tabular data into a harmonized language semantic space using a pre-trained sentence encoder and medical prompts, demonstrating improved resilience to data corruption \cite{ref546}.

Moving beyond simple serialization, and paralleling the move toward universal predictors in scientific domains, a new wave of research aims to build universal tabular predictors that can operate on diverse datasets without task-specific training. The UniPredict framework exemplifies this by training a single LLM on an aggregation of 169 tabular datasets with diverse targets, enabling it to comprehend diverse tabular inputs and predict target variables following instructions \cite{ref547}. This model demonstrated a significant advantage over tree-boosting and neural network baselines, and in low-resource few-shot settings, it achieved over a 100\% performance improvement compared to XGBoost, highlighting the potential of scaling generative models for universal tabular prediction \cite{ref547}. This principle is echoed in the MediTab project, which uses a data engine powered by LLMs to consolidate, align, and refine heterogeneous medical tabular datasets, enabling a pre-trained predictor to infer for arbitrary tabular inputs without fine-tuning and achieving state-of-the-art zero-shot performance on patient outcome prediction \cite{ref548}.

The challenge of heterogeneous tables with varying schemas is directly addressed by methods that treat table understanding as a unified, multimodal problem, akin to the cross-modal alignment strategies seen in molecular modeling. TableGPT is a pioneering framework that introduces the concept of global tabular representations, jointly training an LLM on both table and text modalities to understand and operate on tables using external functional commands for tasks like question answering, data manipulation, and automated prediction \cite{ref549}. This contrasts with purely textual approaches and allows for more complex operations. Another innovative direction is domain transformation, as seen in the TablEye framework, which converts tabular data into images to leverage the wealth of research in few-shot image learning, thereby overcoming the scarcity of shared information among independent tabular datasets \cite{ref550}.

A critical and highly sensitive application domain for tabular ICL is healthcare, particularly the use of electronic health records (EHRs). EHR data is inherently longitudinal, sparse, and knowledge-infused, posing a significant challenge for LLMs. Research has demonstrated that by carefully designing prompting frameworks that account for specific EHR characteristics like units and reference ranges, LLMs can improve prediction performance in key tasks such as mortality, length-of-stay, and 30-day readmission by about 35\% in few-shot settings \cite{ref551}. Specialized agent-based approaches have also emerged, where a predictor agent generates reasoning and a critic agent analyzes incorrect predictions to provide iterative guidance, achieving decent few-shot performance comparable to traditional supervised learning methods in disease prediction \cite{ref552}. The HeLM framework further extends this by grounding multimodal LLMs in individual-specific data, encoding high-dimensional clinical modalities like time-series spirogram data into the LLM's token embedding space to estimate disease risk more effectively than classical machine learning approaches \cite{ref553}.

However, the adaptation of ICL to tabular data is not without its limitations. LLMs can struggle with the specifics of instruction following for tabular prediction, often ignoring instructions or failing to predict specific instances correctly, as revealed by benchmarks like TABLET, which annotate tabular datasets with instructions of varying phrasing and granularity \cite{ref554}. This indicates a need for new capabilities in LLMs to faithfully learn from instructions for structured data, a challenge reminiscent of the need for precise semantic and task-based alignment highlighted in cross-lingual ICL.

Privacy preservation is a paramount concern when applying ICL to tabular data, which often contains highly sensitive personal information. The risk of LLMs leaking or regurgitating private examples from prompts has motivated a significant body of research into privacy-preserving ICL. The DP-ICL framework addresses this by generating differentially private responses through a noisy consensus mechanism among an ensemble of LLM responses based on disjoint exemplar sets, achieving a strong utility-privacy tradeoff for text classification and generation \cite{ref454}. Focusing specifically on the data itself, the DP-TabICL framework investigates the application of local and global differential privacy mechanisms to tabular data prior to serialization and prompting, showing that private ICL can achieve comparable performance to non-LLM baselines, especially under high privacy regimes \cite{ref455}. An alternative strategy is to generate synthetic few-shot demonstrations from the private dataset with formal differential privacy guarantees, which has been shown to be an effective and competitive alternative to non-private ICL and zero-shot solutions \cite{ref476}. The broader security and privacy landscape of adaptation techniques, including ICL, has been systematically compared, revealing that while ICL is vulnerable to membership inference and other attacks, it presents a different risk profile compared to fine-tuning or parameter-efficient methods like LoRA, and there is no single silver bullet for privacy and security \cite{ref441}.

In summary, the adaptation of in-context learning to tabular data is rapidly evolving from simple text serialization to the development of universal, multimodal predictors. The key innovations lie in creating robust natural language representations of structured data, enabling zero-shot generalization across heterogeneous schemas, and integrating privacy-preserving mechanisms like differential privacy directly into the ICL workflow. While challenges remain in instruction faithfulness and bridging the performance gap with specialized tree-based models \cite{ref555}, the ability of LLMs to perform few-shot learning on tabular data without costly and privacy-invasive retraining marks a significant step toward more flexible and accessible data science, extending the bridging capability of ICL from linguistic and scientific domains into the structured data that underpins modern industry.

\subsection{Scientific Discovery and Molecular Modeling}

Extending the bridging capability of ICL from structured tabular data into the frontiers of scientific inquiry, the application of In-Context Learning (ICL) to scientific discovery---particularly within chemistry, material science, and drug discovery---represents a transformative shift in how complex molecular problems are approached. Just as tabular data applications require LLMs to comprehend structured representations, scientific discovery demands that models internalize the language of molecules, reactions, and physical properties. Traditional computational methods often require extensive training data and task-specific model architectures, creating bottlenecks that mirror the heterogeneous schema challenges in tabular prediction. ICL offers a powerful alternative by enabling large language models (LLMs) to perform novel scientific tasks, such as molecular property prediction and cross-modal molecule discovery, directly from a handful of demonstrations embedded in a prompt, without any parameter updates. This is particularly impactful in low-resource scenarios where annotated data for a specific property or reaction is scarce, a challenge that directly connects to the data augmentation and cost-reduction themes that will be expanded in the legal domain analysis.

A central pillar of ICL in this domain is \textbf{molecular property prediction (MPP)} , which parallels the universal prediction objectives seen in tabular data frameworks like UniPredict. The prediction of molecular properties is fundamental to virtual screening and compound optimization. ICL allows models to adapt to new prediction tasks at inference time by understanding textual descriptions of the target property. For instance, a model can be prompted with a few examples of SMILES strings paired with their corresponding absorption, distribution, metabolism, excretion, and toxicity (ADMET) properties and then accurately predict the property for a novel molecule. This paradigm contrasts with the conventional need to train or fine-tune separate models for each new property, an approach embodied by architectures like CLAMP, which are designed to use separate modules for chemical and natural language inputs to boost few-shot learning performance \cite{ref556}. The effectiveness of ICL for MPP is deeply intertwined with the richness of molecular representations, echoing the multimodal fusion strategies in tabular understanding frameworks like TableGPT. While mono-modal methods relying solely on sequences, graphs, or 3D geometries are powerful, multi-modal representation learning, which fuses these disparate views, has been shown to provide a more comprehensive understanding. Deep learning models that convert molecules into SMILES-encoded vectors, ECFP fingerprints, and molecular graphs---processed by Transformers, BiGRUs, and GCNs respectively---have demonstrated that multi-modal fused deep learning (MMFDL) significantly outperforms uni-modal models in accuracy and resilience to data noise \cite{ref557}. The ICL framework naturally complements this by allowing the prompt to interleave different representations, guiding the model to leverage the complementary information encoded within them. Furthermore, the integration of domain knowledge, such as knowledge of molecular similarity and scaffolds, has been shown to enhance model robustness and generalizability, especially in challenging scenarios like activity cliffs \cite{ref558}.

Data scarcity is a chronic bottleneck in scientific discovery that ICL is uniquely positioned to address, much as it reduces annotation costs in legal text analysis. A particularly innovative strategy involves leveraging LLMs for \textbf{retrieval-based pseudo data generation} to overcome the low-resource challenge in cross-modal molecule discovery. A retrieval-based prompting strategy can construct high-quality pseudo data from existing literature or databases, which is then used for domain adaptation. This approach has been shown to outperform traditional methods of cross-modal training while requiring smaller model scales and lower training costs, highlighting the efficiency of using ``artificially-real'' data \cite{ref559}. This method effectively turns a few-shot learning problem into a dramatically more data-rich scenario by expanding the in-context demonstration pool with synthetically generated, yet highly relevant, examples, presaging the clustering-based and generative data augmentation strategies that prove transformative for low-resource legal NLP tasks. This synergy between ICL and data augmentation is similarly observed in the use of domain-adversarial multi-task frameworks that integrate shared knowledge from multiple heterogeneous data domains---such as drug fingerprints, physicochemical properties, and gene expression---to predict novel therapeutic properties of compounds, where pseudo-data and shared representations help bridge gaps in the primary data source \cite{ref560}.

The efficacy of ICL for complex scientific reasoning is substantially amplified by the deliberate construction of \textbf{domain-knowledge-embedded prompts}, a principle that resonates with the structured prompting techniques employed in legal document classification. Rather than relying on naive demonstrations, embedding chemical and physical principles directly into the prompt structure guides LLMs to generate more accurate and relevant responses. This method has proven superior to standard prompt engineering for tasks encompassing the intricate physical-chemical properties of small molecules, their druggability, and the functional attributes of enzymes and crystal materials, as it significantly improves capability, accuracy, and F1 scores while reducing hallucination \cite{ref561}. Such prompts act as a form of soft constraint, aligning the model's latent knowledge with the fundamental laws of the domain. This dovetails with the broader trend of cross-modal representation alignment, where models are explicitly trained to associate the structure of a small molecule with a completely different modality, such as the transcriptional change it induces. In such tasks, ICL can be used to demonstrate the correct alignment between a molecular graph and its perturbational gene expression profile, forcing the model to learn a sophisticated, cross-modal mapping in-context \cite{ref562}.

Beyond MPP, ICL is advancing into more complex territories such as \textbf{automated knowledge extraction from scientific literature} and \textbf{fine-grained chemical entity typing}, tasks that mirror the document analysis challenges in the legal domain. Extracting structured knowledge from the vast, growing body of unstructured text and images in materials science literature is a prime application. Frameworks employing natural language processing (NLP) can automate text and image comprehension for tasks like the precision extraction of composition-structure-processing-property relationships from publications on inorganic glasses \cite{ref563}. In a similar vein, ICL can be used for fine-grained chemical entity typing, a task made challenging by the complex name mentions and graphic representations common in chemistry literature. By providing a few in-context examples that include both the textual mention and a cropped image of the chemical structure, a multimodal model can learn to fuse information from both modalities, as demonstrated by multimodal knowledge representation frameworks that leverage cross-modal attention between text and chemical structure images \cite{ref564}.

The application of ICL to materials discovery is also pushing the boundaries of traditional model capabilities, particularly in extrapolative tasks. The search for materials with properties superior to any currently known is inherently an extrapolation problem, a well-known weak point for models trained to learn probability distributions of data. While rule-based inverse molecular designers like AI-driven combinatorial chemistry can discover such materials by generating all possible molecular structures from fragment combinations \cite{ref565}, ICL offers a meta-learning pathway to acquire extrapolative generalization. By training a model on a diverse set of small extrapolative tasks, it can learn to adapt rapidly to an unseen material domain when given a few examples in its context, thus achieving extrapolative generalization through ``learning to learn'' rather than explicit rule-based programming \cite{ref566}.

Finally, the \texttt{Lo-Hi} benchmark provides a critical, practical lens for evaluating ICL systems in drug discovery, distinguishing between Lead Optimization (Lo) and Hit Identification (Hi) tasks that mirror the real-world process. Applying ICL in this practical setting reveals which modeling strategies are truly robust and guard against the over-optimistic evaluations of current benchmarks that often fail to represent distribution shifts \cite{ref567}. In summary, ICL is not merely a convenience but a fundamental new capability for scientific discovery, enabling flexible, data-efficient, and knowledge-driven exploration in the chemical and material sciences by effectively turning generalist LLMs into specialist scientific reasoners on demand. This capacity for low-resource adaptation and the creative use of synthetic data generation directly anticipates the annotation expansion and human-LLM collaborative paradigms that underpin the transformative impact of ICL in legal text analysis and recommendation systems.

\subsection{Legal Text Analysis and Recommendation Systems}

Building upon the scientific discovery domain where ICL enables knowledge-driven exploration, the legal domain presents a complementary set of challenges where ICL's capacity for data-efficient adaptation proves equally transformative. Legal document analysis encompasses tasks from text classification and clustering to judgment prediction, all traditionally bottlenecked by the labor-intensive process of manual expert annotation. ICL offers a paradigm shift by enabling large language models (LLMs) to perform these tasks with few or no task-specific training examples, thereby reducing reliance on costly labeled data. This section examines the dual role of ICL in legal text analysis and recommendation systems, focusing on its capacity for annotation expansion and low-resource adaptation---themes that directly connect to the annotation and active learning synergies explored in the following subsection.

A primary challenge in legal NLP is the sheer length and structural complexity of legal documents, which can span tens of thousands of words and lack uniform organization. This problem is exemplified in legal judgment prediction, where models must parse extensive case documents to predict outcomes. Research has explored hierarchical frameworks to address this, such as the ``Multi-stage Encoder-based Supervised with-clustering'' (MESc) model, which segments documents, extracts embeddings from LLMs like GPT-Neo and GPT-J, and uses unsupervised clustering to approximate document structure before learning inter-chunk representations \cite{ref568,ref569}. While such approaches often involve some degree of fine-tuning, the underlying principle of leveraging LLMs for document understanding is directly transferable to ICL. ICL-based methods can be employed to classify these large documents without modifying model parameters, and when combined with structured prompting techniques like soft-token tags \cite{ref570}, they can streamline the process of adapting to specific legal classification schemas, such as identifying documents responsive to a discovery request \cite{ref571}.

The most transformative impact of ICL in the legal domain lies in its capacity for data augmentation and the reduction of annotation costs---challenges that directly motivate the active learning synergies discussed in the next subsection. The expense of legal expert annotation drives a critical need for methods that can expand small, expert-curated datasets. ICL-based data augmentation strategies have proven highly effective here. For instance, clustering-based data augmentation, where in-context prompts are used to generate synthetic examples that are then filtered and clustered with existing data, has led to remarkable performance gains of over 15\% in classification accuracy for tasks like analyzing Sustainable Development Goals (SDGs) in legal texts, while expanding the example base by a factor of five \cite{ref572}. Similarly, generative data augmentation frameworks specifically designed for the legal domain, such as DALE, address the limitations of simple rephrasing by leveraging the templatized nature of legal documents. By pre-training on a denoising objective that masks collocated spans of legal text, DALE learns to generate coherent and diverse augmentations with novel contexts, outperforming even general-purpose LLMs on low-resource legal NLP tasks \cite{ref573}. These approaches embody the core ICL promise of using LLMs as cost-effective data labelers, with studies showing that using pseudo-labels from models like GPT-3 can reduce labeling costs by 50\% to 96\% while achieving performance comparable to human-labeled data when combined with a novel framework for mixing human and synthetic labels \cite{ref574}.

This cost-efficiency is further amplified by the development of parameter-efficient learning techniques tailored for the legal domain. Directly fine-tuning massive models is often infeasible and data-inefficient, especially with small legal datasets. Parameter-efficient legal domain adaptation, which uses unsupervised data from public legal forums for pre-training before applying lightweight tuning methods, has been shown to match or exceed the few-shot performance of fully fine-tuned models like LEGAL-BERT on various legal tasks, all while tuning only about 0.1\% of the model's parameters \cite{ref575}. This approach is synergistic with ICL, as the parameter-efficiently adapted model can then serve as a more robust base for in-context demonstrations, or the ICL paradigm itself can be used to guide the adaptation process. Furthermore, mirroring the human-LLM collaborative annotation paradigms that will be detailed in the following subsection, the integration of ICL with active learning (AL) loops creates a powerful synergy for low-resource legal settings. By using an LLM as an annotator within an active learning framework, the model can be queried to label the most informative examples, which are then used for in-context demonstrations or to train a smaller, task-specific student model. This methodology has demonstrated near-state-of-the-art performance with massively reduced data requirements, achieving potential cost savings of at least 42 times compared to human annotation for low-resource languages \cite{ref576}. The concept of ``CoAnnotating,'' where work is allocated between humans and LLMs based on the model's uncertainty, further refines this process, improving performance by up to 21\% over random baselines for various text annotation tasks \cite{ref577}.

In the realm of legal recommendation systems, the challenge is to retrieve relevant case law or literature to support legal reasoning. ICL facilitates this by enabling models to understand the semantic similarity between documents in a zero-shot or few-shot manner, based on natural language descriptions of relevance. While dedicated representation learning methods like averaged fastText word vectors or Poincaré citation embeddings have been highly effective for content-based legal literature recommendations \cite{ref578}, ICL can augment these systems by providing a flexible interface for refining queries and evaluating the contextual relevance of retrieved documents through chain-of-thought prompts. The integration of external knowledge bases with LLMs, as seen in domain-specific models like ChatLaw, which combines vector database retrieval with keyword retrieval to reduce hallucination, represents a path toward more accurate and explainable legal recommendation systems \cite{ref579}. Moreover, the challenge of generating high-quality labeled datasets for training recommenders, a process often hindered by the scarcity of expert annotations \cite{ref580}, can be directly addressed by the ICL-based data augmentation and annotation strategies discussed above, which form the foundation for the broader annotation and active learning synergies explored in the next subsection.

Ultimately, the application of ICL in legal text analysis and recommendation systems is not merely about performing tasks with fewer examples; it is about fundamentally re-engineering the workflow of legal AI. By leveraging LLMs for automated annotation, synthetic data generation, and uncertainty-guided active learning, ICL provides a pathway to democratize access to sophisticated legal NLP tools, enabling low-resource adaptation and significantly reducing the financial and temporal costs associated with the deep domain expertise required for legal text annotation. The combination of these techniques with explainability methods, such as occlusion sensitivity-based sentence extraction \cite{ref568} or automatic explanation generation with tree estimators \cite{ref581}, is critical for building trustworthy systems that can be confidently deployed in the high-stakes legal environment. This emphasis on cost-effective annotation and human-LLM collaboration directly anticipates the comprehensive frameworks for annotation, data augmentation, and active learning synergies that constitute the next frontier in ICL research.

\subsection{Annotation, Data Augmentation, and Active Learning Synergies}

As demonstrated in the legal domain, where ICL-based annotation strategies can reduce labeling costs by over 50\% while enabling five-fold dataset expansion, the high cost of data annotation---especially in specialized domains and low-resource languages---has motivated a powerful synergy between In-Context Learning (ICL), active learning (AL), and data augmentation. Rather than viewing these paradigms in isolation, recent research has demonstrated that integrating them creates a virtuous cycle where ICL enables new forms of low-cost annotation and augmentation, which in turn reduce the labeled data burden for downstream tasks, while AL strategically selects the most informative instances for this process. This subsection surveys these integrative frameworks, building directly on the annotation cost reduction and data augmentation successes discussed in the legal domain, and generalizing these principles across diverse application areas.

A central theme in this synergy is the use of LLMs as active annotators within human-in-the-loop systems, extending the paradigm of LLM-based pseudo-labeling explored in legal text analysis. Traditional AL relies on human experts to label queried instances, but the emergence of instruction-tuned LLMs has opened the possibility of replacing or augmenting human annotators. The framework \textbf{FreeAL} \cite{ref582} exemplifies this shift by revolutionizing traditional AL through an interactive distillation process. In this collaborative framework, an LLM serves as an active annotator that provides coarse-grained labels, while a smaller downstream model filters and refines these annotations, creating a feedback loop that progressively improves the knowledge of both models without any human supervision. This approach demonstrated significant zero-shot performance gains across eight benchmark datasets, illustrating the potential for fully automated annotation pipelines that realize the cost-efficiency promise first quantified in the legal domain.

However, the optimal allocation of annotation work between humans and LLMs remains a nuanced challenge, as the legal domain's experience with domain-specific expertise requirements has shown. The \textbf{CoAnnotating} paradigm \cite{ref577} addresses this by leveraging uncertainty estimation to guide the distribution of annotation tasks---a direct extension of the LLMs-in-the-loop active learning methodology discussed in the previous subsection. By estimating the LLM's confidence in its annotations, the framework dynamically allocates low-uncertainty instances to the LLM and reserves high-uncertainty or ambiguous cases for human annotators. This uncertainty-guided allocation achieves up to 21\% performance improvement over random baselines, demonstrating that strategic human-LLM collaboration can simultaneously maximize quality and minimize costs. This approach is particularly valuable for subjective NLP tasks where label variation is inherent, and the \textbf{Annotator-Centric Active Learning (ACAL)} framework \cite{ref583} further refines this by incorporating annotator selection strategies that approximate the full diversity of human judgments, emphasizing minority perspectives rather than seeking a single majority label.

The integration of AL with ICL for example selection creates another dimension of synergy that addresses the cold-start and budget-constrained scenarios prevalent in specialized domains like law. When operating under a limited annotation budget, the \textbf{AdaICL} algorithm \cite{ref584} proposes a model-adaptive optimization-free approach that identifies examples where the model is uncertain and performs semantic diversity-based selection. This strategy is formulated as a Maximum Coverage problem that dynamically adapts based on model feedback, achieving up to 3x greater budget efficiency than uniform sampling and enabling effective ICL with significantly fewer annotated examples. Similarly, for cold-start scenarios where initial labeled data is scarce, the \textbf{STENCIL} framework \cite{ref585} leverages submodular mutual information and text exemplars to select weakly labeled instances, particularly addressing class imbalance by improving rare-class F1 scores by 17\%--40\% over standard AL methods.

The role of LLMs in low-resource and domain-specific annotation tasks has been empirically investigated across diverse settings, complementing the language-specific and domain adaptation findings from the legal domain. \textbf{LLMs in the Loop} \cite{ref576} demonstrates that integrating LLM annotators into AL loops for low-resource languages can achieve near-state-of-the-art performance with potential cost savings of at least 42.45 times compared to human annotation. However, echoing the legal domain's recognition that deep expertise remains essential for certain tasks, the effectiveness of LLM annotation is not universal. An empirical study on domain-specific tasks \cite{ref586} found that compact models trained with a few hundred expert-annotated examples can outperform GPT-3.5 and achieve comparable performance to GPT-4, suggesting that LLM predictions may serve as effective warmup methods, but human expertise remains indispensable for tasks requiring deep domain knowledge. The \textbf{Active Learning for NLP with Large Language Models} framework \cite{ref587} further investigates this trade-off, proposing a consistency-based mixed annotation strategy where LLMs label most samples while human annotators correct potentially mislabeled instances, achieving similar or better results than fully human-annotated AL on multiple text classification datasets.

Data augmentation synergizes powerfully with ICL and AL to further reduce annotation requirements, building on the generative data augmentation techniques---such as DALE and clustering-based expansion---that proved transformative in the legal domain. The \textbf{CLLM} framework \cite{ref588} leverages LLMs' prior knowledge for tabular data augmentation in the low-data regime, coupling generation with a principled curation mechanism based on learning dynamics and uncertainty metrics to ensure that only high-quality synthetic data is incorporated. This approach demonstrates that LLM-based augmentation, when properly curated, can significantly outperform conventional generators. In the realm of instruction-following data, \textbf{Ensemble-Instruct} \cite{ref589} employs heterogeneous mixtures of LMs with ensembling over multiple outputs to generate high-quality instruction-tuning data, improving performance without requiring very large proprietary models. Furthermore, \textbf{TeaMs-RL} \cite{ref590} uses reinforcement learning with textual operations and rules to directly generate diverse instruction datasets, reducing reliance on external models and human annotators while improving model capabilities.

The synergy extends to explanation-guided and structured data scenarios, mirroring the legal domain's emphasis on explainability methods like occlusion sensitivity-based extraction. The \textbf{ELAD} framework \cite{ref591} incorporates explanation-guided sample selection, where uncertainties in the student model's reasoning steps are exploited to identify challenging samples for annotation by the teacher LLM. This ensures that the knowledge distillation process focuses on examples where the student's reasoning is flawed, maximizing the efficiency of knowledge transfer. For structured prediction tasks, \textbf{partial annotation with self-training} \cite{ref592} reduces annotation costs by selecting only the most informative sub-structures for labeling, while using self-training to propagate labels to the remaining components. The cross-lingual dimension is addressed by \textbf{CREA-ICL} \cite{ref536}, which extracts semantically similar prompts from high-resource languages to improve zero-shot ICL performance in low-resource languages, though its effectiveness varies across classification and generation tasks.

The theoretical and practical foundations of this integration are supported by comprehensive evaluation frameworks. The \textbf{ALE framework} \cite{ref593} provides a reproducible environment for comparing AL strategies, enabling practitioners to make informed decisions about which annotation and augmentation strategies to deploy for specific use cases. This ecosystem of AL, ICL, and data augmentation frameworks collectively enables a paradigm shift from manual annotation at scale to intelligent, human-LLM collaborative annotation that dramatically reduces costs while maintaining or improving quality across diverse domains, from biomedical text processing to legal analysis and beyond. This integrative perspective on annotation and learning efficiency provides the methodological foundation for the emerging research frontiers and open challenges that lie ahead in the continued evolution of ICL.

\section{Emerging Frontiers, Open Challenges, and Future Directions}

\subsection{Advanced Paradigms: Meta-Learning, Reinforcement Learning, and Lifelong Adaptation in Context}

Building upon the theoretical and mechanistic foundations of static in-context learning, a new wave of research is pushing beyond these boundaries toward more dynamic, self-improving, and temporally extended paradigms. While standard ICL operates by conditioning a static, pretrained model on a fixed set of demonstrations, these advanced paradigms---meta-in-context learning, in-context reinforcement learning, and lifelong in-context adaptation---aim to overcome the inherent limitations of static prompting. They enable models to recursively refine their own in-context abilities, make sequential decisions in interactive environments, and accumulate knowledge over non-stationary streams of tasks without catastrophic forgetting. This section explores these frontiers and the synergies between meta-learning, online learning, and context-based meta-reinforcement learning (meta-RL) that underpin them, providing a natural bridge from the theoretical questions of \emph{how} ICL works to the practical challenges of deploying ICL in complex, real-world settings.

A foundational insight driving these advanced paradigms is the recognition of transformers not merely as sequence models but as meta-learners themselves. The capacity of attention mechanisms to infer latent task structures from a sequence of input-output examples is functionally analogous to meta-learning algorithms that learn to learn from a distribution of tasks. This perspective is crystallized in the design of architectures that treat in-context demonstrations as a form of episodic memory for rapid adaptation. The work \cite{ref363} explicitly frames the transformer as a meta-RL agent, where the self-attention mechanism computes a consensus representation over past experiences, effectively minimizing Bayes risk at each layer to enable fast policy adaptation. This view transforms ICL from a passive retrieval process---as often implicitly assumed in static prompting---into an active, context-dependent inference procedure that can be recursively applied to improve its own learning strategies, forming the basis of meta-in-context learning.

The transition from static ICL to in-context reinforcement learning (ICRL) marks a significant expansion into sequential decision-making, directly addressing the types of interactive, feedback-driven scenarios that static prompting cannot handle. In this paradigm, the context is no longer a fixed set of demonstrations but a stream of trajectory-level experience---comprising states, actions, and rewards---from which the model must infer an optimal policy without gradient updates. Context-based meta-RL provides the theoretical and algorithmic scaffolding for ICRL. For instance, methods like \cite{ref310} and \cite{ref318} leverage contrastive learning to train context encoders that produce compact and sufficient task representations from prior trajectories, enabling policies to generalize to novel tasks within a few adaptation steps. The challenge of learning robust task representations is further addressed by \cite{ref311}, which uses Gaussian mixture models to capture both parametric and non-parametric task distributions, demonstrating strong out-of-distribution generalization. These techniques directly inform ICRL by providing the representation learning backbone needed for models to parse interactive experiences from prompts, effectively extending the ICL framework to settings where feedback and sequential reasoning are paramount.

A critical challenge in deploying ICRL---and one that mirrors the theory-practice gap discussed earlier---is the sample efficiency and stability of learning from sequential contexts, particularly when rewards are sparse or environments are non-stationary. Several innovations from meta-RL research offer solutions that dovetail with the theoretical goal of tighter, more realistic generalization bounds. The concept of offline meta-RL, where a policy is meta-trained on a static dataset to adapt quickly at test time, is particularly relevant for large pre-trained models that cannot afford lengthy online interactions. The work \cite{ref594} demonstrates how contextual meta-learning combined with direct online fine-tuning enables rapid adaptation for robotic insertion tasks, achieving high success rates with minimal on-robot trials. To bridge the distribution shift between offline training contexts and those encountered during online adaptation, \cite{ref341} proposes using additional unsupervised online data collection, highlighting a hybrid approach that resonates with the need for pre-trained ICL models to be safely aligned and adapted post-deployment. Moreover, the issue of sparse rewards, a common bottleneck in long-horizon tasks, is tackled by \cite{ref595}, which relabels past experience during meta-training---a technique that could be translated into prompt design strategies where demonstrations are automatically re-annotated to provide richer learning signals for ICRL, thereby improving sample efficiency in a manner that theory must eventually account for.

The synergy between online learning and meta-RL is pivotal for realizing lifelong in-context adaptation, where models must continually learn from a non-stationary, unbounded stream of tasks without catastrophically forgetting previous knowledge. This directly confronts the theoretical challenge of \emph{transience} and the competition between in-context learning (ICL) and in-weights learning (IWL) identified in prior sections. The paradigm of online meta-learning, formalized by \cite{ref596}, merges the fast-adaptation spirit of meta-learning with the sequential, regret-minimizing nature of online learning. This framework has been extended to the fully online, task-boundary-free setting by \cite{ref597}, which does not require resets between tasks, and further augmented with explicit task switch and distribution shift detection mechanisms in \cite{ref598}. These algorithms enable models to reuse the best available knowledge for a current task and adapt their learning rates for out-of-distribution shifts, laying the groundwork for ICL systems that can meta-train their own retrieval and reasoning strategies over a lifetime of use---precisely the type of continual, adaptive capability needed to bridge static prompting and real-world deployment. The theoretical underpinnings are strengthened by \cite{ref599}, which provides dynamic regret guarantees for online meta-learners, characterizing performance in terms of adaptation to environmental changes and addressing the need for more faithful theoretical frameworks highlighted in the previous section.

A crucial emergent phenomenon in these continual settings is the competition between memorizing knowledge in model weights (in-weights learning) and recognizing and adapting to tasks based on context (in-context learning). Lifelong in-context adaptation seeks to strike a deliberate balance, preserving the model's general-purpose rapid adaptation capabilities while still allowing for the accumulation of specialized knowledge. The framework of \cite{ref600} (CoMPS) demonstrates a practical approach by iteratively learning new tasks with RL and then performing completely offline meta-learning on the collected experience, preparing the model for subsequent tasks without revisiting old data. This dual-loop process mirrors an idealized continual ICL system that alternates between applying its in-context abilities to solve a task and then distilling that interaction into improved future prompts or internal representations---a dynamic that theory must formalize to explain when and why ICL persists versus collapsing into IWL. The challenge of catastrophic forgetting is further mitigated by techniques that model task non-stationarity explicitly, such as \cite{ref601}, which learns to track the temporal evolution of latent task parameters, enabling the model to anticipate future task distributions and adapt its context-based inference accordingly.

The integration of skill-based and hierarchical structures represents another crucial direction for scaling ICL to complex, long-horizon tasks in a continual manner, addressing the theoretical need for compositionality-aware generalization. Rather than treating each new context as a monolithic block of demonstrations, advanced systems can decompose tasks into reusable skills conditioned on context. \cite{ref602} extracts a library of skills from offline data and meta-trains a high-level policy to compose these skills for novel, long-horizon target tasks. This decomposable approach is essential for lifelong in-context adaptation, as it allows a model to store a combinatorial library of context-conditioned sub-routines, vastly increasing its representational power and enabling systematic generalization. Similarly, \cite{ref321} proposes a framework that decouples the learning of task contexts and skills using a decoupled Gaussian quantization variational autoencoder, creating discrete codebooks of representative task and skill distributions. These discrete structures are inherently more compatible with the symbolic prompt interfaces of large language models, suggesting future architectures where prompts can directly index and recombine such learned, context-tagged skill modules---a vision that aligns with the data-centric theory of ICL emergence, where training on structured, analogical data improves in-context capabilities. The overarching movement toward these advanced paradigms points to a future where ICL is not a static one-shot prompting technique but a dynamic, self-optimizing process at the heart of generally intelligent, continually adapting systems \cite{ref344}.

\subsection{Bridging the Theory-Practice Gap: Generalization Bounds and Learnability}

Despite the explosive empirical progress in in-context learning (ICL), a significant gulf persists between its practical success and our formal theoretical understanding, a gap that the advanced paradigms discussed in the previous section---meta-in-context learning, in-context reinforcement learning, and lifelong adaptation---bring into sharp relief. Bridging this theory-practice gap is a critical frontier, requiring frameworks that can explain not just \emph{if} ICL works, but \emph{when}, \emph{why}, and \emph{how} it generalizes to novel tasks, makes sequential decisions, and adapts over non-stationary streams of experience. Current theoretical efforts, while illuminating, often rely on simplified settings that abstract away from the complexities of real-world language models and pretraining data. A central open challenge is to develop more faithful and predictive theories that can guide the design of safer, more robust, and more capable in-context learners operating in the dynamic, interactive environments that the shift toward ICRL and lifelong adaptation demands.

One major axis of this challenge is the development of tighter and more realistic generalization bounds for ICL, which are essential to understand how models will behave when deployed in the continual learning scenarios described previously. The foundational work of \cite{ref603} introduced a PAC-based framework for in-context learnability, casting ICL as a two-phase process: an initial pretraining phase on a mixture of latent tasks, followed by an inference phase where the model identifies the task from a prompt. This framework elegantly proved that, under mild assumptions, tasks can be efficiently learned in-context, with the key insight being that ICL is primarily about \emph{task identification} rather than \emph{task learning} from scratch. However, the bounds derived from this framework and related multitask learning perspectives, such as those in \cite{ref76}, are often too loose or rely on assumptions that do not hold in practice---particularly in the dynamic, context-shifting scenarios characteristic of lifelong adaptation and ICRL. A key open problem is refining these bounds to account for the specific structural properties of the transformer architecture and the pretraining data distribution. For instance, \cite{ref46} provides a unified analysis bounding the pretrained model's error by a sum of approximation and generalization errors, but future work must derive compositionality-aware bounds that reflect how models recombine learned primitives, directly addressing the compositional generalization challenges that will be examined in the following section.

The abstraction of the transformer as an algorithm implementing a hypothesis function at inference time, as formalized in \cite{ref2}, provides a powerful lens that resonates with the view of transformers as meta-learners articulated in the advanced paradigms discussion. This work connects the excess risk of ICL to the stability of the algorithm implemented by the transformer, offering a path toward architectural generalization guarantees. A critical direction for future theory is to formally characterize the conditions under which ICL emerges as a stable algorithm versus collapsing into in-weights memorization---the very tension that lifelong in-context adaptation seeks to manage explicitly. The phenomenon of \emph{transience} in ICL, where the capability emerges and then decays as in-weights learning (IWL) takes over, is a subject of intense empirical study \cite{ref55}. A rigorous theoretical account of this competition, perhaps framed as a learning dynamics problem where different circuit modes compete, is needed to explain why models do not always converge to the optimal in-context solution. This ties directly to the need for theory to explain the abrupt phase transitions and learning plateaus observed during training, which are driven by the sequential learning of nested operations necessary for ICL \cite{ref55}, and to predict the conditions under which the meta-learning and online adaptation strategies from the previous section will succeed or fail.

Another significant gap lies in the theoretical understanding of the \emph{mechanisms} that underpin ICL, a question that becomes more pressing as we move from static prompting to the interactive, trajectory-based contexts of ICRL. The provocative hypothesis that transformers implicitly perform gradient descent (GD) during their forward pass \cite{ref25} has generated substantial debate. While this duality offers an elegant explanation for the optimization-like behavior of ICL, its validity in real-world, pre-trained models has been challenged. Studies such as \cite{ref26} and \cite{ref10} highlight critical discrepancies, including differences in sensitivity to demonstration order and a breakdown in the flow of information, termed ``Layer Causality.'' These findings suggest that while a GD-like process may be a useful metaphor, the actual mechanism in pre-trained models is more complex. Bridging this gap requires moving beyond the equivalence drawn in simplified linear attention settings to understand the adaptive, non-linear algorithms learned by deep transformers. The finding that softmax attention adapts to the function landscape, such as its Lipschitzness, in a way that linear attention cannot replicate \cite{ref65} underscores the need for theory that captures the richness of the learned attention patterns. A related perspective models ICL through the lens of contrastive learning, interpreting the self-attention mechanism as a form of gradient descent on a contrastive loss \cite{ref42}, an interpretation that aligns naturally with the contrastive context encoders employed in context-based meta-RL. This suggests a promising avenue: unifying these different theoretical viewpoints---gradient descent, Bayesian model averaging, and contrastive learning---into a single coherent explanation that accounts for both static prompting and the sequential, feedback-driven settings of ICRL.

Finally, a truly comprehensive theory of ICL must be data-centric, explaining how the properties of the pretraining data directly govern the emergence and quality of ICL---a theoretical imperative that directly anticipates the discussion of compositional generalization in the following section. Empirical work has shown that specific data features, such as the presence of long-tail tokens, burstiness, and the challenge of incorporating long-range context, are crucial for fostering ICL \cite{ref41}. A grand challenge for the field is to formalize these observations into a set of necessary and sufficient conditions on the pretraining data distribution that provably lead to the emergence of ICL. For example, \cite{ref190} demonstrates that training on data structured to emphasize analogical reasoning improves ICL. The theory must advance to the point where we can predict, a priori, how a model will learn given a specific data mixture, and conversely, design optimal data curricula to achieve desired in-context capabilities---including the compositional and systematic generalization abilities that the next section will show are not yet guaranteed by ICL alone. Addressing these interconnected theoretical challenges---from tightening generalization bounds and defining transient learning dynamics to resolving the ICL-GD debate and characterizing data-driven emergence---is essential for transforming ICL from a mysterious empirical phenomenon into a well-understood and reliably engineered pillar of machine learning, capable of powering the dynamic, compositional, and lifelong learning systems that lie on the horizon.

\subsection{Compositional and Systematic Generalization in In-Context Learning}

Building directly on the discussion of pretraining data's role in shaping in-context learning (ICL) capabilities, compositional generalization---the capacity to understand and generate novel combinations from known primitives---represents a fundamental hallmark of human intelligence that probes the very nature of the algorithms transformers implement during inference. While ICL has demonstrated remarkable few-shot adaptability, a growing body of evidence suggests that this paradigm does not inherently guarantee systematic compositionality, revealing another dimension of the theory-practice gap where the implicit mechanisms of ICL fall short of the explicit, rule-governed recombination required for out-of-distribution (OOD) generalization. This section analyzes the intricate relationship between ICL and compositional generalization, examining the role of pretraining data, architectural inductive biases, and emerging solutions from modular and skill-based frameworks that seek to impose the systematicity that standard ICL lacks.

The pretraining data distribution is a primary determinant of whether compositional reasoning can emerge and be leveraged in context, directly extending the data-centric theoretical challenges outlined earlier. The theoretical framework presented by \cite{ref108} argues that ICL relies intrinsically on the recombination of compositional operations present in natural language data. This work derives an information-theoretic bound demonstrating that generic next-token prediction gives rise to ICL abilities only when the pretraining distribution contains sufficient compositional structure. This aligns with data-centric analyses showing that properties such as burstiness and long-range coherence drive emergent capabilities. In a controlled setup designed to mimic the compositional nature of language, scaling parameters and data leads to the emergence of ICL, mirroring the behavior of large language models (LLMs). A critical insight is that models prompted to output intermediate steps perform better, providing a mechanistic link between chain-of-thought prompting and the structured decomposition of compositional tasks.

However, empirical investigations reveal that the standard ICL paradigm often fails dramatically on compositional tasks when lexical control is not rigorously maintained, echoing concerns about whether ICL genuinely induces task identification or merely exploits superficial statistical patterns. \cite{ref604} demonstrates that previously reported high performance on benchmarks like COGS is significantly overestimated due to uncontrolled lexical exposure during pretraining. By substituting context-controlled lexical items with novel character sequences or special tokens with novel embeddings, the study reveals a substantial degradation in generalization, with performance inversely scaling with the amount of pretraining data---an instance of ``inverse scaling'' that highlights a fundamental brittleness. This indicates that LLMs often rely on superficial lexical correlations rather than genuine compositional rules, a shortcut that ICL alone does not reliably circumvent, and which mirrors the broader theoretical concern of whether models are performing robust task identification or brittle pattern matching.

The selection and structure of in-context examples themselves critically modulate compositional success, reinforcing the need for a more refined understanding of how demonstrations guide the algorithm implemented by the transformer. \cite{ref605} systematically investigates this through the CoFe test suite, finding that the choice of demonstrations drastically affects OOD generalization. The pivotal factors are that examples should be structurally similar to the test query, diverse from one another, and individually simple in complexity. Strikingly, the study notes that ICL-based compositional generalization is far weaker when applied to fictional words compared to commonly used ones, underscoring that the model's pretrained lexical knowledge can override the incentive to compose from context. Moreover, even powerful pre-trained backbones require the in-context examples to cover the required linguistic structures, suggesting that ICL often functions as a pattern-matching mechanism rather than a true rule induction engine. This echoes findings in \cite{ref606}, which showed that standard sequence-to-sequence models could achieve near-perfect generalization on SCAN when the training distribution was modified in intuitive ways, implying that the capacity for compositionality is latent but highly sensitive to the data regime rather than being an intrinsic property of the architecture or the ICL mechanism.

Current models exhibit clear limitations in systematic, out-of-distribution compositional reasoning, a failure mode that parallels the theoretical challenge of explaining when ICL emerges as a stable algorithm versus collapsing into in-weights memorization. Benchmarks like COGS and SCAN have been instrumental in diagnosing these failures, where models struggle to recombine known parts in novel ways that would be trivial for a symbolic compositional strategy. \cite{ref607} introduced methodologies to construct realistic benchmarks with maximal compound divergence but minimal atom divergence between train and test sets, revealing a strong negative correlation between compound novelty and accuracy across multiple architectures. This failure mode is particularly pronounced in multimodal settings, where compositionality must span visual and linguistic concepts. \cite{ref608} demonstrated that while multimodality can strongly improve generalization in some cases, the compositional generalization of a minimal LSTM-based network was brittle and heavily dependent on factors like color overlaps between objects, falling short of the robust, factor-independent binding required for true systematicity.

To address these shortcomings, a promising direction involves leveraging meta-learning and modular, skill-based decomposition to impose compositional inductive biases into the ICL workflow, effectively engineering the kind of algorithmic stability and systematic recombination that theory suggests is necessary for robust generalization. The hypothesis that forcing models to in-context learn can itself serve as an inductive bias for compositionality was tested in \cite{ref609}. By training a causal Transformer on all possible few-shot learning problems attainable from a dataset---effectively randomizing instance order and labels---the model is incentivized to use earlier examples to generalize to later ones. This procedure led to improved compositional generalization on SCAN, COGS, and GeoQuery, supporting the notion that meta-learning over in-context tasks provides a useful inductive pressure for recombination.

This meta-learning principle extends to explicit architectural and procedural interventions. \cite{ref610} proposes a novel prompting strategy that teaches LLMs how to compose basic skills to solve harder problems. The method is effective because it demonstrates both the constituent skills and their compositional examples within the same context, fostering a synergistic relationship that enables the resolution of unseen problems requiring innovative skill compositions. This approach activates and composes internal competencies acquired during pretraining but not explicitly prompted, effectively bridging the gap between latent knowledge and systematic application. Similarly, \cite{ref611} lays a foundational paradigm with the compositional recursive learner, a framework that reasons about computation by making analogies to previously seen compositional subproblems, demonstrating generalization to more complex tasks than those seen during training. The concept of training a hierarchical curriculum to induce generalization from atomic skills to complex reasoning tasks, as explored in \cite{ref612}, provides a practical training methodology. The finding that atomic skills do not spontaneously compose but can be successfully composed through a hierarchical curriculum underscores the necessity of explicitly teaching the recombination process.

Modular neural architectures offer a structural pathway to compositional ICL, embodying the kind of compositionality-aware design principles that future theoretical bounds must account for. \cite{ref613} presents a neural modular approach that induces a desired sub-task decomposition without strong supervision, allowing for compositional reasoning by linking different tasks through shared modules that handle common routines. This type of architectural design mirrors the ``compositional problem graph'' formalism, where tasks are related by shared subproblems, and modular networks can dynamically rewire to solve novel combinations. In the visual domain, \cite{ref234} demonstrates that the level of compositionality in the token representations---vector, patch, or object slot---governs the trade-offs in compositional generalization. Object-centric tokenizers coupled with cross-attention modules provide the inductive biases crucial for high-fidelity, combinatorial solutions, while non-compositional representations fail on combinatorial tasks but can extend learned composition rules to unseen domains. This highlights that for ICL to be compositionally robust, the underlying representational primitives must align with the true causal factors of variation in the data, a principle rigorously investigated in \cite{ref614}, which derives mild conditions on the support of the training distribution and model architecture sufficient for compositional generalization, reframing it as a property of the data-generating process.

In conclusion, achieving robust compositional generalization through ICL remains an open frontier that crystallizes many of the theoretical challenges discussed previously---from the need to move beyond loose generalization bounds toward compositionality-aware guarantees, to resolving whether ICL implements true rule induction or sophisticated pattern matching. While pretrained LLMs possess vast knowledge and skills, these are not spontaneously recombined in a systematic manner when confronted with novel compositional pressures, much as the transience phenomenon reveals that ICL capabilities do not stably emerge without the right competitive dynamics. The path forward involves a synthesis of meta-learning objectives that directly optimize for OOD generalization \cite{ref615}, data augmentation strategies that correctly sample from the distribution of valid compositions \cite{ref616}, and the development of modular, skill-based architectures that explicitly represent and recompose atomic capabilities \cite{ref610,ref611}. Bridging the gap between the implicit pattern recognition of current ICL and the systematicity of symbolic composition demands a tighter integration of the syntactic and semantic structure of pretraining data with the inductive biases of novel learning frameworks---a challenge that, like the effort to overcome context window limitations explored next, is fundamental to realizing ICL's full potential as a scalable, versatile learning paradigm.

\subsection{Scaling Context Length: Efficiency, Retrieval, and the Lost-in-the-Middle Problem}

A critical frontier for in-context learning (ICL) is overcoming the constraints imposed by fixed context windows, which limit the number of demonstrations and the complexity of tasks that can be presented---a limitation that directly impacts the compositional generalization and systematic reasoning challenges explored in the prior section. Just as modular architectures and skill-based frameworks seek to decompose complex tasks into manageable primitives, innovations in long-context processing must similarly decompose and restructure information to fit within computational boundaries. The core bottleneck is the quadratic computational and memory complexity of the standard self-attention mechanism, which renders processing extremely long sequences computationally prohibitive. This has spurred a wealth of research into efficiency, retrieval, and architectural innovations, all while grappling with a peculiar failure mode known as the ``lost-in-the-middle'' phenomenon. Addressing these intertwined challenges is essential for scaling ICL from few-shot to many-shot and truly long-context regimes, paving the way for the unified, modality-agnostic frameworks discussed in the following section that will demand the ability to process heterogeneous, high-volume multimodal demonstrations.

A primary line of defense against long-context inefficiency is \textbf{context compression}, which seeks to distill the essential information from a long prompt into a shorter, denser representation. Methods like \cite{ref617,ref618} frame this as a token-level pruning problem, using a smaller model to evaluate the information entropy of tokens and removing redundant ones to improve the density of key information in the prompt. Demonstrating a data-driven approach, \cite{ref619} introduces a specialized fine-tuning stage that incorporates paraphrasing, effectively augmenting the model's retrieval capabilities for long contexts without architectural changes. Going further, \cite{ref620} proposes the In-context Autoencoder (ICAE), a lightweight module pretrained to compress long contexts into a small number of compact ``memory slots'' that can be directly conditioned on, achieving 4x compression and significantly reducing latency and GPU memory cost during inference. Similarly, \cite{ref621} combines context compression, retrieval, and parameter-efficient fine-tuning to create a concise offline representation of long documents, enabling a standard model to effectively handle contexts up to 128k tokens while using 30x fewer tokens at inference time. Other approaches target the computational cost of attention directly through \emph{soft prompt compression} \cite{ref622} or by incrementally compressing intermediate key-value (KV) activations into compact sentinel tokens, which reduces both memory and computation for subsequent context processing \cite{ref209}. Systems like \cite{ref623} cleverly avoid redundant computation by precomputing and reusing the attention states of frequently occurring text segments, such as system messages and document contexts, leading to order-of-magnitude reductions in time-to-first-token latency without modifying the model.

To structurally circumvent the quadratic complexity of full attention, many researchers have turned to \textbf{memory-augmented architectures} that mirror the modular decomposition philosophy advocated for compositional generalization. The central idea is to equip a transformer with an external memory module that can store and retrieve information far beyond the model's immediate context window, effectively creating a persistent substrate for the kind of systematic recombination that compositional tasks demand. \cite{ref624} introduces a non-parametric associative memory that can be coupled to any frozen LLM without re-training, dynamically consolidating token representations to enable the model to handle arbitrarily long sequences. \cite{ref625} presents LongMem, a framework with a decoupled network architecture where a frozen LLM acts as a memory encoder and an adaptive residual side-network serves as a memory retriever and reader, enabling the memorization of long history up to 65k tokens and benefiting many-shot ICL. This concept of memory is also central to \cite{ref626} from \cite{ref627}, which incorporates a compressive memory directly into the standard attention mechanism within a single Transformer block, blending masked local attention with long-term linear attention for bounded memory and computation on infinitely long inputs. Alternative attention structures also show promise; \cite{ref628} proposes a hybrid architecture with alternating layers of local and global attention, balancing short-range precision with long-range dependency modeling and achieving superior efficiency. The \cite{ref629} tackles the ``distraction issue'' where models fail to distinguish between relevant and irrelevant keys when the external memory grows large, using a contrastive training procedure to refine the key-value space and successfully extending the effective context length to 256k tokens.

A particularly vexing problem that emerges with long contexts is the \textbf{``lost-in-the-middle'' phenomenon}, first systematically documented in \cite{ref630}. This work demonstrated that language model performance, even on tasks explicitly requiring information retrieval from the context, degrades significantly when the relevant information is positioned in the middle of a long input, with models performing best at the beginning or end. This counter-intuitive failure has profound implications for ICL, where critical demonstrations or task instructions might be placed anywhere within a prompt---a brittleness that echoes the compositional failures where models exploit superficial positional or lexical correlations rather than true rule induction. Subsequent research has both diagnosed and addressed this issue. The study in \cite{ref631} attributes the problem to a ``long-term decay effect'' in Rotary Position Embedding (RoPE) and introduces Multi-scale Positional Encoding (Ms-PoE), a plug-and-play solution that assigns different rescaling ratios to different attention heads, forming a multi-scale context fusion that revitalizes information from all positions. Another effective approach is \cite{ref632} from \cite{ref633}, which combines \emph{reprompting}---repeating the original task instructions periodically throughout the context---with \emph{in-context retrieval}, a two-step process where the model first retrieves relevant passage numbers before answering. This method effectively reduces the distance between the relevant context and the final query, leading to an average 16-point QA accuracy boost, and resonates with the insight from compositional generalization research that the structure and placement of in-context examples critically modulates performance.

Efforts to build robust long-context models from the ground up focus on data and training strategies that instill the systematic information-seeking behavior lacking in standard ICL. The core principle in \cite{ref634} is to combat ``lost-in-the-middle'' by introducing information-intensive training data where the answer requires synthesizing information from multiple specific short segments scattered across a long synthetic context. This forces the model to learn that any position can hold crucial information, resulting in a model robust at retrieving information from its entire 32K window---analogous to how compositional skill-based frameworks force models to attend to all constituent primitives. Multi-agent collaboration represents a creative system-level solution that decomposes the long-context processing task across specialized agents. \cite{ref635} orchestrates a team of LLMs, with a leader delegating information-seeking to members with specific document responsibilities. An inter-member communication mechanism resolves conflicts from hallucinations, enabling the agent team to effectively process 128k-long texts and outperform single models on retrieval and multi-hop QA tasks. Furthermore, scaling the number of demonstrations directly, a goal of many-shot ICL, is facilitated by architectural modifications like \cite{ref208}, which separately encodes demonstrations and uses a rescaled attention mechanism for linear complexity, and \cite{ref636} from \cite{ref637}, which replaces full-attention with a structured version that removes unnecessary inter-demonstration dependencies, making the model invariant to demonstration order and enabling continuous gains with hundreds of examples.

The path forward involves unified frameworks that integrate these strategies, much as the next section envisions unified multimodal architectures that seamlessly blend heterogeneous data streams. \cite{ref638} reformulates existing long-context methods from a memory augmentation perspective across dimensions of management, writing, reading, and injection, proposing an integrated approach that achieves lower perplexity. This holistic view, combined with innovations like the parallel processing paths in \cite{ref639} that discard irrelevant documents early, points toward a future where ICL can efficiently and accurately leverage vast amounts of context across multiple modalities and tasks. Overcoming the ``lost-in-the-middle'' phenomenon and computational bottlenecks is not merely an engineering hurdle but a fundamental step toward realizing ICL's full potential as a scalable, versatile learning paradigm---one capable of the systematic, compositional reasoning discussed previously and the fluid, multimodal integration explored subsequently.

\subsection{Unified Frameworks and Modalities: Toward General-Purpose In-Context Learning Systems}

While early research on in-context learning (ICL) focused predominantly on scaling within a single modality, the trajectory of the field points decisively toward the development of unified, modality-agnostic frameworks that can seamlessly integrate and reason over text, vision, speech, and structured data. Extending the principle of learning from demonstrations to the multimodal realm, the ultimate ambition is to construct general-purpose ICL systems that, much like a human, can fluidly accept demonstrations in any combination of sensory modalities to infer and execute a novel task. This vision moves beyond the current paradigm of siloed, modality-specific ICL models toward a holistic architecture where a single model can, for instance, learn a new visual concept from a handful of labeled images, a new audio classification task from a few sound clips, or a complex cross-modal reasoning task from interleaved video and text examples, all without any parameter updates. Realizing this vision, however, requires overcoming significant challenges in cross-modal representation alignment, transferring ICL capabilities across modalities, designing lightweight adaptation modules, and building systems capable of lifelong, continual learning from streaming multimodal data.

A foundational challenge in building a unified ICL system is the alignment of heterogeneous modalities into a shared representational space. The core operating principle of ICL is that a model can identify and apply the input-output pattern from a sequence of demonstrations. For this to work across modalities, the model must first encode images, audio waveforms, point clouds, and text into a common representational ``language'' where analogous concepts or tasks are represented similarly. Foundational multimodal models have explored various unification strategies. The Meta-Transformer framework, for instance, proposes a frozen encoder that operates on raw data from various modalities mapped into a shared token space, demonstrating unified learning across 12 modalities without paired multimodal data \cite{ref640}. This data tokenizer approach, which converts diverse inputs into a sequence of tokens, is a crucial step. Similarly, the OFA framework unifies a diverse set of cross-modal and unimodal tasks within a simple sequence-to-sequence learning paradigm, treating all tasks as text generation problems \cite{ref641}. This architectural simplification is powerful because it allows ICL demonstrations for any modality to be formatted as a sequence of tokens, making the task of learning a function from context consistent across modalities. However, achieving true alignment requires more than just tokenization; it demands that the semantics of a concept are divorced from its sensory origin. Methods like contrastive learning, as seen in MS-CLIP which shares a transformer encoder across vision and language, can help encode semantic concepts from different modalities more closely in the embedding space, facilitating the transfer of common semantic structure \cite{ref295}.

The transfer of ICL capabilities from one modality to another is a critical, yet nascent, area of research. A model that excels at text-based ICL should ideally be able to leverage that capability when presented with visual or auditory demonstrations. This is closely related to the problem of cross-modal generalization, where a model trained to perform tasks in a source modality must quickly adapt to new tasks in a resource-poor target modality. The concept of meta-alignment is particularly relevant here, as it focuses on aligning the representation spaces of separate source and target modality encoders using strongly and weakly paired cross-modal data, ensuring that the meta-learning ability for rapid task generalization transfers across modalities \cite{ref642}. This principle can be directly applied to ICL: if a model learns to ``learn in context'' from text, meta-alignment can help transfer this ``learning to learn'' capability to the visual domain, enabling visual ICL even with limited labeled examples. The i-Code series of frameworks exemplifies this integrative ambition by flexibly combining vision, speech, and language into unified vector representations, aiming to project different combinations of modalities into a single space for downstream tasks \cite{ref283,ref643}. Such integrative pretraining is essential for building a unified ICL model that can handle any mix of modalities.

To adapt a monolithic, unified model to the specific nuances of different sensory streams and task structures, lightweight tuning modules are indispensable. Full fine-tuning of a massive multimodal model for every new ICL task or modality is computationally prohibitive. Instead, parameter-efficient methods that can be prepended or inserted into a frozen backbone offer a scalable solution. The MultiModal In-conteXt Tuning (M\(^2\)IXT) module is a prime example of this philosophy, designed as a lightweight, prependable module to enhance the ICL capabilities of multimodal unified models \cite{ref218}. By perceiving an expandable context window of multimodal examples, M\(^2\)IXT enables rapid few-shot adaptation across a variety of tasks without updating the core model's parameters. This approach allows a single, powerful generalist model to be efficiently specialized. The \(\pi\)-Tuning method further explores this by aggregating parameters of lightweight, task-specific experts learned from similar tasks, creating a scalable graph of task relationships that can guide transfer learning across both unimodal and multimodal tasks \cite{ref250}. Such frameworks are crucial for a unified ICL system, as they allow the model to dynamically compose and activate the most relevant sub-capabilities for the specific combination of modalities and tasks presented in the prompt.

The final frontier is the development of lifelong, continually learning in-context systems that can accumulate and refine knowledge from a continuous stream of diverse, multimodal inputs. Current ICL models are largely static; their ability to learn in context is a property that emerges from a fixed pretraining distribution. A truly general-purpose system must be able to adapt its in-context learning behavior itself over time, remembering successful strategies for past tasks and applying them to new, related problems without catastrophic forgetting. This requires a move toward meta-in-context learning, where the model's ability to process and learn from new contexts is recursively improved by its own past experiences. OneLLM, which aligns eight modalities to language using a unified framework and a progressive multimodal alignment pipeline, points toward this scalability, but the notion of continual learning from streaming data remains an open challenge \cite{ref644}. Building such a system would involve developing architectures that can dynamically cache and retrieve relevant past in-context experiences, integrating principles from retrieval-augmented generation and external memory. The model would need to decide when to treat new data as a demonstration for a transient task and when to consolidate it as a more permanent component of its underlying knowledge, effectively blurring the lines between in-context learning, in-weights learning, and online adaptation. Overcoming these challenges will pave the way for a new generation of AI that is not just a static pattern matcher but a dynamic, general-purpose learner that grows and adapts alongside the intricate, multimodal world it is designed to understand.


\begin{thebibliography}{644}
\addcontentsline{toc}{section}{References}
\bibitem{ref1} In-context Example Selection with Influences. arXiv:2302.11042v2, 2023. \url{https://arxiv.org/abs/2302.11042v2}
\bibitem{ref2} Transformers as Algorithms Generalization and Stability in In-context Learning. arXiv:2301.07067v2, 2023. \url{https://arxiv.org/abs/2301.07067v2}
\bibitem{ref3} Self-ICL Zero-Shot In-Context Learning with Self-Generated Demonstrations. arXiv:2305.15035v2, 2023. \url{https://arxiv.org/abs/2305.15035v2}
\bibitem{ref4} Universal Self-Adaptive Prompting. arXiv:2305.14926v2, 2023. \url{https://arxiv.org/abs/2305.14926v2}
\bibitem{ref5} Are Human-generated Demonstrations Necessary for In-context Learning. arXiv:2309.14681v4, 2023. \url{https://arxiv.org/abs/2309.14681v4}
\bibitem{ref6} A Data Generation Perspective to the Mechanism of In-Context Learning. arXiv:2402.02212v1, 2024. \url{https://arxiv.org/abs/2402.02212v1}
\bibitem{ref7} Decomposing Label Space, Format and Discrimination Rethinking How LLMs Respond and Solve Tasks via In-Context Learning. arXiv:2404.07546v1, 2024. \url{https://arxiv.org/abs/2404.07546v1}
\bibitem{ref8} Comparable Demonstrations are Important in In-Context Learning A Novel Perspective on Demonstration Selection. arXiv:2312.07476v2, 2023. \url{https://arxiv.org/abs/2312.07476v2}
\bibitem{ref9} Exploring the Relationship between In-Context Learning and Instruction Tuning. arXiv:2311.10367v1, 2023. \url{https://arxiv.org/abs/2311.10367v1}
\bibitem{ref10} In-context Learning and Gradient Descent Revisited. arXiv:2311.07772v4, 2023. \url{https://arxiv.org/abs/2311.07772v4}
\bibitem{ref11} Language Is Not All You Need Aligning Perception with Language Models. arXiv:2302.14045v2, 2023. \url{https://arxiv.org/abs/2302.14045v2}
\bibitem{ref12} The Dawn of LMMs Preliminary Explorations with GPT-4V(ision). arXiv:2309.17421v2, 2023. \url{https://arxiv.org/abs/2309.17421v2}
\bibitem{ref13} VL-ICL Bench The Devil in the Details of Benchmarking Multimodal In-Context Learning. arXiv:2403.13164v1, 2024. \url{https://arxiv.org/abs/2403.13164v1}
\bibitem{ref14} MMICL Empowering Vision-language Model with Multi-Modal In-Context Learning. arXiv:2309.07915v3, 2023. \url{https://arxiv.org/abs/2309.07915v3}
\bibitem{ref15} Many-Shot In-Context Learning. arXiv:2404.11018v1, 2024. \url{https://arxiv.org/abs/2404.11018v1}
\bibitem{ref16} Emergent Abilities in Reduced-Scale Generative Language Models. arXiv:2404.02204v1, 2024. \url{https://arxiv.org/abs/2404.02204v1}
\bibitem{ref17} An Exploration of In-Context Learning for Speech Language Model. arXiv:2310.12477v1, 2023. \url{https://arxiv.org/abs/2310.12477v1}
\bibitem{ref18} Rethinking the Role of Demonstrations What Makes In-Context Learning Work. arXiv:2202.12837v2, 2022. \url{https://arxiv.org/abs/2202.12837v2}
\bibitem{ref19} Dr.ICL Demonstration-Retrieved In-context Learning. arXiv:2305.14128v1, 2023. \url{https://arxiv.org/abs/2305.14128v1}
\bibitem{ref20} Coverage-based Example Selection for In-Context Learning. arXiv:2305.14907v3, 2023. \url{https://arxiv.org/abs/2305.14907v3}
\bibitem{ref21} Let's Learn Step by Step Enhancing In-Context Learning Ability with Curriculum Learning. arXiv:2402.10738v1, 2024. \url{https://arxiv.org/abs/2402.10738v1}
\bibitem{ref22} In-Context Learning Creates Task Vectors. arXiv:2310.15916v1, 2023. \url{https://arxiv.org/abs/2310.15916v1}
\bibitem{ref23} Rectifying Demonstration Shortcut in In-Context Learning. arXiv:2403.09488v3, 2024. \url{https://arxiv.org/abs/2403.09488v3}
\bibitem{ref24} Identifying and Analyzing Task-Encoding Tokens in Large Language Models. arXiv:2401.11323v2, 2024. \url{https://arxiv.org/abs/2401.11323v2}
\bibitem{ref25} Why Can GPT Learn In-Context Language Models Implicitly Perform Gradient Descent as Meta-Optimizers. arXiv:2212.10559v3, 2022. \url{https://arxiv.org/abs/2212.10559v3}
\bibitem{ref26} Revisiting the Hypothesis Do pretrained Transformers Learn In-Context by Gradient Descent. arXiv:2310.08540v4, 2023. \url{https://arxiv.org/abs/2310.08540v4}
\bibitem{ref27} Meta-Learning with Implicit Gradients. arXiv:1909.04630v1, 2019. \url{https://arxiv.org/abs/1909.04630v1}
\bibitem{ref28} OPT-IML Scaling Language Model Instruction Meta Learning through the Lens of Generalization. arXiv:2212.12017v3, 2022. \url{https://arxiv.org/abs/2212.12017v3}
\bibitem{ref29} The Strong Pull of Prior Knowledge in Large Language Models and Its Impact on Emotion Recognition. arXiv:2403.17125v1, 2024. \url{https://arxiv.org/abs/2403.17125v1}
\bibitem{ref30} Few-Shot Parameter-Efficient Fine-Tuning is Better and Cheaper than In-Context Learning. arXiv:2205.05638v2, 2022. \url{https://arxiv.org/abs/2205.05638v2}
\bibitem{ref31} How Does In-Context Learning Help Prompt Tuning. arXiv:2302.11521v1, 2023. \url{https://arxiv.org/abs/2302.11521v1}
\bibitem{ref32} The ICL Consistency Test. arXiv:2312.04945v1, 2023. \url{https://arxiv.org/abs/2312.04945v1}
\bibitem{ref33} On the Relation between Sensitivity and Accuracy in In-context Learning. arXiv:2209.07661v3, 2022. \url{https://arxiv.org/abs/2209.07661v3}
\bibitem{ref34} Data Curation Alone Can Stabilize In-context Learning. arXiv:2212.10378v2, 2022. \url{https://arxiv.org/abs/2212.10378v2}
\bibitem{ref35} In-context Learning with Retrieved Demonstrations for Language Models A Survey. arXiv:2401.11624v5, 2024. \url{https://arxiv.org/abs/2401.11624v5}
\bibitem{ref36} Ambiguity-Aware In-Context Learning with Large Language Models. arXiv:2309.07900v2, 2023. \url{https://arxiv.org/abs/2309.07900v2}
\bibitem{ref37} Mind the instructions a holistic evaluation of consistency and interactions in prompt-based learning. arXiv:2310.13486v1, 2023. \url{https://arxiv.org/abs/2310.13486v1}
\bibitem{ref38} A Study on the Calibration of In-context Learning. arXiv:2312.04021v4, 2023. \url{https://arxiv.org/abs/2312.04021v4}
\bibitem{ref39} On Task Performance and Model Calibration with Supervised and Self-Ensembled In-Context Learning. arXiv:2312.13772v2, 2023. \url{https://arxiv.org/abs/2312.13772v2}
\bibitem{ref40} How Robust are LLMs to In-Context Majority Label Bias. arXiv:2312.16549v1, 2023. \url{https://arxiv.org/abs/2312.16549v1}
\bibitem{ref41} Understanding In-Context Learning via Supportive Pretraining Data. arXiv:2306.15091v1, 2023. \url{https://arxiv.org/abs/2306.15091v1}
\bibitem{ref42} In-context Learning with Transformer Is Really Equivalent to a Contrastive Learning Pattern. arXiv:2310.13220v1, 2023. \url{https://arxiv.org/abs/2310.13220v1}
\bibitem{ref43} Linear Transformers are Versatile In-Context Learners. arXiv:2402.14180v1, 2024. \url{https://arxiv.org/abs/2402.14180v1}
\bibitem{ref44} In-Context Learning of a Linear Transformer Block Benefits of the MLP Component and One-Step GD Initialization. arXiv:2402.14951v1, 2024. \url{https://arxiv.org/abs/2402.14951v1}
\bibitem{ref45} Iterative Forward Tuning Boosts In-context Learning in Language Models. arXiv:2305.13016v2, 2023. \url{https://arxiv.org/abs/2305.13016v2}
\bibitem{ref46} What and How does In-Context Learning Learn Bayesian Model Averaging, Parameterization, and Generalization. arXiv:2305.19420v2, 2023. \url{https://arxiv.org/abs/2305.19420v2}
\bibitem{ref47} In-Context Learning through the Bayesian Prism. arXiv:2306.04891v2, 2023. \url{https://arxiv.org/abs/2306.04891v2}
\bibitem{ref48} Pretraining task diversity and the emergence of non-Bayesian in-context learning for regression. arXiv:2306.15063v2, 2023. \url{https://arxiv.org/abs/2306.15063v2}
\bibitem{ref49} How Many Pretraining Tasks Are Needed for In-Context Learning of Linear Regression. arXiv:2310.08391v2, 2023. \url{https://arxiv.org/abs/2310.08391v2}
\bibitem{ref50} In-Context Learning for MIMO Equalization Using Transformer-Based Sequence Models. arXiv:2311.06101v2, 2023. \url{https://arxiv.org/abs/2311.06101v2}
\bibitem{ref51} In-Context Learning Learns Label Relationships but Is Not Conventional Learning. arXiv:2307.12375v4, 2023. \url{https://arxiv.org/abs/2307.12375v4}
\bibitem{ref52} Data. arXiv:1801.04992v2, 2017. \url{https://arxiv.org/abs/1801.04992v2}
\bibitem{ref53} B-Splines. arXiv:2108.06617v1, 2021. \url{https://arxiv.org/abs/2108.06617v1}
\bibitem{ref54} In-context Learning and Induction Heads. arXiv:2209.11895v1, 2022. \url{https://arxiv.org/abs/2209.11895v1}
\bibitem{ref55} The mechanistic basis of data dependence and abrupt learning in an in-context classification task. arXiv:2312.03002v1, 2023. \url{https://arxiv.org/abs/2312.03002v1}
\bibitem{ref56} What needs to go right for an induction head A mechanistic study of in-context learning circuits and their formation. arXiv:2404.07129v1, 2024. \url{https://arxiv.org/abs/2404.07129v1}
\bibitem{ref57} The Evolution of Statistical Induction Heads In-Context Learning Markov Chains. arXiv:2402.11004v1, 2024. \url{https://arxiv.org/abs/2402.11004v1}
\bibitem{ref58} Identifying Semantic Induction Heads to Understand In-Context Learning. arXiv:2402.13055v1, 2024. \url{https://arxiv.org/abs/2402.13055v1}
\bibitem{ref59} Successor Heads Recurring, Interpretable Attention Heads In The Wild. arXiv:2312.09230v1, 2023. \url{https://arxiv.org/abs/2312.09230v1}
\bibitem{ref60} In-Context Language Learning Architectures and Algorithms. arXiv:2401.12973v2, 2024. \url{https://arxiv.org/abs/2401.12973v2}
\bibitem{ref61} Circuit Component Reuse Across Tasks in Transformer Language Models. arXiv:2310.08744v2, 2023. \url{https://arxiv.org/abs/2310.08744v2}
\bibitem{ref62} In-Context Learning State Vector with Inner and Momentum Optimization. arXiv:2404.11225v1, 2024. \url{https://arxiv.org/abs/2404.11225v1}
\bibitem{ref63} How Do Transformers Learn In-Context Beyond Simple Functions A Case Study on Learning with Representations. arXiv:2310.10616v1, 2023. \url{https://arxiv.org/abs/2310.10616v1}
\bibitem{ref64} Schema-learning and rebinding as mechanisms of in-context learning and emergence. arXiv:2307.01201v1, 2023. \url{https://arxiv.org/abs/2307.01201v1}
\bibitem{ref65} In-Context Learning with Transformers Softmax Attention Adapts to Function Lipschitzness. arXiv:2402.11639v1, 2024. \url{https://arxiv.org/abs/2402.11639v1}
\bibitem{ref66} Label Words are Anchors An Information Flow Perspective for Understanding In-Context Learning. arXiv:2305.14160v4, 2023. \url{https://arxiv.org/abs/2305.14160v4}
\bibitem{ref67} Parallel Structures in Pre-training Data Yield In-Context Learning. arXiv:2402.12530v1, 2024. \url{https://arxiv.org/abs/2402.12530v1}
\bibitem{ref68} Concept-aware Training Improves In-context Learning Ability of Language Models. arXiv:2305.13775v1, 2023. \url{https://arxiv.org/abs/2305.13775v1}
\bibitem{ref69} Pre-training Text-to-Text Transformers for Concept-centric Common Sense. arXiv:2011.07956v2, 2020. \url{https://arxiv.org/abs/2011.07956v2}
\bibitem{ref70} Beyond Scale the Diversity Coefficient as a Data Quality Metric Demonstrates LLMs are Pre-trained on Formally Diverse Data. arXiv:2306.13840v2, 2023. \url{https://arxiv.org/abs/2306.13840v2}
\bibitem{ref71} Skill-it! A Data-Driven Skills Framework for Understanding and Training Language Models. arXiv:2307.14430v1, 2023. \url{https://arxiv.org/abs/2307.14430v1}
\bibitem{ref72} On the Transferability of Pre-trained Language Models A Study from Artificial Datasets. arXiv:2109.03537v2, 2021. \url{https://arxiv.org/abs/2109.03537v2}
\bibitem{ref73} The Transient Nature of Emergent In-Context Learning in Transformers. arXiv:2311.08360v3, 2023. \url{https://arxiv.org/abs/2311.08360v3}
\bibitem{ref74} Breaking through the learning plateaus of in-context learning in Transformer. arXiv:2309.06054v2, 2023. \url{https://arxiv.org/abs/2309.06054v2}
\bibitem{ref75} Human Curriculum Effects Emerge with In-Context Learning in Neural Networks. arXiv:2402.08674v1, 2024. \url{https://arxiv.org/abs/2402.08674v1}
\bibitem{ref76} Generalization Bounds for Meta-Learning via PAC-Bayes and Uniform Stability. arXiv:2102.06589v3, 2021. \url{https://arxiv.org/abs/2102.06589v3}
\bibitem{ref77} A Limitation of the PAC-Bayes Framework. arXiv:2006.13508v3, 2020. \url{https://arxiv.org/abs/2006.13508v3}
\bibitem{ref78} A Unified Model and Dimension for Interactive Estimation. arXiv:2306.06184v1, 2023. \url{https://arxiv.org/abs/2306.06184v1}
\bibitem{ref79} Generalization Bounds of SGLD for Non-convex Learning Two Theoretical Viewpoints. arXiv:1707.05947v1, 2017. \url{https://arxiv.org/abs/1707.05947v1}
\bibitem{ref80} On Lipschitz Continuity and Smoothness of Loss Functions in Learning to Rank. arXiv:1405.0586v3, 2014. \url{https://arxiv.org/abs/1405.0586v3}
\bibitem{ref81} Certified Robust Models with Slack Control and Large Lipschitz Constants. arXiv:2309.06166v1, 2023. \url{https://arxiv.org/abs/2309.06166v1}
\bibitem{ref82} Unifying PAC and Regret Uniform PAC Bounds for Episodic Reinforcement Learning. arXiv:1703.07710v3, 2017. \url{https://arxiv.org/abs/1703.07710v3}
\bibitem{ref83} Uniform-PAC Bounds for Reinforcement Learning with Linear Function Approximation. arXiv:2106.11612v2, 2021. \url{https://arxiv.org/abs/2106.11612v2}
\bibitem{ref84} Convergence of Two-Layer Regression with Nonlinear Units. arXiv:2308.08358v1, 2023. \url{https://arxiv.org/abs/2308.08358v1}
\bibitem{ref85} Rademacher Generalization Bounds for Classifier Chains. arXiv:1807.10166v1, 2018. \url{https://arxiv.org/abs/1807.10166v1}
\bibitem{ref86} Is attention required for ICL Exploring the Relationship Between Model Architecture and In-Context Learning Ability. arXiv:2310.08049v3, 2023. \url{https://arxiv.org/abs/2310.08049v3}
\bibitem{ref87} Mamba Linear-Time Sequence Modeling with Selective State Spaces. arXiv:2312.00752v1, 2023. \url{https://arxiv.org/abs/2312.00752v1}
\bibitem{ref88} Can Mamba Learn How to Learn A Comparative Study on In-Context Learning Tasks. arXiv:2402.04248v2, 2024. \url{https://arxiv.org/abs/2402.04248v2}
\bibitem{ref89} Is Mamba Capable of In-Context Learning. arXiv:2402.03170v2, 2024. \url{https://arxiv.org/abs/2402.03170v2}
\bibitem{ref90} Superiority of Softmax Unveiling the Performance Edge Over Linear Attention. arXiv:2310.11685v1, 2023. \url{https://arxiv.org/abs/2310.11685v1}
\bibitem{ref91} Superiority of Multi-Head Attention in In-Context Linear Regression. arXiv:2401.17426v1, 2024. \url{https://arxiv.org/abs/2401.17426v1}
\bibitem{ref92} Attention is Not All You Need Pure Attention Loses Rank Doubly Exponentially with Depth. arXiv:2103.03404v2, 2021. \url{https://arxiv.org/abs/2103.03404v2}
\bibitem{ref93} Self-attention Networks Localize When QK-eigenspectrum Concentrates. arXiv:2402.02098v1, 2024. \url{https://arxiv.org/abs/2402.02098v1}
\bibitem{ref94} Effects of Parameter Norm Growth During Transformer Training Inductive Bias from Gradient Descent. arXiv:2010.09697v5, 2020. \url{https://arxiv.org/abs/2010.09697v5}
\bibitem{ref95} Adaptively Sparse Transformers. arXiv:1909.00015v2, 2019. \url{https://arxiv.org/abs/1909.00015v2}
\bibitem{ref96} Understanding Multi-Head Attention in Abstractive Summarization. arXiv:1911.03898v1, 2019. \url{https://arxiv.org/abs/1911.03898v1}
\bibitem{ref97} Sparse Modular Activation for Efficient Sequence Modeling. arXiv:2306.11197v4, 2023. \url{https://arxiv.org/abs/2306.11197v4}
\bibitem{ref98} Less is More! A slim architecture for optimal language translation. arXiv:2305.10991v1, 2023. \url{https://arxiv.org/abs/2305.10991v1}
\bibitem{ref99} Max-Margin Token Selection in Attention Mechanism. arXiv:2306.13596v4, 2023. \url{https://arxiv.org/abs/2306.13596v4}
\bibitem{ref100} The Closeness of In-Context Learning and Weight Shifting for Softmax Regression. arXiv:2304.13276v1, 2023. \url{https://arxiv.org/abs/2304.13276v1}
\bibitem{ref101} Towards Understanding the Word Sensitivity of Attention Layers A Study via Random Features. arXiv:2402.02969v1, 2024. \url{https://arxiv.org/abs/2402.02969v1}
\bibitem{ref102} What Can Transformers Learn In-Context A Case Study of Simple Function Classes. arXiv:2208.01066v3, 2022. \url{https://arxiv.org/abs/2208.01066v3}
\bibitem{ref103} What learning algorithm is in-context learning Investigations with linear models. arXiv:2211.15661v3, 2022. \url{https://arxiv.org/abs/2211.15661v3}
\bibitem{ref104} Transformers Learn Higher-Order Optimization Methods for In-Context Learning A Study with Linear Models. arXiv:2310.17086v1, 2023. \url{https://arxiv.org/abs/2310.17086v1}
\bibitem{ref105} How Well Can Transformers Emulate In-context Newton's Method. arXiv:2403.03183v1, 2024. \url{https://arxiv.org/abs/2403.03183v1}
\bibitem{ref106} Can Transformers Learn Sequential Function Classes In Context. arXiv:2312.12655v2, 2023. \url{https://arxiv.org/abs/2312.12655v2}
\bibitem{ref107} Compositional Capabilities of Autoregressive Transformers A Study on Synthetic, Interpretable Tasks. arXiv:2311.12997v2, 2023. \url{https://arxiv.org/abs/2311.12997v2}
\bibitem{ref108} A Theory of Emergent In-Context Learning as Implicit Structure Induction. arXiv:2303.07971v1, 2023. \url{https://arxiv.org/abs/2303.07971v1}
\bibitem{ref109} Characterizing Intrinsic Compositionality in Transformers with Tree Projections. arXiv:2211.01288v2, 2022. \url{https://arxiv.org/abs/2211.01288v2}
\bibitem{ref110} Transformers as Statisticians Provable In-Context Learning with In-Context Algorithm Selection. arXiv:2306.04637v2, 2023. \url{https://arxiv.org/abs/2306.04637v2}
\bibitem{ref111} Implicit regularization of multi-task learning and finetuning in overparameterized neural networks. arXiv:2310.02396v2, 2023. \url{https://arxiv.org/abs/2310.02396v2}
\bibitem{ref112} In-context Learning Generalizes, But Not Always Robustly The Case of Syntax. arXiv:2311.07811v2, 2023. \url{https://arxiv.org/abs/2311.07811v2}
\bibitem{ref113} How Well Does GPT-4V(ision) Adapt to Distribution Shifts A Preliminary Investigation. arXiv:2312.07424v3, 2023. \url{https://arxiv.org/abs/2312.07424v3}
\bibitem{ref114} On the Out-Of-Distribution Generalization of Multimodal Large Language Models. arXiv:2402.06599v1, 2024. \url{https://arxiv.org/abs/2402.06599v1}
\bibitem{ref115} Hidden Covariate Shift A Minimal Assumption For Domain Adaptation. arXiv:1907.12299v1, 2019. \url{https://arxiv.org/abs/1907.12299v1}
\bibitem{ref116} Kernel Robust Bias-Aware Prediction under Covariate Shift. arXiv:1712.10050v1, 2017. \url{https://arxiv.org/abs/1712.10050v1}
\bibitem{ref117} Pretraining Data Mixtures Enable Narrow Model Selection Capabilities in Transformer Models. arXiv:2311.00871v1, 2023. \url{https://arxiv.org/abs/2311.00871v1}
\bibitem{ref118} Explore and Exploit the Diverse Knowledge in Model Zoo for Domain Generalization. arXiv:2306.02595v1, 2023. \url{https://arxiv.org/abs/2306.02595v1}
\bibitem{ref119} Hungry Hungry Hippos Towards Language Modeling with State Space Models. arXiv:2212.14052v3, 2022. \url{https://arxiv.org/abs/2212.14052v3}
\bibitem{ref120} Compositional Attention Disentangling Search and Retrieval. arXiv:2110.09419v2, 2021. \url{https://arxiv.org/abs/2110.09419v2}
\bibitem{ref121} Towards Better Modeling Hierarchical Structure for Self-Attention with Ordered Neurons. arXiv:1909.01562v2, 2019. \url{https://arxiv.org/abs/1909.01562v2}
\bibitem{ref122} Improving In-context Learning via Bidirectional Alignment. arXiv:2312.17055v1, 2023. \url{https://arxiv.org/abs/2312.17055v1}
\bibitem{ref123} Rethinking Query-Key Pairwise Interactions in Vision Transformers. arXiv:2207.00188v2, 2022. \url{https://arxiv.org/abs/2207.00188v2}
\bibitem{ref124} Cross-Architecture Transfer Learning for Linear-Cost Inference Transformers. arXiv:2404.02684v1, 2024. \url{https://arxiv.org/abs/2404.02684v1}
\bibitem{ref125} Ungrammatical-syntax-based In-context Example Selection for Grammatical Error Correction. arXiv:2403.19283v1, 2024. \url{https://arxiv.org/abs/2403.19283v1}
\bibitem{ref126} Retrieval-augmented Multi-modal Chain-of-Thoughts Reasoning for Large Language Models. arXiv:2312.01714v2, 2023. \url{https://arxiv.org/abs/2312.01714v2}
\bibitem{ref127} Retrieve Anything To Augment Large Language Models. arXiv:2310.07554v2, 2023. \url{https://arxiv.org/abs/2310.07554v2}
\bibitem{ref128} PaRaDe Passage Ranking using Demonstrations with Large Language Models. arXiv:2310.14408v1, 2023. \url{https://arxiv.org/abs/2310.14408v1}
\bibitem{ref129} Active Learning Principles for In-Context Learning with Large Language Models. arXiv:2305.14264v2, 2023. \url{https://arxiv.org/abs/2305.14264v2}
\bibitem{ref130} Towards Informative Few-Shot Prompt with Maximum Information Gain for In-Context Learning. arXiv:2310.08923v1, 2023. \url{https://arxiv.org/abs/2310.08923v1}
\bibitem{ref131} Prototypical Calibration for Few-shot Learning of Language Models. arXiv:2205.10183v2, 2022. \url{https://arxiv.org/abs/2205.10183v2}
\bibitem{ref132} RetICL Sequential Retrieval of In-Context Examples with Reinforcement Learning. arXiv:2305.14502v2, 2023. \url{https://arxiv.org/abs/2305.14502v2}
\bibitem{ref133} \$Se\textasciicircum{}2\$ Sequential Example Selection for In-Context Learning. arXiv:2402.13874v2, 2024. \url{https://arxiv.org/abs/2402.13874v2}
\bibitem{ref134} Direct and indirect reinforcement learning. arXiv:1912.10600v2, 2019. \url{https://arxiv.org/abs/1912.10600v2}
\bibitem{ref135} Hierarchical Imitation and Reinforcement Learning. arXiv:1803.00590v2, 2018. \url{https://arxiv.org/abs/1803.00590v2}
\bibitem{ref136} Learning Skills to Navigate without a Master A Sequential Multi-Policy Reinforcement Learning Algorithm. arXiv:2102.00168v3, 2021. \url{https://arxiv.org/abs/2102.00168v3}
\bibitem{ref137} Divide-and-Conquer Monte Carlo Tree Search For Goal-Directed Planning. arXiv:2004.11410v1, 2020. \url{https://arxiv.org/abs/2004.11410v1}
\bibitem{ref138} Coverage-centric Coreset Selection for High Pruning Rates. arXiv:2210.15809v2, 2022. \url{https://arxiv.org/abs/2210.15809v2}
\bibitem{ref139} RD-DPP Rate-Distortion Theory Meets Determinantal Point Process to Diversify Learning Data Samples. arXiv:2304.04137v2, 2023. \url{https://arxiv.org/abs/2304.04137v2}
\bibitem{ref140} A Unified Algorithm Framework for Unsupervised Discovery of Skills based on Determinantal Point Process. arXiv:2212.00211v3, 2022. \url{https://arxiv.org/abs/2212.00211v3}
\bibitem{ref141} Contrastive Demonstration Tuning for Pre-trained Language Models. arXiv:2204.04392v4, 2022. \url{https://arxiv.org/abs/2204.04392v4}
\bibitem{ref142} Objectives Are All You Need Solving Deceptive Problems Without Explicit Diversity Maintenance. arXiv:2311.02283v1, 2023. \url{https://arxiv.org/abs/2311.02283v1}
\bibitem{ref143} Contributing Dimension Structure of Deep Feature for Coreset Selection. arXiv:2401.16193v2, 2024. \url{https://arxiv.org/abs/2401.16193v2}
\bibitem{ref144} Data Diversity Matters for Robust Instruction Tuning. arXiv:2311.14736v2, 2023. \url{https://arxiv.org/abs/2311.14736v2}
\bibitem{ref145} Pattern-Aware Chain-of-Thought Prompting in Large Language Models. arXiv:2404.14812v1, 2024. \url{https://arxiv.org/abs/2404.14812v1}
\bibitem{ref146} Latent Skill Discovery for Chain-of-Thought Reasoning. arXiv:2312.04684v1, 2023. \url{https://arxiv.org/abs/2312.04684v1}
\bibitem{ref147} Synthetic Prompting Generating Chain-of-Thought Demonstrations for Large Language Models. arXiv:2302.00618v1, 2023. \url{https://arxiv.org/abs/2302.00618v1}
\bibitem{ref148} Automatic Prompt Augmentation and Selection with Chain-of-Thought from Labeled Data. arXiv:2302.12822v3, 2023. \url{https://arxiv.org/abs/2302.12822v3}
\bibitem{ref149} Self-Consistency Improves Chain of Thought Reasoning in Language Models. arXiv:2203.11171v4, 2022. \url{https://arxiv.org/abs/2203.11171v4}
\bibitem{ref150} Universal Self-Consistency for Large Language Model Generation. arXiv:2311.17311v1, 2023. \url{https://arxiv.org/abs/2311.17311v1}
\bibitem{ref151} Ask One More Time Self-Agreement Improves Reasoning of Language Models in (Almost) All Scenarios. arXiv:2311.08154v2, 2023. \url{https://arxiv.org/abs/2311.08154v2}
\bibitem{ref152} Progressive-Hint Prompting Improves Reasoning in Large Language Models. arXiv:2304.09797v5, 2023. \url{https://arxiv.org/abs/2304.09797v5}
\bibitem{ref153} Enhancing Chain-of-Thoughts Prompting with Iterative Bootstrapping in Large Language Models. arXiv:2304.11657v3, 2023. \url{https://arxiv.org/abs/2304.11657v3}
\bibitem{ref154} P-Tuning v2 Prompt Tuning Can Be Comparable to Fine-tuning Universally Across Scales and Tasks. arXiv:2110.07602v3, 2021. \url{https://arxiv.org/abs/2110.07602v3}
\bibitem{ref155} Vector-Quantized Input-Contextualized Soft Prompts for Natural Language Understanding. arXiv:2205.11024v2, 2022. \url{https://arxiv.org/abs/2205.11024v2}
\bibitem{ref156} MetaPrompting Learning to Learn Better Prompts. arXiv:2209.11486v4, 2022. \url{https://arxiv.org/abs/2209.11486v4}
\bibitem{ref157} Self-supervised Meta-Prompt Learning with Meta-Gradient Regularization for Few-shot Generalization. arXiv:2303.12314v4, 2023. \url{https://arxiv.org/abs/2303.12314v4}
\bibitem{ref158} Gradient-Regulated Meta-Prompt Learning for Generalizable Vision-Language Models. arXiv:2303.06571v2, 2023. \url{https://arxiv.org/abs/2303.06571v2}
\bibitem{ref159} Residual Prompt Tuning Improving Prompt Tuning with Residual Reparameterization. arXiv:2305.03937v1, 2023. \url{https://arxiv.org/abs/2305.03937v1}
\bibitem{ref160} DePT Decomposed Prompt Tuning for Parameter-Efficient Fine-tuning. arXiv:2309.05173v5, 2023. \url{https://arxiv.org/abs/2309.05173v5}
\bibitem{ref161} Effective Structured Prompting by Meta-Learning and Representative Verbalizer. arXiv:2306.00618v2, 2023. \url{https://arxiv.org/abs/2306.00618v2}
\bibitem{ref162} MPrompt Exploring Multi-level Prompt Tuning for Machine Reading Comprehension. arXiv:2310.18167v1, 2023. \url{https://arxiv.org/abs/2310.18167v1}
\bibitem{ref163} Discourse-Aware Soft Prompting for Text Generation. arXiv:2112.05717v2, 2021. \url{https://arxiv.org/abs/2112.05717v2}
\bibitem{ref164} A Bayesian approach for prompt optimization in pre-trained language models. arXiv:2312.00471v1, 2023. \url{https://arxiv.org/abs/2312.00471v1}
\bibitem{ref165} Hard Prompts Made Easy Gradient-Based Discrete Optimization for Prompt Tuning and Discovery. arXiv:2302.03668v2, 2023. \url{https://arxiv.org/abs/2302.03668v2}
\bibitem{ref166} Automatic Label Sequence Generation for Prompting Sequence-to-sequence Models. arXiv:2209.09401v1, 2022. \url{https://arxiv.org/abs/2209.09401v1}
\bibitem{ref167} Black-Box Tuning for Language-Model-as-a-Service. arXiv:2201.03514v4, 2022. \url{https://arxiv.org/abs/2201.03514v4}
\bibitem{ref168} Black-box Prompt Tuning with Subspace Learning. arXiv:2305.03518v1, 2023. \url{https://arxiv.org/abs/2305.03518v1}
\bibitem{ref169} Promptbreeder Self-Referential Self-Improvement Via Prompt Evolution. arXiv:2309.16797v1, 2023. \url{https://arxiv.org/abs/2309.16797v1}
\bibitem{ref170} PREFER Prompt Ensemble Learning via Feedback-Reflect-Refine. arXiv:2308.12033v1, 2023. \url{https://arxiv.org/abs/2308.12033v1}
\bibitem{ref171} Dialogue for Prompting a Policy-Gradient-Based Discrete Prompt Generation for Few-shot Learning. arXiv:2308.07272v2, 2023. \url{https://arxiv.org/abs/2308.07272v2}
\bibitem{ref172} PRompt Optimization in Multi-Step Tasks (PROMST) Integrating Human Feedback and Preference Alignment. arXiv:2402.08702v2, 2024. \url{https://arxiv.org/abs/2402.08702v2}
\bibitem{ref173} Strings from the Library of Babel Random Sampling as a Strong Baseline for Prompt Optimisation. arXiv:2311.09569v2, 2023. \url{https://arxiv.org/abs/2311.09569v2}
\bibitem{ref174} Evaluating the Robustness of Discrete Prompts. arXiv:2302.05619v1, 2023. \url{https://arxiv.org/abs/2302.05619v1}
\bibitem{ref175} Flatness-Aware Prompt Selection Improves Accuracy and Sample Efficiency. arXiv:2305.10713v2, 2023. \url{https://arxiv.org/abs/2305.10713v2}
\bibitem{ref176} Larger language models do in-context learning differently. arXiv:2303.03846v2, 2023. \url{https://arxiv.org/abs/2303.03846v2}
\bibitem{ref177} Mitigating Label Biases for In-context Learning. arXiv:2305.19148v3, 2023. \url{https://arxiv.org/abs/2305.19148v3}
\bibitem{ref178} Batch Calibration Rethinking Calibration for In-Context Learning and Prompt Engineering. arXiv:2309.17249v2, 2023. \url{https://arxiv.org/abs/2309.17249v2}
\bibitem{ref179} Generative Calibration for In-context Learning. arXiv:2310.10266v1, 2023. \url{https://arxiv.org/abs/2310.10266v1}
\bibitem{ref180} Addressing Order Sensitivity of In-Context Demonstration Examples in Causal Language Models. arXiv:2402.15637v1, 2024. \url{https://arxiv.org/abs/2402.15637v1}
\bibitem{ref181} Premise Order Matters in Reasoning with Large Language Models. arXiv:2402.08939v2, 2024. \url{https://arxiv.org/abs/2402.08939v2}
\bibitem{ref182} Large Language Models Sensitivity to The Order of Options in Multiple-Choice Questions. arXiv:2308.11483v1, 2023. \url{https://arxiv.org/abs/2308.11483v1}
\bibitem{ref183} Position-Aware Parameter Efficient Fine-Tuning Approach for Reducing Positional Bias in LLMs. arXiv:2404.01430v1, 2024. \url{https://arxiv.org/abs/2404.01430v1}
\bibitem{ref184} Instruction Position Matters in Sequence Generation with Large Language Models. arXiv:2308.12097v1, 2023. \url{https://arxiv.org/abs/2308.12097v1}
\bibitem{ref185} Found in the Middle Permutation Self-Consistency Improves Listwise Ranking in Large Language Models. arXiv:2310.07712v2, 2023. \url{https://arxiv.org/abs/2310.07712v2}
\bibitem{ref186} Are We Falling in a Middle-Intelligence Trap An Analysis and Mitigation of the Reversal Curse. arXiv:2311.07468v2, 2023. \url{https://arxiv.org/abs/2311.07468v2}
\bibitem{ref187} Causal Prompting Debiasing Large Language Model Prompting based on Front-Door Adjustment. arXiv:2403.02738v1, 2024. \url{https://arxiv.org/abs/2403.02738v1}
\bibitem{ref188} MetaICL Learning to Learn In Context. arXiv:2110.15943v2, 2021. \url{https://arxiv.org/abs/2110.15943v2}
\bibitem{ref189} How does Multi-Task Training Affect Transformer In-Context Capabilities Investigations with Function Classes. arXiv:2404.03558v1, 2024. \url{https://arxiv.org/abs/2404.03558v1}
\bibitem{ref190} Concept-aware Data Construction Improves In-context Learning of Language Models. arXiv:2403.09703v1, 2024. \url{https://arxiv.org/abs/2403.09703v1}
\bibitem{ref191} Bridging Multi-Task Learning and Meta-Learning Towards Efficient Training and Effective Adaptation. arXiv:2106.09017v1, 2021. \url{https://arxiv.org/abs/2106.09017v1}
\bibitem{ref192} Towards Task Sampler Learning for Meta-Learning. arXiv:2307.08924v3, 2023. \url{https://arxiv.org/abs/2307.08924v3}
\bibitem{ref193} In-context Learning Distillation Transferring Few-shot Learning Ability of Pre-trained Language Models. arXiv:2212.10670v1, 2022. \url{https://arxiv.org/abs/2212.10670v1}
\bibitem{ref194} LESS Selecting Influential Data for Targeted Instruction Tuning. arXiv:2402.04333v2, 2024. \url{https://arxiv.org/abs/2402.04333v2}
\bibitem{ref195} Multimodal Instruction Tuning with Conditional Mixture of LoRA. arXiv:2402.15896v1, 2024. \url{https://arxiv.org/abs/2402.15896v1}
\bibitem{ref196} Prompt Tuning based Adapter for Vision-Language Model Adaption. arXiv:2303.15234v1, 2023. \url{https://arxiv.org/abs/2303.15234v1}
\bibitem{ref197} Batch-ICL Effective, Efficient, and Order-Agnostic In-Context Learning. arXiv:2401.06469v2, 2024. \url{https://arxiv.org/abs/2401.06469v2}
\bibitem{ref198} Fine-tune Language Models to Approximate Unbiased In-context Learning. arXiv:2310.03331v1, 2023. \url{https://arxiv.org/abs/2310.03331v1}
\bibitem{ref199} A Survey on Curriculum Learning. arXiv:2010.13166v2, 2020. \url{https://arxiv.org/abs/2010.13166v2}
\bibitem{ref200} Curriculum Learning A Survey. arXiv:2101.10382v3, 2021. \url{https://arxiv.org/abs/2101.10382v3}
\bibitem{ref201} A Probabilistic Interpretation of Self-Paced Learning with Applications to Reinforcement Learning. arXiv:2102.13176v3, 2021. \url{https://arxiv.org/abs/2102.13176v3}
\bibitem{ref202} Let the Model Decide its Curriculum for Multitask Learning. arXiv:2205.09898v2, 2022. \url{https://arxiv.org/abs/2205.09898v2}
\bibitem{ref203} Curriculum Learning with Adam The Devil Is in the Wrong Details. arXiv:2308.12202v1, 2023. \url{https://arxiv.org/abs/2308.12202v1}
\bibitem{ref204} In-sample Curriculum Learning by Sequence Completion for Natural Language Generation. arXiv:2211.11297v2, 2022. \url{https://arxiv.org/abs/2211.11297v2}
\bibitem{ref205} Batch Prompting Efficient Inference with Large Language Model APIs. arXiv:2301.08721v2, 2023. \url{https://arxiv.org/abs/2301.08721v2}
\bibitem{ref206} \$k\$NN Prompting Beyond-Context Learning with Calibration-Free Nearest Neighbor Inference. arXiv:2303.13824v1, 2023. \url{https://arxiv.org/abs/2303.13824v1}
\bibitem{ref207} Compressed Context Memory For Online Language Model Interaction. arXiv:2312.03414v2, 2023. \url{https://arxiv.org/abs/2312.03414v2}
\bibitem{ref208} Structured Prompting Scaling In-Context Learning to 1,000 Examples. arXiv:2212.06713v1, 2022. \url{https://arxiv.org/abs/2212.06713v1}
\bibitem{ref209} Context Compression for Auto-regressive Transformers with Sentinel Tokens. arXiv:2310.08152v2, 2023. \url{https://arxiv.org/abs/2310.08152v2}
\bibitem{ref210} Model Tells You What to Discard Adaptive KV Cache Compression for LLMs. arXiv:2310.01801v3, 2023. \url{https://arxiv.org/abs/2310.01801v3}
\bibitem{ref211} XC-Cache Cross-Attending to Cached Context for Efficient LLM Inference. arXiv:2404.15420v1, 2024. \url{https://arxiv.org/abs/2404.15420v1}
\bibitem{ref212} Language Models are Few-Shot Learners. arXiv:2005.14165v4, 2020. \url{https://arxiv.org/abs/2005.14165v4}
\bibitem{ref213} Foundations and Trends in Multimodal Machine Learning Principles, Challenges, and Open Questions. arXiv:2209.03430v2, 2022. \url{https://arxiv.org/abs/2209.03430v2}
\bibitem{ref214} Multimodal Machine Learning A Survey and Taxonomy. arXiv:1705.09406v2, 2017. \url{https://arxiv.org/abs/1705.09406v2}
\bibitem{ref215} SpeechGPT Empowering Large Language Models with Intrinsic Cross-Modal Conversational Abilities. arXiv:2305.11000v2, 2023. \url{https://arxiv.org/abs/2305.11000v2}
\bibitem{ref216} SVIT Scaling up Visual Instruction Tuning. arXiv:2307.04087v3, 2023. \url{https://arxiv.org/abs/2307.04087v3}
\bibitem{ref217} MMICT Boosting Multi-Modal Fine-Tuning with In-Context Examples. arXiv:2312.06363v2, 2023. \url{https://arxiv.org/abs/2312.06363v2}
\bibitem{ref218} Lightweight In-Context Tuning for Multimodal Unified Models. arXiv:2310.05109v1, 2023. \url{https://arxiv.org/abs/2310.05109v1}
\bibitem{ref219} UNIMO-3 Multi-granularity Interaction for Vision-Language Representation Learning. arXiv:2305.13697v1, 2023. \url{https://arxiv.org/abs/2305.13697v1}
\bibitem{ref220} What Makes Multimodal In-Context Learning Work. arXiv:2404.15736v2, 2024. \url{https://arxiv.org/abs/2404.15736v2}
\bibitem{ref221} Understanding and Improving In-Context Learning on Vision-language Models. arXiv:2311.18021v1, 2023. \url{https://arxiv.org/abs/2311.18021v1}
\bibitem{ref222} Multimodal Representation Learning by Alternating Unimodal Adaptation. arXiv:2311.10707v2, 2023. \url{https://arxiv.org/abs/2311.10707v2}
\bibitem{ref223} Hybrid Contrastive Learning of Tri-Modal Representation for Multimodal Sentiment Analysis. arXiv:2109.01797v1, 2021. \url{https://arxiv.org/abs/2109.01797v1}
\bibitem{ref224} Deeply Coupled Cross-Modal Prompt Learning. arXiv:2305.17903v3, 2023. \url{https://arxiv.org/abs/2305.17903v3}
\bibitem{ref225} Token-Level Contrastive Learning with Modality-Aware Prompting for Multimodal Intent Recognition. arXiv:2312.14667v1, 2023. \url{https://arxiv.org/abs/2312.14667v1}
\bibitem{ref226} Can MLLMs Perform Text-to-Image In-Context Learning. arXiv:2402.01293v2, 2024. \url{https://arxiv.org/abs/2402.01293v2}
\bibitem{ref227} High-Modality Multimodal Transformer Quantifying Modality \& Interaction Heterogeneity for High-Modality Representation Learning. arXiv:2203.01311v4, 2022. \url{https://arxiv.org/abs/2203.01311v4}
\bibitem{ref228} IMProv Inpainting-based Multimodal Prompting for Computer Vision Tasks. arXiv:2312.01771v1, 2023. \url{https://arxiv.org/abs/2312.01771v1}
\bibitem{ref229} Instruct Me More! Random Prompting for Visual In-Context Learning. arXiv:2311.03648v1, 2023. \url{https://arxiv.org/abs/2311.03648v1}
\bibitem{ref230} Exploring Effective Factors for Improving Visual In-Context Learning. arXiv:2304.04748v1, 2023. \url{https://arxiv.org/abs/2304.04748v1}
\bibitem{ref231} In-Context Learning Unlocked for Diffusion Models. arXiv:2305.01115v2, 2023. \url{https://arxiv.org/abs/2305.01115v2}
\bibitem{ref232} Context Diffusion In-Context Aware Image Generation. arXiv:2312.03584v1, 2023. \url{https://arxiv.org/abs/2312.03584v1}
\bibitem{ref233} InstructDiffusion A Generalist Modeling Interface for Vision Tasks. arXiv:2309.03895v1, 2023. \url{https://arxiv.org/abs/2309.03895v1}
\bibitem{ref234} Im-Promptu In-Context Composition from Image Prompts. arXiv:2305.17262v3, 2023. \url{https://arxiv.org/abs/2305.17262v3}
\bibitem{ref235} All in an Aggregated Image for In-Image Learning. arXiv:2402.17971v2, 2024. \url{https://arxiv.org/abs/2402.17971v2}
\bibitem{ref236} Visual In-Context Learning for Large Vision-Language Models. arXiv:2402.11574v1, 2024. \url{https://arxiv.org/abs/2402.11574v1}
\bibitem{ref237} MaPLe Multi-modal Prompt Learning. arXiv:2210.03117v3, 2022. \url{https://arxiv.org/abs/2210.03117v3}
\bibitem{ref238} APLe Token-Wise Adaptive for Multi-Modal Prompt Learning. arXiv:2401.06827v2, 2024. \url{https://arxiv.org/abs/2401.06827v2}
\bibitem{ref239} Progressive Multi-modal Conditional Prompt Tuning. arXiv:2404.11864v2, 2024. \url{https://arxiv.org/abs/2404.11864v2}
\bibitem{ref240} APoLLo Unified Adapter and Prompt Learning for Vision Language Models. arXiv:2312.01564v1, 2023. \url{https://arxiv.org/abs/2312.01564v1}
\bibitem{ref241} Scaffolding Coordinates to Promote Vision-Language Coordination in Large Multi-Modal Models. arXiv:2402.12058v1, 2024. \url{https://arxiv.org/abs/2402.12058v1}
\bibitem{ref242} BBA Bi-Modal Behavioral Alignment for Reasoning with Large Vision-Language Models. arXiv:2402.13577v1, 2024. \url{https://arxiv.org/abs/2402.13577v1}
\bibitem{ref243} Tackling Vision Language Tasks Through Learning Inner Monologues. arXiv:2308.09970v1, 2023. \url{https://arxiv.org/abs/2308.09970v1}
\bibitem{ref244} Lost in Translation When GPT-4V(ision) Can't See Eye to Eye with Text. A Vision-Language-Consistency Analysis of VLLMs and Beyond. arXiv:2310.12520v1, 2023. \url{https://arxiv.org/abs/2310.12520v1}
\bibitem{ref245} Manipulating the Label Space for In-Context Classification. arXiv:2312.00351v2, 2023. \url{https://arxiv.org/abs/2312.00351v2}
\bibitem{ref246} CaMML Context-Aware Multimodal Learner for Large Models. arXiv:2401.03149v2, 2024. \url{https://arxiv.org/abs/2401.03149v2}
\bibitem{ref247} UniAdapter Unified Parameter-Efficient Transfer Learning for Cross-modal Modeling. arXiv:2302.06605v2, 2023. \url{https://arxiv.org/abs/2302.06605v2}
\bibitem{ref248} Cross-Modal Adapter Parameter-Efficient Transfer Learning Approach for Vision-Language Models. arXiv:2404.12588v1, 2024. \url{https://arxiv.org/abs/2404.12588v1}
\bibitem{ref249} MuDPT Multi-modal Deep-symphysis Prompt Tuning for Large Pre-trained Vision-Language Models. arXiv:2306.11400v1, 2023. \url{https://arxiv.org/abs/2306.11400v1}
\bibitem{ref250} \$$\pi$\$-Tuning Transferring Multimodal Foundation Models with Optimal Multi-task Interpolation. arXiv:2304.14381v3, 2023. \url{https://arxiv.org/abs/2304.14381v3}
\bibitem{ref251} MoPE Parameter-Efficient and Scalable Multimodal Fusion via Mixture of Prompt Experts. arXiv:2403.10568v1, 2024. \url{https://arxiv.org/abs/2403.10568v1}
\bibitem{ref252} Tuning LayerNorm in Attention Towards Efficient Multi-Modal LLM Finetuning. arXiv:2312.11420v1, 2023. \url{https://arxiv.org/abs/2312.11420v1}
\bibitem{ref253} MultiWay-Adapater Adapting large-scale multi-modal models for scalable image-text retrieval. arXiv:2309.01516v3, 2023. \url{https://arxiv.org/abs/2309.01516v3}
\bibitem{ref254} MmAP Multi-modal Alignment Prompt for Cross-domain Multi-task Learning. arXiv:2312.08636v1, 2023. \url{https://arxiv.org/abs/2312.08636v1}
\bibitem{ref255} 3D-VisTA Pre-trained Transformer for 3D Vision and Text Alignment. arXiv:2308.04352v1, 2023. \url{https://arxiv.org/abs/2308.04352v1}
\bibitem{ref256} Towards Unified Representation of Multi-Modal Pre-training for 3D Understanding via Differentiable Rendering. arXiv:2404.13619v1, 2024. \url{https://arxiv.org/abs/2404.13619v1}
\bibitem{ref257} JM3D \& JM3D-LLM Elevating 3D Understanding with Joint Multi-modal Cues. arXiv:2310.09503v3, 2023. \url{https://arxiv.org/abs/2310.09503v3}
\bibitem{ref258} ViT-Lens Initiating Omni-Modal Exploration through 3D Insights. arXiv:2308.10185v2, 2023. \url{https://arxiv.org/abs/2308.10185v2}
\bibitem{ref259} Autoencoders as Cross-Modal Teachers Can Pretrained 2D Image Transformers Help 3D Representation Learning. arXiv:2212.08320v2, 2022. \url{https://arxiv.org/abs/2212.08320v2}
\bibitem{ref260} 3DMIT 3D Multi-modal Instruction Tuning for Scene Understanding. arXiv:2401.03201v2, 2024. \url{https://arxiv.org/abs/2401.03201v2}
\bibitem{ref261} SUGAR Pre-training 3D Visual Representations for Robotics. arXiv:2404.01491v1, 2024. \url{https://arxiv.org/abs/2404.01491v1}
\bibitem{ref262} Learning Viewpoint-Agnostic Visual Representations by Recovering Tokens in 3D Space. arXiv:2206.11895v4, 2022. \url{https://arxiv.org/abs/2206.11895v4}
\bibitem{ref263} Grounding Symbols in Multi-Modal Instructions. arXiv:1706.00355v1, 2017. \url{https://arxiv.org/abs/1706.00355v1}
\bibitem{ref264} All in One Exploring Unified Video-Language Pre-training. arXiv:2203.07303v1, 2022. \url{https://arxiv.org/abs/2203.07303v1}
\bibitem{ref265} MERLOT Multimodal Neural Script Knowledge Models. arXiv:2106.02636v3, 2021. \url{https://arxiv.org/abs/2106.02636v3}
\bibitem{ref266} With a Little Help from my Temporal Context Multimodal Egocentric Action Recognition. arXiv:2111.01024v1, 2021. \url{https://arxiv.org/abs/2111.01024v1}
\bibitem{ref267} Summarize the Past to Predict the Future Natural Language Descriptions of Context Boost Multimodal Object Interaction Anticipation. arXiv:2301.09209v4, 2023. \url{https://arxiv.org/abs/2301.09209v4}
\bibitem{ref268} Language as the Medium Multimodal Video Classification through text only. arXiv:2309.10783v1, 2023. \url{https://arxiv.org/abs/2309.10783v1}
\bibitem{ref269} Adapting Short-Term Transformers for Action Detection in Untrimmed Videos. arXiv:2312.01897v2, 2023. \url{https://arxiv.org/abs/2312.01897v2}
\bibitem{ref270} Interaction-Aware Prompting for Zero-Shot Spatio-Temporal Action Detection. arXiv:2304.04688v4, 2023. \url{https://arxiv.org/abs/2304.04688v4}
\bibitem{ref271} Enhancing Video Transformers for Action Understanding with VLM-aided Training. arXiv:2403.16128v1, 2024. \url{https://arxiv.org/abs/2403.16128v1}
\bibitem{ref272} Language-based Action Concept Spaces Improve Video Self-Supervised Learning. arXiv:2307.10922v3, 2023. \url{https://arxiv.org/abs/2307.10922v3}
\bibitem{ref273} Long-Short Temporal Contrastive Learning of Video Transformers. arXiv:2106.09212v3, 2021. \url{https://arxiv.org/abs/2106.09212v3}
\bibitem{ref274} Can Temporal Information Help with Contrastive Self-Supervised Learning. arXiv:2011.13046v1, 2020. \url{https://arxiv.org/abs/2011.13046v1}
\bibitem{ref275} TCGL Temporal Contrastive Graph for Self-supervised Video Representation Learning. arXiv:2112.03587v3, 2021. \url{https://arxiv.org/abs/2112.03587v3}
\bibitem{ref276} Video Representation Learning with Visual Tempo Consistency. arXiv:2006.15489v2, 2020. \url{https://arxiv.org/abs/2006.15489v2}
\bibitem{ref277} End-to-end Multi-modal Video Temporal Grounding. arXiv:2107.05624v2, 2021. \url{https://arxiv.org/abs/2107.05624v2}
\bibitem{ref278} VLTinT Visual-Linguistic Transformer-in-Transformer for Coherent Video Paragraph Captioning. arXiv:2211.15103v2, 2022. \url{https://arxiv.org/abs/2211.15103v2}
\bibitem{ref279} Vision-Language Integration in Multimodal Video Transformers (Partially) Aligns with the Brain. arXiv:2311.07766v1, 2023. \url{https://arxiv.org/abs/2311.07766v1}
\bibitem{ref280} Divert More Attention to Vision-Language Object Tracking. arXiv:2307.10046v1, 2023. \url{https://arxiv.org/abs/2307.10046v1}
\bibitem{ref281} CLAP Learning Audio Concepts From Natural Language Supervision. arXiv:2206.04769v1, 2022. \url{https://arxiv.org/abs/2206.04769v1}
\bibitem{ref282} CTAL Pre-training Cross-modal Transformer for Audio-and-Language Representations. arXiv:2109.00181v1, 2021. \url{https://arxiv.org/abs/2109.00181v1}
\bibitem{ref283} i-Code An Integrative and Composable Multimodal Learning Framework. arXiv:2205.01818v2, 2022. \url{https://arxiv.org/abs/2205.01818v2}
\bibitem{ref284} VATLM Visual-Audio-Text Pre-Training with Unified Masked Prediction for Speech Representation Learning. arXiv:2211.11275v2, 2022. \url{https://arxiv.org/abs/2211.11275v2}
\bibitem{ref285} Towards More Unified In-context Visual Understanding. arXiv:2312.02520v2, 2023. \url{https://arxiv.org/abs/2312.02520v2}
\bibitem{ref286} Accommodating Audio Modality in CLIP for Multimodal Processing. arXiv:2303.06591v1, 2023. \url{https://arxiv.org/abs/2303.06591v1}
\bibitem{ref287} VALOR Vision-Audio-Language Omni-Perception Pretraining Model and Dataset. arXiv:2304.08345v1, 2023. \url{https://arxiv.org/abs/2304.08345v1}
\bibitem{ref288} CoAVT A Cognition-Inspired Unified Audio-Visual-Text Pre-Training Model for Multimodal Processing. arXiv:2401.12264v2, 2024. \url{https://arxiv.org/abs/2401.12264v2}
\bibitem{ref289} u-HuBERT Unified Mixed-Modal Speech Pretraining And Zero-Shot Transfer to Unlabeled Modality. arXiv:2207.07036v2, 2022. \url{https://arxiv.org/abs/2207.07036v2}
\bibitem{ref290} Fine-grained Audio-Visual Joint Representations for Multimodal Large Language Models. arXiv:2310.05863v2, 2023. \url{https://arxiv.org/abs/2310.05863v2}
\bibitem{ref291} Audio-Text Models Do Not Yet Leverage Natural Language. arXiv:2303.10667v1, 2023. \url{https://arxiv.org/abs/2303.10667v1}
\bibitem{ref292} ViT-Lens Towards Omni-modal Representations. arXiv:2311.16081v2, 2023. \url{https://arxiv.org/abs/2311.16081v2}
\bibitem{ref293} Connecting Multi-modal Contrastive Representations. arXiv:2305.14381v2, 2023. \url{https://arxiv.org/abs/2305.14381v2}
\bibitem{ref294} Extending Multi-modal Contrastive Representations. arXiv:2310.08884v1, 2023. \url{https://arxiv.org/abs/2310.08884v1}
\bibitem{ref295} Learning Visual Representation from Modality-Shared Contrastive Language-Image Pre-training. arXiv:2207.12661v1, 2022. \url{https://arxiv.org/abs/2207.12661v1}
\bibitem{ref296} Variational Mixture-of-Experts Autoencoders for Multi-Modal Deep Generative Models. arXiv:1911.03393v1, 2019. \url{https://arxiv.org/abs/1911.03393v1}
\bibitem{ref297} Omni-SMoLA Boosting Generalist Multimodal Models with Soft Mixture of Low-rank Experts. arXiv:2312.00968v2, 2023. \url{https://arxiv.org/abs/2312.00968v2}
\bibitem{ref298} MULTI Multimodal Understanding Leaderboard with Text and Images. arXiv:2402.03173v2, 2024. \url{https://arxiv.org/abs/2402.03173v2}
\bibitem{ref299} BLINK Multimodal Large Language Models Can See but Not Perceive. arXiv:2404.12390v2, 2024. \url{https://arxiv.org/abs/2404.12390v2}
\bibitem{ref300} Q-Bench A Benchmark for General-Purpose Foundation Models on Low-level Vision. arXiv:2309.14181v3, 2023. \url{https://arxiv.org/abs/2309.14181v3}
\bibitem{ref301} A Benchmark for Multi-modal Foundation Models on Low-level Vision from Single Images to Pairs. arXiv:2402.07116v1, 2024. \url{https://arxiv.org/abs/2402.07116v1}
\bibitem{ref302} Muffin or Chihuahua Challenging Large Vision-Language Models with Multipanel VQA. arXiv:2401.15847v2, 2024. \url{https://arxiv.org/abs/2401.15847v2}
\bibitem{ref303} IllusionVQA A Challenging Optical Illusion Dataset for Vision Language Models. arXiv:2403.15952v2, 2024. \url{https://arxiv.org/abs/2403.15952v2}
\bibitem{ref304} ReForm-Eval Evaluating Large Vision Language Models via Unified Re-Formulation of Task-Oriented Benchmarks. arXiv:2310.02569v2, 2023. \url{https://arxiv.org/abs/2310.02569v2}
\bibitem{ref305} Uncertainty-Aware Evaluation for Vision-Language Models. arXiv:2402.14418v2, 2024. \url{https://arxiv.org/abs/2402.14418v2}
\bibitem{ref306} Reassessing Evaluation Practices in Visual Question Answering A Case Study on Out-of-Distribution Generalization. arXiv:2205.12191v2, 2022. \url{https://arxiv.org/abs/2205.12191v2}
\bibitem{ref307} Large Language Models can Implement Policy Iteration. arXiv:2210.03821v2, 2022. \url{https://arxiv.org/abs/2210.03821v2}
\bibitem{ref308} Human-Timescale Adaptation in an Open-Ended Task Space. arXiv:2301.07608v1, 2023. \url{https://arxiv.org/abs/2301.07608v1}
\bibitem{ref309} RL\$\textasciicircum{}3\$ Boosting Meta Reinforcement Learning via RL inside RL\$\textasciicircum{}2\$. arXiv:2306.15909v4, 2023. \url{https://arxiv.org/abs/2306.15909v4}
\bibitem{ref310} Towards Effective Context for Meta-Reinforcement Learning an Approach based on Contrastive Learning. arXiv:2009.13891v3, 2020. \url{https://arxiv.org/abs/2009.13891v3}
\bibitem{ref311} Meta-Reinforcement Learning Based on Self-Supervised Task Representation Learning. arXiv:2305.00286v1, 2023. \url{https://arxiv.org/abs/2305.00286v1}
\bibitem{ref312} Decomposed Mutual Information Optimization for Generalized Context in Meta-Reinforcement Learning. arXiv:2210.04209v1, 2022. \url{https://arxiv.org/abs/2210.04209v1}
\bibitem{ref313} Meta-Learning for Contextual Bandit Exploration. arXiv:1901.08159v1, 2019. \url{https://arxiv.org/abs/1901.08159v1}
\bibitem{ref314} Meta-Reinforcement Learning in Broad and Non-Parametric Environments. arXiv:2108.03718v2, 2021. \url{https://arxiv.org/abs/2108.03718v2}
\bibitem{ref315} Guided Meta-Policy Search. arXiv:1904.00956v2, 2019. \url{https://arxiv.org/abs/1904.00956v2}
\bibitem{ref316} Concurrent Meta Reinforcement Learning. arXiv:1903.02710v1, 2019. \url{https://arxiv.org/abs/1903.02710v1}
\bibitem{ref317} Enhanced Meta Reinforcement Learning using Demonstrations in Sparse Reward Environments. arXiv:2209.13048v1, 2022. \url{https://arxiv.org/abs/2209.13048v1}
\bibitem{ref318} Improving Context-Based Meta-Reinforcement Learning with Self-Supervised Trajectory Contrastive Learning. arXiv:2103.06386v1, 2021. \url{https://arxiv.org/abs/2103.06386v1}
\bibitem{ref319} Unsupervised Curricula for Visual Meta-Reinforcement Learning. arXiv:1912.04226v1, 2019. \url{https://arxiv.org/abs/1912.04226v1}
\bibitem{ref320} Self-Paced Contextual Reinforcement Learning. arXiv:1910.02826v1, 2019. \url{https://arxiv.org/abs/1910.02826v1}
\bibitem{ref321} Decoupling Meta-Reinforcement Learning with Gaussian Task Contexts and Skills. arXiv:2312.06518v1, 2023. \url{https://arxiv.org/abs/2312.06518v1}
\bibitem{ref322} Adaptive Asynchronous Control Using Meta-learned Neural Ordinary Differential Equations. arXiv:2207.12062v5, 2022. \url{https://arxiv.org/abs/2207.12062v5}
\bibitem{ref323} Meta-Gradients in Non-Stationary Environments. arXiv:2209.06159v1, 2022. \url{https://arxiv.org/abs/2209.06159v1}
\bibitem{ref324} Continuous Adaptation via Meta-Learning in Nonstationary and Competitive Environments. arXiv:1710.03641v2, 2017. \url{https://arxiv.org/abs/1710.03641v2}
\bibitem{ref325} Meta-Thompson Sampling. arXiv:2102.06129v2, 2021. \url{https://arxiv.org/abs/2102.06129v2}
\bibitem{ref326} Meta-Learning Bandit Policies by Gradient Ascent. arXiv:2006.05094v2, 2020. \url{https://arxiv.org/abs/2006.05094v2}
\bibitem{ref327} AutoML for Contextual Bandits. arXiv:1909.03212v2, 2019. \url{https://arxiv.org/abs/1909.03212v2}
\bibitem{ref328} Meta-learners' learning dynamics are unlike learners'. arXiv:1905.01320v1, 2019. \url{https://arxiv.org/abs/1905.01320v1}
\bibitem{ref329} Meta-in-context learning in large language models. arXiv:2305.12907v1, 2023. \url{https://arxiv.org/abs/2305.12907v1}
\bibitem{ref330} Hierarchical Adaptive Contextual Bandits for Resource Constraint based Recommendation. arXiv:2004.01136v2, 2020. \url{https://arxiv.org/abs/2004.01136v2}
\bibitem{ref331} Off-policy Bandits with Deficient Support. arXiv:2006.09438v1, 2020. \url{https://arxiv.org/abs/2006.09438v1}
\bibitem{ref332} Estimation Considerations in Contextual Bandits. arXiv:1711.07077v4, 2017. \url{https://arxiv.org/abs/1711.07077v4}
\bibitem{ref333} Efficient Off-Policy Meta-Reinforcement Learning via Probabilistic Context Variables. arXiv:1903.08254v1, 2019. \url{https://arxiv.org/abs/1903.08254v1}
\bibitem{ref334} On Context Distribution Shift in Task Representation Learning for Offline Meta RL. arXiv:2304.00354v2, 2023. \url{https://arxiv.org/abs/2304.00354v2}
\bibitem{ref335} Context Shift Reduction for Offline Meta-Reinforcement Learning. arXiv:2311.03695v1, 2023. \url{https://arxiv.org/abs/2311.03695v1}
\bibitem{ref336} Robust Task Representations for Offline Meta-Reinforcement Learning via Contrastive Learning. arXiv:2206.10442v1, 2022. \url{https://arxiv.org/abs/2206.10442v1}
\bibitem{ref337} Disentangling Policy from Offline Task Representation Learning via Adversarial Data Augmentation. arXiv:2403.07261v1, 2024. \url{https://arxiv.org/abs/2403.07261v1}
\bibitem{ref338} FOCAL Efficient Fully-Offline Meta-Reinforcement Learning via Distance Metric Learning and Behavior Regularization. arXiv:2010.01112v4, 2020. \url{https://arxiv.org/abs/2010.01112v4}
\bibitem{ref339} Provably Improved Context-Based Offline Meta-RL with Attention and Contrastive Learning. arXiv:2102.10774v2, 2021. \url{https://arxiv.org/abs/2102.10774v2}
\bibitem{ref340} Generalizable Task Representation Learning for Offline Meta-Reinforcement Learning with Data Limitations. arXiv:2312.15909v1, 2023. \url{https://arxiv.org/abs/2312.15909v1}
\bibitem{ref341} Offline Meta-Reinforcement Learning with Online Self-Supervision. arXiv:2107.03974v4, 2021. \url{https://arxiv.org/abs/2107.03974v4}
\bibitem{ref342} Model-Based Offline Meta-Reinforcement Learning with Regularization. arXiv:2202.02929v2, 2022. \url{https://arxiv.org/abs/2202.02929v2}
\bibitem{ref343} Bridging Distributionally Robust Learning and Offline RL An Approach to Mitigate Distribution Shift and Partial Data Coverage. arXiv:2310.18434v1, 2023. \url{https://arxiv.org/abs/2310.18434v1}
\bibitem{ref344} Towards an Information Theoretic Framework of Context-Based Offline Meta-Reinforcement Learning. arXiv:2402.02429v1, 2024. \url{https://arxiv.org/abs/2402.02429v1}
\bibitem{ref345} Variational Empowerment as Representation Learning for Goal-Based Reinforcement Learning. arXiv:2106.01404v1, 2021. \url{https://arxiv.org/abs/2106.01404v1}
\bibitem{ref346} Unleash Model Potential Bootstrapped Meta Self-supervised Learning. arXiv:2308.14267v1, 2023. \url{https://arxiv.org/abs/2308.14267v1}
\bibitem{ref347} Robust Imitation of a Few Demonstrations with a Backwards Model. arXiv:2210.09337v1, 2022. \url{https://arxiv.org/abs/2210.09337v1}
\bibitem{ref348} Discriminator-Weighted Offline Imitation Learning from Suboptimal Demonstrations. arXiv:2207.10050v1, 2022. \url{https://arxiv.org/abs/2207.10050v1}
\bibitem{ref349} TRAIL Near-Optimal Imitation Learning with Suboptimal Data. arXiv:2110.14770v1, 2021. \url{https://arxiv.org/abs/2110.14770v1}
\bibitem{ref350} SubIQ Inverse Soft-Q Learning for Offline Imitation with Suboptimal Demonstrations. arXiv:2402.13147v1, 2024. \url{https://arxiv.org/abs/2402.13147v1}
\bibitem{ref351} Learning to Generalize Across Long-Horizon Tasks from Human Demonstrations. arXiv:2003.06085v2, 2020. \url{https://arxiv.org/abs/2003.06085v2}
\bibitem{ref352} GO-DICE Goal-Conditioned Option-Aware Offline Imitation Learning via Stationary Distribution Correction Estimation. arXiv:2312.10802v1, 2023. \url{https://arxiv.org/abs/2312.10802v1}
\bibitem{ref353} Model-based trajectory stitching for improved behavioural cloning and its applications. arXiv:2212.04280v1, 2022. \url{https://arxiv.org/abs/2212.04280v1}
\bibitem{ref354} Model-based Trajectory Stitching for Improved Offline Reinforcement Learning. arXiv:2211.11603v1, 2022. \url{https://arxiv.org/abs/2211.11603v1}
\bibitem{ref355} Context-Former Stitching via Latent Conditioned Sequence Modeling. arXiv:2401.16452v2, 2024. \url{https://arxiv.org/abs/2401.16452v2}
\bibitem{ref356} Diffusion Model-Augmented Behavioral Cloning. arXiv:2302.13335v3, 2023. \url{https://arxiv.org/abs/2302.13335v3}
\bibitem{ref357} Cold Diffusion on the Replay Buffer Learning to Plan from Known Good States. arXiv:2310.13914v1, 2023. \url{https://arxiv.org/abs/2310.13914v1}
\bibitem{ref358} Learning to Discern Imitating Heterogeneous Human Demonstrations with Preference and Representation Learning. arXiv:2310.14196v1, 2023. \url{https://arxiv.org/abs/2310.14196v1}
\bibitem{ref359} Offline Learning from Demonstrations and Unlabeled Experience. arXiv:2011.13885v1, 2020. \url{https://arxiv.org/abs/2011.13885v1}
\bibitem{ref360} Get Back Here Robust Imitation by Return-to-Distribution Planning. arXiv:2305.01400v1, 2023. \url{https://arxiv.org/abs/2305.01400v1}
\bibitem{ref361} CCIL Context-conditioned imitation learning for urban driving. arXiv:2305.02649v1, 2023. \url{https://arxiv.org/abs/2305.02649v1}
\bibitem{ref362} SWBT Similarity Weighted Behavior Transformer with the Imperfect Demonstration for Robotic Manipulation. arXiv:2401.08957v1, 2024. \url{https://arxiv.org/abs/2401.08957v1}
\bibitem{ref363} Transformers are Meta-Reinforcement Learners. arXiv:2206.06614v1, 2022. \url{https://arxiv.org/abs/2206.06614v1}
\bibitem{ref364} Meta-learning of Sequential Strategies. arXiv:1905.03030v2, 2019. \url{https://arxiv.org/abs/1905.03030v2}
\bibitem{ref365} Memory-Based Meta-Learning on Non-Stationary Distributions. arXiv:2302.03067v2, 2023. \url{https://arxiv.org/abs/2302.03067v2}
\bibitem{ref366} Transformers as Decision Makers Provable In-Context Reinforcement Learning via Supervised Pretraining. arXiv:2310.08566v1, 2023. \url{https://arxiv.org/abs/2310.08566v1}
\bibitem{ref367} Emergent Agentic Transformer from Chain of Hindsight Experience. arXiv:2305.16554v1, 2023. \url{https://arxiv.org/abs/2305.16554v1}
\bibitem{ref368} Trainable Transformer in Transformer. arXiv:2307.01189v2, 2023. \url{https://arxiv.org/abs/2307.01189v2}
\bibitem{ref369} TransDreamer Reinforcement Learning with Transformer World Models. arXiv:2202.09481v1, 2022. \url{https://arxiv.org/abs/2202.09481v1}
\bibitem{ref370} Self-Paced Multi-Task Learning. arXiv:1604.01474v2, 2016. \url{https://arxiv.org/abs/1604.01474v2}
\bibitem{ref371} Self-Paced Deep Reinforcement Learning. arXiv:2004.11812v5, 2020. \url{https://arxiv.org/abs/2004.11812v5}
\bibitem{ref372} Self-Paced Context Evaluation for Contextual Reinforcement Learning. arXiv:2106.05110v1, 2021. \url{https://arxiv.org/abs/2106.05110v1}
\bibitem{ref373} Curriculum Reinforcement Learning using Optimal Transport via Gradual Domain Adaptation. arXiv:2210.10195v1, 2022. \url{https://arxiv.org/abs/2210.10195v1}
\bibitem{ref374} On the Benefit of Optimal Transport for Curriculum Reinforcement Learning. arXiv:2309.14091v1, 2023. \url{https://arxiv.org/abs/2309.14091v1}
\bibitem{ref375} Meta Automatic Curriculum Learning. arXiv:2011.08463v3, 2020. \url{https://arxiv.org/abs/2011.08463v3}
\bibitem{ref376} Trying AGAIN instead of Trying Longer Prior Learning for Automatic Curriculum Learning. arXiv:2004.03168v1, 2020. \url{https://arxiv.org/abs/2004.03168v1}
\bibitem{ref377} Unsupervised Meta-Learning for Reinforcement Learning. arXiv:1806.04640v3, 2018. \url{https://arxiv.org/abs/1806.04640v3}
\bibitem{ref378} Learning Curricula in Open-Ended Worlds. arXiv:2312.03126v2, 2023. \url{https://arxiv.org/abs/2312.03126v2}
\bibitem{ref379} CLUTR Curriculum Learning via Unsupervised Task Representation Learning. arXiv:2210.10243v2, 2022. \url{https://arxiv.org/abs/2210.10243v2}
\bibitem{ref380} Reward-Machine-Guided, Self-Paced Reinforcement Learning. arXiv:2305.16505v1, 2023. \url{https://arxiv.org/abs/2305.16505v1}
\bibitem{ref381} Hacking Task Confounder in Meta-Learning. arXiv:2312.05771v3, 2023. \url{https://arxiv.org/abs/2312.05771v3}
\bibitem{ref382} Multi-skill Mobile Manipulation for Object Rearrangement. arXiv:2209.02778v1, 2022. \url{https://arxiv.org/abs/2209.02778v1}
\bibitem{ref383} Bottom-Up Skill Discovery from Unsegmented Demonstrations for Long-Horizon Robot Manipulation. arXiv:2109.13841v2, 2021. \url{https://arxiv.org/abs/2109.13841v2}
\bibitem{ref384} Language-guided Skill Learning with Temporal Variational Inference. arXiv:2402.16354v1, 2024. \url{https://arxiv.org/abs/2402.16354v1}
\bibitem{ref385} Skill Preferences Learning to Extract and Execute Robotic Skills from Human Feedback. arXiv:2108.05382v1, 2021. \url{https://arxiv.org/abs/2108.05382v1}
\bibitem{ref386} Meta-Learning Parameterized Skills. arXiv:2206.03597v4, 2022. \url{https://arxiv.org/abs/2206.03597v4}
\bibitem{ref387} Generalizing to New Tasks via One-Shot Compositional Subgoals. arXiv:2205.07716v2, 2022. \url{https://arxiv.org/abs/2205.07716v2}
\bibitem{ref388} OCEAN Online Task Inference for Compositional Tasks with Context Adaptation. arXiv:2008.07087v1, 2020. \url{https://arxiv.org/abs/2008.07087v1}
\bibitem{ref389} LLM Augmented Hierarchical Agents. arXiv:2311.05596v1, 2023. \url{https://arxiv.org/abs/2311.05596v1}
\bibitem{ref390} Acquiring Diverse Skills using Curriculum Reinforcement Learning with Mixture of Experts. arXiv:2403.06966v1, 2024. \url{https://arxiv.org/abs/2403.06966v1}
\bibitem{ref391} SkillS Adaptive Skill Sequencing for Efficient Temporally-Extended Exploration. arXiv:2211.13743v3, 2022. \url{https://arxiv.org/abs/2211.13743v3}
\bibitem{ref392} Robust Policy Learning via Offline Skill Diffusion. arXiv:2403.00225v2, 2024. \url{https://arxiv.org/abs/2403.00225v2}
\bibitem{ref393} SemTra A Semantic Skill Translator for Cross-Domain Zero-Shot Policy Adaptation. arXiv:2402.07418v1, 2024. \url{https://arxiv.org/abs/2402.07418v1}
\bibitem{ref394} Model Primitive Hierarchical Lifelong Reinforcement Learning. arXiv:1903.01567v1, 2019. \url{https://arxiv.org/abs/1903.01567v1}
\bibitem{ref395} Language as an Abstraction for Hierarchical Deep Reinforcement Learning. arXiv:1906.07343v2, 2019. \url{https://arxiv.org/abs/1906.07343v2}
\bibitem{ref396} LEAGUE Guided Skill Learning and Abstraction for Long-Horizon Manipulation. arXiv:2210.12631v2, 2022. \url{https://arxiv.org/abs/2210.12631v2}
\bibitem{ref397} Sim-to-Real Transfer in Deep Reinforcement Learning for Robotics a Survey. arXiv:2009.13303v2, 2020. \url{https://arxiv.org/abs/2009.13303v2}
\bibitem{ref398} Robot Learning from Randomized Simulations A Review. arXiv:2111.00956v2, 2021. \url{https://arxiv.org/abs/2111.00956v2}
\bibitem{ref399} Solving Rubik's Cube with a Robot Hand. arXiv:1910.07113v1, 2019. \url{https://arxiv.org/abs/1910.07113v1}
\bibitem{ref400} Data-efficient Domain Randomization with Bayesian Optimization. arXiv:2003.02471v4, 2020. \url{https://arxiv.org/abs/2003.02471v4}
\bibitem{ref401} Continual Domain Randomization. arXiv:2403.12193v1, 2024. \url{https://arxiv.org/abs/2403.12193v1}
\bibitem{ref402} Domain Randomization via Entropy Maximization. arXiv:2311.01885v2, 2023. \url{https://arxiv.org/abs/2311.01885v2}
\bibitem{ref403} MELD Meta-Reinforcement Learning from Images via Latent State Models. arXiv:2010.13957v2, 2020. \url{https://arxiv.org/abs/2010.13957v2}
\bibitem{ref404} Prototypical context-aware dynamics generalization for high-dimensional model-based reinforcement learning. arXiv:2211.12774v1, 2022. \url{https://arxiv.org/abs/2211.12774v1}
\bibitem{ref405} TWIST Teacher-Student World Model Distillation for Efficient Sim-to-Real Transfer. arXiv:2311.03622v1, 2023. \url{https://arxiv.org/abs/2311.03622v1}
\bibitem{ref406} Privileged Knowledge Distillation for Sim-to-Real Policy Generalization. arXiv:2305.18464v1, 2023. \url{https://arxiv.org/abs/2305.18464v1}
\bibitem{ref407} Bi-directional Domain Adaptation for Sim2Real Transfer of Embodied Navigation Agents. arXiv:2011.12421v2, 2020. \url{https://arxiv.org/abs/2011.12421v2}
\bibitem{ref408} Pre-training Contextualized World Models with In-the-wild Videos for Reinforcement Learning. arXiv:2305.18499v2, 2023. \url{https://arxiv.org/abs/2305.18499v2}
\bibitem{ref409} Dreaming of Many Worlds Learning Contextual World Models Aids Zero-Shot Generalization. arXiv:2403.10967v1, 2024. \url{https://arxiv.org/abs/2403.10967v1}
\bibitem{ref410} Task Aware Dreamer for Task Generalization in Reinforcement Learning. arXiv:2303.05092v3, 2023. \url{https://arxiv.org/abs/2303.05092v3}
\bibitem{ref411} Out-of-Variable Generalization for Discriminative Models. arXiv:2304.07896v3, 2023. \url{https://arxiv.org/abs/2304.07896v3}
\bibitem{ref412} Adapting to Unseen Environments through Explicit Representation of Context. arXiv:2002.05640v2, 2020. \url{https://arxiv.org/abs/2002.05640v2}
\bibitem{ref413} Goal Exploration via Adaptive Skill Distribution for Goal-Conditioned Reinforcement Learning. arXiv:2404.12999v1, 2024. \url{https://arxiv.org/abs/2404.12999v1}
\bibitem{ref414} Open-Ended Learning Leads to Generally Capable Agents. arXiv:2107.12808v2, 2021. \url{https://arxiv.org/abs/2107.12808v2}
\bibitem{ref415} One-Shot Learning from a Demonstration with Hierarchical Latent Language. arXiv:2203.04806v1, 2022. \url{https://arxiv.org/abs/2203.04806v1}
\bibitem{ref416} Transferable Reinforcement Learning via Generalized Occupancy Models. arXiv:2403.06328v1, 2024. \url{https://arxiv.org/abs/2403.06328v1}
\bibitem{ref417} Unveiling the Tapestry the Interplay of Generalization and Forgetting in Continual Learning. arXiv:2211.11174v5, 2022. \url{https://arxiv.org/abs/2211.11174v5}
\bibitem{ref418} Calibration in Deep Learning A Survey of the State-of-the-Art. arXiv:2308.01222v2, 2023. \url{https://arxiv.org/abs/2308.01222v2}
\bibitem{ref419} Overthinking the Truth Understanding how Language Models Process False Demonstrations. arXiv:2307.09476v3, 2023. \url{https://arxiv.org/abs/2307.09476v3}
\bibitem{ref420} Improving Calibration through the Relationship with Adversarial Robustness. arXiv:2006.16375v2, 2020. \url{https://arxiv.org/abs/2006.16375v2}
\bibitem{ref421} Calibration Attack A Framework For Adversarial Attacks Targeting Calibration. arXiv:2401.02718v1, 2024. \url{https://arxiv.org/abs/2401.02718v1}
\bibitem{ref422} Calibration Meets Explanation A Simple and Effective Approach for Model Confidence Estimates. arXiv:2211.03041v1, 2022. \url{https://arxiv.org/abs/2211.03041v1}
\bibitem{ref423} How are Prompts Different in Terms of Sensitivity. arXiv:2311.07230v1, 2023. \url{https://arxiv.org/abs/2311.07230v1}
\bibitem{ref424} Ground-Truth Labels Matter A Deeper Look into Input-Label Demonstrations. arXiv:2205.12685v2, 2022. \url{https://arxiv.org/abs/2205.12685v2}
\bibitem{ref425} Measuring Inductive Biases of In-Context Learning with Underspecified Demonstrations. arXiv:2305.13299v1, 2023. \url{https://arxiv.org/abs/2305.13299v1}
\bibitem{ref426} Adversarial Examples in Modern Machine Learning A Review. arXiv:1911.05268v2, 2019. \url{https://arxiv.org/abs/1911.05268v2}
\bibitem{ref427} Adversarial Learning in Statistical Classification A Comprehensive Review of Defenses Against Attacks. arXiv:1904.06292v3, 2019. \url{https://arxiv.org/abs/1904.06292v3}
\bibitem{ref428} Adversarial Demonstration Attacks on Large Language Models. arXiv:2305.14950v2, 2023. \url{https://arxiv.org/abs/2305.14950v2}
\bibitem{ref429} Hijacking Large Language Models via Adversarial In-Context Learning. arXiv:2311.09948v1, 2023. \url{https://arxiv.org/abs/2311.09948v1}
\bibitem{ref430} Prompt Backdoors in Visual Prompt Learning. arXiv:2310.07632v1, 2023. \url{https://arxiv.org/abs/2310.07632v1}
\bibitem{ref431} Extreme Miscalibration and the Illusion of Adversarial Robustness. arXiv:2402.17509v1, 2024. \url{https://arxiv.org/abs/2402.17509v1}
\bibitem{ref432} An LLM can Fool Itself A Prompt-Based Adversarial Attack. arXiv:2310.13345v1, 2023. \url{https://arxiv.org/abs/2310.13345v1}
\bibitem{ref433} RobustBench a standardized adversarial robustness benchmark. arXiv:2010.09670v3, 2020. \url{https://arxiv.org/abs/2010.09670v3}
\bibitem{ref434} Reliable evaluation of adversarial robustness with an ensemble of diverse parameter-free attacks. arXiv:2003.01690v2, 2020. \url{https://arxiv.org/abs/2003.01690v2}
\bibitem{ref435} BadCLIP Dual-Embedding Guided Backdoor Attack on Multimodal Contrastive Learning. arXiv:2311.12075v3, 2023. \url{https://arxiv.org/abs/2311.12075v3}
\bibitem{ref436} Confidence-Calibrated Adversarial Training Generalizing to Unseen Attacks. arXiv:1910.06259v4, 2019. \url{https://arxiv.org/abs/1910.06259v4}
\bibitem{ref437} Membership Inference Attacks against Machine Learning Models. arXiv:1610.05820v2, 2016. \url{https://arxiv.org/abs/1610.05820v2}
\bibitem{ref438} Towards Demystifying Membership Inference Attacks. arXiv:1807.09173v2, 2018. \url{https://arxiv.org/abs/1807.09173v2}
\bibitem{ref439} Quantifying Privacy Risks of Prompts in Visual Prompt Learning. arXiv:2310.11970v1, 2023. \url{https://arxiv.org/abs/2310.11970v1}
\bibitem{ref440} Does Prompt-Tuning Language Model Ensure Privacy. arXiv:2304.03472v2, 2023. \url{https://arxiv.org/abs/2304.03472v2}
\bibitem{ref441} Last One Standing A Comparative Analysis of Security and Privacy of Soft Prompt Tuning, LoRA, and In-Context Learning. arXiv:2310.11397v1, 2023. \url{https://arxiv.org/abs/2310.11397v1}
\bibitem{ref442} Analyzing Leakage of Personally Identifiable Information in Language Models. arXiv:2302.00539v4, 2023. \url{https://arxiv.org/abs/2302.00539v4}
\bibitem{ref443} Evaluations of Machine Learning Privacy Defenses are Misleading. arXiv:2404.17399v1, 2024. \url{https://arxiv.org/abs/2404.17399v1}
\bibitem{ref444} Systematic Evaluation of Privacy Risks of Machine Learning Models. arXiv:2003.10595v2, 2020. \url{https://arxiv.org/abs/2003.10595v2}
\bibitem{ref445} Label-Only Membership Inference Attacks. arXiv:2007.14321v3, 2020. \url{https://arxiv.org/abs/2007.14321v3}
\bibitem{ref446} User Inference Attacks on Large Language Models. arXiv:2310.09266v2, 2023. \url{https://arxiv.org/abs/2310.09266v2}
\bibitem{ref447} Privacy Backdoors Enhancing Membership Inference through Poisoning Pre-trained Models. arXiv:2404.01231v1, 2024. \url{https://arxiv.org/abs/2404.01231v1}
\bibitem{ref448} PreCurious How Innocent Pre-Trained Language Models Turn into Privacy Traps. arXiv:2403.09562v1, 2024. \url{https://arxiv.org/abs/2403.09562v1}
\bibitem{ref449} Information Leakage Detection through Approximate Bayes-optimal Prediction. arXiv:2401.14283v1, 2024. \url{https://arxiv.org/abs/2401.14283v1}
\bibitem{ref450} A New Approach to Adaptive Data Analysis and Learning via Maximal Leakage. arXiv:1903.01777v1, 2019. \url{https://arxiv.org/abs/1903.01777v1}
\bibitem{ref451} Can LLMs Keep a Secret Testing Privacy Implications of Language Models via Contextual Integrity Theory. arXiv:2310.17884v1, 2023. \url{https://arxiv.org/abs/2310.17884v1}
\bibitem{ref452} Understanding Practical Membership Privacy of Deep Learning. arXiv:2402.06674v1, 2024. \url{https://arxiv.org/abs/2402.06674v1}
\bibitem{ref453} On the Privacy Risks of Model Explanations. arXiv:1907.00164v6, 2019. \url{https://arxiv.org/abs/1907.00164v6}
\bibitem{ref454} Privacy-Preserving In-Context Learning for Large Language Models. arXiv:2305.01639v2, 2023. \url{https://arxiv.org/abs/2305.01639v2}
\bibitem{ref455} DP-TabICL In-Context Learning with Differentially Private Tabular Data. arXiv:2403.05681v1, 2024. \url{https://arxiv.org/abs/2403.05681v1}
\bibitem{ref456} DP-UTIL Comprehensive Utility Analysis of Differential Privacy in Machine Learning. arXiv:2112.12998v1, 2021. \url{https://arxiv.org/abs/2112.12998v1}
\bibitem{ref457} Research Challenges in Designing Differentially Private Text Generation Mechanisms. arXiv:2012.05403v1, 2020. \url{https://arxiv.org/abs/2012.05403v1}
\bibitem{ref458} Synthetic Text Generation with Differential Privacy A Simple and Practical Recipe. arXiv:2210.14348v3, 2022. \url{https://arxiv.org/abs/2210.14348v3}
\bibitem{ref459} How reparametrization trick broke differentially-private text representation learning. arXiv:2202.12138v2, 2022. \url{https://arxiv.org/abs/2202.12138v2}
\bibitem{ref460} DP-Rewrite Towards Reproducibility and Transparency in Differentially Private Text Rewriting. arXiv:2208.10400v1, 2022. \url{https://arxiv.org/abs/2208.10400v1}
\bibitem{ref461} Budget Recycling Differential Privacy. arXiv:2403.11445v3, 2024. \url{https://arxiv.org/abs/2403.11445v3}
\bibitem{ref462} Differentially Empirical Risk Minimization under the Fairness Lens. arXiv:2106.02674v2, 2021. \url{https://arxiv.org/abs/2106.02674v2}
\bibitem{ref463} Spending Privacy Budget Fairly and Wisely. arXiv:2204.12903v1, 2022. \url{https://arxiv.org/abs/2204.12903v1}
\bibitem{ref464} Towards Fair Classifiers Without Sensitive Attributes Exploring Biases in Related Features. arXiv:2104.14537v4, 2021. \url{https://arxiv.org/abs/2104.14537v4}
\bibitem{ref465} When Fairness Meets Privacy Fair Classification with Semi-Private Sensitive Attributes. arXiv:2207.08336v2, 2022. \url{https://arxiv.org/abs/2207.08336v2}
\bibitem{ref466} Fairness without the sensitive attribute via Causal Variational Autoencoder. arXiv:2109.04999v1, 2021. \url{https://arxiv.org/abs/2109.04999v1}
\bibitem{ref467} Fair Classification via Domain Adaptation A Dual Adversarial Learning Approach. arXiv:2206.03656v2, 2022. \url{https://arxiv.org/abs/2206.03656v2}
\bibitem{ref468} On the Impact of Multi-dimensional Local Differential Privacy on Fairness. arXiv:2312.04404v3, 2023. \url{https://arxiv.org/abs/2312.04404v3}
\bibitem{ref469} (Local) Differential Privacy has NO Disparate Impact on Fairness. arXiv:2304.12845v2, 2023. \url{https://arxiv.org/abs/2304.12845v2}
\bibitem{ref470} De-amplifying Bias from Differential Privacy in Language Model Fine-tuning. arXiv:2402.04489v1, 2024. \url{https://arxiv.org/abs/2402.04489v1}
\bibitem{ref471} Fair NLP Models with Differentially Private Text Encoders. arXiv:2205.06135v1, 2022. \url{https://arxiv.org/abs/2205.06135v1}
\bibitem{ref472} Privacy for Fairness Information Obfuscation for Fair Representation Learning with Local Differential Privacy. arXiv:2402.10473v1, 2024. \url{https://arxiv.org/abs/2402.10473v1}
\bibitem{ref473} Differentially Private and Fair Classification via Calibrated Functional Mechanism. arXiv:2001.04958v2, 2020. \url{https://arxiv.org/abs/2001.04958v2}
\bibitem{ref474} Stochastic Differentially Private and Fair Learning. arXiv:2210.08781v2, 2022. \url{https://arxiv.org/abs/2210.08781v2}
\bibitem{ref475} Learning with Impartiality to Walk on the Pareto Frontier of Fairness, Privacy, and Utility. arXiv:2302.09183v1, 2023. \url{https://arxiv.org/abs/2302.09183v1}
\bibitem{ref476} Privacy-Preserving In-Context Learning with Differentially Private Few-Shot Generation. arXiv:2309.11765v2, 2023. \url{https://arxiv.org/abs/2309.11765v2}
\bibitem{ref477} On the Privacy Risks of Algorithmic Fairness. arXiv:2011.03731v4, 2020. \url{https://arxiv.org/abs/2011.03731v4}
\bibitem{ref478} Neither Private Nor Fair Impact of Data Imbalance on Utility and Fairness in Differential Privacy. arXiv:2009.06389v3, 2020. \url{https://arxiv.org/abs/2009.06389v3}
\bibitem{ref479} An Empirical Analysis of Fairness Notions under Differential Privacy. arXiv:2302.02910v1, 2023. \url{https://arxiv.org/abs/2302.02910v1}
\bibitem{ref480} Unlocking Accuracy and Fairness in Differentially Private Image Classification. arXiv:2308.10888v1, 2023. \url{https://arxiv.org/abs/2308.10888v1}
\bibitem{ref481} Privacy-preserving Fine-tuning of Large Language Models through Flatness. arXiv:2403.04124v1, 2024. \url{https://arxiv.org/abs/2403.04124v1}
\bibitem{ref482} Beyond Incompatibility Trade-offs between Mutually Exclusive Fairness Criteria in Machine Learning and Law. arXiv:2212.00469v3, 2022. \url{https://arxiv.org/abs/2212.00469v3}
\bibitem{ref483} HappyMap A Generalized Multi-calibration Method. arXiv:2303.04379v1, 2023. \url{https://arxiv.org/abs/2303.04379v1}
\bibitem{ref484} Function Composition in Trustworthy Machine Learning Implementation Choices, Insights, and Questions. arXiv:2302.09190v1, 2023. \url{https://arxiv.org/abs/2302.09190v1}
\bibitem{ref485} Fault Tolerance of Neural Networks in Adversarial Settings. arXiv:1910.13875v2, 2019. \url{https://arxiv.org/abs/1910.13875v2}
\bibitem{ref486} ComplAI Theory of A Unified Framework for Multi-factor Assessment of Black-Box Supervised Machine Learning Models. arXiv:2212.14599v1, 2022. \url{https://arxiv.org/abs/2212.14599v1}
\bibitem{ref487} Auditing and Generating Synthetic Data with Controllable Trust Trade-offs. arXiv:2304.10819v3, 2023. \url{https://arxiv.org/abs/2304.10819v3}
\bibitem{ref488} Towards a Responsible AI Development Lifecycle Lessons From Information Security. arXiv:2203.02958v1, 2022. \url{https://arxiv.org/abs/2203.02958v1}
\bibitem{ref489} Private Prediction Sets. arXiv:2102.06202v3, 2021. \url{https://arxiv.org/abs/2102.06202v3}
\bibitem{ref490} Federated Conformal Predictors for Distributed Uncertainty Quantification. arXiv:2305.17564v2, 2023. \url{https://arxiv.org/abs/2305.17564v2}
\bibitem{ref491} Efficient Conformal Prediction under Data Heterogeneity. arXiv:2312.15799v1, 2023. \url{https://arxiv.org/abs/2312.15799v1}
\bibitem{ref492} Trust, but Verify Using Self-Supervised Probing to Improve Trustworthiness. arXiv:2302.02628v1, 2023. \url{https://arxiv.org/abs/2302.02628v1}
\bibitem{ref493} Calibrate to Interpret. arXiv:2207.03324v1, 2022. \url{https://arxiv.org/abs/2207.03324v1}
\bibitem{ref494} An Empirical Study of Accuracy, Fairness, Explainability, Distributional Robustness, and Adversarial Robustness. arXiv:2109.14653v1, 2021. \url{https://arxiv.org/abs/2109.14653v1}
\bibitem{ref495} Information Theoretic Evaluation of Privacy-Leakage, Interpretability, and Transferability for Trustworthy AI. arXiv:2106.06046v5, 2021. \url{https://arxiv.org/abs/2106.06046v5}
\bibitem{ref496} Trustworthy AI A Computational Perspective. arXiv:2107.06641v3, 2021. \url{https://arxiv.org/abs/2107.06641v3}
\bibitem{ref497} Connecting the Dots in Trustworthy Artificial Intelligence From AI Principles, Ethics, and Key Requirements to Responsible AI Systems and Regulation. arXiv:2305.02231v2, 2023. \url{https://arxiv.org/abs/2305.02231v2}
\bibitem{ref498} Natural Language Processing for Information Extraction. arXiv:1807.02383v1, 2018. \url{https://arxiv.org/abs/1807.02383v1}
\bibitem{ref499} Information Extraction in Low-Resource Scenarios Survey and Perspective. arXiv:2202.08063v5, 2022. \url{https://arxiv.org/abs/2202.08063v5}
\bibitem{ref500} Zero-Shot Information Extraction via Chatting with ChatGPT. arXiv:2302.10205v1, 2023. \url{https://arxiv.org/abs/2302.10205v1}
\bibitem{ref501} C-ICL Contrastive In-context Learning for Information Extraction. arXiv:2402.11254v1, 2024. \url{https://arxiv.org/abs/2402.11254v1}
\bibitem{ref502} Guideline Learning for In-context Information Extraction. arXiv:2310.05066v2, 2023. \url{https://arxiv.org/abs/2310.05066v2}
\bibitem{ref503} Unified Structure Generation for Universal Information Extraction. arXiv:2203.12277v1, 2022. \url{https://arxiv.org/abs/2203.12277v1}
\bibitem{ref504} Retrieval-Augmented Code Generation for Universal Information Extraction. arXiv:2311.02962v1, 2023. \url{https://arxiv.org/abs/2311.02962v1}
\bibitem{ref505} STAR Boosting Low-Resource Information Extraction by Structure-to-Text Data Generation with Large Language Models. arXiv:2305.15090v3, 2023. \url{https://arxiv.org/abs/2305.15090v3}
\bibitem{ref506} ProgGen Generating Named Entity Recognition Datasets Step-by-step with Self-Reflexive Large Language Models. arXiv:2403.11103v1, 2024. \url{https://arxiv.org/abs/2403.11103v1}
\bibitem{ref507} Enhancing Few-shot Text-to-SQL Capabilities of Large Language Models A Study on Prompt Design Strategies. arXiv:2305.12586v1, 2023. \url{https://arxiv.org/abs/2305.12586v1}
\bibitem{ref508} Structured Prediction as Translation between Augmented Natural Languages. arXiv:2101.05779v3, 2021. \url{https://arxiv.org/abs/2101.05779v3}
\bibitem{ref509} Retrieval-Augmented Generation-based Relation Extraction. arXiv:2404.13397v1, 2024. \url{https://arxiv.org/abs/2404.13397v1}
\bibitem{ref510} Large Language Models for Generative Information Extraction A Survey. arXiv:2312.17617v1, 2023. \url{https://arxiv.org/abs/2312.17617v1}
\bibitem{ref511} Exploring the In-context Learning Ability of Large Language Model for Biomedical Concept Linking. arXiv:2307.01137v1, 2023. \url{https://arxiv.org/abs/2307.01137v1}
\bibitem{ref512} Biomedical knowledge graph-enhanced prompt generation for large language models. arXiv:2311.17330v1, 2023. \url{https://arxiv.org/abs/2311.17330v1}
\bibitem{ref513} LLMs in Biomedicine A study on clinical Named Entity Recognition. arXiv:2404.07376v1, 2024. \url{https://arxiv.org/abs/2404.07376v1}
\bibitem{ref514} MedCAT -- Medical Concept Annotation Tool. arXiv:1912.10166v1, 2019. \url{https://arxiv.org/abs/1912.10166v1}
\bibitem{ref515} Distilling Large Language Models for Biomedical Knowledge Extraction A Case Study on Adverse Drug Events. arXiv:2307.06439v1, 2023. \url{https://arxiv.org/abs/2307.06439v1}
\bibitem{ref516} Knowledge-augmented Graph Neural Networks with Concept-aware Attention for Adverse Drug Event Detection. arXiv:2301.10451v2, 2023. \url{https://arxiv.org/abs/2301.10451v2}
\bibitem{ref517} A Systematic Literature Review of Automated ICD Coding and Classification Systems using Discharge Summaries. arXiv:2107.10652v1, 2021. \url{https://arxiv.org/abs/2107.10652v1}
\bibitem{ref518} TransICD Transformer Based Code-wise Attention Model for Explainable ICD Coding. arXiv:2104.10652v1, 2021. \url{https://arxiv.org/abs/2104.10652v1}
\bibitem{ref519} From Extreme Multi-label to Multi-class A Hierarchical Approach for Automated ICD-10 Coding Using Phrase-level Attention. arXiv:2102.09136v2, 2021. \url{https://arxiv.org/abs/2102.09136v2}
\bibitem{ref520} CoRelation Boosting Automatic ICD Coding Through Contextualized Code Relation Learning. arXiv:2402.15700v1, 2024. \url{https://arxiv.org/abs/2402.15700v1}
\bibitem{ref521} Graph-Based Retriever Captures the Long Tail of Biomedical Knowledge. arXiv:2402.12352v1, 2024. \url{https://arxiv.org/abs/2402.12352v1}
\bibitem{ref522} LLMs Accelerate Annotation for Medical Information Extraction. arXiv:2312.02296v1, 2023. \url{https://arxiv.org/abs/2312.02296v1}
\bibitem{ref523} Diverse Demonstrations Improve In-context Compositional Generalization. arXiv:2212.06800v3, 2022. \url{https://arxiv.org/abs/2212.06800v3}
\bibitem{ref524} Adapt and Decompose Efficient Generalization of Text-to-SQL via Domain Adapted Least-To-Most Prompting. arXiv:2308.02582v3, 2023. \url{https://arxiv.org/abs/2308.02582v3}
\bibitem{ref525} Selective Demonstrations for Cross-domain Text-to-SQL. arXiv:2310.06302v1, 2023. \url{https://arxiv.org/abs/2310.06302v1}
\bibitem{ref526} Improving Demonstration Diversity by Human-Free Fusing for Text-to-SQL. arXiv:2402.10663v1, 2024. \url{https://arxiv.org/abs/2402.10663v1}
\bibitem{ref527} Importance of Synthesizing High-quality Data for Text-to-SQL Parsing. arXiv:2212.08785v1, 2022. \url{https://arxiv.org/abs/2212.08785v1}
\bibitem{ref528} Leveraging Code to Improve In-context Learning for Semantic Parsing. arXiv:2311.09519v2, 2023. \url{https://arxiv.org/abs/2311.09519v2}
\bibitem{ref529} Semantic Decomposition of Question and SQL for Text-to-SQL Parsing. arXiv:2310.13575v1, 2023. \url{https://arxiv.org/abs/2310.13575v1}
\bibitem{ref530} DocuT5 Seq2seq SQL Generation with Table Documentation. arXiv:2211.06193v1, 2022. \url{https://arxiv.org/abs/2211.06193v1}
\bibitem{ref531} Schema-Aware Multi-Task Learning for Complex Text-to-SQL. arXiv:2403.09706v1, 2024. \url{https://arxiv.org/abs/2403.09706v1}
\bibitem{ref532} XRICL Cross-lingual Retrieval-Augmented In-Context Learning for Cross-lingual Text-to-SQL Semantic Parsing. arXiv:2210.13693v1, 2022. \url{https://arxiv.org/abs/2210.13693v1}
\bibitem{ref533} Conversational Text-to-SQL An Odyssey into State-of-the-Art and Challenges Ahead. arXiv:2302.11054v1, 2023. \url{https://arxiv.org/abs/2302.11054v1}
\bibitem{ref534} Going Beyond Word Matching Syntax Improves In-context Example Selection for Machine Translation. arXiv:2403.19285v1, 2024. \url{https://arxiv.org/abs/2403.19285v1}
\bibitem{ref535} Cross-Lingual Retrieval Augmented Prompt for Low-Resource Languages. arXiv:2212.09651v4, 2022. \url{https://arxiv.org/abs/2212.09651v4}
\bibitem{ref536} From Classification to Generation Insights into Crosslingual Retrieval Augmented ICL. arXiv:2311.06595v3, 2023. \url{https://arxiv.org/abs/2311.06595v3}
\bibitem{ref537} Crosslingual Retrieval Augmented In-context Learning for Bangla. arXiv:2311.00587v2, 2023. \url{https://arxiv.org/abs/2311.00587v2}
\bibitem{ref538} Multilingual LLMs are Better Cross-lingual In-context Learners with Alignment. arXiv:2305.05940v3, 2023. \url{https://arxiv.org/abs/2305.05940v3}
\bibitem{ref539} LLMs Are Few-Shot In-Context Low-Resource Language Learners. arXiv:2403.16512v2, 2024. \url{https://arxiv.org/abs/2403.16512v2}
\bibitem{ref540} Enhancing Cross-lingual Sentence Embedding for Low-resource Languages with Word Alignment. arXiv:2404.02490v1, 2024. \url{https://arxiv.org/abs/2404.02490v1}
\bibitem{ref541} Cross-lingual Retrieval for Iterative Self-Supervised Training. arXiv:2006.09526v2, 2020. \url{https://arxiv.org/abs/2006.09526v2}
\bibitem{ref542} Are Structural Concepts Universal in Transformer Language Models Towards Interpretable Cross-Lingual Generalization. arXiv:2310.12794v2, 2023. \url{https://arxiv.org/abs/2310.12794v2}
\bibitem{ref543} Democratizing LLMs for Low-Resource Languages by Leveraging their English Dominant Abilities with Linguistically-Diverse Prompts. arXiv:2306.11372v1, 2023. \url{https://arxiv.org/abs/2306.11372v1}
\bibitem{ref544} Embeddings for Tabular Data A Survey. arXiv:2302.11777v1, 2023. \url{https://arxiv.org/abs/2302.11777v1}
\bibitem{ref545} TabText A Flexible and Contextual Approach to Tabular Data Representation. arXiv:2206.10381v4, 2022. \url{https://arxiv.org/abs/2206.10381v4}
\bibitem{ref546} P-Transformer A Prompt-based Multimodal Transformer Architecture For Medical Tabular Data. arXiv:2303.17408v3, 2023. \url{https://arxiv.org/abs/2303.17408v3}
\bibitem{ref547} UniPredict Large Language Models are Universal Tabular Classifiers. arXiv:2310.03266v2, 2023. \url{https://arxiv.org/abs/2310.03266v2}
\bibitem{ref548} MediTab Scaling Medical Tabular Data Predictors via Data Consolidation, Enrichment, and Refinement. arXiv:2305.12081v2, 2023. \url{https://arxiv.org/abs/2305.12081v2}
\bibitem{ref549} TableGPT Towards Unifying Tables, Nature Language and Commands into One GPT. arXiv:2307.08674v3, 2023. \url{https://arxiv.org/abs/2307.08674v3}
\bibitem{ref550} TablEye Seeing small Tables through the Lens of Images. arXiv:2307.02491v1, 2023. \url{https://arxiv.org/abs/2307.02491v1}
\bibitem{ref551} Prompting Large Language Models for Zero-Shot Clinical Prediction with Structured Longitudinal Electronic Health Record Data. arXiv:2402.01713v2, 2024. \url{https://arxiv.org/abs/2402.01713v2}
\bibitem{ref552} LLMs-based Few-Shot Disease Predictions using EHR A Novel Approach Combining Predictive Agent Reasoning and Critical Agent Instruction. arXiv:2403.15464v1, 2024. \url{https://arxiv.org/abs/2403.15464v1}
\bibitem{ref553} Multimodal LLMs for health grounded in individual-specific data. arXiv:2307.09018v2, 2023. \url{https://arxiv.org/abs/2307.09018v2}
\bibitem{ref554} TABLET Learning From Instructions For Tabular Data. arXiv:2304.13188v1, 2023. \url{https://arxiv.org/abs/2304.13188v1}
\bibitem{ref555} Trompt Towards a Better Deep Neural Network for Tabular Data. arXiv:2305.18446v2, 2023. \url{https://arxiv.org/abs/2305.18446v2}
\bibitem{ref556} Enhancing Activity Prediction Models in Drug Discovery with the Ability to Understand Human Language. arXiv:2303.03363v2, 2023. \url{https://arxiv.org/abs/2303.03363v2}
\bibitem{ref557} Integrating Chemical Language and Molecular Graph in Multimodal Fused Deep Learning for Drug Property Prediction. arXiv:2312.17495v1, 2023. \url{https://arxiv.org/abs/2312.17495v1}
\bibitem{ref558} From molecules to scaffolds to functional groups building context-dependent molecular representation via multi-channel learning. arXiv:2311.02798v1, 2023. \url{https://arxiv.org/abs/2311.02798v1}
\bibitem{ref559} From Artificially Real to Real Leveraging Pseudo Data from Large Language Models for Low-Resource Molecule Discovery. arXiv:2309.05203v3, 2023. \url{https://arxiv.org/abs/2309.05203v3}
\bibitem{ref560} Domain-Adversarial Multi-Task Framework for Novel Therapeutic Property Prediction of Compounds. arXiv:1810.00867v1, 2018. \url{https://arxiv.org/abs/1810.00867v1}
\bibitem{ref561} Integrating Chemistry Knowledge in Large Language Models via Prompt Engineering. arXiv:2404.14467v1, 2024. \url{https://arxiv.org/abs/2404.14467v1}
\bibitem{ref562} Cross-modal representation alignment of molecular structure and perturbation-induced transcriptional profiles. arXiv:1911.10241v2, 2019. \url{https://arxiv.org/abs/1911.10241v2}
\bibitem{ref563} Looking Through Glass Knowledge Discovery from Materials Science Literature using Natural Language Processing. arXiv:2101.01508v1, 2021. \url{https://arxiv.org/abs/2101.01508v1}
\bibitem{ref564} Fine-Grained Chemical Entity Typing with Multimodal Knowledge Representation. arXiv:2108.12899v1, 2021. \url{https://arxiv.org/abs/2108.12899v1}
\bibitem{ref565} Materials Discovery with Extreme Properties via AI-Driven Combinatorial Chemistry. arXiv:2303.11833v1, 2023. \url{https://arxiv.org/abs/2303.11833v1}
\bibitem{ref566} Advancing Extrapolative Predictions of Material Properties through Learning to Learn. arXiv:2404.08657v1, 2024. \url{https://arxiv.org/abs/2404.08657v1}
\bibitem{ref567} Lo-Hi Practical ML Drug Discovery Benchmark. arXiv:2310.06399v1, 2023. \url{https://arxiv.org/abs/2310.06399v1}
\bibitem{ref568} A Hierarchical Neural Framework for Classification and its Explanation in Large Unstructured Legal Documents. arXiv:2309.10563v2, 2023. \url{https://arxiv.org/abs/2309.10563v2}
\bibitem{ref569} Exploring Large Language Models and Hierarchical Frameworks for Classification of Large Unstructured Legal Documents. arXiv:2403.06872v1, 2024. \url{https://arxiv.org/abs/2403.06872v1}
\bibitem{ref570} ICL Markup Structuring In-Context Learning using Soft-Token Tags. arXiv:2312.07405v1, 2023. \url{https://arxiv.org/abs/2312.07405v1}
\bibitem{ref571} A Framework for Explainable Text Classification in Legal Document Review. arXiv:1912.09501v1, 2019. \url{https://arxiv.org/abs/1912.09501v1}
\bibitem{ref572} Text clustering applied to data augmentation in legal contexts. arXiv:2404.08683v1, 2024. \url{https://arxiv.org/abs/2404.08683v1}
\bibitem{ref573} DALE Generative Data Augmentation for Low-Resource Legal NLP. arXiv:2310.15799v1, 2023. \url{https://arxiv.org/abs/2310.15799v1}
\bibitem{ref574} Want To Reduce Labeling Cost GPT-3 Can Help. arXiv:2108.13487v1, 2021. \url{https://arxiv.org/abs/2108.13487v1}
\bibitem{ref575} Parameter-Efficient Legal Domain Adaptation. arXiv:2210.13712v2, 2022. \url{https://arxiv.org/abs/2210.13712v2}
\bibitem{ref576} LLMs in the Loop Leveraging Large Language Model Annotations for Active Learning in Low-Resource Languages. arXiv:2404.02261v1, 2024. \url{https://arxiv.org/abs/2404.02261v1}
\bibitem{ref577} CoAnnotating Uncertainty-Guided Work Allocation between Human and Large Language Models for Data Annotation. arXiv:2310.15638v1, 2023. \url{https://arxiv.org/abs/2310.15638v1}
\bibitem{ref578} Evaluating Document Representations for Content-based Legal Literature Recommendations. arXiv:2104.13841v1, 2021. \url{https://arxiv.org/abs/2104.13841v1}
\bibitem{ref579} ChatLaw Open-Source Legal Large Language Model with Integrated External Knowledge Bases. arXiv:2306.16092v1, 2023. \url{https://arxiv.org/abs/2306.16092v1}
\bibitem{ref580} Datasets for Portuguese Legal Semantic Textual Similarity Comparing weak supervision and an annotation process approaches. arXiv:2306.00007v1, 2023. \url{https://arxiv.org/abs/2306.00007v1}
\bibitem{ref581} Automatic explanation of the classification of Spanish legal judgments in jurisdiction-dependent law categories with tree estimators. arXiv:2404.00437v1, 2024. \url{https://arxiv.org/abs/2404.00437v1}
\bibitem{ref582} FreeAL Towards Human-Free Active Learning in the Era of Large Language Models. arXiv:2311.15614v1, 2023. \url{https://arxiv.org/abs/2311.15614v1}
\bibitem{ref583} Annotator-Centric Active Learning for Subjective NLP Tasks. arXiv:2404.15720v1, 2024. \url{https://arxiv.org/abs/2404.15720v1}
\bibitem{ref584} Which Examples to Annotate for In-Context Learning Towards Effective and Efficient Selection. arXiv:2310.20046v1, 2023. \url{https://arxiv.org/abs/2310.20046v1}
\bibitem{ref585} STENCIL Submodular Mutual Information Based Weak Supervision for Cold-Start Active Learning. arXiv:2402.13468v1, 2024. \url{https://arxiv.org/abs/2402.13468v1}
\bibitem{ref586} Human Still Wins over LLM An Empirical Study of Active Learning on Domain-Specific Annotation Tasks. arXiv:2311.09825v1, 2023. \url{https://arxiv.org/abs/2311.09825v1}
\bibitem{ref587} Active Learning for NLP with Large Language Models. arXiv:2401.07367v1, 2024. \url{https://arxiv.org/abs/2401.07367v1}
\bibitem{ref588} Curated LLM Synergy of LLMs and Data Curation for tabular augmentation in ultra low-data regimes. arXiv:2312.12112v2, 2023. \url{https://arxiv.org/abs/2312.12112v2}
\bibitem{ref589} Ensemble-Instruct Generating Instruction-Tuning Data with a Heterogeneous Mixture of LMs. arXiv:2310.13961v1, 2023. \url{https://arxiv.org/abs/2310.13961v1}
\bibitem{ref590} TeaMs-RL Teaching LLMs to Teach Themselves Better Instructions via Reinforcement Learning. arXiv:2403.08694v1, 2024. \url{https://arxiv.org/abs/2403.08694v1}
\bibitem{ref591} ELAD Explanation-Guided Large Language Models Active Distillation. arXiv:2402.13098v1, 2024. \url{https://arxiv.org/abs/2402.13098v1}
\bibitem{ref592} Data-efficient Active Learning for Structured Prediction with Partial Annotation and Self-Training. arXiv:2305.12634v2, 2023. \url{https://arxiv.org/abs/2305.12634v2}
\bibitem{ref593} ALE A Simulation-Based Active Learning Evaluation Framework for the Parameter-Driven Comparison of Query Strategies for NLP. arXiv:2308.02537v1, 2023. \url{https://arxiv.org/abs/2308.02537v1}
\bibitem{ref594} Offline Meta-Reinforcement Learning for Industrial Insertion. arXiv:2110.04276v4, 2021. \url{https://arxiv.org/abs/2110.04276v4}
\bibitem{ref595} Hindsight Task Relabelling Experience Replay for Sparse Reward Meta-RL. arXiv:2112.00901v1, 2021. \url{https://arxiv.org/abs/2112.00901v1}
\bibitem{ref596} Online Meta-Learning. arXiv:1902.08438v4, 2019. \url{https://arxiv.org/abs/1902.08438v4}
\bibitem{ref597} Fully Online Meta-Learning Without Task Boundaries. arXiv:2202.00263v3, 2022. \url{https://arxiv.org/abs/2202.00263v3}
\bibitem{ref598} Algorithm Design for Online Meta-Learning with Task Boundary Detection. arXiv:2302.00857v1, 2023. \url{https://arxiv.org/abs/2302.00857v1}
\bibitem{ref599} Dynamic Regret Analysis for Online Meta-Learning. arXiv:2109.14375v1, 2021. \url{https://arxiv.org/abs/2109.14375v1}
\bibitem{ref600} CoMPS Continual Meta Policy Search. arXiv:2112.04467v1, 2021. \url{https://arxiv.org/abs/2112.04467v1}
\bibitem{ref601} Meta-Reinforcement Learning by Tracking Task Non-stationarity. arXiv:2105.08834v1, 2021. \url{https://arxiv.org/abs/2105.08834v1}
\bibitem{ref602} Skill-based Meta-Reinforcement Learning. arXiv:2204.11828v1, 2022. \url{https://arxiv.org/abs/2204.11828v1}
\bibitem{ref603} The Learnability of In-Context Learning. arXiv:2303.07895v1, 2023. \url{https://arxiv.org/abs/2303.07895v1}
\bibitem{ref604} Uncontrolled Lexical Exposure Leads to Overestimation of Compositional Generalization in Pretrained Models. arXiv:2212.10769v1, 2022. \url{https://arxiv.org/abs/2212.10769v1}
\bibitem{ref605} How Do In-Context Examples Affect Compositional Generalization. arXiv:2305.04835v3, 2023. \url{https://arxiv.org/abs/2305.04835v3}
\bibitem{ref606} Revisiting the Compositional Generalization Abilities of Neural Sequence Models. arXiv:2203.07402v1, 2022. \url{https://arxiv.org/abs/2203.07402v1}
\bibitem{ref607} Measuring Compositional Generalization A Comprehensive Method on Realistic Data. arXiv:1912.09713v2, 2019. \url{https://arxiv.org/abs/1912.09713v2}
\bibitem{ref608} Generalization in Multimodal Language Learning from Simulation. arXiv:2108.02319v1, 2021. \url{https://arxiv.org/abs/2108.02319v1}
\bibitem{ref609} Towards Understanding the Relationship between In-context Learning and Compositional Generalization. arXiv:2403.11834v1, 2024. \url{https://arxiv.org/abs/2403.11834v1}
\bibitem{ref610} Skills-in-Context Prompting Unlocking Compositionality in Large Language Models. arXiv:2308.00304v2, 2023. \url{https://arxiv.org/abs/2308.00304v2}
\bibitem{ref611} Automatically Composing Representation Transformations as a Means for Generalization. arXiv:1807.04640v2, 2018. \url{https://arxiv.org/abs/1807.04640v2}
\bibitem{ref612} Laying the Foundation First Investigating the Generalization from Atomic Skills to Complex Reasoning Tasks. arXiv:2403.09479v1, 2024. \url{https://arxiv.org/abs/2403.09479v1}
\bibitem{ref613} Explainable Neural Computation via Stack Neural Module Networks. arXiv:1807.08556v3, 2018. \url{https://arxiv.org/abs/1807.08556v3}
\bibitem{ref614} Compositional Generalization from First Principles. arXiv:2307.05596v1, 2023. \url{https://arxiv.org/abs/2307.05596v1}
\bibitem{ref615} Meta-Learning to Compositionally Generalize. arXiv:2106.04252v2, 2021. \url{https://arxiv.org/abs/2106.04252v2}
\bibitem{ref616} Simple and effective data augmentation for compositional generalization. arXiv:2401.09815v1, 2024. \url{https://arxiv.org/abs/2401.09815v1}
\bibitem{ref617} LongLLMLingua Accelerating and Enhancing LLMs in Long Context Scenarios via Prompt Compression. arXiv:2310.06839v1, 2023. \url{https://arxiv.org/abs/2310.06839v1}
\bibitem{ref618} LLMLingua-2 Data Distillation for Efficient and Faithful Task-Agnostic Prompt Compression. arXiv:2403.12968v1, 2024. \url{https://arxiv.org/abs/2403.12968v1}
\bibitem{ref619} Training With Paraphrasing the Original Text'' Improves Long-Context Performance. arXiv:2312.11193v8, 2023. \url{https://arxiv.org/abs/2312.11193v8}
\bibitem{ref620} In-context Autoencoder for Context Compression in a Large Language Model. arXiv:2307.06945v3, 2023. \url{https://arxiv.org/abs/2307.06945v3}
\bibitem{ref621} LLoCO Learning Long Contexts Offline. arXiv:2404.07979v1, 2024. \url{https://arxiv.org/abs/2404.07979v1}
\bibitem{ref622} Adapting LLMs for Efficient Context Processing through Soft Prompt Compression. arXiv:2404.04997v2, 2024. \url{https://arxiv.org/abs/2404.04997v2}
\bibitem{ref623} Prompt Cache Modular Attention Reuse for Low-Latency Inference. arXiv:2311.04934v2, 2023. \url{https://arxiv.org/abs/2311.04934v2}
\bibitem{ref624} CAMELoT Towards Large Language Models with Training-Free Consolidated Associative Memory. arXiv:2402.13449v1, 2024. \url{https://arxiv.org/abs/2402.13449v1}
\bibitem{ref625} Augmenting Language Models with Long-Term Memory. arXiv:2306.07174v1, 2023. \url{https://arxiv.org/abs/2306.07174v1}
\bibitem{ref626} Invertible Attention. arXiv:2106.09003v2, 2021. \url{https://arxiv.org/abs/2106.09003v2}
\bibitem{ref627} Leave No Context Behind Efficient Infinite Context Transformers with Infini-attention. arXiv:2404.07143v1, 2024. \url{https://arxiv.org/abs/2404.07143v1}
\bibitem{ref628} Zebra Extending Context Window with Layerwise Grouped Local-Global Attention. arXiv:2312.08618v1, 2023. \url{https://arxiv.org/abs/2312.08618v1}
\bibitem{ref629} Focused Transformer Contrastive Training for Context Scaling. arXiv:2307.03170v2, 2023. \url{https://arxiv.org/abs/2307.03170v2}
\bibitem{ref630} Lost in the Middle How Language Models Use Long Contexts. arXiv:2307.03172v3, 2023. \url{https://arxiv.org/abs/2307.03172v3}
\bibitem{ref631} Found in the Middle How Language Models Use Long Contexts Better via Plug-and-Play Positional Encoding. arXiv:2403.04797v1, 2024. \url{https://arxiv.org/abs/2403.04797v1}
\bibitem{ref632} Research Re search \& Re-search. arXiv:2403.13705v1, 2024. \url{https://arxiv.org/abs/2403.13705v1}
\bibitem{ref633} Can't Remember Details in Long Documents You Need Some R\&R. arXiv:2403.05004v1, 2024. \url{https://arxiv.org/abs/2403.05004v1}
\bibitem{ref634} Make Your LLM Fully Utilize the Context. arXiv:2404.16811v2, 2024. \url{https://arxiv.org/abs/2404.16811v2}
\bibitem{ref635} LongAgent Scaling Language Models to 128k Context through Multi-Agent Collaboration. arXiv:2402.11550v2, 2024. \url{https://arxiv.org/abs/2402.11550v2}
\bibitem{ref636} SHACL A Description Logic in Disguise. arXiv:2108.06096v3, 2021. \url{https://arxiv.org/abs/2108.06096v3}
\bibitem{ref637} Scaling In-Context Demonstrations with Structured Attention. arXiv:2307.02690v1, 2023. \url{https://arxiv.org/abs/2307.02690v1}
\bibitem{ref638} UniMem Towards a Unified View of Long-Context Large Language Models. arXiv:2402.03009v1, 2024. \url{https://arxiv.org/abs/2402.03009v1}
\bibitem{ref639} Superposition Prompting Improving and Accelerating Retrieval-Augmented Generation. arXiv:2404.06910v1, 2024. \url{https://arxiv.org/abs/2404.06910v1}
\bibitem{ref640} Meta-Transformer A Unified Framework for Multimodal Learning. arXiv:2307.10802v1, 2023. \url{https://arxiv.org/abs/2307.10802v1}
\bibitem{ref641} OFA Unifying Architectures, Tasks, and Modalities Through a Simple Sequence-to-Sequence Learning Framework. arXiv:2202.03052v2, 2022. \url{https://arxiv.org/abs/2202.03052v2}
\bibitem{ref642} Cross-Modal Generalization Learning in Low Resource Modalities via Meta-Alignment. arXiv:2012.02813v1, 2020. \url{https://arxiv.org/abs/2012.02813v1}
\bibitem{ref643} i-Code V2 An Autoregressive Generation Framework over Vision, Language, and Speech Data. arXiv:2305.12311v1, 2023. \url{https://arxiv.org/abs/2305.12311v1}
\bibitem{ref644} OneLLM One Framework to Align All Modalities with Language. arXiv:2312.03700v1, 2023. \url{https://arxiv.org/abs/2312.03700v1}
\end{thebibliography}

\end{document}
